\PassOptionsToPackage{unicode}{hyperref}
\PassOptionsToPackage{hyphens}{url}
\PassOptionsToPackage{dvipsnames,svgnames,x11names}{xcolor}
\documentclass[american,11pt,letterpaper]{article}
\usepackage{xcolor}
\usepackage[margin=1in]{geometry}
\usepackage{amsmath,amssymb,amsthm,mathtools}
\usepackage{iftex}
\ifPDFTeX
  \usepackage[T1]{fontenc}
  \usepackage[utf8]{inputenc}
  \usepackage{textcomp}
\else
  \usepackage{fontspec}
  \defaultfontfeatures{Scale=MatchLowercase}
  \defaultfontfeatures[\rmfamily]{Ligatures=TeX,Scale=1}
\fi
\usepackage{lmodern}
\ifPDFTeX\else
  \setmainfont[]{Times New Roman}
\fi
\IfFileExists{upquote.sty}{\usepackage{upquote}}{}
\IfFileExists{microtype.sty}{%
  \usepackage[]{microtype}
  \UseMicrotypeSet[protrusion]{basicmath}
}{}
\usepackage{longtable,booktabs,array}
\usepackage{calc}
\usepackage{etoolbox}
\makeatletter
\patchcmd\longtable{\par}{\if@noskipsec\mbox{}\fi\par}{}{}
\makeatother
\IfFileExists{footnotehyper.sty}{\usepackage{footnotehyper}}{\usepackage{footnote}}
\makesavenoteenv{longtable}
\ifLuaTeX
  \usepackage[bidi=basic,shorthands=off]{babel}
\else
  \usepackage[bidi=default,shorthands=off]{babel}
\fi
\ifPDFTeX\else
  \babelfont{rm}[]{Times New Roman}
\fi
\ifLuaTeX
  \usepackage{selnolig}
\fi
\usepackage{fancyhdr}
\usepackage{enumitem}
\usepackage{caption}
\usepackage{bm}
\usepackage{float}
\usepackage{graphicx}
\usepackage{appendix}
\usepackage[hyphens]{xurl}
\usepackage{ragged2e}

\pagestyle{fancy}
\fancyhf{}
\fancyfoot[C]{\thepage}
\renewcommand{\headrulewidth}{0pt}

\setlength{\parindent}{0pt}
\setlength{\parskip}{6pt plus 2pt minus 1pt}
\setlist{nosep,leftmargin=*}
\captionsetup{font=small,labelfont=bf}
\setlength{\emergencystretch}{3em}
\sloppy
\hyphenation{im-ple-men-ta-tion op-er-a-tion-al-i-za-tion}
\renewcommand{\familydefault}{\rmdefault}
\AtBeginEnvironment{longtable}{\small}
\urlstyle{same}
\usepackage{hyperref}
\usepackage{bookmark}

\hypersetup{
  pdftitle={Super Information Theory: The Coherence Conservation Law Unifying the Wave Function, Gravity, and Time (Version 81)},
  pdfauthor={Micah Blumberg},
  pdflang={en-US},
  colorlinks=true,
  linkcolor={blue},
  filecolor={Maroon},
  citecolor={blue},
  urlcolor={blue}
}

\newtheorem{result}{Result}
\newtheorem{prediction}{Prediction}
\newtheorem{theorem}{Theorem}[section]
\newtheorem{lemma}[theorem]{Lemma}
\newtheorem{proposition}[theorem]{Proposition}
\newtheorem{corollary}[theorem]{Corollary}
\newtheorem{definition}[theorem]{Definition}
\newtheorem{remark}[theorem]{Remark}
\newcommand{\Rcoh}{R_{\mathrm{coh}}}
\newcommand{\Scoh}{S_{\mathrm{coh}}}
\newcommand{\Ecoh}{E_{\mathrm{coh}}}
\newcommand{\Edec}{E_{\mathrm{dec}}}
\newcommand{\bra}[1]{\langle #1 \rvert}
\newcommand{\ket}[1]{\lvert #1 \rangle}
\newcommand{\braket}[2]{\langle #1 \mid #2 \rangle}
\newcommand{\blockquote}[1]{\begin{quote}#1\end{quote}}
\newcommand{\SITConditionalFigure}[4]{%
  \begin{figure}[htbp]
  \centering
  \IfFileExists{figures/#1.pdf}{%
    \includegraphics[width=0.92\textwidth]{figures/#1.pdf}%
  }{%
    \IfFileExists{../../simulations/figures/#1.pdf}{%
      \includegraphics[width=0.92\textwidth]{../../simulations/figures/#1.pdf}%
    }{%
      \fbox{\begin{minipage}{0.90\textwidth}
      \vspace{0.5em}
      \textbf{Numerical Verification Figure Panel (\texttt{figures/#1.pdf}):}\\[0.35em]
      #4
      \vspace{0.5em}
      \end{minipage}}%
    }%
  }%
  \caption{#2}
  \label{#3}
  \end{figure}%
}
\pdfstringdefDisableCommands{%
  \def\mathbb#1{#1}%
  \def\mathbf#1{#1}%
  \def\mathcal#1{#1}%
  \def\mathrm#1{#1}%
  \def\texttt#1{#1}%
  \def\emph#1{#1}%
  \def\Rcoh{R_coh}%
  \def\Scoh{S_coh}%
  \def\Ecoh{E_coh}%
  \def\Edec{E_dec}%
  \def\bra#1{<#1|}%
  \def\ket#1{|#1>}%
  \def\braket#1#2{<#1|#2>}%
}
\author{}
\date{}

\begin{document}


\begin{center}
{\LARGE\bfseries Super Information Theory\par}
\vspace{0.35em}
{\large\itshape The Coherence Conservation Law Unifying the Wave Function, Gravity, and Time\\(Version 81)\par}
\vspace{0.8em}
{\small\textbf{Author: Micah Blumberg}\par}
{\small Self Aware Networks Research Institute\par}
{\small\url{https://selfawarenetworks.com}\par}
{\small Original Publication Date: February 9, 2025\par}
{\small Version 81 Updated on November 19, 2026\par}
\end{center}
\vspace{0.65em}

% --- BEGIN INPUT sections/00_abstract_and_note.tex ---
\begin{abstract}
We formulate \emph{Super Information Theory} (SIT) as a generally covariant effective field theory in mostly-plus Lorentzian signature $(-,+,+,+)$ in which physical reality is described at low and mesoscopic energies ($E \ll \Lambda_{\mathrm{UV}}$) by two primitive dynamical informational fields alongside spacetime geometry, gauge connections, and Standard Model matter: a complex coherence order-parameter field $\psi(x) = |\psi(x)| e^{i\theta(x)}$ with dimensionless coherence ratio $\Rcoh(x) := |\psi(x)|/\psi_0 \ge 0$, and a real temporal-density scalar $\rho_t(x)$ governing local operational clock rates. In its authoritative consolidated November 2026 formulation (\texttt{M\_AUTH\_SIT\_EFT\_v1}, canon \texttt{v1\_1}), designated \textbf{Version~81} to preserve accurate version chronology after the July~2026 compact corrective \emph{Version~80} lineage (Paper~002, Zenodo DOI \texttt{10.5281/zenodo.21693867}, July~29, 2026; Version~80 Draft~13, July~26, 2026), dynamics are governed by the single covariant action
\begin{align}
S_{\mathrm{AUTH}}[g,\rho_t,\psi,A_\mu,\Psi_{\mathrm{SM}}]
&= \int d^4x\,\sqrt{-g}\Biggl[
  \frac{R}{16\pi G}
  - \frac{\kappa_t}{2}\,g^{\mu\nu}(\partial_\mu\rho_t)(\partial_\nu\rho_t)
  - \frac{\kappa_c}{2}\,g^{\mu\nu}(D_\mu\psi)^*(D_\nu\psi)
  \nonumber\\
&\qquad
  - V_{\mathrm{tot}}(\rho_t,|\psi|)
  - \frac{1}{4}\,f_{\mathrm{EM}}(\Rcoh)\,F_{\mu\nu}F^{\mu\nu}
  + f_{\mathrm{SM}}(\Rcoh)\,\mathcal{L}_{\mathrm{matter}}
\Biggr],
\label{eq:abstract_master_action}
\end{align}
with gauge-covariant derivative $D_\mu\psi = (\nabla_\mu - i\hat{q}A_\mu)\psi$, strictly vacuum-positive constitutive dressings
\[
f_{\mathrm{EM}}(\Rcoh) := 1 + c_{\mathrm{EM}}\Rcoh^2 \ge 1,
\qquad
f_{\mathrm{SM}}(\Rcoh) := 1 + \frac{\Rcoh^2}{2 M_{\mathrm{sym}}^2} \ge 1,
\]
and globally non-negative $\mathbb{Z}_2$-lifted scalar potential $V_{\mathrm{tot}}(\rho_t,|\psi|) = V(\rho_t) + U(|\psi|) + U_{\mathrm{link}}(\rho_t,|\psi|) \ge 0$ attaining zero vacuum energy uniquely at $(\rho_0, v_0)$ while lifting $(-\rho_0, v_0)$ by $2 m_{\mathrm{link}}^2 \rho_0^2 > 0$. Because $f_{\mathrm{EM}} \ge 1$ and $f_{\mathrm{SM}} \ge 1$ hold globally, the continuum effective theory is free of zero-coherence $1/f$ Hamiltonian singularities and propagates exactly $7$ ghost-free bosonic polarizations in pure-SIT vacuum ($8$ modes when coupled to a representative neutral matter scalar in curved $3+1$D ADM phase space of $36 - 2\times 10 = 16$ physical dimensions). Stage-6 action-response diagnostics (\texttt{Stage6FunctionalResponse.\allowbreak{}ActionIdentityDiagnostics}, $6$ Lean~4.30 theorems) prove that the unchanged multiplicative portal $+\frac{|\psi|^2}{\Lambda_{\mathrm{sym}}^2}\mathcal{L}_{\mathrm{matter}}$ does not algebraically match a four-dimensional conformal Jordan metric $\tilde{g}_{\mu\nu} = A^2 g_{\mu\nu}$ unless $f_{\mathrm{SM}} = 1$; accordingly, Version~81 explicitly separates the November~2026 multiplicative action \texttt{M\_AUTH\_SIT\_EFT\_v1} from the July~2026 Version~80 conformal two-real-scalar benchmark (Paper~002), qualifying Jordan-frame identities $\Phi_J + \Psi_J = 2\Phi_E$ and $|\gamma_{\mathrm{PPN}} - 1| \le 2\alpha_{\mathrm{conf}}[A]\,S(r)/(1 + \alpha_{\mathrm{conf}}[A]\,S(r))$ as conditional on adopting a consistently normalized conformal extension rather than automatic consequences of \texttt{M\_AUTH\_SIT\_EFT\_v1}.

Invariance under the global $U(1)$ phase subgroup implies a conserved gauge-invariant coherence current
\begin{equation}
J^\mu_{\mathrm{coh}}
= \frac{i\kappa_c}{2}\bigl(\psi(D^\mu\psi)^* - \psi^*(D^\mu\psi)\bigr)
= \kappa_c |\psi|^2\bigl(\nabla^\mu\theta - \hat{q}A^\mu\bigr),
\label{eq:abstract_coherence_current}
\end{equation}
reducing to $\kappa_c|\psi|^2\nabla^\mu\theta$ when $A_\mu \to 0$. Directly from the unchanged continuum action \texttt{M\_AUTH\_SIT\_EFT\_v1}, Stage~6 (\texttt{UnchangedActionWeakFieldObservable}, $5$ Lean~4.30 theorems) performs a \emph{conditional weak-field E\"ot-Wash non-exclusion calculation}: under a declared pressureless-dust worldline closure ($q_{\mathrm{eff}} = 1$, contrasted with homogeneous free Klein--Gordon oscillation where $\langle\mathcal{L}_m\rangle = 0$ gives $q_{\mathrm{eff}} = 0$, neither deriving microscopic composite-material charges), canonical benchmark \texttt{BM\_PHYS\_EFT\_1} yields light-mode range $\lambda_1 \approx 8.056\times 10^{-31}\,\mathrm{m}$ ($\alpha_1 \approx 4.745\times 10^{-7}$), giving negligible force ratio at Lee~\emph{et~al.}'s $52\,\mu\mathrm{m}$ separation ($\log_{10} R_{\mathrm{slab}} \approx -2.80\times 10^{25}$), while a deliberately tuned $m_1 = 5.112\,\mathrm{meV}$ uniform-slab illustration gives $2.80\%$ normal-force contrast (distinct from patterned rotating-mass E\"ot-Wash torques, whose full run-by-run reproduction from Lee's dissertation Chapter~11 Table~11.4 and Appendix~A remains an open next step). Version~81 incorporates both the Stage-6/v80.3 repairs---disambiguating the local clock response $\alpha_\rho$ ($y_{\mathrm{clock}} = \alpha_\rho r$, $r = \delta\rho_t/\rho_0$) from the gravitational-potential redshift slope $\alpha_\Phi := d\ln\nu/d(\Phi/c^2) = \alpha_\rho\mathcal{T}$, repairing linearized scalar equations in Appendix~E (with diagonal coordinate congruence $D = \mathrm{diag}(1, \psi_0)$ to Stage~6), withdrawing the flat-connection $A_\mu = (\hbar/e)\partial_\mu\theta$ magnetism construction and legacy CRT $10^{-8}$ figure, removing unsupported laboratory forecast amplitudes and universal falsification criteria, anchoring six primary experimental provenance records in \texttt{primary\_measurements\_manifest\_v2.json} (including Takamoto~\emph{et~al.}'s Tokyo Skytree combined measurement $(+1.4 \pm 9.1)\times 10^{-5}$, NIST/BACON's \emph{Phys.\ Rev.\ Lett.}\ \textbf{137}, 033201 optical frequency-ratio measurement uncertainties $\le 3.2\times 10^{-18}$ across phase-stabilized optical fiber links, and the corrected Lee~\emph{et~al.}\ 18-fold/120-fold patterned platinum--epoxy apparatus description), and quarantining the legacy $3\times 10^{-8}$ clock target, $29.4\,\mu\mathrm{Gal}$ BEC ceiling, and hand-coded \texttt{np.piecewise} Yukawa sketch---and the bidirectional mathematical and experimental controls recovered from the July~2026 intervening papers (retracting Version~79's false universal CPTP purity-decay theorem via amplitude-damping and complete-dephasing counterexamples in favor of unital CPTP purity non-increase and general CPTP relative-entropy contraction, separating scalar field amplitudes from quantum purity/condensate-fraction/visibility proxies via open calibration kernels, and integrating Paper~002/003/007/014/055 mass-calibration, form-factor identifiability, matched-preparation residual, and implication-boundary guardrails).

Beyond classical and perturbative effective field theory below $\Lambda_{\mathrm{UV}} := \min(M_{\mathrm{Pl}}, \Lambda_{\mathrm{EM}}, \Lambda_{\mathrm{sym}})$, Version~81 incorporates a machine-checked formalization in Lean~4.30 (verified strictly under standard foundational axioms \texttt{[propext, Classical.\allowbreak{}choice, Quot.\allowbreak{}sound]} with zero \texttt{sorry}, zero \texttt{admit}, and zero custom axioms across $660$ frozen Stage~1--5 \texttt{.lean} files, $6$ frozen Stage-6 modules in \texttt{Stage6FunctionalResponse} auditing $30/30$ theorems in \texttt{SITStage6Audit}, and $3$ publication-branch addendum modules in \texttt{Stage6Addendum}---periodic lattice Fourier eigenvalues [$3$], sharp $\frac{3\beta}{4}$ nonnegative-field coercivity allowing zero corners with both equality witnesses [$10$], and canonical action-response propagator inverse $(p^2 K + H)^{-1} = V\,\mathrm{diag}(1/(p^2+m_a^2))\,V^\top$ with coordinate-rescaling equivalence and direct instantiation on \texttt{AuthModelParams} [$5$]---totaling $669$ \texttt{.lean} files and $48$ Stage-6/Addendum audited theorems in Appendix~\ref{app:epistemic-ledger-lean}) together with SymPy/NumPy dual-engine symbolic verification and exploratory finite-volume 4D $SU(2)$ Markov-chain Monte Carlo consistency benchmarks. Each formal theorem addresses a specific, bounded mathematical question; aggregate module or theorem counts never substitute for an unproved matter-source derivation, an interacting quantum continuum limit, or an observational confirmation. On a non-perturbative 4D hypercubic Euclidean Wilson lattice with compact Haar measure, unified Model~L3 (\texttt{sitAMGMGatedOSAction}) replaces the constant plaquette baseline of Model~L1 by the four-corner arithmetic-mean counterterm $\beta\,\mathrm{amgmCoherenceGap} \ge 0$ (differing from Model~L1 by $\beta\,\mathrm{amgmCoherenceGap} + \kappa_t\sum_{x,\mu}(\rho(x+e_\mu)-\rho(x))^2$, isolating the pure AM--GM gap at $\kappa_t = 0$), simultaneously proving $0 < Z_{\mathrm{L3}} \le 1$, exact zero-matter ($\delta_0 \otimes \mu_{\mathrm{temp}}$ and conditional $\psi = 0$) $SU(2)$ Haar disorder for vanishing gates $f_{\mathrm{gate}}(\rho,0) = 0$ (with zero single-plaquette expectation for $l \ge 2$), identical conditional gauge law $P(U \mid \rho_t,\psi)$, and coupled $(U,\psi,\rho_t)$ Osterwalder--Schrader $2L_t$ reflection positivity across compact-ball and unbounded/Gaussian site priors. Stage~6 plus the addenda complete the \emph{classical} lattice-to-continuum functional, spectral, and nonnegative-coercivity analysis directly on $\mathrm{Site}\,l = (\mathrm{Fin}\,4 \to \mathrm{ZMod}\,l)$: proving exact finite-lattice corner Hessian cancellation of nonlinear second derivatives $j_i$, proving that periodic lattice Fourier modes satisfy the exact eigenvalue equation $(H_{\mathrm{lat}} u)(x) = a^2\lambda_Q(k;a)\,u(x)$ ($a = 1/l$) obeying $0 \le 3\beta|k|^2/(4f_0) - \lambda_Q(k;a) \le \beta a^2|k|^4/(16f_0)$, proving $O(1/l^2)$ convergence of the rescaled 4-corner energy $l^{-2}\beta\,\mathrm{amgmCoherenceGap}$ to $\frac{3\beta}{8}\int_{[0,1]^4}\frac{|\nabla f|^2}{f}\,du$ for every conventional $C^4$ (\texttt{ContDiff} $\mathbb{R}\;4$) strictly positive 1-periodic field $f$, proving the sharp zero-division nonnegative-field coercivity bound $Q_l(F) \ge \frac{3\beta}{4} a^2\sum_{x,\mu}(\sqrt{F(x+e_\mu)}-\sqrt{F(x)})^2$ allowing zero corners (strictly strengthening the valid weaker $\frac{\beta}{2}$ bound from the sharp $1/12$ six-pair corner inequality, attained identically on both a nontrivial single-site spike $F(0)=s>0$ for $l \ge 2$ and a nontrivial strictly positive parity-alternating checkerboard $F(x)\in\{A^2,B^2\}$ with $A,B>0, A \ne B$ for even $l$, while noting that together with a uniform bound on $a^4\sum_x F(x)$ this yields strong $L^2(\mathbb{T}^4)$ compactness of $\sqrt{F_l}$ with weak-derivative $H^1(\mathbb{T}^4)$ control and hence strong $L^1(\mathbb{T}^4)$ compactness of $F_l$ for future variational $\Gamma$-convergence), and proving exact finite-lattice embedding into Stage-5 Model~L3 at $\kappa_t = 0$ (explicitly qualifying $\kappa_t = 0$ isolation and affine $1 + c|\psi|^2$ positivity vs.\ vanishing Haar-disorder gates). All open mathematical and observational frontiers---ultraviolet quantum gravity at $E \ge \Lambda_{\mathrm{UV}}$ (\texttt{OPEN-1}), quantum Wheeler--DeWitt ordering and Sobolev/Regge--Teitelboim curved Cauchy PDE (\texttt{OPEN-2}), interacting quantum continuum-measure convergence, OS transfer-operator Hamiltonian construction, unmodified 4D pure Yang--Mills mass gap, and unbounded 3D $L^2(\mathbb{R}^3)$ PDEs (\texttt{OPEN-3}), the conditionally isolated Riemann spectral bridge (\texttt{OPEN-4}), massless Weyl determinant phases, three-generation soliton homotopy, and $P \ne NP$ (\texttt{OPEN-5}), and the microphysical action-to-apparatus bridge of composite scalar charges, visible $U(1)$ photon dictionary, laboratory $J_{\mathrm{prep}}$, and patterned rotating-mass E\"ot-Wash likelihoods (\texttt{OPEN-6})---are demarcated explicitly as unsolved open problems (\texttt{NOT\_SOLVED}) under the five-tier epistemic ledger \texttt{QG49}. Version~81 is explicitly presented as a working research monograph rather than a claim of experimental confirmation or final scientific completion.
\end{abstract}

\subsection*{Author's Note on Version 81, Version Chronology, and Model Continuity (November 19, 2026)}
\addcontentsline{toc}{subsection}{Author's Note on Version 81, Version Chronology, and Model Continuity (November 19, 2026)}

\textbf{Version Chronology and Why This Monograph Is Designated Version~81.} Accurate public chronology separates three distinct stages of the SIT monograph lineage rather than collapsing them into a single undifferentiated version number:
\begin{enumerate}[leftmargin=2em]
  \item \textbf{Version~79 (December~18, 2025):} The prior full-length monograph baseline whose conceptual architecture and 2017--2025 priority documentation are preserved for archival continuity.
  \item \textbf{July~2026 Version~80 Compact Corrective Lineage (Papers~002, 048, and Specialized July--October~2026 Portfolio):} The public Version~80 designation is already occupied by the July~2026 compact corrective lineage: Paper~048 (\emph{Causal Fermion Systems and Super Information Theory}, July~28, 2026) cites its frozen comparison source as \emph{Version~80 Draft~13} dated July~26, 2026, and Paper~002 (\emph{Super Information Theory: Coherence, Time Density, Informational Dissipation, and Gravity}, 26~pages, published July~29, 2026, Zenodo DOI \texttt{10.5281/zenodo.21693867}) is explicitly cataloged as the correction-integrated Version~80 successor to Version~79. Across the same July--October~2026 window, specialized companion notes developed bounded sub-models and controls: Paper~003 (\emph{Level~0 SIT}, Penrose--Onsager condensate-source contract), Paper~007 (\emph{Super Dark Time}), Paper~014 (\emph{Designing a Coherence-Gravity Test with Bose--Einstein Condensates}), Paper~028 (\emph{Cross-Disciplinary SIT Translation Guide}), Paper~033 (finite 2D abelian $U(1)$ lattice toy model), Paper~034 (\emph{SuperTimePosition and Gravitational Phase-Lock}), Paper~055 (\emph{QGTCD Constraints Note}), Papers~062--063 (\emph{Physics in Review} comparisons), and Paper~114 (\emph{Learning to Feel}: initial edition October~17, 2026; current main paper v1.2 dated October~20, 2026, with technical supplement intentionally retaining its v1.1 label; pursuing an independent neuroscience track whose neural synchrony observables are not identified with canonical SIT spacetime fields).
  \item \textbf{November~2026 Version~81 Full Monograph and Machine-Checked Formalization Series:} Because Version~80 is already taken by the July~2026 compact corrective lineage, this November~2026 full-length monograph---advancing the frozen November internal working drafts v80.2 (Stage~5.2, November~16, 2026, commit \texttt{8955eab}) and v80.3 (Stage~6 plus three Lean~4.30 addenda, Hydrogen house-style alignment, and magnetism/Part~III/lattice consistency repairs, November~18--19, 2026, commit \texttt{30a71e4}) while integrating bidirectional mathematical and experimental controls recovered from the July~2026 intervening papers---is designated \textbf{Super Information Theory: Version~81} (accompanied by the November~2026 Companion Papers~I--III, Open Frontiers, and standalone lattice paper). Version~81 is explicitly a working research monograph rather than a claim of final scientific confirmation.
\end{enumerate}

\textbf{Model-Continuity Contract: July~2026 Version~80 Benchmark vs.\ November~2026 Version~81 Authoritative Action.} The July~2026 compact corrective lineage (Paper~002/003) and the November~2026 authoritative continuum/lattice formalization (\texttt{M\_AUTH\_SIT\_EFT\_v1} and 4D lattice Models~L1--L3) share research questions, but their actions, field content, and matter couplings are \emph{not} interchangeable (Table~\ref{tab:july-v80-vs-nov-v81}). Every result carried between models must declare its original action, destination action, background, matter sector, observable, and exact relation type (exact field redefinition, consistent truncation, controlled approximation, or shared quadratic response calculation); any formula lacking such a bridge remains attached strictly to its originating model. In particular, Stage~6's proof that the multiplicative portal $+f_{\mathrm{SM}}(\Rcoh)\mathcal{L}_{\mathrm{matter}}$ does not algebraically match a conformal metric unless $f_{\mathrm{SM}} = 1$ does \emph{not} invalidate Paper~002's explicitly conformal July benchmark, and conversely Paper~002's conformal scalar-tensor formulas do \emph{not} automatically become predictions of \texttt{M\_AUTH\_SIT\_EFT\_v1}.

\begin{table}[htbp]
\centering
\caption{Model-continuity crosswalk separating the July~2026 Version~80 compact benchmark (Paper~002/003) from the November~2026 Version~81 authoritative continuum action (\texttt{M\_AUTH\_SIT\_EFT\_v1}) and 4D lattice Models~L1--L3.}
\label{tab:july-v80-vs-nov-v81}
\vspace{0.35em}
\footnotesize
\begin{tabular}{|p{2.7cm}|p{5.8cm}|p{6.2cm}|}
\hline
\textbf{Structural Element} & \textbf{July~2026 Version~80 Benchmark (Paper~002 / Paper~003)} & \textbf{November~2026 Version~81 Authoritative Action (\texttt{M\_AUTH\_SIT\_EFT\_v1} \& Models~L1--L3)} \\
\hline
Scalar field content & Two neutral real scalar amplitudes $(\tau, \sigma)$ of mass dimension $[M]^1$ (plus optional shift-symmetric global phase $\theta$ with kinetic stiffness $K(\sigma)$). & Real gauge-singlet temporal scalar $\rho_t$ and charged complex scalar $\psi = |\psi|e^{i\theta}$ ($\Rcoh = |\psi|/\psi_0 \ge 0$). \\
\hline
Matter coupling & Explicit conformal physical metric $\tilde{g}_{\mu\nu} = \Omega^2(\tau,\sigma)g_{\mu\nu}$ ($\Omega > 0$); Paper~003 adds a declared phenomenological dimension-4 Penrose--Onsager condensate source for $\sigma$. & Multiplicative portal $+f_{\mathrm{SM}}(\Rcoh)\mathcal{L}_{\mathrm{matter}}$ ($f_{\mathrm{SM}} = 1 + \Rcoh^2/(2M_{\mathrm{sym}}^2) \ge 1$); proved non-conformal in 4D unless $f_{\mathrm{SM}} = 1$ (Stage~6); composite/condensate sources open (\texttt{OPEN-6A}/\texttt{6C}). \\
\hline
Gauge / EM sector & Independent physical Maxwell connection $A_\mu$ with dressing $B_F$ and standard unbroken visible photon. & Internal $U(1)$ connection $A_\mu$ in $D_\mu\psi = (\nabla_\mu - i\hat{q}A_\mu)\psi$ with Proca mass $m_A^2 = \kappa_c\hat{q}^2 v_0^2/f_{\mathrm{EM}0}$ when $\hat{q}v_0 \ne 0$; visible $U(1)_{\mathrm{EM}}$ dictionary open (\texttt{OPEN-6B}). \\
\hline
Lattice / non-perturbative sector & Separate 2D finite abelian $U(1)$ toy model (Paper~033, finite readout-fibre theorem). & 4D hypercubic Euclidean Haar Wilson-coherence lattice Models~L1--L3 ($G \subset U(n+1)$, $2L_t$ OS positivity, classical $C^4$ $O(1/l^2)$ limit, sharp $3\beta/4$ nonnegative coercivity). \\
\hline
Shared / transferable content & Unital-CPTP vs.\ relative-entropy repair, operational amplitude-vs-purity separation, mass-calibration/form-factor cautions, matched-preparation controls, QGTCD counterexamples. & Shared two-mode linear response $(p^2 K + H)^{-1}$ on matched quadratic truncations; distinct non-linear actions remain non-interchangeable unless an explicit truncation/redefinition is proved. \\
\hline
\end{tabular}
\end{table}

Seven major upgrades distinguish Version~81 from Version~79 and earlier working drafts:
\begin{enumerate}[leftmargin=2em]
  \item \textbf{Canonical $(-,+,+,+)$ Effective Field Theory Lock (\texttt{M\_AUTH\_SIT\_EFT\_v1}, canon \texttt{v1\_1}):} Every gravitational, scalar, gauge, and stress-energy equation in the core theory is unified in mostly-plus metric signature $(-,+,+,+)$ with canonical kinetic signs ($-\frac{\kappa_t}{2}g^{\mu\nu}\partial_\mu\rho_t\partial_\nu\rho_t$, $-\frac{\kappa_c}{2}g^{\mu\nu}(D_\mu\psi)^*(D_\nu\psi)$, $-\frac{1}{4}f_{\mathrm{EM}}F_{\mu\nu}F^{\mu\nu}$), vacuum-positive dressings $f_{\mathrm{EM}}(\Rcoh) \ge 1$ and $f_{\mathrm{SM}}(\Rcoh) \ge 1$, $\mathbb{Z}_2$-lifted non-negative linkage potential $V_{\mathrm{tot}} \ge 0$, conditional Jordan-frame null-cone invariance $\Phi_J + \Psi_J = 2\Phi_E$ (explicitly distinguished from multiplicative $f_{\mathrm{SM}}\mathcal{L}_{\mathrm{matter}}$), and explicit ultraviolet cutoff crosswalk $\Lambda_{\mathrm{UV}} = \min(M_{\mathrm{Pl}}, \Lambda_{\mathrm{EM}}, \Lambda_{\mathrm{sym}})$. All historical sign flips, $1/2$-vs-$1/4$ Maxwell prefactors, wrong-sign ghost kinetic branches (\texttt{C04}, \texttt{C05}, \texttt{C07}), unconstructed supersymmetric factorizations (\texttt{C06}), naive row-normalized stochastic clock averaging (\texttt{C08}), and one-loop Sakharov induced gravity (\texttt{C12}) are retired and documented transparently.
  \item \textbf{Infinite Sobolev Coherence Tower and Curved $3+1$D ADM Closure (Part~IV, Sections~51--52):} We present the machine-checked $\ell^2$ weighted Sobolev hierarchy $H^2(\mu) \subsetneq H^1(\mu) \subsetneq \ell^2(\mathbb{N})$ and strong Galerkin convergence $\|u - P_N u\|_{\ell^2} \to 0$, paired with the November 2026 perturbative diagnosis (\emph{The Hawking--Page transition at weak coupling}, \texttt{hexagon:2611.00175v1}, \url{https://hexagonmath.org/pdf/2611.00175v1}) showing why sharp spatial mode cutoffs $P_N$ induce boost-breaking $-\frac{4}{3}g^2$ Coulomb and quadratic scalar mass counterterms in 4D gauge loops. We then establish nonlinear curved $3+1$D ADM canonical closure ($36 - 2\times 10 = 16$ phase-space dimensions, $13\times 13$ DeWitt--Cholesky Legendre inversion, and local pointwise formal 3-jet first-class Dirac--Bergmann closure modulo explicit spatial divergences $\partial_k B^k$ of $\{\mathcal{H}_\perp,\mathcal{H}_\perp\}$, $\{\mathcal{D},\mathcal{H}_\perp\}$, $\{\mathcal{D},\mathcal{D}\}$, $\{\mathcal{H}_\perp,G_{\mathrm{Gauss}}\}$, and $\{\mathcal{D},G_{\mathrm{Gauss}}\}$, including unconstrained 6-component symmetric $3\times 3$ off-diagonal $q_{ij}$ spatial-diffeomorphism / super-Hamiltonian closure in \texttt{Stage5\allowbreak{}Extensions.\allowbreak{}Off\allowbreak{}Diagonal\allowbreak{}Curved\allowbreak{}ADMClosure} [$227$ declarations] specialized to the physical Sylvester Riemannian positive-definite cone \texttt{IsSPD3} in \texttt{Stage5\allowbreak{}Extensions.\allowbreak{}Matter\allowbreak{}And\allowbreak{}ADMDomain\allowbreak{}Refinements} [$66$ declarations], while demarcating integrated constraint algebras on Sobolev/Banach spaces with Regge--Teitelboim asymptotic boundary conditions under \texttt{OPEN-2B}).
  \item \textbf{Non-Perturbative 4D Haar Wilson-Coherence Lattice Resolution (Section~49 and Part~IV, Sections~53--54):} Replacing sharp spatial cutoffs and the 2025 continuum $f\to 0$ heuristic (\texttt{C09}) with a compact-group 4D hypercubic Euclidean Haar Wilson lattice resolves the $\Pi^2/(2f)$ singularity non-perturbatively ($0 < Z \le 1$ even at $f_{\mathrm{gate}}(\rho,0) = 0$). Explicitly distinguishing unshifted Model~L1 (\texttt{sit\allowbreak{}Wilson\allowbreak{}Coherence\allowbreak{}Action}, frozen Stage~4, which exhibits exact zero-matter Haar disorder yet encounters a reduced two-state gate-kernel obstruction $\det K = m_{11} - 1 < 0$ for $m_{11} = e^{-\beta} < 1$ when isolating the unshifted cross-slab factor $e^{-\beta g_0 g_1}$, \texttt{unshifted\_\allowbreak{}gate\_\allowbreak{}two\_\allowbreak{}by\_\allowbreak{}two\_\allowbreak{}det\_\allowbreak{}neg}, motivating an arithmetic baseline) and ungated-temporal Model~L2 (\texttt{sitOSCoupledAction} in \texttt{Stage5\allowbreak{}Extensions.\allowbreak{}Coupled\allowbreak{}OSReflection\allowbreak{}Positivity} [$153$ declarations]), Stage~5.1 constructs unified Model~L3 (\texttt{sitAMGMGatedOSAction} / \texttt{sitAMGMGatedOSLaw} in \texttt{Stage5\allowbreak{}Extensions.\allowbreak{}Unified\allowbreak{}AMGMGated\allowbreak{}OSPositivity} [$98$ declarations]), differing from Model~L1 by the gauge-link-independent sum $\beta\,\mathrm{amgmCoherenceGap} + \kappa_t\sum_{x,\mu}(\rho(x+e_\mu)-\rho(x))^2 \ge 0$ (isolating the pure AM--GM gap at $\kappa_t = 0$ or relative to a reference action carrying the same clock-hopping term) and simultaneously proving within one single measure: exact $SU(2)$ Haar disorder in the concentrated zero-matter specialization ($\mu_{\mathrm{site}} = \delta_0 \otimes \mu_{\mathrm{temp}}$, \texttt{sit\allowbreak{}AMGMGated\allowbreak{}OSLaw\_\allowbreak{}zero\allowbreak{}Matter\_\allowbreak{}fst\_\allowbreak{}eq\_\allowbreak{}pi\_\allowbreak{}haar}, alongside conditional $P_{\mathrm{L3}}(U\mid\rho_t,\psi=0)$, with zero single-plaquette expectation for $l \ge 2$), identical conditional link law $P_{\mathrm{L3}}(U\mid\rho_t,\psi) = P_{\mathrm{L1}}(U\mid\rho_t,\psi)$ (preserving the $48 M \beta L_f \sum_x \|\psi(x)\|^2$ activation bound and uniform-background $O(a^2)$ gauge-sector recovery of \texttt{M\_AUTH\_SIT\_EFT\_v1}, while inducing a radial scalar kinetic shift $\Delta S_{\mathrm{AMGM}} = \frac{3\beta}{8a^2}\int d^4x\,\frac{(\partial_\mu f)^2}{f} + O(a^0)$ whose three continuum/EFT scaling regimes are proved in \texttt{amgm\_\allowbreak{}three\_\allowbreak{}eft\_\allowbreak{}scaling\_\allowbreak{}regimes}), and coupled $(U,\psi,\rho_t)$ Osterwalder--Schrader $2L_t$ reflection positivity extended from compact-ball priors $B_R$ to unbounded/Gaussian unitary-invariant priors via dominated-convergence exhaustion. In the thermal sector, heavy-coherence coupling induces first-order Polyakov coexistence in the effective holonomy potential with exact cubic barrier $(4/27)c u_*^3 = 1/16 = 0.0625$ (cross-checked by an independent 2-variable zero-mode effective-potential sampler \texttt{heavy\allowbreak{}Coherence\allowbreak{}Polyakov\allowbreak{}Joint} at $\kappa = 2.0$ yielding $\Delta V_{\mathrm{MCMC}} = 0.0610 \pm 0.0021$, $0.71\sigma$, \texttt{is\_\allowbreak{}microscopic\_\allowbreak{}4d\_\allowbreak{}lattice\_\allowbreak{}extraction = false}, separate from exploratory finite-volume 4D $SU(2)$ Monte Carlo consistency benchmarks on $4^4$, $6^4$, $4^3\times 4$, $6^3\times 4$, and $4^3\times 8$ lattices at $\kappa = 1.58$, where the microscopic-to-effective parameter map remains analytical/open), alongside a dated synthesis of the 2025 SIT holonomy corpus and the October--November 2026 constructive Yang--Mills milestones.
  \item \textbf{Informational Bundle Holonomy for Matter and Chiral Standard Model Interface (Section~55):} Separating frozen Stage-4 algebraic/conditional interfaces (\texttt{Stage4Matter/*}) from Stage-5.1/5.2 constructive derivations (\texttt{Stage5\allowbreak{}Extensions.\allowbreak{}Matter\allowbreak{}And\allowbreak{}ADMDomain\allowbreak{}Refinements} [$66$ declarations]) and Tier-2 continuum analytics, Lean~4.30 proves $4\times 4$ Dirac Clifford and noncommutative magnetic commutator $[\Pi_i,\Pi_j] = i q\hbar\varepsilon_{ijk}B_k$ forcing $g_{\mathrm{tree}} = 2$ (\texttt{dirac\_\allowbreak{}clifford\_\allowbreak{}pauli\_\allowbreak{}derives\_\allowbreak{}g\_\allowbreak{}eq\_\allowbreak{}two}) with $O(\Lambda_{\mathrm{UV}}^{-2})$ dipole suppression, circular Coulomb \texttt{HasDerivAt} force balance plus $2\pi n$ phase holonomy deriving Bohr radii $r_n$ and $E_n = E_1/n^2$ (\texttt{coulomb\_\allowbreak{}holonomy\_\allowbreak{}derives\_\allowbreak{}bohr\_\allowbreak{}spectrum}), and strictly non-zero finite-matrix Neuberger overlap chirality $D_{\mathrm{ov}} = \frac{1}{a}(I - \gamma_5 H_W) \ne 0$ (\texttt{neuberger\allowbreak{}Overlap\_\allowbreak{}witness\_\allowbreak{}nontrivial}) with injective $SU(3)_c\times SU(2)_L\times U(1)_Y \hookrightarrow U(6)$ embedding and massive fermion determinant floor $m_f^{2N} \le \det(D^\dagger D + m_f^2 I)$ (\texttt{massive\allowbreak{}Fermion\allowbreak{}Det\allowbreak{}Weight\_\allowbreak{}constructs\_\allowbreak{}sm\_\allowbreak{}interface}), while demarcating continuum unbounded 3D $L^2(\mathbb{R}^3)$ Dirac/Schr\"odinger PDEs, Brillouin-zone single-pole doubler dispersion, and massless chiral determinant phases as analytical/open.
  \item \textbf{Five-Tier Epistemic Discipline and Four-Layer Verification (\texttt{QG49}, Appendix~AB):} Every claim in this monograph is classified into five explicit epistemic tiers separating (Tier~1) Lean~4.30 kernel-compiled formal theorem types under \texttt{[propext, Classical.\allowbreak{}choice, Quot.\allowbreak{}sound]}---reported at the level of individual mathematical questions rather than conflating aggregate theorem counts with unsolved physical derivations---(Tier~2) analytical derivations, dual-engine SymPy/NumPy checks, and exploratory finite-volume MCMC benchmarks, (Tier~3) conditional formal bridges, (Tier~4) retired historical branches and quarantined legacy observational approximations, and (Tier~5) unsolved open problems (\texttt{NOT\_SOLVED}, \texttt{OPEN-1}--\texttt{OPEN-6}), backed by a four-layer audit distinguishing static source hygiene, symbolic-numerical Python execution, Lean kernel type-checking, and human mathematical-physics specification-to-prose review.
  \item \textbf{Stage-6 Classical $C^4$ Continuum Limit, Periodic Fourier Eigenvectors, Sharp $\frac{3\beta}{4}$ Nonnegative Coercivity, Canonical Propagator Instantiation, and Provenance Repairs:} Incorporating all $6$ frozen Stage-6 Lean~4.30 modules (\texttt{Stage6FunctionalResponse/*}, $30/30$ audited theorems in \texttt{SITStage6Audit}) plus $3$ publication-branch addendum modules (\texttt{Stage6Addendum.\allowbreak{}PeriodicLatticeFourierEigenvalue} [$3$], \texttt{NonnegativeFieldCoercivity} [$10$], and \texttt{ActionResponsePropagatorIdentity} [$5$]; $669$ total \texttt{.lean} files and $48$ Stage-6/Addendum audited theorems), Version~81 (a)~proves $O(1/l^2)$ convergence of the exact 4-corner AM--GM lattice energy to $\frac{3\beta}{8}\int_{[0,1]^4}\frac{|\nabla f|^2}{f}\,du$ for conventional $C^4$ strictly positive periodic fields together with exact Stage-5 finite-lattice embedding at $\kappa_t = 0$, exact periodic lattice Fourier-mode eigenvalue equations $(H_{\mathrm{lat}} u)(x) = a^2\lambda_Q(k;a)\,u(x)$, and the sharp nonnegative-field coercivity bound $Q_l(F) \ge \frac{3\beta}{4} a^2\sum_{x,\mu}(\sqrt{F(x+e_\mu)}-\sqrt{F(x)})^2$ valid even at zero corners with both exact equality witnesses (nontrivial single-site spike $s > 0$ for $l \ge 2$ and nontrivial strictly positive checkerboard $A, B > 0, A \ne B$ for even $l$, strengthening the valid weaker $\beta/2$ bound and yielding strong $L^2(\mathbb{T}^4)$ compactness of $\sqrt{F_l}$ with $H^1(\mathbb{T}^4)$ control and strong $L^1(\mathbb{T}^4)$ compactness of $F_l$ when paired with a uniform bound on $a^4\sum_x F(x)$; also presented independently of SIT physics in \emph{A Classical Continuum Limit for a Four-Corner Lattice Energy}); (b)~proves that the unchanged multiplicative portal $+\frac{|\psi|^2}{\Lambda_{\mathrm{sym}}^2}\mathcal{L}_{\mathrm{matter}}$ does not match a conformal Jordan metric unless $f_{\mathrm{SM}} = 1$, proves $(p^2 K + H)^{-1} = V\,\mathrm{diag}(1/(p^2+m_a^2))\,V^\top$ together with its $D = \mathrm{diag}(1,\psi_0)$ coordinate-rescaling invariance and direct instantiation on \texttt{AuthModelParams} proving $0 < m_0^2 \le m_1^2$ and automatic $p^2 + m_a^2 \ne 0$ for all $p^2 \ge 0$, and qualifies all Jordan-frame Cassini margins accordingly; (c)~executes a conditional weak-field uniform-slab E\"ot-Wash non-exclusion calculation directly from unchanged \texttt{M\_AUTH\_SIT\_EFT\_v1} while documenting the corrected Lee~\emph{et~al.}\ 18-fold/120-fold patterned Pt--epoxy apparatus on BK7 supports in \texttt{primary\_measurements\_manifest\_v2.json} and leaving patterned torque modeling, composite material charges $\rho_{\mathrm{Pt}} q_{\mathrm{Pt}} - \rho_{\mathrm{epoxy}} q_{\mathrm{epoxy}}$, visible $U(1)$ photon coupling, and laboratory $J_{\mathrm{prep}}$ open (\texttt{OPEN-6}); and (d)~replaces legacy unsupported forecast amplitudes and universal falsification criteria with an explicit missing-derivation checklist, withdraws the flat-connection $A_\mu = (\hbar/e)\partial_\mu\theta$ magnetism construction and CRT $10^{-8}$ figure, repairs Appendix~E's linearized scalar action/stress-energy derivations and coordinate congruence, unifies the $\alpha_\rho$ / $\mathcal{T}$ / $\alpha_\Phi = \alpha_\rho\mathcal{T}$ clock dictionary, and quarantines legacy $3\times 10^{-8}$, $29.4\,\mu\mathrm{Gal}$, and \texttt{np.piecewise} approximations.
  \item \textbf{Bidirectional Corrections and Controls Imported from the July~2026 Intervening Papers (Version~81 Addition):} Importing the strongest mathematical and methodological safeguards developed across the July~2026 portfolio (Papers~002, 003, 007, 014, 028, and~055), Version~81 (a)~retracts Version~79's false claim that quantum purity decreases monotonically under arbitrary CPTP maps (explicitly recording Paper~002's amplitude-damping and complete-dephasing counterexamples, restricting universal purity non-increase to \emph{unital} CPTP channels, and grounding general open-system contraction in the quantum relative-entropy data-processing inequality $D(\Phi(\rho)\|\Phi(\omega)) \le D(\rho\|\omega)$, Section~\ref{sec:CoherenceConservationArrow}); (b)~separates classical EFT field amplitudes $(\Rcoh, \rho_t)$ cleanly from operational quantum statistics (purity, visibility, Penrose--Onsager condensate fraction, and point-process event rates) via open calibration-kernel contracts (Section~\ref{sec:PrimitiveFieldDefinitions}); (c)~integrates Paper~003/014's Penrose--Onsager condensate-source boundary and matched-preparation residual controls $\Delta a_{\mathrm{res}} = (a_{\mathrm{high-C}} - a_{\mathrm{low-C}}) - \Delta a_{\mathrm{standard}}$ (Section~\ref{sec:smoking-gun-bec}); and (d)~codifies Paper~002/007's mass-calibration and finite-source form-factor identifiability cautions alongside Paper~055's four standing QGTCD non-implication guardrails (Section~\ref{sec:experimental-predictions}).
\end{enumerate}

\tableofcontents
\newpage
% --- END INPUT sections/00_abstract_and_note.tex ---
% --- BEGIN INPUT sections/01_part1_sec01_to_03.tex ---
\part{Foundations of Super Information Theory}

\section{Introduction}
\label{sec:introduction}

\subsection{Motivation and Theoretical Scope}

The foundational postulate of Super Information Theory (SIT) is that ``information'' is not an abstract Platonic symbol but an operationally defined property of physical systems, quantified through the local coincidence of independent physical events. A coincidence event occurs whenever two or more physical processes intersect within a defined spatiotemporal neighborhood, from the simultaneous arrival of particles in a detector to the synchronized firing of cortical neurons. The fundamental fields of SIT---the complex coherence field $\psi(x)$ and the real time-density field $\rho_t(x)$---provide the formal mathematical description of the structure and dynamics of these coincidence processes.

The unification of general relativity and quantum mechanics remains a central challenge of theoretical physics. Super Information Theory proposes a covariant framework in which quantum phase coherence is not treated merely as a passive statistical descriptor on a fixed background, but as an active order-parameter field whose gradients and modulus couple dynamically to spacetime curvature and local proper-time accumulation. Physical phenomena across gravitational, electromagnetic, and thermodynamic regimes emerge from the coupled variational dynamics of $(\psi, \rho_t, g_{\mu\nu}, A_\mu)$ alongside Standard Model matter, without introducing unobservable ontological substrates.

Crucially, the core framework developed in this monograph does not claim to constitute an unconditional ultraviolet-complete quantum theory of gravity at trans-Planckian energies. All continuum field equations and perturbative loop results are derived and interpreted within a controlled Wilsonian effective field theory (EFT) regime valid below an explicit ultraviolet cutoff $\Lambda_{\mathrm{UV}} := \min(M_{\mathrm{Pl}}, \Lambda_{\mathrm{EM}}, \Lambda_{\mathrm{sym}})$; whether SIT admits a regulator-independent non-perturbative 4D gravitational UV fixed point above $\Lambda_{\mathrm{UV}}$ remains an explicit open problem (\texttt{OPEN-1}, Appendix~\ref{app:epistemic-ledger-lean}).

Within this well-defined effective regime, gravity, electromagnetism, and quantum interference represent distinct geometric aspects of a unified informational field structure:
\begin{itemize}
  \item \textbf{Gravity} is sourced by the stress--energy of matter together with the coupled dynamics of the time-density scalar $\rho_t(x)$ and the gauge-invariant coherence modulus $\Rcoh(x) := |\psi(x)|/\psi_0$.
  \item \textbf{Electromagnetism} is realized through a $U(1)$ principal bundle connection $A_\mu(x)$ coupled minimally to the coherence phase $\theta(x) \equiv \arg(\psi(x))$ via $D_\mu\psi = (\nabla_\mu - i\hat{q}A_\mu)\psi$, reducing on flat pure-gauge patches ($F_{\mu\nu}=0$) to Aharonov--Bohm phase holonomy while retaining two physical transverse polarizations (plus a longitudinal Proca mode when $v_0 > 0$) globally.
\end{itemize}

The central objective of this monograph is to establish this unified canonical action (\texttt{M\_AUTH\_SIT\_EFT\_v1}), prove its mathematical stability, ghost freedom, curved $3+1$D ADM constraint closure, and non-perturbative 4D lattice gauge consistency, demonstrate its exact reduction to General Relativity, Maxwell/Proca theory, quantum mechanics, and kinetic theory in tested limits, and derive concrete laboratory falsifiers in optical-clock metrology, cold-atom interferometry, and $^{87}\mathrm{Rb}$ Bose--Einstein condensates.

\subsection{Foundational Insight: From Neuroscience to a Unified Physical Theory}

The conceptual architecture of Super Information Theory grew out of a synthesis uniting two inverse scaling laws across neuroscience and fundamental physics. This synthesis provided the bridge for extending an oscillatory theory of mind (\emph{Self-Aware Networks}) into a covariant field theory of spacetime and coherence, beginning in 2022 with \emph{Quantum Gradient Time Crystal Dilation (QGTCD)}.

In systems neuroscience, the \textbf{$1/f$ power law} (spectral ``pink noise'') characterizes cortical population activity in which oscillatory amplitude scales inversely with temporal frequency: slower, large-scale rhythms coordinate broad spatial ensembles, whereas faster, localized bursts carry rapid state updates. In fundamental wave mechanics, wavelength $\lambda$ and period $T$ likewise scale inversely with frequency $f$. Synthesizing these observations yields an organizing physical principle: at fixed informational energy budget, a coherent wave configuration trades off spatial/structural \textbf{Magnitude} against dynamical \textbf{Frequency}.

Visualized through the physiological contrast between sustained tonic synchronization and rapid phasic signaling, this trade-off motivates an informational reading of Einstein's mass--energy equivalence $E = mc^2$ across distinct morphologies of coherent organization:
\begin{itemize}
  \item \textbf{Radiative Energy ($E$):} Coherent excitations organized as high-frequency, propagating, low-amplitude modes (analogous to phasic bursts).
  \item \textbf{Localized Mass ($m$):} Coherent excitations organized as persistent, phase-locked, spatially localized standing structures---informational solitons or ``time crystals'' (analogous to tonic attractors)---whose coherence modulus thickens local time density $\rho_t$.
  \item \textbf{Decoherent Vacuum Baseline:} The reference state $(\rho_t = \rho_0,\, \Rcoh = R_0)$ relative to which gradients of coherence and time density source measurable curvature and clock shifts.
\end{itemize}

\subsection{Relation to Causal Fermion Systems and Bounded Comparison Scope}
\label{sec:cfs-convergence}

An instructive point of comparison in foundational physics is the relationship between Super Information Theory (SIT) and Causal Fermion Systems (CFS), developed by Felix Finster and collaborators from operator algebras and a causal action principle. Following the bounded four-lens audit established in Paper~048 (\emph{Causal Fermion Systems and Super Information Theory}, July~28, 2026), Version~81 strictly separates dated conceptual parallels from mathematical equivalence or empirical validation:
\begin{itemize}
  \item \textbf{Dated Conceptual Parallels (Narrative Level Only):} SIT's 2017 premise framing information around relational ``coincidence patterns'' and internal distributed observation (NAPOT) shares high-level vocabulary with later 2024--2025 CFS expositions describing spacetime as a ``web of correlations'' and intrinsic ``effective dynamical collapse.'' Chronological overlap in conceptual framing is \emph{not} a mathematical embedding, a proof of operational equivalence, or an independent empirical validation of SIT.
  \item \textbf{Bounded Effective-Equation Overlap and Null Control (Paper~048):} Under Paper~048's frozen comparison against July~2026 Version~80 Draft~13, both a regulated constant-coupling CFS continuum regime and a specialized neutral-Dirac constant-background SIT submodel reduce to standard Einstein--Dirac equations, matching a direct Einstein--Dirac null control with \emph{zero discriminatory value}. Extending that constant-background truncation check to the November~2026 multiplicative action \texttt{M\_AUTH\_SIT\_EFT\_v1} remains a separate open calculation.
\end{itemize}

Table~\ref{tab:sot-cfs-map} summarizes the conceptual chronology, while Appendices~\ref{app:cfs-chronology} and~\ref{app:chronology} preserve the dated historical timeline alongside the explicitly open formalization obligations.

\begin{table}[htbp]
\centering
\caption{Conceptual comparison between high-level framing elements of Super Information Theory (SIT) and Causal Fermion Systems (CFS), with first public disclosure dates (bounded narrative level only; not a proof of mathematical equivalence or independent validation, cf.\ Paper~048).}
\label{tab:sot-cfs-map}
\vspace{0.4em}
\begin{tabular}{|p{3.0cm}|p{5.8cm}|p{5.8cm}|}
\hline
\textbf{Conceptual Pillar} & \textbf{Super Information Theory (Blumberg)} & \textbf{Causal Fermion Systems (Finster et al.)} \\
\hline
\textbf{Relational Substrate} & Reality framed via relational \textbf{``coincidence patterns''} of events (\emph{First Public Appearance: 2017}). & Spacetime reframed as an informational \textbf{``web of correlations''} of a quantum many-body system (\emph{First Public Appearance: 2025}). \\
\hline
\textbf{Measurement Framing} & Observation modeled conceptually as \textbf{distributed internal phase alignment} (``inner screens'', \emph{First Public Appearance: 2017}). & Measurement studied via \textbf{``effective dynamical collapse''} from the causal action (\emph{First Public Appearance: 2024}). \\
\hline
\end{tabular}
\end{table}

Where CFS emphasizes abstract operator measures, SIT formulates an explicit sub-Planckian scalar--tensor--gauge EFT (\texttt{M\_AUTH\_SIT\_EFT\_v1}) and 4D lattice gauge models, whose concrete empirical confrontation awaits completion of the action-to-apparatus bridges (\texttt{OPEN-6}).

\subsection{Relation to Prior Work and Technical Comparison with Unification Frameworks}
\label{sec:relation-prior-work}
\label{sec:SITTechnicalComparison}

\subsubsection{From Static Bits to Dynamical Coherence}

Information theory entered physics in modern quantitative form through Claude Shannon's 1948 measure of symbolic uncertainty, $H = -\sum_i p_i \log p_i$, defined over discrete probability distributions on a timeless register. While foundational to communication engineering, combinatorial symbol counting says nothing about continuous phase relationships inside physical media. John Wheeler's \emph{``It from Bit''} shifted attention toward binary physical interrogations, and Wojciech Zurek's environment-induced decoherence program demonstrated how open-system coupling selects pointer states by exporting phase correlations into an environment.

Super Information Theory embeds these motivations into a local covariant field theory by introducing two dynamical EFT variables: a complex order parameter $\psi(x)$ with amplitude ratio $\Rcoh(x) = |\psi(x)|/\psi_0$ and a real scalar $\rho_t(x)$ (kept distinct from quantum purity and event-rate proxies via open calibration kernels, Section~\ref{sec:PrimitiveFieldDefinitions}). Closed-system global $U(1)$ invariance yields the conserved Noether current $J^\mu_{\mathrm{coh}}$ (distinct from operational open-system purity, which is non-increasing specifically under unital CPTP maps while general CPTP channels contract quantum relative entropy, Section~\ref{sec:arrow-of-time}), and thermodynamic-gravity analogies are investigated as coarse-grained effective comparisons rather than identities.

\begin{table}[htbp]
\centering
\caption{Technical comparison of Super Information Theory (SIT) with established gravitational and quantum-informational paradigms.}
\label{tab:sit-framework-comparison}
\vspace{0.4em}
\small
\begin{tabular}{|p{3.2cm}|p{4.8cm}|p{6.8cm}|}
\hline
\textbf{Framework} & \textbf{Core Mechanism / Field Content} & \textbf{How SIT Differs and Extends} \\
\hline
\textbf{Entropic Gravity (Verlinde, Padmanabhan)} & Gravity from entropy gradients across holographic screens; no local bulk action for coherence. & SIT supplies a local covariant bulk action (\texttt{M\_AUTH\_SIT\_EFT\_v1}) for $(\rho_t, \psi)$ with hyperbolic wave equations, exact Bianchi--Noether conservation, and prospective clock/BEC test channels (\texttt{OPEN-6}). \\
\hline
\textbf{Holography \& AdS/CFT ('t Hooft, Susskind, Maldacena)} & Bulk AdS gravity dual to boundary CFT operator algebra; relies on asymptotic boundary structure. & SIT operates directly in four-dimensional Lorentzian bulk spacetimes of arbitrary geometry without requiring anti-de Sitter asymptotics. \\
\hline
\textbf{Tensor Networks \& QEC (MERA, Almheiri--Harlow)} & Geometry from discrete entanglement graphs and boundary quantum error-correcting codes. & SIT formulates continuous bulk EFT fields below $\Lambda_{\mathrm{UV}}$ paired with a gauge-invariant 4D Euclidean Haar Wilson lattice at the non-perturbative gauge level. \\
\hline
\textbf{Loop Quantum Gravity \& Causal Sets} & Quantized $SU(2)$ holonomy spin networks or locally finite partially ordered causal sets. & SIT couples continuum $3+1$D ADM geometry to $(\rho_t, \psi)$ in the infrared and employs compact $U(1)/SU(N)$ bundle holonomy on 4D Wilson lattices for gauge-matter dynamics. \\
\hline
\textbf{Brans--Dicke \& Symmetron Scalar--Tensor Gravity} & Real scalar coupled to Ricci curvature or matter trace with density-dependent screening. & SIT pairs a real time-density scalar $\rho_t$ with a $U(1)$-charged complex coherence order parameter $\psi$, with vacuum-positive dressings $f_{\mathrm{EM}} \ge 1$ and $f_{\mathrm{SM}} \ge 1$ (keeping conformal-frame Shapiro relations explicitly conditional on a conformal matter-response derivation, Section~\ref{sec:ppn-shapiro-exact}). \\
\hline
\end{tabular}
\end{table}

\section{Providing a Micro-Mechanism for Informational Gravity Frameworks}
\label{sec:entropic_RTI_SIT}

Two influential conceptual traditions that framed gravitation and quantum measurement in informational terms---Verlinde's entropic gravity and Cramer--Kastner's Relativistic Transactional Interpretation (RTI)---acquire a concrete local field-theoretic realization inside Super Information Theory through the coupled dynamics of $\Rcoh(x)$ and $\rho_t(x)$.

\subsection*{RTI's Handshake as Phase Alignment}

In transactional language, a completed quantum event corresponds to phase matching between offer and confirmation amplitudes. Within SIT, local phase alignment is measured by the vanishing of relative phase mismatch $\Delta\theta(x) = \arg(\psi^{\mathrm{ret}}) - \arg(\psi^{\mathrm{adv}})$, maximizing local interference visibility and rendering the coherence current $J^\mu_{\mathrm{coh}}$ stationary across the interaction region. Rather than invoking retrocausal signaling, SIT tracks how constructive phase alignment elevates $\Rcoh(x)$ and, via the linkage potential $U_{\mathrm{link}}(\rho_t, |\psi|)$, shifts the local time-density scalar $\rho_t(x)$ in strictly causal, hyperbolic fashion.

\subsection*{Entropic Gravity as a Coarse-Grained Limit}

In entropic gravity, displacing a test mass $m$ by $\Delta x$ near a holographic screen of temperature $T$ changes screen entropy by $\Delta S = (2\pi k_{\mathrm{B}}/\hbar) m c \Delta x$, yielding $F\Delta x = T\Delta S$. In SIT, a localized rise in $\Rcoh$ shifts $\rho_t$ through $U_{\mathrm{link}}$ and dresses matter stress--energy via $f_{\mathrm{SM}}(\Rcoh)$, generating a screened Yukawa addition to the Einstein-frame Newtonian potential $\Phi_E(r) = -G f_{\mathrm{SM}0} M / r$. Coarse-graining over microscopic phase modes reproduces thermodynamic force balance without sacrificing local stress--energy conservation $\nabla^\mu T^{\mathrm{tot}}_{\mu\nu} = 0$.

\subsection*{Quantum Coherence as Information Density and Quasicrystal Echoes}

Regions whose constituent modes share a common phase exhibit elevated $\Rcoh$ and lowered linkage potential energy. At fixed mode energy, the inverse amplitude--frequency trade-off implies that phase-locked, high-amplitude configurations oscillate at lower effective coordinate rates, corresponding operationally to gravitational redshift. Across scales, constructive and destructive interference domains in $(\Rcoh, \rho_t)$ play a geometric role analogous to quasiperiodic projections in optical and plasmonic lattices: coherent condensations deepen gravitational wells while decoherent expanses relax toward the vacuum baseline.

\subsection{Operational and Mathematical Definitions of Primitive Fields}
\label{sec:PrimitiveFieldDefinitions}
\label{sec:math}

To anchor empirical evaluation, we define SIT's primitive fields hierarchically from their mathematical status in the canonical action to their laboratory measurement proxies.

\paragraph{(1) The Complex Coherence Field $\psi(x)$.}
The primary complex field $\psi(x) = |\psi(x)| e^{i\theta(x)}$ is a dynamical complex scalar field of mass dimension $[M]^1$ (in natural units $\hbar = c = 1$) on the four-dimensional Lorentzian manifold $(\mathcal{M}^4, g_{\mu\nu})$, transforming in a representation of local $U(1)$ gauge symmetry $\psi(x) \mapsto e^{i\alpha(x)}\psi(x)$ with covariant derivative $D_\mu\psi = (\nabla_\mu - i\hat{q}A_\mu)\psi$. We emphasize that $\psi(x)$ is \textbf{not} a single-particle quantum-mechanical Schr\"odinger wavefunction and $|\psi(x)|^2$ is \textbf{not} a Born probability density; rather, $\psi(x)$ is a macroscopic covariant order parameter in the effective-field-theory sense (analogous to the Ginzburg--Landau condensate order parameter in superconductivity or superfluidity), quantifying coherent structural order available to dress matter and gauge couplings.

\paragraph{(2) The Dimensionless Coherence Ratio $\Rcoh(x)$ vs.\ Operational Quantum Statistics.}
\begin{itemize}
  \item \textbf{Mathematical Definition:} Relative to the characteristic coherence reference scale $\psi_0 > 0$ ($[\psi_0] = [M]^1$), the dimensionless gauge-invariant coherence ratio is
  \begin{equation}
  \Rcoh(x) := \frac{|\psi(x)|}{\psi_0} \ge 0.
  \label{eq:def_Rcoh_canonical}
  \end{equation}
  (In normalized dimensionless units where $\psi_0 = 1$, this coincides with the shorthand $\Rcoh \equiv |\psi|$ used in phenomenological sections.) All gauge-invariant potentials and constitutive dressings depend on $\psi$ exclusively through $\Rcoh^2 = |\psi|^2/\psi_0^2$.
  \item \textbf{Operational Separation Contract (Imported from Paper~002 \& Paper~003):} Following the July~2026 operational discipline (Paper~002, Section~4.2), the classical field amplitude $\Rcoh(x) = |\psi(x)|/\psi_0$ is \emph{not} identically equal to quantum purity, fringe visibility, mutual information, or Bose--Einstein condensate fraction. Different experimental domains measure distinct declared statistics $C_M(x)$---such as normalized register purity for a finite-dimensional density matrix $\rho(x)$ of dimension $d > 1$,
  \begin{equation}
  \Rcoh^{\mathrm{proxy}}(x) = \frac{\mathrm{Tr}[\rho^2(x)] - 1/d}{1 - 1/d},
  \label{eq:purity_proxy_def}
  \end{equation}
  interferometric visibility, ultraviolet-regulated mutual information, or the Penrose--Onsager condensate fraction $f_{\mathrm{BEC}}$ of Paper~003. Any state-to-field link takes the form of an explicit, testable calibration kernel $C_M(x) = K_M[\psi,\rho_t](x) + \eta_M(x)$ (\texttt{OPEN-6C}); until $K_M$ and the corresponding laboratory source $J_{\mathrm{prep}}$ are derived and tested, $\Rcoh(x)$ remains a covariant EFT field amplitude rather than a universal quantum-information resource measure.
\end{itemize}

\paragraph{(3) The Time-Density Scalar Field $\rho_t(x)$.}
\begin{itemize}
  \item \textbf{Mathematical Definition:} $\rho_t(x)$ is a real gauge-singlet Lorentz scalar field of physical dimension $[T^{-1}]$ (mass dimension $[M]^1$ in natural units), possessing its own positive-norm kinetic term $-\frac{\kappa_t}{2}g^{\mu\nu}\partial_\mu\rho_t\partial_\nu\rho_t$ ($\kappa_t > 0$) and coupled to $|\psi|$ via the portal/linkage potential $U_{\mathrm{link}}(\rho_t, |\psi|)$.
  \item \textbf{Operational Measurement Contract (Paper~002 \& Paper~007):} Conceptually, a \emph{coincidence event} denotes an operational state transition resolved within a spacetime cell---such as photodetector clicks, two-photon Hong--Ou--Mandel coincidences, or atomic-clock transitions---motivating an operational event-rate proxy along a timelike worldline with proper time $\tau$:
  \begin{equation}
  \rho_t^{\mathrm{proxy}}(x) := \lim_{\Delta\tau \to 0} \frac{\Delta N_{\mathrm{events}}}{\Delta\tau}.
  \label{eq:def_rhot_operational}
  \end{equation}
  As formalized in Paper~002 (Section~4.1) and Paper~007, an event count per unit proper time and an atomic transition-frequency shift are distinct laboratory observables; linking $\rho_t(x)$ to either requires an explicit calibrated response function $R_{\mathcal{O},\epsilon}$ (accounting for exposure, dead time, and coupling dictionary \texttt{OPEN-6B}) rather than asserting an uncalibrated identity.
\end{itemize}

\begin{table}[htbp]
\centering
\caption{Summary of primitive SIT informational fields, strictly separated operational laboratory proxies (linked only via open calibration kernels under \texttt{OPEN-6}), and canonical dimensions ($\hbar = c = 1$).}
\label{tab:primitive-fields-summary}
\vspace{0.4em}
\footnotesize
\begin{tabular}{|p{2.9cm}|p{3.9cm}|p{4.3cm}|p{3.3cm}|}
\hline
\textbf{Field / Ratio} & \textbf{Canonical EFT Role} & \textbf{Distinct Operational Proxy (Requires \texttt{OPEN-6} Kernel)} & \textbf{Mass Dim.\ / Type} \\
\hline
$\psi(x)$ & Macroscopic $U(1)$ order parameter $|\psi|e^{i\theta}$ & Condensate or phase-order statistic (e.g.\ Penrose--Onsager mode, Paper~003) & $[M]^1$, complex $U(1)$ scalar \\
$\Rcoh(x) = |\psi|/\psi_0$ & Gauge-invariant EFT amplitude ratio & Normalized purity $\frac{\mathrm{Tr}[\rho^2]-1/d}{1-1/d}$, fringe visibility, or $f_{\mathrm{BEC}}$ (not bare identities) & $[M]^0$, real gauge invariant \\
$\rho_t(x)$ & Dynamical real temporal-density scalar & Atomic clock shift ($\alpha_\rho \delta\rho_t/\rho_0$) or calibrated event rate $\Delta N_{\mathrm{events}}/\Delta\tau$ & $[M]^1$ ($[T^{-1}]$), real scalar \\
\hline
\end{tabular}
\end{table}

For clarity of evaluation, Sections~1--8 establish the canonical physical EFT framework of Super Information Theory (\texttt{M\_AUTH\_SIT\_EFT\_v1}), including its action, covariant field equations, screening consistency, conditional PPN/Shapiro structure, and observational interfaces; Sections~9--25 develop conservation laws, canonical quantization, corrected open-system thermodynamics (unital CPTP vs.\ relative-entropy contraction), stability, and disambiguated optical-clock relations; Part~II (Sections~26--31) specifies experimental channels, identifiability cautions, and missing-derivation checklists; Part~III (Sections~32--50) examines interdisciplinary, neuroscientific, cosmological, and carefully scoped frontier extensions; and Part~IV (Sections~51--55) presents the machine-checked Lean~4.30 formalization across infinite Sobolev towers, curved $3+1$D ADM gravity, 4D non-perturbative Wilson-coherence lattice gauge theory, and bundle holonomy for matter.

\section{Operational Definitions and the Informational Fields}
\label{sec:operational-definitions}

A central methodological commitment of Super Information Theory is that ``information'' is treated as an operationally instantiated physical structure rather than a disembodied metaphysical substance. Its physical realization is mediated by concrete degrees of freedom whose interactions generate measurable events. Accordingly, $\Rcoh(x)$ and $\rho_t(x)$ do not postulate unobservable Platonic entities acting outside physical law; rather, they furnish covariant field variables in \texttt{M\_AUTH\_SIT\_EFT\_v1} whose connection to concrete laboratory observables (atomic clock frequency shifts, torsion-balance torques, interferometric visibilities, or condensate fractions) is governed explicitly by the open calibration and source-to-apparatus contracts of \texttt{OPEN-6} (Paper~002, Section~4; Section~\ref{sec:experimental-predictions}).

No claim is made that abstract informational notation is ontologically prior to physical degrees of freedom, nor that renaming a laboratory measurement (such as purity or clock rate) by a field symbol establishes a new dynamical field. Instead, following the calibration checklist of Paper~002 (Section~4.4), the field interpretation becomes empirically informative only when an explicit state-to-field map compresses multiple measurements, transfers out of sample across apparatuses or preparations, and predicts a reversible residual beyond established general-relativistic, electromagnetic, thermal, and mechanical systematics.
% --- END INPUT sections/01_part1_sec01_to_03.tex ---
% --- BEGIN INPUT sections/02_part1_sec04_to_08.tex ---
\section{The Gravitational Sector}
\label{sec:gravitational-sector}

\subsection{Derivation of the Modified Einstein Equations}

We derive the gravitational field equations by varying the canonical SIT action $S_{\mathrm{AUTH}}$ (Eq.~\eqref{eq:canonical_action_sec5}) with respect to the inverse metric $g^{\mu\nu}$ in mostly-plus Lorentzian signature $(-,+,+,+)$. Decomposing the total action into five additive covariant pieces,
\begin{equation}
S_{\mathrm{AUTH}} = S_{\mathrm{EH}} + S_{\rho_t} + S_{\psi} + S_{\mathrm{pot}} + S_{\mathrm{EM}}^{\mathrm{coh}} + S_{\mathrm{SM}}^{\mathrm{coh}},
\label{eq:action_sector_split}
\end{equation}
variation of the Einstein--Hilbert sector $S_{\mathrm{EH}} = \int d^4x\,\sqrt{-g}\,R/(16\pi G)$ yields the Einstein tensor:
\begin{equation}
\frac{\delta S_{\mathrm{EH}}}{\delta g^{\mu\nu}}
= \frac{\sqrt{-g}}{16\pi G}\left(R_{\mu\nu} - \frac{1}{2}g_{\mu\nu}R\right)
= \frac{\sqrt{-g}}{2\kappa_G}\,G_{\mu\nu},
\qquad
\kappa_G := 8\pi G > 0.
\label{eq:eh_variation}
\end{equation}
Defining the Hilbert stress--energy tensor of each non-gravitational sector by the canonical variational convention
\[
T_{\mu\nu}^{(X)} := -\frac{2}{\sqrt{-g}}\frac{\delta S_X}{\delta g^{\mu\nu}},
\]
stationarity $\delta S_{\mathrm{AUTH}} / \delta g^{\mu\nu} = 0$ yields the modified Einstein field equations:
\begin{equation}
G_{\mu\nu}
= \kappa_G\,T_{\mu\nu}^{\mathrm{tot}}
= 8\pi G\Bigl(
  T_{\mu\nu}^{(\rho_t)}
  + T_{\mu\nu}^{(\psi)}
  + T_{\mu\nu}^{(\mathrm{EM})}
  + T_{\mu\nu}^{(\mathrm{pot})}
  + T_{\mu\nu}^{(\mathrm{matter})}
\Bigr).
\label{eq:Einstein-like}
\end{equation}
Each contribution is manifestly symmetric, gauge-invariant, and ghost-free under $\kappa_t > 0$, $\kappa_c > 0$, and $f_{\mathrm{EM}}(\Rcoh) \ge 1$:
\begin{enumerate}[leftmargin=2em]
  \item \textbf{Time-Density Scalar Sector ($T_{\mu\nu}^{(\rho_t)}$):} Varying $-\frac{\kappa_t}{2}\int d^4x\,\sqrt{-g}\,g^{\alpha\beta}\partial_\alpha\rho_t\partial_\beta\rho_t$ gives
  \begin{equation}
  T_{\mu\nu}^{(\rho_t)}
  = \kappa_t\left(\partial_\mu\rho_t\,\partial_\nu\rho_t - \frac{1}{2}g_{\mu\nu}\,g^{\alpha\beta}\partial_\alpha\rho_t\,\partial_\beta\rho_t\right),
  \label{eq:T_rhot_sec4}
  \end{equation}
  whose Minkowski-frame energy density $T_{00}^{(\rho_t)} = \frac{\kappa_t}{2}\bigl[(\partial_t\rho_t)^2 + (\nabla\rho_t)^2\bigr] \ge 0$ is strictly positive-definite for $\kappa_t > 0$.
  \item \textbf{Complex Coherence Sector ($T_{\mu\nu}^{(\psi)}$):} Varying $-\frac{\kappa_c}{2}\int d^4x\,\sqrt{-g}\,g^{\alpha\beta}(D_\alpha\psi)^*(D_\beta\psi)$ gives the symmetrized gauge-invariant stress--energy tensor
  \begin{equation}
  T_{\mu\nu}^{(\psi)}
  = \kappa_c\left(\mathrm{Re}\bigl[(D_\mu\psi)^*(D_\nu\psi)\bigr] - \frac{1}{2}g_{\mu\nu}\,g^{\alpha\beta}(D_\alpha\psi)^*(D_\beta\psi)\right),
  \label{eq:T_psi_sec4}
  \end{equation}
  with $D_\mu\psi = (\nabla_\mu - i\hat{q}A_\mu)\psi$, sourcing curvature directly from coherence amplitude gradients and gauge-covariant phase supercurrents.
  \item \textbf{Coherence-Dressed Maxwell/Proca Sector ($T_{\mu\nu}^{(\mathrm{EM})}$):} Varying $-\frac{1}{4}\int d^4x\,\sqrt{-g}\,f_{\mathrm{EM}}(\Rcoh)F_{\alpha\beta}F^{\alpha\beta}$ yields
  \begin{equation}
  T_{\mu\nu}^{(\mathrm{EM})}
  = f_{\mathrm{EM}}(\Rcoh)\left(F_{\mu\lambda}F_\nu^{\ \lambda} - \frac{1}{4}g_{\mu\nu}F_{\alpha\beta}F^{\alpha\beta}\right),
  \qquad
  f_{\mathrm{EM}}(\Rcoh) = 1 + c_{\mathrm{EM}}\Rcoh^2 \ge 1.
  \label{eq:T_EM_sec4}
  \end{equation}
  \item \textbf{Potential and Portal Linkage Sector ($T_{\mu\nu}^{(\mathrm{pot})}$):} Since $V_{\mathrm{tot}}(\rho_t, |\psi|) = V(\rho_t) + U(|\psi|) + U_{\mathrm{link}}(\rho_t,|\psi|) \ge 0$ depends on the metric exclusively through $\sqrt{-g}$, its variation contributes
  \begin{equation}
  T_{\mu\nu}^{(\mathrm{pot})} = -g_{\mu\nu}\,V_{\mathrm{tot}}(\rho_t, |\psi|),
  \label{eq:T_pot_sec4}
  \end{equation}
  yielding $T_{00}^{(\mathrm{pot})} = +V_{\mathrm{tot}}(\rho_t, |\psi|) \ge 0$ in $(-,+,+,+)$ and vanishing identically at the equilibrium vacuum $(\rho_0, v_0)$.
  \item \textbf{Coherence-Dressed Matter Sector ($T_{\mu\nu}^{(\mathrm{matter})}$):} Varying $+\int d^4x\,\sqrt{-g}\,f_{\mathrm{SM}}(\Rcoh)\mathcal{L}_{\mathrm{matter}}$ gives
  \begin{equation}
  T_{\mu\nu}^{(\mathrm{matter})}
  = f_{\mathrm{SM}}(\Rcoh)\,T_{\mu\nu}^{\mathrm{SM}},
  \qquad
  f_{\mathrm{SM}}(\Rcoh) = 1 + \frac{\Rcoh^2}{2 M_{\mathrm{sym}}^2} \ge 1,
  \label{eq:T_matt_sec4}
  \end{equation}
  where $T_{\mu\nu}^{\mathrm{SM}} := -\frac{2}{\sqrt{-g}}\frac{\delta(\sqrt{-g}\mathcal{L}_{\mathrm{matter}})}{\delta g^{\mu\nu}}$ is the canonical Standard Model stress--energy tensor. This coherence dressing modulates the effective gravitational coupling of matter in unscreened coherent regimes, grounding the equivalence-principle tests of Sections~\ref{sec:smoking-gun-bec} and~\ref{sec:experimental-predictions}.
\end{enumerate}

\subsection{Summary of Operational Definitions}

The operational constructs $\rho_t$ and $\Rcoh$ translate across physical and biological domains without ambiguity:
\begin{itemize}
  \item \textbf{Coincidence as a Bit:} Rather than an abstract binary voltage, an operational bit corresponds to a registered spatiotemporal confluence of events, elevating local $\Rcoh$.
  \item \textbf{Self-Aware Networks:} Oscillatory synchrony and phase-wave differentials in neural networks realize mesoscopic high-coherence states measurable via EEG, MEG, and multi-electrode arrays (Section~\ref{sec:neuroscience-implications}).
  \item \textbf{Coherence as an Order Parameter:} In quantum optics and cold-atom physics, operational coherence proxies track interferometric visibility and condensate fraction (kept strictly distinct from the classical EFT modulus $\Rcoh = |\psi|/\psi_0$ pending \texttt{OPEN-6C}, Section~\ref{sec:PrimitiveFieldDefinitions}); in finite-dimensional open quantum systems, monotonicity is governed by unital CPTP purity non-increase and general CPTP relative-entropy contraction (Result~\ref{thm:CoherenceMonotonicity}).
\end{itemize}

\subsection{Coherence--Time Coupling and State-Dependent Gravitational Effects}

While classical General Relativity couples geometry universally to total $T_{\mu\nu}$ regardless of internal phase coherence, SIT introduces an explicit, screened state-dependent dressing alongside scalar kinetic and linkage stress--energy:
\begin{enumerate}[leftmargin=2em]
  \item \textbf{State-Dependent Sourcing:} Regions of elevated macroscopic phase alignment ($\Rcoh > R_0$) dress matter stress--energy via $f_{\mathrm{SM}}(\Rcoh)$ and excite $(\rho_t, |\psi|)$ gradients.
  \item \textbf{Covariant Energy Redistribution:} Total energy--momentum remains strictly conserved on-shell ($\nabla^\mu T_{\mu\nu}^{\mathrm{tot}} = 0$) by diffeomorphism invariance; work expended to prepare or destroy macroscopic coherence exchanges energy conservatively among matter, electromagnetic, and $(\rho_t, \psi)$ sectors.
  \item \textbf{Clock-Rate Modulation via Linkage:} Spatial variations in $|\psi|$ pull $\rho_t$ through the portal potential $U_{\mathrm{link}}(\rho_t, |\psi|)$, linking local clock rates to coherence profiles without introducing extra temporal dimensions.
\end{enumerate}

\paragraph{SIT Gravity Postulates (GP1--GP7).}
\begin{itemize}
  \item \textbf{GP1 (Field Content):} Gravitational and electromagnetic deviations from GR+Maxwell are mediated by $\psi(x) = |\psi(x)|e^{i\theta(x)}$ ($\Rcoh = |\psi|/\psi_0$) and $\rho_t(x)$ on $(\mathcal{M}^4, g_{\mu\nu})$.
  \item \textbf{GP2 (Coherence Conservation):} Global $U(1)$ phase invariance yields the conserved Noether current $\nabla_\mu J^\mu_{\mathrm{coh}} = 0$ (Section~\ref{sec:coherence-conservation}).
  \item \textbf{GP3 (Coherence--Time Linkage):} Along the portal valley of $U_{\mathrm{link}}(\rho_t, |\psi|)$, $\rho_t - \rho_0 = \frac{\sqrt{\lambda_{\mathrm{portal}}}}{m_{\mathrm{link}}}(|\psi|^2 - v_0^2)$. In the separate exploratory register-level phenomenology of Section~\ref{sec:QIP} (not derived from unchanged \texttt{M\_AUTH\_SIT\_EFT\_v1}, \texttt{OPEN-6C}), any proposed state-dependent proper-time ansatz takes the form $d\tau/dt = \rho_g(x)[1 + \gamma\,\mathcal{C}_{\mathrm{state}}[\rho] + O(\gamma^2)]$ using a declared \emph{nonlinear} state functional $\mathcal{C}_{\mathrm{state}}[\rho]$ (or a state-specific linear witness $\hat{W}_\Psi$), explicitly replacing the retired impossible operator $\hat{\mathcal{R}}_{\mathrm{coh}}$.
  \item \textbf{GP4 (Curvature Source):} Metric variation yields $G_{\mu\nu} = 8\pi G\,T_{\mu\nu}^{\mathrm{tot}}$, reducing in the static weak-field limit around $(\rho_0, v_0)$ to a Yukawa-screened Poisson law.
  \item \textbf{GP5 (GR Recovery and Coupled Consistency):} At the constant broken-phase vacuum $(\rho_t, |\psi|) = (\rho_0, v_0)$ with $D_\mu\psi = 0$, $V_{\mathrm{tot}} = 0$ and $\nabla V_{\mathrm{tot}} = 0$, reducing the metric equation \emph{algebraically} to $G_{\mu\nu} = 8\pi G f_{\mathrm{SM}}(R_0) T_{\mu\nu}^{\mathrm{SM}} + 8\pi G f_{\mathrm{EM}}(R_0) T_{\mu\nu}^{\mathrm{EM},0}$; this yields an exact coupled solution when the scalar and gauge equations also vanish (such as source-free vacuum $\mathcal{L}_{\mathrm{matter}} = 0, F_{\mu\nu} = 0$, whereas $\mathcal{L}_{\mathrm{matter}} = -\rho_m < 0$ or $\hat{q}v_0 F_{\mu\nu} \ne 0$ excites scalar/Proca equations at $v_0 > 0$), while broader exact sourced GR+Maxwell ($f_{\mathrm{SM}} = f_{\mathrm{EM}} = 1, m_A = 0$) holds in the separate unbroken parameter limit $v_0 = 0$, and heavy-scalar decoupling provides a controlled asymptotic approximation.
  \item \textbf{GP6 (Operational Dictionary):} Normalized operational coherence proxies approaching $1$ mark maximal phase alignment, while proxies approaching $0$ mark complete local decoherence (distinct from the unbounded classical field ratio $\Rcoh = |\psi|/\psi_0 \ge 0$).
  \item \textbf{GP7 (Laboratory Channels):} State-dependent free-fall comparisons in $^{87}\mathrm{Rb}$ BECs and differential optical-lattice clock redshift measurements provide prospective falsification channels once microscopic source and detector bridges (\texttt{OPEN-6}) are derived.
\end{itemize}

\subsection{Relativistic Clarification: Internal vs.\ External Coherence Perspectives}

External laboratory observers see ultrarelativistic wavepackets cycle rapidly in coordinate time ($\Delta t$ small), whereas in the particle's comoving proper frame ($\Delta\tau = \gamma^{-1}\Delta t$) internal phase evolution is Lorentz-dilated, preserving $\Delta E\,\Delta t \ge \hbar/2$. Because $\rho_t(x)$ is a genuine Lorentz scalar rather than a coordinate time axis, SIT requires no multitemporal signature $(-, -, +, +)$ (which would introduce fatal Ostrogradsky/kleinian ghosts, purged as branch \texttt{C05} under \texttt{ODQ-2}): proper causality is governed everywhere by a single Lorentzian metric $g_{\mu\nu}$ of signature $(-,+,+,+)$.

In phenomenological regimes where $\rho_t$ tracks $\Rcoh$ monotonically around $\rho_0$, one may write the smooth adiabatic representation $\rho_t(x) \approx \rho_0\exp[\alpha\Rcoh(x)] = \rho_0[1 + \alpha\Rcoh(x) + O(\alpha^2)]$, or equivalently via the quadratic portal minimum of \texttt{M\_AUTH\_SIT\_EFT\_v1}, $\rho_t - \rho_0 = (\sqrt{\lambda_{\mathrm{portal}}}/m_{\mathrm{link}})(|\psi|^2 - v_0^2)$. In both parameterizations, constant fields at a stationary vacuum recover standard General Relativity with constant effective gravitational coupling.

\subsection{Determining the Functional Form Linking Coherence--Decoherence to Local Time Density}

Integrating deterministic oscillatory synchronization (\emph{SuperTimePosition}), local difference dissipation (\emph{Micah's New Law of Thermodynamics}), and interference-driven clock modulation (\emph{Super Dark Time}) motivates mesoscopic analogies between oscillator ensembles and the continuum EFT fields $(\Rcoh, \rho_t)$ (distinct from canonical action derivations).

\subsection{Synchronization and Informational Coherence}

For an ensemble of $N$ coupled phase oscillators with instantaneous phases $\theta_j(t)$, the Kuramoto complex order parameter
\begin{equation}
r e^{i\Theta} = \frac{1}{N}\sum_{j=1}^N e^{i\theta_j}
\label{eq:kuramoto_order_sec4}
\end{equation}
provides a discrete mesoscopic analogy for a bounded phase-alignment amplitude $r \in [0,1]$ and collective phase $\Theta$.

\subsection{Wave-Driven Dissipation and Equilibrium Dynamics (Heuristic Oscillator Analogy)}

In open oscillator networks exchanging phase signals, pairwise phase differences $\Delta\phi_{ij} = \phi_i - \phi_j$ satisfy $\sin(\Delta\phi_{ij}) = -\sin(\Delta\phi_{ji})$, so the antisymmetric double sum $\sum_{i,j}\sin(\Delta\phi_{ij})$ vanishes identically; earlier drafts writing $d\rho_t/dt = \alpha_{\mathrm{sync}}\sum_{i,j}\sin(\Delta\phi_{ij})$ therefore yielded $d\rho_t/dt \equiv 0$. A non-vanishing heuristic synchronization-rate proxy instead uses the symmetric phase-alignment defect:
\begin{equation}
\frac{d\rho_t^{\mathrm{osc}}}{dt} = -\alpha_{\mathrm{sync}}\frac{1}{N^2}\sum_{i,j=1}^N \bigl[1 - \cos(\Delta\phi_{ij})\bigr] = -\alpha_{\mathrm{sync}}\bigl(1 - r^2\bigr) \le 0,
\end{equation}
which vanishes precisely at full phase lock ($r = 1$) and is explicitly marked as an exploratory oscillator-network proxy rather than an Euler--Lagrange equation of \texttt{M\_AUTH\_SIT\_EFT\_v1}.

\subsection{Quantum--Gravitational Interference Interpretation}

For a superposition $\sum_n C_n e^{i\phi_n}$ of local modes, interference cross-terms modulate the local coherence expectation and induced time-density shift according to
\begin{equation}
\rho_t(\Rcoh) = \rho_0 + \gamma_{\mathrm{int}}\,\mathrm{Re}\!\left[\sum_{n,m} C_n C_m^* e^{i(\phi_n - \phi_m)}\right],
\end{equation}
linking constructive phase alignment ($\Delta\phi_{nm} \to 0$) to positive $\delta\rho_t$ shifts and destructive dephasing to relaxation back toward $\rho_0$.

\subsection{Empirical Validation Pathways}

Three complementary measurement modalities probe this linkage directly (contingent on the Stage-6 microphysical action-to-apparatus dictionaries in Section~\ref{sec:stage6-action-diagnostics} and \texttt{OPEN-6}): (i) optical-lattice atomic clock comparisons (testing terrestrial gravitational redshift at $(+1.4 \pm 9.1)\times 10^{-5}$ in the Tokyo Skytree campaign and achieving optical frequency-ratio measurement uncertainties at or below $3.2\times 10^{-18}$ in the NIST/BACON report; legacy $3\times 10^{-8}$ targets are quarantined as hypothetical future sensitivities); (ii) cold-atom Mach--Zehnder and fountain interferometers ($\sim 10^{-3}\,\mathrm{rad}$ phase resolution); and (iii) cluster-scale weak and strong gravitational lensing tomography ($\sim 1\%$ shear/delay precision).

\subsection{Information as Active Substrate}

Within this covariant architecture, spatial gradients $\nabla_i\Rcoh$ and covariant phase currents $D_\mu\psi$ enter $T_{\mu\nu}^{\mathrm{tot}}$ on equal mathematical footing with conventional matter fields, elevating structured correlation from passive bookkeeping to a dynamically tested field sector.

\subsection{Coherence as Informational Teamwork}

Self-interaction in $U(|\psi|) = \frac{\lambda_\psi}{4}(|\psi|^2 - v_0^2)^2$ ($\lambda_\psi > 0$) stabilizes collective condensate formation much like a Ginzburg--Landau superfluid or phased antenna array: phase-aligned modes reinforce a coherent expectation value $|\psi| = v_0$, whereas incoherent thermal excitations average out in phase without shifting the macroscopic order parameter.

\subsection{Integration into the Mathematical Framework}

Across continuum PDEs, discrete oscillator lattices, and $U(1)$ principal bundles, resonance, geometric holonomy, and coincidence counting refer to projections of the same gauge-coupled complex scalar $\psi(x)$ and real scalar $\rho_t(x)$, whose canonical effective field theory we now state in complete precision.

\section{What the Theory Is: Mathematical Framework (\texttt{M\_AUTH\_SIT\_EFT\_v1})}
\label{sec:framework}
\label{sec:quantitative-coupling}

\subsection{Primitive Fields, Spacetime Conventions, and Dimensional Ledger}

The authoritative continuum effective field theory \texttt{M\_AUTH\_SIT\_EFT\_v1} (canon \texttt{v1\_1}, locked across \texttt{ODQ-1} through \texttt{ODQ-7}) is defined on a four-dimensional time-oriented Lorentzian manifold $(\mathcal{M}^4, g_{\mu\nu})$ in mostly-plus metric signature $(-,+,+,+)$ with Minkowski reference $\eta_{\mu\nu} = \mathrm{diag}(-1,+1,+1,+1)$ and natural units $\hbar = c = 1$ ($[M] = \mathrm{GeV}$). All field and parameter mass dimensions are locked strictly as follows:
\begin{itemize}
  \item \textbf{Dimension $[M]^{-2}$:} Gravitational coupling $\kappa_G := 8\pi G = M_{\mathrm{Pl}}^{-2} > 0$.
  \item \textbf{Dimension $[M]^0$ (dimensionless):} Metric tensor $g_{\mu\nu}$, dimensionless coherence ratio $\Rcoh := |\psi|/\psi_0 \ge 0$, vacuum coherence ratio $R_0 := v_0/\psi_0 \ge 0$, electromagnetic gating coefficient $c_{\mathrm{EM}} \ge 0$, dimensionless symmetron scale parameter $M_{\mathrm{sym}} > 0$, quartic self-couplings $\lambda_\rho > 0$ and $\lambda_\psi > 0$, portal coupling $\lambda_{\mathrm{portal}} \ge 0$, $U(1)$ charge coupling $\hat{q} \ge 0$, and positive kinetic normalizations $\kappa_t > 0$, $\kappa_c > 0$.
  \item \textbf{Dimension $[M]^1$ ($\mathrm{GeV}$):} Real temporal-density scalar $\rho_t$, complex coherence scalar $\psi = u + iv = |\psi|e^{i\theta}$, $U(1)$ gauge connection $A_\mu$ with curvature $F_{\mu\nu} = \partial_\mu A_\nu - \partial_\nu A_\mu$ and covariant derivative $D_\mu\psi = (\partial_\mu - i\hat{q}A_\mu)\psi$, together with reference scales $\psi_0 > 0$, $\rho_0 > 0$, $v_0 = \psi_0 R_0 \ge 0$, $m_{\mathrm{link}} > 0$, $M_{\mathrm{Pl}}$, $\Lambda_{\mathrm{EM}}$, $\Lambda_{\mathrm{sym}}$, and $\Lambda_{\mathrm{UV}}$.
  \item \textbf{Dimension $[M]^4$ Lagrangian densities:} $\mathcal{L}_{\mathrm{EH}}$, $\mathcal{L}_{\mathrm{coh}}$, $\mathcal{L}_{\mathrm{EM}}^{\mathrm{coh}}$, and $\mathcal{L}_{\mathrm{SM}}^{\mathrm{coh}}$.
  \item \textbf{Dimension $[M]^6$ EFT operators:} $-\frac{c_{\mathrm{EM}}}{4\psi_0^2}|\psi|^2 F_{\mu\nu}F^{\mu\nu} = -\frac{1}{4\Lambda_{\mathrm{EM}}^2}|\psi|^2 F_{\mu\nu}F^{\mu\nu}$ and $+\frac{|\psi|^2}{2\psi_0^2 M_{\mathrm{sym}}^2}\mathcal{L}_{\mathrm{matter}} = +\frac{|\psi|^2}{\Lambda_{\mathrm{sym}}^2}\mathcal{L}_{\mathrm{matter}}$.
\end{itemize}

Operationally, a reference clock of proper frequency $\omega_0$ in the baseline environment $\rho_t = \rho_0$ ticks locally at $\omega(x) = \omega_0(\rho_0/\rho_t(x))$, yielding fractional shift $\Delta\omega/\omega \approx -\delta\rho_t/\rho_0$ to first order. Local $U(1)$ gauge transformations act via $\psi(x) \mapsto e^{i\hat{q}\chi(x)}\psi(x)$ and $A_\mu(x) \mapsto A_\mu(x) + \partial_\mu\chi(x)$, leaving $|\psi|$, $\Rcoh$, $\rho_t$, $F_{\mu\nu}$, and $|D_\mu\psi|^2$ strictly invariant.

\subsection{The Canonical Master Action and Vacuum-Positive Couplings}

The authoritative effective action is
\begin{equation}
S_{\mathrm{AUTH}}[g,\rho_t,\psi,A_\mu,\Psi_{\mathrm{SM}}]
= \int d^4x\,\sqrt{-g}\Bigl[
  \mathcal{L}_{\mathrm{EH}}
  + \mathcal{L}_{\mathrm{coh}}
  + \mathcal{L}_{\mathrm{EM}}^{\mathrm{coh}}
  + \mathcal{L}_{\mathrm{SM}}^{\mathrm{coh}}
\Bigr],
\label{eq:canonical_action_sec5}
\end{equation}
where the four covariant sectors are defined unambiguously in signature $(-,+,+,+)$ by
\begin{align}
\mathcal{L}_{\mathrm{EH}} &= \frac{R}{16\pi G} = \frac{R}{2\kappa_G} = \frac{M_{\mathrm{Pl}}^2}{2}\,R, \label{eq:L_EH_def}\\
\mathcal{L}_{\mathrm{coh}} &= -\frac{\kappa_t}{2}\,g^{\mu\nu}(\partial_\mu\rho_t)(\partial_\nu\rho_t)
  - \frac{\kappa_c}{2}\,g^{\mu\nu}(D_\mu\psi)^*(D_\nu\psi)
  - V(\rho_t) - U(|\psi|) - U_{\mathrm{link}}(\rho_t,\psi), \label{eq:L_coh_def}\\
\mathcal{L}_{\mathrm{EM}}^{\mathrm{coh}} &= -\frac{1}{4}\,f_{\mathrm{EM}}(\Rcoh)\,F_{\mu\nu}F^{\mu\nu},
  \qquad f_{\mathrm{EM}}(\Rcoh) := 1 + c_{\mathrm{EM}}\Rcoh^2 \ge 1, \label{eq:L_EM_def}\\
\mathcal{L}_{\mathrm{SM}}^{\mathrm{coh}} &= +f_{\mathrm{SM}}(\Rcoh)\,\mathcal{L}_{\mathrm{matter}},
  \qquad f_{\mathrm{SM}}(\Rcoh) := 1 + \frac{\Rcoh^2}{2 M_{\mathrm{sym}}^2} \ge 1. \label{eq:L_SM_def}
\end{align}

\paragraph{Potentials, $\mathbb{Z}_2$ Degeneracy Lifting, and Positive-Definite Vacuum Hessian.}
The dimension-$\le 4$ polynomial scalar potential block $V_{\mathrm{tot}}(\rho_t,\psi) := V(\rho_t) + U(|\psi|) + U_{\mathrm{link}}(\rho_t,\psi)$ is locked as
\begin{align}
V(\rho_t) &= \frac{\lambda_\rho}{4}\bigl(\rho_t^2 - \rho_0^2\bigr)^2,
\qquad
U(|\psi|) = \frac{\lambda_\psi}{4}\bigl(|\psi|^2 - v_0^2\bigr)^2,
\qquad v_0 = \psi_0 R_0 \ge 0, \label{eq:V_U_canonical}\\
U_{\mathrm{link}}(\rho_t,\psi) &= \frac{1}{2}\Bigl[m_{\mathrm{link}}(\rho_t - \rho_0) - \sqrt{\lambda_{\mathrm{portal}}}\,\bigl(|\psi|^2 - v_0^2\bigr)\Bigr]^2.
\label{eq:U_link_canonical}
\end{align}
Because $V_{\mathrm{tot}}$ is a sum of three real squares with $\lambda_\rho, \lambda_\psi, m_{\mathrm{link}} > 0$, it satisfies $V_{\mathrm{tot}}(\rho_t,\psi) \ge 0$ globally on $\mathbb{R}\times\mathbb{C}$ (\texttt{potentialsTotal\_globally\_nonneg}), vanishing with zero gradient uniquely on the physical vacuum circle $\rho_t = \rho_0$, $|\psi| = v_0$ (\texttt{vacuumEquilibrium\_zeroCosmologicalConstant}). Crucially, the linear-in-$(\rho_t - \rho_0)$ portal term explicitly breaks the unphysical negative-clock $\mathbb{Z}_2$ mirror branch at $(-\rho_0, v_0)$:
\begin{equation}
V_{\mathrm{tot}}(-\rho_0, v_0) = 2 m_{\mathrm{link}}^2 \rho_0^2 > 0.
\label{eq:Z2_lifting_gap}
\end{equation}
Expanding $V_{\mathrm{tot}}$ around $(\rho_0, v_0)$ in radial fluctuations $\delta\Phi = (\delta\rho_t, \delta|\psi|)^\top$ with kinetic matrix $K = \mathrm{diag}(\kappa_t, \kappa_c) \succ 0$ yields the symmetric $2\times 2$ potential-curvature Hessian matrix $H$:
\begin{equation}
H = \begin{pmatrix}
2\lambda_\rho\rho_0^2 + m_{\mathrm{link}}^2 & -2 m_{\mathrm{link}}\sqrt{\lambda_{\mathrm{portal}}}\,v_0 \\[0.4em]
-2 m_{\mathrm{link}}\sqrt{\lambda_{\mathrm{portal}}}\,v_0 & 2 v_0^2\bigl(\lambda_\psi + 2\lambda_{\mathrm{portal}}\bigr)
\end{pmatrix},
\label{eq:radial_Hessian_sec5}
\end{equation}
with trace $\mathrm{tr}(H) > 0$ and determinant
\begin{equation}
\det(H) = 2 v_0^2\Bigl(2\lambda_\rho\lambda_\psi\rho_0^2 + \lambda_\psi m_{\mathrm{link}}^2 + 4\lambda_\rho\rho_0^2\lambda_{\mathrm{portal}}\Bigr).
\label{eq:det_H_positive}
\end{equation}
For $v_0 > 0$ and $\kappa_t, \kappa_c > 0$, $\det(H) > 0$ with Rayleigh margin $\lambda_{\min}(H) \ge \det(H)/\mathrm{tr}(H) > 0$. Because $K = \mathrm{diag}(\kappa_t, \kappa_c) \succ 0$ can carry nontrivial kinetic normalizations, the \emph{physical squared eigenmasses} $(m_1^2, m_2^2)$ and canonical orthogonal mixing angle $\vartheta$ are the eigenvalues and rotation angle of the canonically normalized mass matrix $M_{\mathrm{phys}}^2 := K^{-1/2} H K^{-1/2}$ (equivalently solving the generalized eigenproblem $H u_a = m_a^2 K u_a$, Section~\ref{sec:stage6-action-diagnostics} and Appendix~E.5), which are strictly positive by Sylvester's law of inertia ($\lambda_{\min}(M_{\mathrm{phys}}^2) \ge \lambda_{\min}(H)/\max(\kappa_t, \kappa_c) > 0$) and coincide with the raw eigenvalues of $H$ when $K = I_2$ ($\kappa_t = \kappa_c = 1$). In the stiff-link decoupling limit $m_{\mathrm{link}} \to \infty$, the Schur complement $(H_{22} - H_{12}^2/H_{11})/\kappa_c \to 2\lambda_\psi v_0^2/\kappa_c > 0$ remains smooth and positive.

\paragraph{Elimination of the Zero-Coherence $1/f$ Singularity.}
Both constitutive dressings $f_{\mathrm{EM}}(\Rcoh) \ge 1$ and $f_{\mathrm{SM}}(\Rcoh) \ge 1$ are bounded away from zero across all field configurations, including $\Rcoh \to 0$. In decoherent empty space ($\Rcoh \to 0$), $f_{\mathrm{EM}}(0) = f_{\mathrm{SM}}(0) = 1$, recovering standard Maxwell electrodynamics without dividing canonical electric momenta $\Pi^2/(2f_{\mathrm{EM}})$ by zero (the singularity that required quarantining the 2025 degenerate-vacuum draft \texttt{C09} from continuum EFT under \texttt{ODQ-3}, resolved non-perturbatively on the 4D Haar Wilson lattice in Sections~\ref{sec:mass-gap} and~\ref{sec:stage4-lattice}).

\subsubsection{Classification as a Two-Scalar--Vector--Tensor EFT and Cutoff Crosswalk}

Using CODATA 2022 physical constants ($\hbar = 1.054571817\times 10^{-34}\,\mathrm{J\,s}$, $c = 299792458\,\mathrm{m\,s^{-1}}$, $G = 6.6743\times 10^{-11}\,\mathrm{m^3\,kg^{-1}\,s^{-2}}$, $1\,\mathrm{GeV} = 1.602176634\times 10^{-10}\,\mathrm{J}$), the three Planck mass conventions in the literature and the two dimension-6 operator cutoffs satisfy exact identities (\texttt{ADV\_D1\_01}):
\begin{align}
M_{\mathrm{Pl,unred}} &= \sqrt{\frac{\hbar c}{G}}\,\frac{c^2}{\mathrm{GeV}} = 1.220890128209864\times 10^{19}\,\mathrm{GeV}, \nonumber\\
M_{\mathrm{Pl}} &:= \kappa_G^{-1/2} = (8\pi G)^{-1/2} = \frac{M_{\mathrm{Pl,unred}}}{\sqrt{8\pi}} = 2.4353234593382036\times 10^{18}\,\mathrm{GeV}, \label{eq:Planck_scales_crosswalk}\\
M_{\mathrm{Pl,QG31}} &:= (2\kappa_G)^{-1/2} = (16\pi G)^{-1/2} = \frac{M_{\mathrm{Pl}}}{\sqrt{2}} = 1.7220337324807252\times 10^{18}\,\mathrm{GeV}, \nonumber
\end{align}
paired with dimension-6 gauge-gating and symmetron scales
\begin{equation}
\Lambda_{\mathrm{EM}} := \frac{\psi_0}{\sqrt{c_{\mathrm{EM}}}} \quad (+\infty\ \text{if } c_{\mathrm{EM}}=0),
\qquad
\Lambda_{\mathrm{sym}} := \sqrt{2}\,\psi_0 M_{\mathrm{sym}},
\qquad
\Lambda_{\mathrm{UV}} := \min\bigl(M_{\mathrm{Pl}},\,\Lambda_{\mathrm{EM}},\,\Lambda_{\mathrm{sym}}\bigr).
\label{eq:Lambda_UV_ceiling}
\end{equation}
For any $\Rcoh > 0$ and $c_{\mathrm{EM}} > 0$, the vacuum-dressing deviations obey the exact structural identity
\begin{equation}
\left(\frac{\Lambda_{\mathrm{sym}}}{\Lambda_{\mathrm{EM}}}\right)^2
= 2 c_{\mathrm{EM}} M_{\mathrm{sym}}^2
= \frac{c_{\mathrm{EM}}\Rcoh^2}{\Rcoh^2 / (2 M_{\mathrm{sym}}^2)}
= \frac{f_{\mathrm{EM}}(\Rcoh) - 1}{f_{\mathrm{SM}}(\Rcoh) - 1},
\label{eq:cutoff_ratio_identity}
\end{equation}
guaranteeing radiative stability across electromagnetic and matter-dressing channels whenever $M_{\mathrm{sym}} \gg 1$ and $c_{\mathrm{EM}} \ll 1$ (e.g.\ canonical benchmark \texttt{BM\_PHYS\_EFT\_1} in Appendix~\ref{app:epistemic-ledger-lean}, where $\psi_0 = 2\times 10^{16}\,\mathrm{GeV}$, $M_{\mathrm{sym}} = 50$, $c_{\mathrm{EM}} = 10^{-4}$ yield $\Lambda_{\mathrm{EM}} = 2.0\times 10^{18}\,\mathrm{GeV}$ and $\Lambda_{\mathrm{sym}} = \Lambda_{\mathrm{UV}} = 1.4142\times 10^{18}\,\mathrm{GeV}$).

\subsection{Canonical Stress--Energy vs.\ Separately Postulated Phenomenological Mass-Density Ansatz}

In unchanged \texttt{M\_AUTH\_SIT\_EFT\_v1}, gravitational sourcing is governed strictly by the variational stress--energy tensor $T_{\mu\nu}^{\mathrm{tot}}$ of Eq.~\eqref{eq:covariant_einstein_sec5}, which vanishes identically at the constant vacuum $(\rho_0, v_0, A_\mu = 0)$ where all gradients and $V_{\mathrm{tot}}(\rho_0, v_0)$ vanish. Separately, earlier phenomenological SIT notes (Section~\ref{sec:informational_energy_equivalence}) explored an auxiliary mass-density ansatz
\begin{equation}
\rho_{\mathrm{info}}^{\mathrm{phen}}(x) := \frac{\varepsilon_{\mathrm{SIT}}^{\mathrm{phen}}(x)}{c^2} = \frac{\zeta}{c^2}\,\Rcoh(x)\,\bigl[\rho_t(x)\bigr]^2,
\qquad
M_{\mathrm{info}}^{\mathrm{phen}} = \int_{\mathcal{V}} \rho_{\mathrm{info}}^{\mathrm{phen}}(x)\,d^3x,
\label{eq:informational_mass_def}
\end{equation}
where $\rho_{\mathrm{info}}^{\mathrm{phen}}$ is a \emph{mass density} (not a total mass without volume integration). Two strict boundaries apply:
\begin{enumerate}[leftmargin=2em]
  \item \textbf{Separate Specification (Not Derived from Unchanged \texttt{M\_AUTH\_SIT\_EFT\_v1}):} Because the canonical vacuum has $T_{\mu\nu}^{(\rho_t)} + T_{\mu\nu}^{(\psi)} + T_{\mu\nu}^{(\mathrm{pot})} = 0$ at $(\rho_0, v_0)$ whereas $\zeta R_0 \rho_0^2 > 0$, Eq.~\eqref{eq:informational_mass_def} is \emph{not} a consequence of unchanged \texttt{M\_AUTH\_SIT\_EFT\_v1}; connecting it to the canonical stress--energy tensor would require an explicit vacuum-subtraction scheme or an additional covariant operator with proved stress--energy conservation (Section~\ref{sec:informational_energy_equivalence}).
  \item \textbf{Standard Model Rest Mass Preserved in Decoherent Limits:} Ordinary Standard Model rest masses carried inside $\mathcal{L}_{\mathrm{matter}}$ never disappear when $\Rcoh \to 0$; at $v_0 \to 0$, $f_{\mathrm{SM}}(0) = 1$ and gravitational sourcing reverts smoothly to standard General Relativity coupled to $T_{\mu\nu}^{\mathrm{SM}}$.
\end{enumerate}

\subsection{Field Equations, Bianchi--Noether Identity, and Dynamical Sources}
\label{sec:stress-energy}

\subsubsection{Covariant Euler--Lagrange Field Equations}

Varying $S_{\mathrm{AUTH}}$ with respect to $g^{\mu\nu}$, $\rho_t$, $\psi^*$, and $A_\nu$ yields the closed coupled second-order field system of \texttt{M\_AUTH\_SIT\_EFT\_v1}:
\begin{enumerate}[leftmargin=2em]
  \item \textbf{Einstein Equation:}
  \begin{equation}
  G_{\mu\nu} = \kappa_G\Bigl(T_{\mu\nu}^{(\rho_t)} + T_{\mu\nu}^{(\psi)} + T_{\mu\nu}^{(\mathrm{EM})} + T_{\mu\nu}^{(\mathrm{pot})} + f_{\mathrm{SM}}(\Rcoh)\,T_{\mu\nu}^{\mathrm{SM}}\Bigr).
  \label{eq:covariant_einstein_sec5}
  \end{equation}
  \item \textbf{Temporal-Density Scalar Equation of Motion:}
  \begin{equation}
  \kappa_t\,\Box_g\rho_t
  - \Bigl[\lambda_\rho\rho_t(\rho_t^2 - \rho_0^2) + m_{\mathrm{link}}\bigl(m_{\mathrm{link}}(\rho_t - \rho_0) - \sqrt{\lambda_{\mathrm{portal}}}\,(|\psi|^2 - v_0^2)\bigr)\Bigr] = 0,
  \label{eq:rhot_eom_canonical}
  \end{equation}
  where $\Box_g\rho_t := \frac{1}{\sqrt{-g}}\partial_\mu(\sqrt{-g}\,g^{\mu\nu}\partial_\nu\rho_t)$ reduces in flat Minkowski space to $(-\partial_t^2 + \nabla^2)\rho_t$.
  \item \textbf{Complex Coherence Scalar Equation of Motion:}
  \begin{align}
  \kappa_c\,D_\mu D^\mu\psi
  &- \Biggl[
    \lambda_\psi\bigl(|\psi|^2 - v_0^2\bigr)
    - 2\sqrt{\lambda_{\mathrm{portal}}}\,\Bigl(m_{\mathrm{link}}(\rho_t - \rho_0) - \sqrt{\lambda_{\mathrm{portal}}}\,\bigl(|\psi|^2 - v_0^2\bigr)\Bigr)
    \nonumber\\
  &\qquad\quad
    + \frac{c_{\mathrm{EM}}}{2\psi_0^2}\,F_{\alpha\beta}F^{\alpha\beta}
    - \frac{1}{\psi_0^2 M_{\mathrm{sym}}^2}\,\mathcal{L}_{\mathrm{matter}}
  \Biggr]\psi = 0.
  \label{eq:Rcoh-eom}
  \end{align}
  \item \textbf{Coherence-Gated Maxwell/Proca Equation:}
  \begin{equation}
  \nabla_\mu\bigl(f_{\mathrm{EM}}(\Rcoh)\,F^{\nu\mu}\bigr)
  = \hat{q}\,J^\nu_{\mathrm{coh}} + J^\nu_{\mathrm{EM,SM}},
  \qquad
  J^\nu_{\mathrm{coh}} = \kappa_c|\psi|^2\bigl(\partial^\nu\theta - \hat{q}A^\nu\bigr).
  \label{eq:maxwell_proca_eom_canonical}
  \end{equation}
\end{enumerate}

\subsubsection{Contracted Bianchi--Noether Conservation Identity}

Although individual dressed sectors exchange energy--momentum proportionally to $\nabla_\nu f_{\mathrm{EM}}(\Rcoh)$ and $\nabla_\nu f_{\mathrm{SM}}(\Rcoh)$, off-shell diffeomorphism and $U(1)$ gauge invariance enforce exact derivative cancellation across Eqs.~\eqref{eq:rhot_eom_canonical}--\eqref{eq:maxwell_proca_eom_canonical} (verified symbolically in 4D and formally in Lean's \texttt{Jet1Plus1D} module), yielding exact on-shell conservation of the total stress--energy tensor:
\begin{equation}
\nabla^\mu G_{\mu\nu} \equiv 0
\implies
\nabla^\mu T_{\mu\nu}^{\mathrm{tot}} = 0.
\label{eq:bianchi_noether_total}
\end{equation}

\subsection{Physical Classification: Scalar--Tensor Effective Field Theory}

Within the standard scalar--tensor EFT taxonomy, $S_{\mathrm{AUTH}}$ retains a canonical Einstein-frame Ricci term $\frac{M_{\mathrm{Pl}}^2}{2}R$ while coupling the gauge-invariant coherence modulus $\Rcoh^2$ quadratically to $F_{\mu\nu}F^{\mu\nu}$ and $\mathcal{L}_{\mathrm{matter}}$ and linking $\rho_t$ to $|\psi|^2$ via $U_{\mathrm{link}}$. When $(\rho_t, |\psi|)$ freezes to its vacuum expectation $(\rho_0, v_0)$, all scalar gradients vanish, $V_{\mathrm{tot}}(\rho_0, v_0) = 0$, and General Relativity is recovered with rescaled constants $G_{\mathrm{eff}} = G f_{\mathrm{SM}}(R_0)$ and $\alpha_{\mathrm{EM,eff}} = \alpha_{\mathrm{bare}}/f_{\mathrm{EM}}(R_0)$ (conditional on identifying $A_\mu$ with the visible electromagnetic sector, \texttt{OPEN-6B}).

\subsection{Electromagnetic Holonomy from the Coherence Phase}

Under local $U(1)$ gauge transformations $\psi \mapsto e^{i\hat{q}\chi}\psi$, $A_\mu \mapsto A_\mu + \partial_\mu\chi$, the field strength $F_{\mu\nu} = \partial_\mu A_\nu - \partial_\nu A_\mu$ is the curvature two-form of a principal $U(1)$ bundle over $\mathcal{M}^4$. On any simply connected pure-gauge patch where $F_{\mu\nu} = 0$, one may set $A_\mu = \hat{q}^{-1}\partial_\mu\theta$, recovering Aharonov--Bohm phase holonomy $\exp\bigl(i\hat{q}\oint_\gamma A_\mu dx^\mu\bigr)$; globally, whenever $F_{\mu\nu} \ne 0$, $A_\mu$ is a genuine dynamical connection not globally reducible to the gradient of a single-valued scalar phase (Sections~\ref{sec:magnetism-holonomy} and~\ref{sec:stage4-matter-holonomy}).

\section{Why the Theory Is Safe: Screening Mechanism, Shapiro Null-Cone Invariance, and the Newtonian Limit}
\label{sec:screening-mechanism}

To reconcile state-dependent coherence couplings with stringent Solar System tests (Cassini radioscience, Lunar Laser Ranging, E\"ot-Wash torsion balances), \texttt{M\_AUTH\_SIT\_EFT\_v1} incorporates an exact quadratic $\mathbb{Z}_2$-symmetric matter dressing $f_{\mathrm{SM}}(\Rcoh) = 1 + \Rcoh^2/(2 M_{\mathrm{sym}}^2)$ and isotropic electromagnetic gating $f_{\mathrm{EM}}(\Rcoh) = 1 + c_{\mathrm{EM}}\Rcoh^2$. We analyze coupled two-field environmental symmetron screening---distinguishing exact two-field stationarity and stability conditions from decoupled or stiff-portal single-field approximations---and derive consistent conditional conformal-frame post-Newtonian bounds.

\subsection{Environmental Ground-State Selection, Coupled Stationarity, Two-Field Stability, and Fixed-Equilibrium Stiffening}

For non-relativistic dust with $\mathcal{L}_{\mathrm{matter}} = -\rho_m \le 0$, the matter-dressed Lagrangian $+f_{\mathrm{SM}}(\Rcoh)\mathcal{L}_{\mathrm{matter}} = -\rho_m - \frac{\rho_m}{2\psi_0^2 M_{\mathrm{sym}}^2}|\psi|^2$ adds a positive quadratic density term $+\frac{\rho_m}{2\psi_0^2 M_{\mathrm{sym}}^2}|\psi|^2$ to the effective potential for $(\rho_t, |\psi|)$:
\begin{equation}
V_{\mathrm{eff}}(\rho_t, |\psi|;\,\rho_m)
= V_{\mathrm{tot}}(\rho_t, |\psi|) + \frac{\rho_m}{2\psi_0^2 M_{\mathrm{sym}}^2}\,|\psi|^2.
\label{eq:symmetron_Veff_sec6}
\end{equation}
Three distinct statements govern this effective potential, and care is required when passing from the fixed-background Hessian to nonlinear two-field interior equilibria:
\begin{enumerate}[leftmargin=2em]
  \item \textbf{Fixed-Equilibrium Quadratic Stiffening (\texttt{ADV\_D1\_03}, Exact Algebraic Result):} Evaluated around a fixed reference equilibrium, static dust $\rho_m \ge 0$ contributes a strictly non-negative curvature stiffening $\Delta H_{22}(\rho_m) = +\frac{\rho_m}{\psi_0^2 M_{\mathrm{sym}}^2} \ge 0$ (\texttt{symmetronDustShift\_nonnegative\_stabilizing}), raising the effective radial stiffness inside dense media.
  \item \textbf{Coupled Nonlinear Stationarity Equations (Correcting the Portal-Valley Approximation):} A static homogeneous interior configuration $(\rho_{\mathrm{in}}^*, v_{\mathrm{in}}^*)$ must satisfy \emph{both} coupled stationarity equations $\partial_{\rho_t} V_{\mathrm{eff}} = 0$ and $\partial_{|\psi|} V_{\mathrm{eff}} = 0$ simultaneously. Crucially, minimizing the portal term $U_{\mathrm{link}}$ alone by setting $\rho_{\mathrm{in}}^{\mathrm{naive}} = \rho_0 - (\sqrt{\lambda_{\mathrm{portal}}}/m_{\mathrm{link}})v_0^2$ at $|\psi| = 0$ does \emph{not} in general solve $\partial_{\rho_t} V_{\mathrm{eff}} = 0$, because $V'(\rho_{\mathrm{in}}^{\mathrm{naive}}) = \lambda_\rho\rho_{\mathrm{in}}^{\mathrm{naive}}\bigl((\rho_{\mathrm{in}}^{\mathrm{naive}})^2 - \rho_0^2\bigr) \ne 0$ (for example, at $\lambda_\rho = \lambda_\psi = \lambda_{\mathrm{portal}} = \rho_0 = v_0 = \psi_0 = M_{\mathrm{sym}} = 1$ and $m_{\mathrm{link}} = 2$, $\rho_{\mathrm{in}}^{\mathrm{naive}} = 1/2$ yields $\partial_{\rho_t} V_{\mathrm{eff}}(1/2, 0) = -3/8 \ne 0$, with no dust $\rho_t$-derivative to cancel it). Instead, at the $\mathbb{Z}_2$/$U(1)$-symmetric interior locus $\psi = u + iv = 0$ (where $\partial_u V_{\mathrm{eff}} = \partial_v V_{\mathrm{eff}} = 0$ identically), exact stationarity requires $\rho_{\mathrm{in}}^*$ to solve the coupled cubic equation
  \begin{equation}
  \lambda_\rho\,\rho_{\mathrm{in}}^*\bigl((\rho_{\mathrm{in}}^*)^2 - \rho_0^2\bigr)
  + m_{\mathrm{link}}\Bigl(m_{\mathrm{link}}(\rho_{\mathrm{in}}^* - \rho_0) + \sqrt{\lambda_{\mathrm{portal}}}\,v_0^2\Bigr) = 0.
  \label{eq:exact_coupled_interior_cubic}
  \end{equation}
  \item \textbf{Full Two-Field Hessian Stability Criterion (Strict Minimum vs.\ Saddle and Degenerate Boundary Counterexamples):} Because $V_{\mathrm{eff}}$ is a multi-field potential in $(\rho_t, u, v)$ (with $|\psi|^2 = u^2 + v^2$), radial curvature alone is \emph{not} sufficient for local stability. At any stationary point $(\rho_{\mathrm{in}}^*, 0, 0)$ solving Eq.~\eqref{eq:exact_coupled_interior_cubic}, the mixed derivatives $\partial_{\rho_t}\partial_u V_{\mathrm{eff}} = \partial_{\rho_t}\partial_v V_{\mathrm{eff}} = \partial_u\partial_v V_{\mathrm{eff}} = 0$ vanish identically by $U(1)$ symmetry, so the full Hessian is diagonal:
  \begin{equation}
  H_{\mathrm{sym}}(\rho_{\mathrm{in}}^*, 0; \rho_m) = \mathrm{diag}\!\Bigl(H_{\rho\rho}(\rho_{\mathrm{in}}^*),\; H_{rr}(\rho_{\mathrm{in}}^*; \rho_m),\; H_{rr}(\rho_{\mathrm{in}}^*; \rho_m)\Bigr),
  \label{eq:exact_two_field_sym_hessian}
  \end{equation}
  with temporal-field curvature
  \begin{equation}
  H_{\rho\rho}(\rho_{\mathrm{in}}^*) := \partial_{\rho_t}^2 V_{\mathrm{eff}}(\rho_{\mathrm{in}}^*, 0) = \lambda_\rho\bigl(3(\rho_{\mathrm{in}}^*)^2 - \rho_0^2\bigr) + m_{\mathrm{link}}^2
  \label{eq:exact_H_rhorho}
  \end{equation}
  and radial curvature $H_{rr}(\rho_{\mathrm{in}}^*; \rho_m) := \partial_u^2 V_{\mathrm{eff}}(\rho_{\mathrm{in}}^*, 0; \rho_m) = (\rho_m - \rho_{m,\mathrm{crit}})/(\psi_0^2 M_{\mathrm{sym}}^2)$, where the exact coupled radial threshold is
  \begin{align}
  \rho_{m,\mathrm{crit}}
  &:= \psi_0^2 M_{\mathrm{sym}}^2\left[\lambda_\psi v_0^2 + 2\sqrt{\lambda_{\mathrm{portal}}}\Bigl(m_{\mathrm{link}}(\rho_{\mathrm{in}}^* - \rho_0) + \sqrt{\lambda_{\mathrm{portal}}}\,v_0^2\Bigr)\right] \nonumber\\
  &= \psi_0^2 M_{\mathrm{sym}}^2\left[\lambda_\psi v_0^2 - \frac{2\sqrt{\lambda_{\mathrm{portal}}}\,\lambda_\rho}{m_{\mathrm{link}}}\,\rho_{\mathrm{in}}^*\bigl((\rho_{\mathrm{in}}^*)^2 - \rho_0^2\bigr)\right].
  \label{eq:exact_coupled_screening_threshold}
  \end{align}
  Therefore, a \textbf{sufficient condition} for $(\rho_{\mathrm{in}}^*, 0)$ to be a strict local minimum is that \emph{both} eigenvalues are strictly positive: $H_{\rho\rho}(\rho_{\mathrm{in}}^*) > 0$ \emph{and} $\rho_m > \rho_{m,\mathrm{crit}}$ ($H_{rr} > 0$), while $H_{\rho\rho} \ge 0$ and $\rho_m \ge \rho_{m,\mathrm{crit}}$ are \textbf{necessary}. Earlier wording claiming local stability ``iff $\partial_{|\psi|}^2 V_{\mathrm{eff}} \ge 0$'' is refuted by two explicit exact counterexamples:
  \begin{itemize}
    \item \emph{Temporal-direction saddle counterexample ($H_{rr} > 0$ yet $H_{\rho\rho} < 0$):} Choosing $\lambda_\rho = \lambda_\psi = \rho_0 = v_0 = \psi_0 = M_{\mathrm{sym}} = 1$, $m_{\mathrm{link}} = 1/2$, $\lambda_{\mathrm{portal}} = (27/32)^2$, and $\rho_m = 2$, the exact stationary point $(\rho_{\mathrm{in}}^*, 0) = (1/4, 0)$ satisfies Eq.~\eqref{eq:exact_coupled_interior_cubic} identically and passes the radial check with $H_{rr} = 107/512 > 0$, yet $H_{\rho\rho} = -9/16 < 0$ (radial-plane Hessian $\mathrm{diag}(-9/16,\; 107/512)$), making $(1/4, 0)$ a \emph{saddle point} rather than a local minimum.
    \item \emph{Degenerate zero-eigenvalue instability counterexample ($H_{\rho\rho} > 0, H_{rr} = 0$):} At the equality boundary $\rho_m = \rho_{m,\mathrm{crit}}$ ($H_{rr} = 0$) with $H_{\rho\rho} > 0$, minimizing over $\delta\rho_t = (m_{\mathrm{link}}\sqrt{\lambda_{\mathrm{portal}}}/H_{\rho\rho})|\psi|^2 + O(|\psi|^4)$ along the coupled valley shifts the effective $|\psi|^4$ coefficient to $\lambda_{\mathrm{eff}}^{(4)} = \frac{\lambda_\psi}{4} + \frac{\lambda_{\mathrm{portal}}}{2} - \frac{m_{\mathrm{link}}^2\lambda_{\mathrm{portal}}}{2 H_{\rho\rho}}$. At $\lambda_\rho = \rho_0 = v_0 = m_{\mathrm{link}} = \psi_0 = M_{\mathrm{sym}} = 1$, $\lambda_{\mathrm{portal}} = (7/8)^2$, $\lambda_\psi = 1/10$, $\rho_{\mathrm{in}}^* = 1/2$, and $\rho_m = \rho_{m,\mathrm{crit}} = 121/160$, the Hessian is positive semidefinite $\mathrm{diag}(3/4,\; 0)$, yet $\lambda_{\mathrm{eff}}^{(4)} = -197/1920 < 0$, so $(1/2, 0)$ is \emph{unstable} along the coupled quartic valley.
  \end{itemize}
  By contrast, at the unit-parameter $m_{\mathrm{link}} = 2$ benchmark above, the exact real root is $\rho_{\mathrm{in}}^* \approx 0.59607$, the radial threshold is $\rho_{m,\mathrm{crit}} \approx 1.38429$ (exceeding the naive decoupled value $1.0$), and the temporal curvature is strictly positive, $H_{\rho\rho}(\rho_{\mathrm{in}}^*) \approx 4.06590 > 0$, so $(\rho_{\mathrm{in}}^*, 0)$ is indeed a strict local minimum for every $\rho_m > 1.38429$. Only in the controlled \emph{stiff-portal asymptotic regime} $m_{\mathrm{link}}^2 \gg \lambda_\rho\rho_0^2$ (which also ensures $H_{\rho\rho} > 0$ automatically) with $\sqrt{\lambda_{\mathrm{portal}}}\,v_0^2 \ll m_{\mathrm{link}}\rho_0$ does Eq.~\eqref{eq:exact_coupled_interior_cubic} reduce to $\rho_{\mathrm{in}}^* = \rho_0 - (\sqrt{\lambda_{\mathrm{portal}}}/m_{\mathrm{link}})v_0^2 + O(\lambda_\rho\rho_0^2/m_{\mathrm{link}}^2)$ and Eq.~\eqref{eq:exact_coupled_screening_threshold} approach $\lambda_\psi v_0^2\psi_0^2 M_{\mathrm{sym}}^2$:
  \begin{itemize}
    \item \emph{High-density screened interior ($\rho_m > \rho_{m,\mathrm{crit}}$ with $H_{\rho\rho}(\rho_{\mathrm{in}}^*) > 0$):} The symmetric stationary point $(\rho_{\mathrm{in}}^*, 0)$ solving Eq.~\eqref{eq:exact_coupled_interior_cubic} is a strict local minimum, pinning $|\psi| = 0$ so the linear scalar-matter coupling $f'_{\mathrm{SM}}(0) = 0$ vanishes in the bulk interior.
    \item \emph{Low-density unscreened vacuum ($\rho_m \ll \rho_{m,\mathrm{crit}}$):} The coupled ground state relaxes toward $(\rho_0, v_0)$ ($v_0 > 0$), permitting state-dependent laboratory responses in ultra-high-vacuum chambers.
  \end{itemize}
\end{enumerate}

\paragraph{Falsifiability of the Screening Window.}
Viability requires a single non-empty parameter domain that simultaneously satisfies Cassini and torsion-balance bounds in high-density/macroscopic bodies (solving the full coupled two-field equations) while preserving calculable laboratory sensitivity in unscreened vacuum; if no parameter region satisfies both constraints concurrently using total conserved $T_{\mu\nu}^{\mathrm{tot}}$, that screening realization is excluded.

\subsection{Linearization and Scope of Decoupled Single-Field Approximations}

Linearizing the full two-scalar system $(\delta\rho_t, \delta|\psi|)$ around a coupled homogeneous equilibrium $(\rho_{\mathrm{bg}}, v_{\mathrm{bg}})$ yields the coupled Helmholtz/Yukawa system $(K\nabla^2 - H_{\mathrm{eff}}(\rho_{\mathrm{bg}}))\,\delta\Phi = J_{\mathrm{src}}$; its radial source entry vanishes identically when $v_{\mathrm{bg}} \to 0$. Passing to canonically normalized modes $\hat{\Phi} = K^{1/2}\delta\Phi$ and diagonalizing the positive-definite physical mass matrix $M_{\mathrm{phys}}^2 = K^{-1/2} H K^{-1/2}$ ($H u_a = m_a^2 K u_a$, Eq.~\eqref{eq:radial_Hessian_sec5}) with mixing angle $\vartheta$ and physical eigenmasses $(m_1, m_2) > 0$, a static point mass $M$ generates the normalized two-Yukawa screening profile
\begin{equation}
S(r) := \sin^2\vartheta\,e^{-m_1 r} + \cos^2\vartheta\,e^{-m_2 r} \in (0, 1],
\qquad
S(0) = 1,
\qquad
\lim_{r\to\infty} S(r) = 0.
\label{eq:two_yukawa_S_profile}
\end{equation}
Reducing this two-field system to a single-field radial Helmholtz approximation $(\nabla^2 - m_{\mathrm{eff}}^2(\rho_{\mathrm{bg}}))\,\delta|\psi| = \frac{v_{\mathrm{bg}}}{\kappa_c\psi_0^2 M_{\mathrm{sym}}^2}\,\delta\rho_m$ requires \emph{both} heavy temporal decoupling ($m_{\mathrm{link}}^2 \gg \lambda_\rho\rho_0^2$, $\lambda_{\mathrm{portal}} \to 0$) \emph{and} negligible induced kinetic mixing ($4\kappa_t\lambda_{\mathrm{portal}} v_0^2 / (\kappa_c m_{\mathrm{link}}^2) \ll 1$, since eliminating $\rho_t - \rho_0 \approx (\sqrt{\lambda_{\mathrm{portal}}}/m_{\mathrm{link}})(|\psi|^2 - v_0^2)$ in $-\frac{\kappa_t}{2}(\partial\rho_t)^2$ adds $-\frac{2\kappa_t\lambda_{\mathrm{portal}}|\psi|^2}{m_{\mathrm{link}}^2}(\partial|\psi|)^2$).

\subsection{Thin-Shell Suppression Estimate (Imported Single-Field Conformal Benchmark)}

As an \emph{imported heuristic benchmark} from single-field conformal symmetron literature (valid only in the decoupled single-field conformal regime, not yet derived for the full coupled multiplicative two-field action \texttt{M\_AUTH\_SIT\_EFT\_v1}), an extended spherical body of radius $R_{\mathrm{body}}$, mass $M_{\mathrm{body}}$, and surface Newtonian potential $\Phi_N = G M_{\mathrm{body}} / R_{\mathrm{body}}$ whose interior pins $v_{\mathrm{in}} \approx 0$ exhibits thin-shell scaling
\begin{equation}
\epsilon_{\mathrm{th}} \simeq \frac{|v_{\mathrm{vac}} - v_{\mathrm{in}}|}{6\,\alpha_{\mathrm{eff,body}}\,M_{\mathrm{Pl}}\,\Phi_N},
\label{eq:thin_shell_factor}
\end{equation}
suggesting exterior scalar-charge suppression ($\epsilon_{\mathrm{th}} \ll 1$) when $\Phi_N$ is large (Sun, Earth) versus $\epsilon_{\mathrm{th}} \gtrsim 1$ in microscopic laboratory traps ($\Phi_N \ll 10^{-20}$). Obtaining rigorous exterior charges for unchanged \texttt{M\_AUTH\_SIT\_EFT\_v1} requires solving the coupled two-field nonlinear spherical boundary-value problem for $(\rho_t(r), |\psi(r)|)$ directly (\texttt{OPEN-6A}/\texttt{OPEN-6D}).

\subsection{Planar Large-Radius Decoupled Exterior Profile Approximation}

In the decoupled single-field ($\lambda_{\mathrm{portal}} \to 0$), planar/large-radius ($m_\psi R_{\mathrm{body}} \gg 1$, neglecting the spherical curvature term $-(2/r)\partial_r|\psi|$) approximation outside a screened boundary ($v_{\mathrm{in}} = 0$ at $r = R_{\mathrm{body}}$), the coherence amplitude recovers toward its exterior vacuum value via the 1D $\phi^4$ kink
\begin{equation}
\Rcoh(r) \simeq R_0 \tanh\!\left[\sqrt{\frac{\lambda_\psi v_0^2}{2\kappa_c}}\,(r - R_{\mathrm{body}})\right],
\qquad r \gtrsim R_{\mathrm{body}},
\label{eq:radial_tanh_profile}
\end{equation}
which receives both spherical $O(1/(m_\psi R_{\mathrm{body}}))$ and coupled two-field $O(\lambda_{\mathrm{portal}}\lambda_\rho\rho_0^2 / (\lambda_\psi m_{\mathrm{link}}^2))$ corrections whenever $\rho_t(r)$ varies simultaneously from $\rho_{\mathrm{in}}^*$ to $\rho_0$.

\subsection{Exact Shapiro Null-Cone Invariance and Consistently Normalized Conditional Conformal-Frame PPN Bounds}
\label{sec:ppn-shapiro-exact}

We now state the Parameterized Post-Newtonian (PPN) potential identities for a \emph{consistently specified conditional conformal extension} together with the mandatory Stage-6 action-equivalence qualification (\texttt{Stage6FunctionalResponse.\allowbreak{}ActionIdentityDiagnostics}, Section~\ref{sec:stage6-action-diagnostics}). Writing the static spherically symmetric Einstein-frame metric in isotropic coordinates as
\[
ds_E^2 = -\bigl(1 + 2\Phi_E(r)\bigr)dt^2 + \bigl(1 - 2\Psi_E(r)\bigr)d\mathbf{x}^2,
\qquad
\Phi_E(r) = \Psi_E(r) = -\frac{G M_E}{r},
\]
suppose one adopts a conditional universal conformal Jordan metric
\begin{equation}
g_{\mu\nu}^{(J)} := A^2(\Rcoh)\,g_{\mu\nu}^{(E)},
\qquad
S_{\mathrm{pp}} = -m\int d\tau_J = -m\int A(\Rcoh)\,d\tau_E,
\label{eq:conformal_metric_general_A}
\end{equation}
with smooth positive conformal factor $A(\Rcoh) > 0$. Deriving \emph{both} the static point-particle scalar source ($\nabla^2\delta|\psi| - m_{\mathrm{eff}}^2\delta|\psi| = \frac{M_E}{\kappa_c}\frac{A'(v_0)}{A(v_0)}\delta^{(3)}(\mathbf{r})$, where $M_E = m A(v_0)$) \emph{and} the weak-field Jordan metric perturbation $g_{\mu\nu}^{(J)}/A_0^2 = [1 + 2(A'(v_0)/A(v_0))\delta|\psi|]g_{\mu\nu}^{(E)}$ consistently from the same $A(\Rcoh)$ yields the modified Jordan-frame potentials
\begin{equation}
\Phi_J(r) = \Phi_E(r)\bigl[1 + \alpha_{\mathrm{conf}}[A]\,S(r)\bigr],
\qquad
\Psi_J(r) = \Phi_E(r)\bigl[1 - \alpha_{\mathrm{conf}}[A]\,S(r)\bigr],
\label{eq:jordan_potentials_exact}
\end{equation}
where $S(r) \in [0,1]$ is the normalized Yukawa profile~\eqref{eq:two_yukawa_S_profile} and the dimensionless conformal scalar-strength coefficient is uniquely determined by $A$:
\begin{equation}
\alpha_{\mathrm{conf}}[A]
:= \frac{2 M_{\mathrm{Pl}}^2}{\kappa_c}\left(\frac{A'(v_0)}{A(v_0)}\right)^2 \ge 0.
\label{eq:alpha_conf_general}
\end{equation}
\textbf{Resolution of Earlier Conformal Normalization Mismatch (Factor of $1/4$):} Earlier drafts simultaneously wrote $g_{\mu\nu}^{(J)} = f_{\mathrm{SM}}(\Rcoh)^{-1} g_{\mu\nu}^{(E)}$ (which implies $A = f_{\mathrm{SM}}^{-1/2}$) while printing $\alpha_{\mathrm{sym}} := \frac{2 M_{\mathrm{Pl}}^2}{\kappa_c}\bigl(\frac{f'_{\mathrm{SM}}(v_0)}{f_{\mathrm{SM}}(R_0)}\bigr)^2$ (which corresponds to $A = f_{\mathrm{SM}}^{\pm 1}$). Version~81 resolves that mismatch explicitly by separating the two distinct conditional conformal models:
\begin{itemize}
  \item \textbf{Conformal Model C1 ($g_{\mu\nu}^{(J)} = f_{\mathrm{SM}}(\Rcoh)^{-1} g_{\mu\nu}^{(E)}$, i.e.\ $A = f_{\mathrm{SM}}^{-1/2}$):} Because $A'/A = -\frac{1}{2}f'_{\mathrm{SM}}/f_{\mathrm{SM}}$, consistent source and test-body derivation gives
  \begin{equation}
  \alpha_{\mathrm{conf}}^{(\mathrm{C1})} = \frac{M_{\mathrm{Pl}}^2}{2\kappa_c}\left(\frac{f'_{\mathrm{SM}}(v_0)}{f_{\mathrm{SM}}(R_0)}\right)^2 = \frac{1}{4}\,\alpha_{\mathrm{sym}}.
  \label{eq:alpha_conf_C1}
  \end{equation}
  \item \textbf{Conformal Model C2 ($g_{\mu\nu}^{(J)} = f_{\mathrm{SM}}(\Rcoh)^{\pm 2} g_{\mu\nu}^{(E)}$, i.e.\ $A = f_{\mathrm{SM}}^{\pm 1}$):} Because $A'/A = \pm f'_{\mathrm{SM}}/f_{\mathrm{SM}}$, consistent derivation gives
  \begin{equation}
  \alpha_{\mathrm{conf}}^{(\mathrm{C2})} = \frac{2 M_{\mathrm{Pl}}^2}{\kappa_c}\left(\frac{f'_{\mathrm{SM}}(v_0)}{f_{\mathrm{SM}}(R_0)}\right)^2 = \alpha_{\mathrm{sym}},
  \label{eq:alpha_sym_exact}
  \end{equation}
  which algebraically matches the unchanged-action dust ($q_{\mathrm{eff}} = 1$) two-mode exchange coefficient $\alpha_{\mathrm{sym}}$ of Section~\ref{sec:stage6-action-diagnostics} (where the worldline action couples linearly to $f_{\mathrm{SM}}$ rather than $f_{\mathrm{SM}}^{1/2}$).
\end{itemize}
\textbf{Stage-6 Action-Equivalence Qualification:} Stage~6 proves in Lean~4.30 (\texttt{multiplicative\_\allowbreak{}vs\_\allowbreak{}conformal\_\allowbreak{}scalar\_\allowbreak{}coeff\_\allowbreak{}forces\_\allowbreak{}unit}) that the unchanged multiplicative portal $+\sqrt{-g}\,f_{\mathrm{SM}}(\Rcoh)\mathcal{L}_{\mathrm{matter}}$ in Eq.~\eqref{eq:abstract_master_action} is \emph{not} algebraically identical to \emph{either} conformal extension C1 or C2 unless $f_{\mathrm{SM}} = 1$. Accordingly, the two conformal-frame PPN results below hold strictly \emph{conditional} on explicitly adopting one chosen conformal extension (C1 or C2)---rather than as unconditional predictions of unchanged \texttt{M\_AUTH\_SIT\_EFT\_v1}:
\begin{enumerate}[leftmargin=2em]
  \item \textbf{Exact Shapiro Null-Cone Invariance (Conformal Frame):} Because the Maxwell kinetic term $-\frac{1}{4}f_{\mathrm{EM}}(\Rcoh)\sqrt{-g}g^{\mu\alpha}g^{\nu\beta}F_{\mu\nu}F_{\alpha\beta}$ in $d=4$ spacetime dimensions is Weyl-invariant under $g_{\mu\nu}^{(J)} = A^2 g_{\mu\nu}^{(E)}$ and $f_{\mathrm{EM}}(\Rcoh)$ is an isotropic scalar function, geometric-optics light rays follow the null cones $ds^2 = 0$ identically, and under Eq.~\eqref{eq:jordan_potentials_exact} the sum of the Jordan-frame potentials satisfies the exact cancellation
  \begin{equation}
  \Phi_J(r) + \Psi_J(r) = 2\Phi_E(r),
  \label{eq:shapiro_null_cone_invariance}
  \end{equation}
  eliminating anomalous tree-level birefringence and locking gravitational light bending to $2\Phi_E(r)$.
  \item \textbf{Conditional Eddington--Cassini PPN $\gamma$ Bound:} Formalized algebraically in Lean~4.30 (\texttt{SitQg.Stage2.QG15To19.ppnGammaBound\_of\_alpha\_and\_screening}) for any abstract conformal strength $\alpha_{\mathrm{conf}} \ge 0$ in Eq.~\eqref{eq:jordan_potentials_exact}, the ratio $\gamma_{\mathrm{PPN}}(r) := \Psi_J(r)/\Phi_J(r)$ satisfies
  \begin{equation}
  \gamma_{\mathrm{PPN}}(r) = \frac{1 - \alpha_{\mathrm{conf}} S(r)}{1 + \alpha_{\mathrm{conf}} S(r)},
  \qquad
  \bigl|\gamma_{\mathrm{PPN}}(r) - 1\bigr|
  = \frac{2\alpha_{\mathrm{conf}} S(r)}{1 + \alpha_{\mathrm{conf}} S(r)}
  \le \frac{2\alpha_{\mathrm{conf}}}{1 + \alpha_{\mathrm{conf}}} < 2.
  \label{eq:gamma_ppn_exact_bound}
  \end{equation}
  Hence $|\gamma_{\mathrm{PPN}} - 1| < 2.3 \times 10^{-5}$ holds even unscreened ($S(r) = 1$) whenever $\alpha_{\mathrm{conf}} < 1.15 \times 10^{-5}$. Evaluating the \texttt{BM\_PHYS\_EFT\_1} parameter tuple strictly as conditional conformal-extension illustrations gives $\alpha_{\mathrm{conf}}^{(\mathrm{C2})} = \alpha_{\mathrm{sym}} = 4.7446\times 10^{-7}$ ($|\gamma_{\mathrm{PPN}}-1|_{S=1} = 9.4893\times 10^{-7}$) under Model~C2 ($A = f_{\mathrm{SM}}^{\pm 1}$), or $\alpha_{\mathrm{conf}}^{(\mathrm{C1})} = \frac{1}{4}\alpha_{\mathrm{sym}} = 1.1862\times 10^{-7}$ ($|\gamma_{\mathrm{PPN}}-1|_{S=1} = 2.3723\times 10^{-7}$) under Model~C1 ($A = f_{\mathrm{SM}}^{-1/2}$), while in the unbroken symmetric limit $v_0 = 0$ (\texttt{BM\_V0\_SYMMETRIC\_GR}), $\alpha_{\mathrm{conf}} = 0$, $f_{\mathrm{EM}0} = f_{\mathrm{SM}0} = 1$, and $\gamma_{\mathrm{PPN}} = \beta_{\mathrm{PPN}} = 1$ identically. Neither conformal illustration replaces deriving the gravitational response directly from unchanged \texttt{M\_AUTH\_SIT\_EFT\_v1} (\texttt{OPEN-6D}; see Section~\ref{sec:stage6-action-diagnostics} for the direct unchanged-action E\"ot-Wash channel).
\end{enumerate}

\subsection{Consistency of Parameter Scales}

Across both high-scale EFT benchmarks (\texttt{BM\_PHYS\_EFT\_1}: $\psi_0 = \rho_0 = m_{\mathrm{link}} = 2\times 10^{16}\,\mathrm{GeV}$, $R_0 = 0.01$, $M_{\mathrm{sym}} = 50$, where the heavy scalar eigenmasses $m_1 = 2.4495\times 10^{14}\,\mathrm{GeV}$ [range $\lambda_1 = 8.056\times 10^{-31}\,\mathrm{m}$] and $m_2 = 2.8285\times 10^{16}\,\mathrm{GeV}$ exponentially suppress macroscopic fifth forces under the declared dust closure of Section~\ref{sec:stage6-action-diagnostics}) and light-scalar thin-shell regimes ($\mu \sim 10^{-3}\,\mathrm{eV}$, $M \sim 10^5\text{--}10^6\,\mathrm{GeV}$, $\epsilon_{\mathrm{th}} \ll 1$ for Solar System bodies with $\Phi_N \sim 10^{-9}\text{--}10^{-6}$ while $\epsilon_{\mathrm{th}} \gtrsim 1$ for laboratory traps with $\Phi_N \ll 10^{-20}$), \texttt{M\_AUTH\_SIT\_EFT\_v1} admits broad, radiatively stable EFT parameter domains whose direct laboratory falsifiability awaits completion of the microphysical action-to-apparatus bridges (\texttt{OPEN-6}).

\section{What the Theory Predicts: Map to Observables}
\label{sec:map-to-observables}

Super Information Theory organizes its prospective laboratory and astrophysical observable channels into five interfaces, with strict Version~81 separation between the one conditional weak-field channel derived directly from unchanged \texttt{M\_AUTH\_SIT\_EFT\_v1} and future channels requiring completion of \texttt{OPEN-6}:
\begin{itemize}
  \item \textbf{Sub-Millimeter Torsion-Balance Fifth-Force Tests (Stage-6 Conditional Non-Exclusion Channel):} Under a declared pressureless-dust worldline closure ($q_{\mathrm{eff}} = 1$, contrasted with $q_{\mathrm{eff}} = 0$ for homogeneous free Klein--Gordon oscillation), unchanged \texttt{M\_AUTH\_SIT\_EFT\_v1} maps directly onto two-mode Yukawa exchange with range $\lambda_1 = 8.056\times 10^{-31}\,\mathrm{m}$ at \texttt{BM\_PHYS\_EFT\_1} (exponentially indistinguishable from Newtonian gravity at the $52\,\mu\mathrm{m}$ E\"ot-Wash separation, Section~\ref{sec:stage6-action-diagnostics}); reproducing Lee~\emph{et~al.}'s patterned rotating-mass torque likelihood from Lee's dissertation Appendix~A with microscopic test-material scalar charges remains \texttt{OPEN-6A}/\texttt{OPEN-6D}.
  \item \textbf{Optical-Clock Redshift Metrology (Conditional on Visible $U(1)$ Dictionary, \texttt{OPEN-6B}):} Separating the local time-density response $\alpha_\rho$ from the gravitational-potential transfer factor $\mathcal{T} := (\rho_0^{-1}\partial\rho_t/\partial(\Phi/c^2))_{\mathrm{lab}}$ and potential redshift slope $\alpha_\Phi := d\ln\nu/d(\Phi/c^2) = \alpha_\rho\mathcal{T}$ (repairing the Version~79/80.2 symbol collision that reused $\alpha_{\mathrm{eff}}$ for both), a potential-driven clock shift obeys without double-counting $\mathcal{T}$:
  \begin{equation}
  \frac{\Delta\nu}{\nu} = \alpha_\rho\,\frac{\delta\rho_t}{\rho_0} = \alpha_\rho\,\mathcal{T}\,\frac{\Delta\Phi}{c^2} = \alpha_\Phi\,\frac{\Delta\Phi}{c^2}.
  \label{eq:clock_shift_map_sec7}
  \end{equation}
  The published Tokyo Skytree optical-lattice clock test (Takamoto~\emph{et~al.}, 2020) reported $\alpha_{\mathrm{RS}} = (+1.4 \pm 9.1)\times 10^{-5}$, while NIST (2026) achieved $3.2\times 10^{-18}$ optical-clock ratio uncertainty; the legacy $3\times 10^{-8}$ figure cited in earlier drafts is formally quarantined in Version~81 as a \emph{hypothetical future sensitivity benchmark}, not an empirical $95\%$~CL bound (Section~\ref{sec:atomic-clocks} and Appendix~\ref{appendix:ClockConstraints}).
  \item \textbf{Cold-Atom Interferometry (Conditional on $J_{\mathrm{prep}}$, \texttt{OPEN-6C}):} Controlled spatial gradients of $\Rcoh$ and $\rho_t$ motivate exploratory phase-shift and deflection sensitivity studies (Sections~\ref{sec:applications-predictions} and~\ref{sec:informational-torque}), contingent on deriving the microscopic state-preparation source $J_{\mathrm{prep}}$.
  \item \textbf{Bose--Einstein Condensate vs.\ Thermal Cloud Equivalence Test (Hypothetical Sensitivity Benchmark, \texttt{OPEN-6A}/\texttt{OPEN-6C}):} Comparing a pure $^{87}\mathrm{Rb}$ BEC against an equal-mass thermal cloud tests coherence-dependent gravitational coupling (subject to Paper~003/014's condensate-source and matched-preparation systematic controls); the legacy $|\delta a| \le 29.4\,\mu\mathrm{Gal}$ figure derived from $|\alpha_\Phi| \le 3\times 10^{-8}$ at $\chi = 1$ is explicitly relabeled in Version~81 as a \emph{hypothetical sensitivity target} rather than an experimentally anchored ceiling (Section~\ref{sec:smoking-gun-bec} and Appendix~\ref{app:BEC-benchmark}).
  \item \textbf{Astrophysical Lensing Tomography (Level-1 Exploratory Regime):} Slow cluster-scale coherence gradients provide exploratory templates for Euclid, JWST, and Rubin Observatory weak/strong lensing surveys (Sections~\ref{sec:part2-astro-probes} and~\ref{sec:black-holes-voids}).
\end{itemize}

\subsection{Parameter Summary, Vacuum Recovery Endpoints, and Observational Scope}

The canonical parameter tuple $\{\kappa_G, \kappa_t, \kappa_c, \psi_0, \rho_0, v_0, m_{\mathrm{link}}, \lambda_\rho, \lambda_\psi, \lambda_{\mathrm{portal}}, c_{\mathrm{EM}}, M_{\mathrm{sym}}, \hat{q}\}$ satisfies strict kinetic positivity, bounded-below scalar potential ($V_{\mathrm{tot}} \ge 0$), and principal-symbol hyperbolic wave propagation ($c_{\mathrm{char}} = 1$). Regarding GR/Maxwell recovery and Solar System safety, Version~81 enforces three explicit distinctions (correcting earlier wording that claimed exact GR/Maxwell recovery as $(\rho_t, |\psi|) \to (\rho_0, 0)$ at $v_0 > 0$):
\begin{enumerate}[leftmargin=2em]
  \item \textbf{Why $(\rho_0, 0)$ Is Not a Broken-Phase Vacuum ($v_0 > 0$):} Evaluating Eq.~\eqref{eq:U_link_canonical} at $(\rho_t, |\psi|) = (\rho_0, 0)$ when $v_0 > 0$ gives $V_{\mathrm{tot}}(\rho_0, 0) = (\frac{\lambda_\psi}{4} + \frac{\lambda_{\mathrm{portal}}}{2})v_0^4 > 0$ and $\partial_{\rho_t} V_{\mathrm{tot}}(\rho_0, 0) = m_{\mathrm{link}}\sqrt{\lambda_{\mathrm{portal}}}\,v_0^2 \ne 0$; thus $(\rho_0, 0)$ is neither stationary nor zero-energy in the broken phase.
  \item \textbf{Algebraic Metric Reduction vs.\ Full Coupled Solution Conditions:} At the broken-phase vacuum $(\rho_t, |\psi|) = (\rho_0, v_0)$ with $D_\mu\psi = 0$ ($V_{\mathrm{tot}} = 0$, $\nabla V_{\mathrm{tot}} = 0$), the metric equation reduces \emph{algebraically} to GR with constant $G_{\mathrm{eff}} = G f_{\mathrm{SM}}(R_0)$ and $f_{\mathrm{EM}}(R_0)$, forming an exact solution of the full coupled system whenever the scalar and gauge equations are simultaneously satisfied (such as source-free vacuum $\mathcal{L}_{\mathrm{matter}} = 0, F_{\mu\nu} = 0$; whereas at $v_0 > 0$, neutral dust $\mathcal{L}_{\mathrm{matter}} = -\rho_m < 0$ produces a nonzero scalar residual $-v_0\rho_m/(\psi_0^2 M_{\mathrm{sym}}^2) \ne 0$ and $[D_\mu, D_\nu]\psi = -i\hat{q}F_{\mu\nu}\psi = 0$ forces $F_{\mu\nu} = 0$ when $\hat{q}v_0 \ne 0$, Appendix~D.2). Broader exact sourced GR+Maxwell ($f_{\mathrm{EM}} = f_{\mathrm{SM}} = 1$, $m_A = 0$) holds in the separate unbroken parameter limit $v_0 = 0$ (\texttt{BM\_V0\_SYMMETRIC\_GR}), while heavy-scalar decoupling ($m_1, m_2 \gg E_{\mathrm{obs}}$) provides a controlled asymptotic approximation.
  \item \textbf{Observational Safety Requires a Derived Response:} Positivity of the parameter tuple alone does \emph{not} establish Solar System PPN safety; empirical viability requires deriving the gravitational and electromagnetic response from the stated action and testing it against observational data (\texttt{OPEN-6D}, Section~\ref{sec:stage6-action-diagnostics}).
\end{enumerate}

\section{Action Principle and Field Equations}
\label{sec:action-principle}
\label{sec:core-action}

\subsection{Symmetries and Minimal Canonical Model Selection (Not Uniqueness)}

The canonical action $S_{\mathrm{AUTH}}$ is a \emph{specific minimal ghost-free ansatz}---not the unique action permitted by symmetry---chosen within the class of dimension-$\le 4$ operators (plus leading $\mathbb{Z}_2$-even gauge-invariant dimension-6 dressings $f_{\mathrm{EM}}, f_{\mathrm{SM}}$) satisfying four baseline requirements: (1) four-dimensional diffeomorphism invariance; (2) local $U(1)$ gauge invariance $\psi \mapsto e^{i\hat{q}\chi}\psi$, $A_\mu \mapsto A_\mu + \partial_\mu\chi$; (3) discrete $C, P, T$ invariance and absence of bare phase-dependent potentials; and (4) strict second-order hyperbolicity excluding higher-derivative Ostrogradsky ghosts such as $(\Box\rho_t)^2$ or $(\nabla_\mu\nabla_\nu\Rcoh)^2$ (\texttt{C07} purged under \texttt{ODQ-2}). Because $U_{\mathrm{link}}$ explicitly breaks $\rho_t \mapsto -\rho_t$ symmetry via $m_{\mathrm{link}}(\rho_t - \rho_0)$, symmetry and second-order hyperbolicity alone still permit additional dimension-$\le 4$ operators---including non-minimal curvature couplings ($\xi_1 R\rho_t$, $\xi_2 R\rho_t^2$, $\xi_3 R|\psi|^2$), odd polynomial terms ($\rho_t$, $\rho_t^3$, $|\psi|^2\rho_t$), and independent quartic cross-terms $|\psi|^2\rho_t^2$---which are set to zero in \texttt{M\_AUTH\_SIT\_EFT\_v1} by explicit canonical model selection rather than claimed symmetry uniqueness.

\subsection{Symmetries and Gauge Structure}

Over $(\mathcal{M}^4, g_{\mu\nu})$ in signature $(-,+,+,+)$, every scalar potential and constitutive coupling depends on $\psi$ exclusively through $|\psi|^2 = \psi^*\psi = \psi_0^2\Rcoh^2$, while spacetime derivatives of $\psi$ enter solely via $D_\mu\psi = (\nabla_\mu - i\hat{q}A_\mu)\psi$. Time-density $\rho_t$ transforms as a neutral real Lorentz scalar.

\subsection{Unified Canonical SIT Action and Lagrangian}

Collecting Eqs.~\eqref{eq:L_EH_def}--\eqref{eq:U_link_canonical}, the unified action in explicit component form is
\begin{align}
S_{\mathrm{AUTH}}
&= \int d^4x\,\sqrt{-g}\Biggl[
  \frac{R}{16\pi G}
  - \frac{\kappa_t}{2}\,g^{\mu\nu}\nabla_\mu\rho_t\nabla_\nu\rho_t
  - \frac{\kappa_c}{2}\,g^{\mu\nu}(D_\mu\psi)^*(D_\nu\psi)
  \nonumber\\
&\qquad\quad
  - V(\rho_t) - U(|\psi|) - U_{\mathrm{link}}(\rho_t,\psi)
  - \frac{1}{4}\,f_{\mathrm{EM}}(\Rcoh)\,F_{\mu\nu}F^{\mu\nu}
  + f_{\mathrm{SM}}(\Rcoh)\,\mathcal{L}_{\mathrm{matter}}
\Biggr].
\label{eq:SIT_action}
\end{align}

\subsection{Exclusion of Pathological and Superseded Operators}

To maintain complete mathematical consistency and prevent any ambiguity with pre-Version-80 drafts, four operator classes are strictly excluded from $S_{\mathrm{AUTH}}$:
\begin{itemize}
  \item \textbf{Wrong-Sign Ghost Kinetic Terms (\texttt{C04} Purged):} In mostly-plus signature $(-,+,+,+)$, a scalar kinetic term $+\frac{1}{2}g^{\mu\nu}\partial_\mu\rho_t\partial_\nu\rho_t = -\frac{1}{2}(\partial_t\rho_t)^2 + \frac{1}{2}(\nabla\rho_t)^2$ carries negative kinetic energy and produces catastrophic vacuum instability. All legacy $+\frac{1}{2}g^{\mu\nu}\partial_\mu\rho_t\partial_\nu\rho_t$ sign conventions from early drafts are permanently replaced by $-\frac{\kappa_t}{2}g^{\mu\nu}\partial_\mu\rho_t\partial_\nu\rho_t$ with $\kappa_t > 0$.
  \item \textbf{Higher-Derivative Ostrogradsky Operators (\texttt{C07} Purged):} Operators such as $(\Box_g\rho_t)^2$ or $(\nabla_\mu\nabla_\nu\Rcoh)^2$ introduce fourth-order equations of motion with unbounded-below Ostrogradsky Hamiltonians and are excluded from the tree-level action (appearing only as Planck/cutoff-suppressed perturbative EFT counterterms treated via order reduction below $\Lambda_{\mathrm{UV}}$).
  \item \textbf{Non-Gauge-Invariant Phase Potentials:} Any bare dependence on $\theta = \arg\psi$ outside $D_\mu\psi$ violates local $U(1)$ gauge invariance and charge conservation.
  \item \textbf{Zero-Floor Gauge Dressings in Continuum EFT (\texttt{C09} Quarantined):} Dressings with $f_{\mathrm{EM}}(\rho_t, 0) = 0$ make the continuum canonical Hamiltonian $\frac{1}{2f_{\mathrm{EM}}}\Pi^2$ singular at $|\psi| = 0$; continuum \texttt{M\_AUTH\_SIT\_EFT\_v1} enforces $f_{\mathrm{EM}}(\Rcoh) = 1 + c_{\mathrm{EM}}\Rcoh^2 \ge 1$.
\end{itemize}

\subsection{Motivation for the Canonical Action's Form}

Under diffeomorphism invariance, $U(1)$ gauge symmetry, $\mathbb{Z}_2$ vacuum selection, and second-order hyperbolicity, Eq.~\eqref{eq:SIT_action} is the minimal ghost-free covariant effective action coupling spacetime curvature, temporal clock-rate density $\rho_t$, complex phase-coherence order $\psi$, Abelian gauge curvature $F_{\mu\nu}$, and Standard Model matter $\mathcal{L}_{\mathrm{matter}}$.

\subsection{Derivation of the Field Equations and Total Stress--Energy Tensor}

\subsubsection{Variational Principle}

Requiring $\delta S_{\mathrm{AUTH}} = 0$ under arbitrary compactly supported variations $(\delta g^{\mu\nu}, \delta\rho_t, \delta\psi^*, \delta A_\nu)$ on $(\mathcal{M}^4, g_{\mu\nu})$ yields $G_{\mu\nu} = 8\pi G\,T_{\mu\nu}^{\mathrm{tot}}$ with $T_{\mu\nu} = -2\frac{\delta\mathcal{L}_{\mathrm{non\text{-}EH}}}{\delta g^{\mu\nu}} + g_{\mu\nu}\mathcal{L}_{\mathrm{non\text{-}EH}}$.

\subsubsection{Contribution from the Time-Density and Coherence Scalars}

Using $\delta\sqrt{-g} = -\frac{1}{2}\sqrt{-g}\,g_{\mu\nu}\delta g^{\mu\nu}$ on $\mathcal{L}_{\mathrm{coh}}$ gives
\begin{align}
T_{\mu\nu}^{(\rho_t)} + T_{\mu\nu}^{(\psi)} + T_{\mu\nu}^{(\mathrm{pot})}
&= \kappa_t\,\partial_\mu\rho_t\partial_\nu\rho_t
  + \kappa_c\,\mathrm{Re}\bigl[(D_\mu\psi)^*(D_\nu\psi)\bigr]
  \nonumber\\
&\quad
  - g_{\mu\nu}\left[
    \frac{\kappa_t}{2}\,g^{\alpha\beta}\partial_\alpha\rho_t\partial_\beta\rho_t
    + \frac{\kappa_c}{2}\,g^{\alpha\beta}(D_\alpha\psi)^*(D_\beta\psi)
    + V_{\mathrm{tot}}(\rho_t,|\psi|)
  \right].
\label{eq:T_scalar_combined_sec8}
\end{align}

\subsubsection{Contribution from Coherence-Dressed Gauge and Matter Sectors}

Using $\frac{\delta(F_{\alpha\beta}F^{\alpha\beta})}{\delta g^{\mu\nu}} = 2 F_{\mu\lambda}F_\nu^{\ \lambda}$ on $\mathcal{L}_{\mathrm{EM}}^{\mathrm{coh}} = -\frac{1}{4}f_{\mathrm{EM}}(\Rcoh)F_{\alpha\beta}F^{\alpha\beta}$ (with canonical Maxwell factor $-1/4$, repairing historical $-1/2$ typos \texttt{T1}) and varying $+f_{\mathrm{SM}}(\Rcoh)\mathcal{L}_{\mathrm{matter}}$ yields
\begin{equation}
T_{\mu\nu}^{(\mathrm{EM})} + T_{\mu\nu}^{(\mathrm{matter})}
= f_{\mathrm{EM}}(\Rcoh)\left(F_{\mu\lambda}F_\nu^{\ \lambda} - \frac{1}{4}g_{\mu\nu}F_{\alpha\beta}F^{\alpha\beta}\right)
  + f_{\mathrm{SM}}(\Rcoh)\,T_{\mu\nu}^{\mathrm{SM}}.
\label{eq:T_dressed_combined_sec8}
\end{equation}

\subsubsection{Physical Interpretation of Coherence Effects and Low-Energy GR Recovery}

In regions of elevated macroscopic coherence ($\Rcoh > R_0$), $(\rho_t, |\psi|)$ gradients and dressings $(f_{\mathrm{EM}}, f_{\mathrm{SM}})$ contribute positive-energy gravitational sourcing; in decohered or screened regions where $(\rho_t, |\psi|) \to (\rho_0, 0)$, $V_{\mathrm{tot}}(\rho_0, 0) = \frac{\lambda_\psi}{4}v_0^4 + \frac{\lambda_{\mathrm{portal}}}{2}v_0^4$ acts as a non-stationary positive potential offset for $v_0 > 0$, while at the broken-phase vacuum $(\rho_0, v_0)$ we have $V_{\mathrm{tot}}(\rho_0, v_0) = 0$ identically (giving exact source-free vacuum GR, controlled $O(m_{1,2}^{-2})$ sourced GR recovery in the heavy-scalar limit, and exact sourced GR+Maxwell in the unbroken limit $v_0 = 0$).

\subsubsection{Summary of Canonical Field Equations}

For ease of reference across Parts~I--IV, the complete canonical Euler--Lagrange equations of \texttt{M\_AUTH\_SIT\_EFT\_v1} in $(-,+,+,+)$ (aligned directly with $\delta S_{\mathrm{AUTH}}/\delta\psi^* = 0$ in Eq.~\eqref{eq:appD_full_variation}, where the kinetic term $-\frac{\kappa_c}{2}|D\psi|^2$ produces $+\frac{\kappa_c}{2}D_\mu D^\mu\psi$ against undoubled $\partial/\partial\psi^*$ derivatives, multiplying by $2$ to give Eq.~\eqref{eq:Rcoh-eom}) are:
\begin{align}
G_{\mu\nu} &= 8\pi G\,\Bigl(T_{\mu\nu}^{(\rho_t)} + T_{\mu\nu}^{(\psi)} + T_{\mu\nu}^{(\mathrm{EM})} + T_{\mu\nu}^{(\mathrm{pot})} + f_{\mathrm{SM}}(\Rcoh)\,T_{\mu\nu}^{\mathrm{SM}}\Bigr), \label{eq:summary_einstein_sec8}\\
\kappa_t\,\Box_g\rho_t &= \frac{\partial V_{\mathrm{tot}}}{\partial\rho_t}, \label{eq:summary_rhot_sec8}\\
\frac{\kappa_c}{2}\,D_\mu D^\mu\psi &= \frac{\partial V_{\mathrm{tot}}}{\partial\psi^*}
  + \frac{1}{4}\frac{\partial f_{\mathrm{EM}}}{\partial\psi^*}\,F_{\alpha\beta}F^{\alpha\beta}
  - \frac{\partial f_{\mathrm{SM}}}{\partial\psi^*}\,\mathcal{L}_{\mathrm{matter}}, \label{eq:summary_psi_sec8}\\
\nabla_\mu\bigl(f_{\mathrm{EM}}(\Rcoh)\,F^{\nu\mu}\bigr) &= \hat{q}\,\kappa_c|\psi|^2\bigl(\partial^\nu\theta - \hat{q}A^\nu\bigr) + J^\nu_{\mathrm{EM,SM}}. \label{eq:summary_maxwell_sec8}
\end{align}
Notice the canonical minus sign in $\kappa_t\Box_g\rho_t - \partial_{\rho_t}V_{\mathrm{tot}} = 0$ (repairing historical sign flip \texttt{R1}/\texttt{R2}): around the minimum $(\rho_0, v_0)$ where $m_\rho^2 = \partial_{\rho_t}^2 V_{\mathrm{tot}} > 0$, flat-space perturbations satisfy $(-\partial_t^2 + \nabla^2)\delta\rho_t - (m_\rho^2/\kappa_t)\delta\rho_t = 0$, i.e.\ $\partial_t^2\delta\rho_t - \nabla^2\delta\rho_t + (m_\rho^2/\kappa_t)\delta\rho_t = 0$, guaranteeing stable real dispersion $\omega^2 = \mathbf{k}^2 + m_\rho^2/\kappa_t > 0$.
% --- END INPUT sections/02_part1_sec04_to_08.tex ---
% --- BEGIN INPUT sections/03_part1_sec09_to_15.tex ---
\section{Symmetries, Noether Charge Conservation, and Operational Coherence Boundaries}
\label{sec:coherence-conservation}
\label{sec:symmetry-noether-moved}
\label{app:NoetherCurrents}

This section derives the exact Noether currents associated with general covariance and the internal $U(1)$ phase symmetry of the canonical SIT action (\texttt{M\_AUTH\_SIT\_EFT\_v1}), and strictly separates proved on-shell Noether charge conservation from operational quantum basis coherence, density-matrix purity, entanglement, and neural synchronization (following the operational separation contract of Section~\ref{sec:PrimitiveFieldDefinitions} and Paper~002).

\subsection{General Covariance and Energy--Momentum Conservation}

Because $S_{\mathrm{AUTH}} = \int d^4x\,\sqrt{-g}\,\mathcal{L}_{\mathrm{AUTH}}$ is invariant under smooth diffeomorphisms $x^\mu \mapsto x^\mu + \xi^\mu(x)$, Noether's second theorem (equivalently the contracted Bianchi identity $\nabla_\mu G^{\mu\nu} \equiv 0$) implies exact on-shell covariant conservation of the total stress--energy tensor:
\[
\nabla_\mu T^{\mu\nu}_{\mathrm{tot}} = 0,
\qquad
T^{\mu\nu}_{\mathrm{tot}} = T^{\mu\nu}_{(\rho_t)} + T^{\mu\nu}_{(\psi)} + T^{\mu\nu}_{(\mathrm{EM})} + T^{\mu\nu}_{(\mathrm{pot})} + f_{\mathrm{SM}}(\Rcoh)\,T^{\mu\nu}_{\mathrm{SM}}.
\]

\subsection{U(1) Phase Symmetry and the Gauge-Invariant Coherence Current}

The canonical action~\eqref{eq:canonical_action_sec5} is invariant under local $U(1)$ gauge transformations
\begin{equation}
\psi(x) \to e^{i\hat{q}\chi(x)}\psi(x),
\qquad
A_\mu(x) \to A_\mu(x) + \partial_\mu\chi(x),
\label{eq:u1_gauge_transform_sec9}
\end{equation}
implemented through $D_\mu\psi = (\nabla_\mu - i\hat{q}A_\mu)\psi$ inside the $(-,+,+,+)$ kinetic density $\mathcal{L}_\psi = -\frac{\kappa_c}{2}g^{\mu\nu}(D_\mu\psi)^*(D_\nu\psi) - U(|\psi|)$. Invariance under the global subgroup ($\chi = \mathrm{const}$) yields the gauge-invariant conserved coherence current
\begin{equation}
J^\mu_{\mathrm{coh}}
= \frac{i\kappa_c}{2}\Bigl(\psi(D^\mu\psi)^* - \psi^*(D^\mu\psi)\Bigr),
\qquad
\nabla_\mu J^\mu_{\mathrm{coh}} = 0 \quad (\text{on-shell}).
\label{eq:Jcoh_covariant_sec9}
\end{equation}
Writing $\psi(x) = |\psi(x)| e^{i\theta(x)}$ casts Eq.~\eqref{eq:Jcoh_covariant_sec9} into its explicit London superflow form:
\begin{equation}
J^\mu_{\mathrm{coh}} = \kappa_c |\psi|^2 \bigl(\nabla^\mu\theta - \hat{q}A^\mu\bigr),
\label{eq:Jcoh_london_sec9}
\end{equation}
which reduces in the zero-field limit $A_\mu \to 0$ to $J^\mu_{\mathrm{coh}} = \kappa_c |\psi|^2 \nabla^\mu\theta$ and sources the coherence-gated Maxwell/Proca equation~\eqref{eq:maxwell_proca_eom_canonical}.

\paragraph{Exact Scope of Noether Charge Conservation vs.\ Operational Coherence Measures:}
The conservation law $\nabla_\mu J^\mu_{\mathrm{coh}} = 0$ is an exact on-shell theorem strictly for the global $U(1)$ Noether charge $Q_{\mathrm{coh}} = \int_\Sigma J^0_{\mathrm{coh}}\sqrt{q}\,d^3x$ of the complex EFT field $\psi$ in \texttt{M\_AUTH\_SIT\_EFT\_v1}. It does \emph{not} imply conservation or monotonicity of arbitrary operational coherence statistics:
\begin{itemize}
    \item \textbf{What Noether Conservation Establishes:} In a closed classical solution of \texttt{M\_AUTH\_SIT\_EFT\_v1} (with appropriate boundary falloff), total $U(1)$ Noether charge $Q_{\mathrm{coh}}$ is conserved, and local charge changes are balanced by spatial flux $\nabla\cdot\mathbf{J}_{\mathrm{coh}}$.
    \item \textbf{What Noether Conservation Does \emph{Not} Establish:} Conservation of $Q_{\mathrm{coh}}$ does \emph{not} prove conservation of fixed-basis quantum off-diagonal coherence, density-matrix purity, multipartite entanglement, or neural cross-frequency synchrony. Under unitary quantum evolution, fixed-basis coherence is \emph{not} conserved: for example, a single-qubit Hadamard unitary maps $|0\rangle \mapsto |+\rangle$, changing computational-basis off-diagonal coherence $2|\rho_{01}|$ from $0$ to $1$ while purity $\mathrm{Tr}(\rho^2) = 1$ stays constant. Similarly, under non-unital open-system CPTP maps (Section~\ref{sec:arrow-of-time}), even purity can strictly increase.
\end{itemize}

\subsection{Why Uncertainty Relations Do Not Establish a Universal Coherence Budget}
\label{subsec:uncertainty-vs-coherence}

Earlier drafts of SIT proposed a universal conservation-of-coherence rule across complementary observables $(A, B)$ of the schematic form $\mathcal{C}_A \mathcal{C}_B \gtrsim \chi > 0$, attempting to derive the Heisenberg--Robertson uncertainty principle ($\Delta x\,\Delta p \ge \hbar/2$) and neural cross-frequency trade-offs as instances of a single conserved coherence budget. Version~81 explicitly \emph{retracts} any claim of a universal conserved coherence budget or universal positive product lower bound across complementary bases:
\begin{enumerate}[leftmargin=2em]
  \item \textbf{Counterexample to Universal Lower Bound $\mathcal{C}_A \mathcal{C}_B \ge \chi > 0$:} For a single qubit in the maximally mixed state $\rho = I/2$, off-diagonal coherence vanishes identically in \emph{every} orthonormal basis ($\mathcal{C}_Z = \mathcal{C}_X = \mathcal{C}_Y = 0$), so $\mathcal{C}_A \mathcal{C}_B = 0$ for every complementary pair, refuting any universal positive lower bound $\chi > 0$ on basis coherence measures.
  \item \textbf{Counterexample to Unitary Basis-Coherence Conservation:} A unitary transformation (such as the Hadamard gate $|0\rangle \mapsto |+\rangle$) changes fixed-basis coherence freely while leaving purity invariant, showing that unitary quantum dynamics conserves inner products and von~Neumann entropy rather than a single fixed-basis coherence functional.
  \item \textbf{Distinction Between Variance Uncertainty and Coherence Conservation:} Standard quantum mechanics derives Robertson--Schr\"odinger variance bounds ($\Delta A\,\Delta B \ge \frac{1}{2}|\langle[A,B]\rangle|$) and entropic uncertainty relations from non-commutativity of observables on Hilbert space. Those established theorems constrain simultaneous sharpness in incompatible bases, but do not follow from $U(1)$ Noether charge conservation of a classical EFT field $\psi$, nor do they derive neural band-power trade-offs (which belong to separate empirical neuroscience modeling, Section~\ref{sec:neuroscience-implications}).
\end{enumerate}

The motivating conceptual research question---how basis-dependent coherence, entanglement, and measurement-induced disturbance trade off inside specified quantum protocols, and whether cross-frequency synchrony exhibits empirical trade-offs in neural recordings---remains part of the interdisciplinary research program, kept cleanly separated from the proved Noether current $\nabla_\mu J^\mu_{\mathrm{coh}} = 0$ of \texttt{M\_AUTH\_SIT\_EFT\_v1}.

\subsection{Time-Density Field and Internal Symmetries}

For $\rho_t(x)$, if $V(\rho_t)$ is invariant under shifts or certain scaling transformations, similar Noether currents arise.

\paragraph{Example: Shift Symmetry in $\rho_t$.}
If $V(\rho_t)$ is constant, consider infinitesimal shifts $\rho_t \to \rho_t + \beta$:
\[
J^\mu_{(\rho_t)} = \frac{\partial \mathcal{L}_{\rho_t}}{\partial (\partial_\mu \rho_t)} \delta \rho_t = \partial^\mu \rho_t\, \beta
\]
Conservation:
\[
\nabla_\mu J^\mu_{(\rho_t)} = \square \rho_t = 0
\]
if the equation of motion holds and $V'(\rho_t) = 0$.

\paragraph{Remarks:}
If $V'(\rho_t) \neq 0$, this current is explicitly broken as well.

\subsection{Electromagnetic Gauge Invariance}

If matter and electromagnetic fields are present, standard $U(1)$ gauge invariance applies:
\[
A_\mu \to A_\mu + \partial_\mu \chi, \qquad \psi \to e^{ie\chi/\hbar} \psi
\]
Noether's theorem yields the standard electromagnetic current:
\[
J^\mu_{\mathrm{EM}} = \frac{\partial \mathcal{L}_{\mathrm{matter}}}{\partial (\partial_\mu \psi)} (i e / \hbar) \psi + \text{c.c.}
\]
This is unaffected by $R_{\mathrm{coh}}$ and $\rho_t$ unless the coupling functions $f_1, f_2$ break $U(1)$ invariance, which is not permitted if standard EM is to be recovered.

\subsection{Summary Table of Symmetries and Currents}
\begin{center}
\small
\begin{tabular}{|p{0.16\textwidth}|p{0.17\textwidth}|p{0.33\textwidth}|p{0.22\textwidth}|}
\hline
\textbf{Field/Sector} & \textbf{Symmetry} & \textbf{Noether Current} & \textbf{Conservation Condition} \\
\hline
Metric, all fields & Diffeomorphism & $T^{\mu\nu}_{\text{total}}$ & Always (on-shell) \\
Complex Field $\psi$ & Global U(1) phase rotation & $J^\mu_{\mathrm{coh}} = \frac{i\kappa_c}{2}(\psi\nabla^\mu\psi^* - \psi^*\nabla^\mu\psi)$ & Holds if action depends only on $|\psi|$ \\
$\rho_t$ & Shift symmetry & $J^\mu_{(\rho_t)} \propto \nabla^\mu \rho_t$ & Broken if potential $V(\rho_t)$ is non-trivial \\
EM/matter & Local U(1) gauge & $J^\mu_{\mathrm{EM}}$ & Holds by construction \\
\hline
\end{tabular}
\end{center}

\noindent
In the presence of explicit symmetry-breaking potentials $U(R_{\mathrm{coh}})$, $V(\rho_t)$, or non-invariant coupling functions $f_1, f_2$, the corresponding currents are only approximately conserved. SIT thus recovers standard physical conservation laws in the appropriate limits, and introduces generalized coherence/entropy currents whose exact conservation is symmetry-dependent.



\section{Quantum Operator Formalism for the Informational Fields}
\label{sec:quantum-operator-formalism}
\label{sec:path-int}

Having established the classical action, symmetries, and post-Newtonian consistency of \texttt{M\_AUTH\_SIT\_EFT\_v1}, we formulate its canonical quantization as a perturbative quantum effective field theory valid at energies $E \ll \Lambda_{\mathrm{UV}} := \min(M_{\mathrm{Pl}}, \Lambda_{\mathrm{EM}}, \Lambda_{\mathrm{sym}})$.

\subsection{Canonical Quantization and Equal-Time Commutation Relations}

Promoting $\psi(x)$ and $\rho_t(x)$ around Minkowski space $\eta_{\mu\nu} = \mathrm{diag}(-1,+1,+1,+1)$ to operator-valued distributions $\hat{\psi}(\mathbf{x},t)$, $\hat{\psi}^\dagger(\mathbf{x},t)$, and $\hat{\rho}_t(\mathbf{x},t)$, the $(-,+,+,+)$ kinetic Lagrangian $-\frac{\kappa_t}{2}\eta^{\mu\nu}\partial_\mu\rho_t\partial_\nu\rho_t - \frac{\kappa_c}{2}\eta^{\mu\nu}(D_\mu\psi)^\dagger(D_\nu\psi) = +\frac{\kappa_t}{2}\dot{\rho}_t^2 + \frac{\kappa_c}{2}|D_0\psi|^2 - \dots$ defines positive-norm canonical conjugate momenta:
\begin{align}
\hat{\Pi}_\rho(\mathbf{x},t) &:= \frac{\partial\mathcal{L}_{\mathrm{AUTH}}}{\partial(\partial_0\hat{\rho}_t)} = +\kappa_t\,\partial_0\hat{\rho}_t(\mathbf{x},t), \label{eq:canon_mom_rho}\\
\hat{\Pi}_\psi(\mathbf{x},t) &:= \frac{\partial\mathcal{L}_{\mathrm{AUTH}}}{\partial(\partial_0\hat{\psi})} = +\frac{\kappa_c}{2}\,(D_0\hat{\psi})^\dagger(\mathbf{x},t), \label{eq:canon_mom_psi}\\
\hat{\Pi}_{\psi^\dagger}(\mathbf{x},t) &:= \frac{\partial\mathcal{L}_{\mathrm{AUTH}}}{\partial(\partial_0\hat{\psi}^\dagger)} = +\frac{\kappa_c}{2}\,D_0\hat{\psi}(\mathbf{x},t). \label{eq:canon_mom_psidag}
\end{align}
We impose standard equal-time canonical commutation relations ($\hbar = 1$, restoring $\hbar$ explicitly where informative):
\begin{align}
\bigl[\hat{\psi}(\mathbf{x},t),\,\hat{\Pi}_\psi(\mathbf{y},t)\bigr] &= i\hbar\,\delta^{(3)}(\mathbf{x}-\mathbf{y}),
\qquad
\bigl[\hat{\psi}^\dagger(\mathbf{x},t),\,\hat{\Pi}_{\psi^\dagger}(\mathbf{y},t)\bigr] = i\hbar\,\delta^{(3)}(\mathbf{x}-\mathbf{y}), \nonumber\\
\bigl[\hat{\rho}_t(\mathbf{x},t),\,\hat{\Pi}_\rho(\mathbf{y},t)\bigr] &= i\hbar\,\delta^{(3)}(\mathbf{x}-\mathbf{y}),
\label{eq:equal_time_ccr}
\end{align}
with all remaining equal-time commutators vanishing.

\subsection{Bounded-Below Linearized Fock/KLMN Hamiltonian and Retirement of Unconstructed SUSY (\texttt{C06})}

Expanding \texttt{M\_AUTH\_SIT\_EFT\_v1} to quadratic order around the stable vacuum $(\eta_{\mu\nu}, \rho_0, v_0, A_\mu = 0)$ in ADM scalar--vector--tensor variables reveals how gauge constraints eliminate every unphysical negative-norm mode (\texttt{SITQG09LinearizedFluctuationAndObstruction}, \texttt{SITTrancheBAuthoritativeModelAndQuantum}):
\begin{itemize}
  \item \textbf{Constraint Elimination of Conformal and Temporal-Gauge Ghosts:} Although the unconstrained metric trace carries a conformal kinetic contribution $-6(\partial_t\Psi)^2$, the linearized ADM Hamiltonian constraint $\Delta\Psi = +\frac{\kappa_G}{2}T_{00}^{(2)}$ renders $\Psi$ non-propagating, while the $U(1)$ Gauss constraint $\nabla\cdot\mathbf{\Pi} - \hat{q}(u\pi_v - v\pi_u) = 0$ eliminates $A_0$.
  \item \textbf{Bounded-Below Physical Fock/KLMN Spectrum:} Across the remaining $7$ physical bosonic polarizations in pure-SIT vacuum ($2$ transverse-traceless gravitons $h_{+,\times}$ with $c_T = 1$, $2$ transverse photons/Proca modes, $1$ longitudinal Proca mode of mass $m_A = \sqrt{\hat{q}^2\kappa_c v_0^2 / f_{\mathrm{EM}0}}$ for $v_0 > 0$ or Goldstone phase $\theta$ for $v_0 = 0$, and $2$ massive scalar eigenstates of $M_{\mathrm{phys}}^2 = K^{-1/2} H K^{-1/2}$ with $\det(H) > 0$), every kinetic weight ($\frac{1}{8\kappa_G}$, $\frac{\kappa_t}{2}$, $\frac{\kappa_c}{2}$, $\frac{f_{\mathrm{EM}0}}{2}$) and every physical mass-squared eigenvalue is strictly non-negative. Under ultraviolet mode regularization $\Lambda_{\mathrm{reg}} \le \Lambda_{\mathrm{UV}}$ with normal ordering (or Kato--Lions--Lax--Milgram--Nelson quadratic-form perturbation for small interactions), the physical quantum Hamiltonian is bounded below:
  \begin{equation}
  \hat{H}_{\mathrm{phys}} \ge E_0 > -\infty.
  \label{eq:klmn_bounded_below}
  \end{equation}
  \item \textbf{Explicit Retirement of Unconstructed Supersymmetry (\texttt{C06}, \texttt{ODQ-4}):} Earlier exploratory drafts postulated an unconstructed supersymmetric factorization $\hat{H}_{\mathrm{SIT}} = \hat{A}^\dagger\hat{A}$ to argue for energy non-negativity. Under \texttt{QG49} Tier-4 hygiene, \texttt{C06} is permanently retired: \texttt{M\_AUTH\_SIT\_EFT\_v1} contains no superpartner fields and requires no supersymmetry, because Eq.~\eqref{eq:klmn_bounded_below} follows directly from $\kappa_G, \kappa_t, \kappa_c > 0$, $f_{\mathrm{EM}} \ge 1$, $V_{\mathrm{tot}} \ge 0$, and $\det(H) > 0$.
\end{itemize}

\subsection{Consistency of Heisenberg, Path-Integral, and Classical Formulations}

\paragraph{(a) Heisenberg Evolution.} Any physical observable $\hat{O}$ evolves via $i\hbar\,d\hat{O}/dt = [\hat{O}, \hat{H}_{\mathrm{phys}}]$, reproducing the operator field equations upon inserting Eq.~\eqref{eq:equal_time_ccr}.

\paragraph{(b) Regulated EFT Path Integral.} Below $\Lambda_{\mathrm{UV}}$, transition amplitudes admit the Faddeev--Popov gauge-fixed path-integral representation
\begin{equation}
\braket{\Psi_f}{e^{-i\hat{H}_{\mathrm{phys}}(t_f-t_i)/\hbar}\Psi_i}
= \int_{\Lambda_{\mathrm{reg}} \le \Lambda_{\mathrm{UV}}} \mathcal{D}[g_{\mu\nu}]\,\mathcal{D}[\rho_t]\,\mathcal{D}[\psi]\,\mathcal{D}[\psi^*]\,\mathcal{D}[A_\mu]\,\mu_{\mathrm{FP}}\,
  \exp\!\left(\frac{i}{\hbar}S_{\mathrm{AUTH}}\right),
\label{eq:eft_path_integral}
\end{equation}
while non-perturbative gauge-coherence path integrals are constructed rigorously on a compact-group 4D hypercubic Euclidean Haar Wilson lattice in Part~IV (Section~\ref{sec:stage4-lattice}).

\paragraph{(c) Semiclassical Recovery.} Both coherent-state expectation values $\bra{\Psi}\delta S_{\mathrm{AUTH}}/\delta\psi^*\ket{\Psi} = 0$ and the $\hbar \to 0$ stationary-phase saddle $\delta S_{\mathrm{AUTH}} = 0$ recover the exact classical Euler--Lagrange equations~\eqref{eq:summary_einstein_sec8}--\eqref{eq:summary_maxwell_sec8}:
\[
\bra{\Psi}\Bigl(\frac{\kappa_c}{2} D_\mu D^\mu\hat{\psi} - \frac{\partial V_{\mathrm{tot}}}{\partial\hat{\psi}^*} + \dots\Bigr)\ket{\Psi} = 0
\;\Longrightarrow\;
\frac{\kappa_c}{2} D_\mu D^\mu\psi - \frac{\partial V_{\mathrm{tot}}}{\partial\psi^*} + \dots = 0.
\]

\section{SIT 3.0: Exploratory State-Dependent Clock Hypothesis and Impossibility of a Linear Coherence Operator}
\label{sec:SIT3}
\label{sec:QIP}

This section preserves the exploratory ``SIT~3.0'' research question---whether quantum coherence or entanglement in a localized register could enter an effective proper-time rate---while correcting a fatal mathematical inconsistency in earlier drafts (which postulated a single fixed Hermitian operator $\hat{\mathcal{R}}_{\mathrm{coh}}$ with mutually incompatible properties). This hypothesis is a \emph{separate phenomenological extension}, not derived from unchanged \texttt{M\_AUTH\_SIT\_EFT\_v1}.

\paragraph{Mathematical Impossibility of the Earlier Linear Operator $\hat{\mathcal{R}}_{\mathrm{coh}}$.}
Earlier SIT drafts simultaneously required a fixed, state-independent Hermitian operator $\hat{\mathcal{R}}_{\mathrm{coh}}$ to satisfy: (a)~$\mathrm{Tr}(\rho\,\hat{\mathcal{R}}_{\mathrm{coh}}) = 2|\rho_{01}|$ on single-qubit density matrices, (b)~$\hat{\mathcal{R}}_{\mathrm{coh}}|\psi_1\rangle\otimes\cdots\otimes|\psi_N\rangle = 0$ on every $N$-partite product state, and (c)~positive eigenvalues $N_c > 0$ and $N_w > 0$ on GHZ and $W$ states. Two elementary proofs establish that no such linear operator exists:
\begin{enumerate}[leftmargin=2em]
  \item \emph{Violation of Linearity Under Mixing:} Both $|+\rangle = (|0\rangle + |1\rangle)/\sqrt{2}$ and $|-\rangle = (|0\rangle - |1\rangle)/\sqrt{2}$ satisfy $2|\rho_{01}| = 1$. Their equal statistical mixture is $\rho = \frac{1}{2}|+\rangle\langle +| + \frac{1}{2}|-\rangle\langle -| = I/2$, which has $2|\rho_{01}| = 0$. Because the expectation value $\mathrm{Tr}(\rho\hat{O})$ of any fixed linear observable $\hat{O}$ is affine in $\rho$, having expectation $1$ on both $|+\rangle$ and $|-\rangle$ forces expectation $1$ (not $0$) on $I/2$. Thus $2|\rho_{01}|$ is a \emph{nonlinear} convex state functional, never the expectation value of a single state-independent Hermitian operator.
  \item \emph{Computational Product Basis Forces $\hat{\mathcal{R}}_{\mathrm{coh}} \equiv 0$:} The $2^N$ computational product states $\{|i_1\rangle\otimes\cdots\otimes|i_N\rangle : i_k \in \{0,1\}\}$ form a complete orthonormal basis of $(\mathbb{C}^2)^{\otimes N}$. Any linear operator that annihilates every product state therefore vanishes identically ($\hat{\mathcal{R}}_{\mathrm{coh}} = 0$) and must also annihilate every superposition thereof, including $(|0\rangle^{\otimes N} + |1\rangle^{\otimes N})/\sqrt{2}$ and $|W_N\rangle$; it cannot possess nonzero GHZ or $W$ eigenvalues.
\end{enumerate}

\paragraph{Corrected Exploratory Formulation (Nonlinear State Functional vs.\ Restricted Linear Witness).}
To keep the state-dependent-clock research question mathematically well posed, Version~81 replaces the impossible linear operator $\hat{\mathcal{R}}_{\mathrm{coh}}$ with one of two distinct, mathematically consistent choices (neither yet derived from \texttt{M\_AUTH\_SIT\_EFT\_v1}):
\begin{itemize}
  \item \textbf{Option A (Declared Nonlinear State Functional $\mathcal{C}_{\mathrm{state}}[\rho]$):} A state-dependent proper-time ansatz takes the form
  \begin{equation}
  \frac{d\tau}{dt} \;=\; \rho_g(x)\,\exp\!\Bigl[\gamma\,\mathcal{C}_{\mathrm{state}}[\rho]\Bigr]
  \;=\; \rho_g(x)\Bigl[1 + \gamma\,\mathcal{C}_{\mathrm{state}}[\rho] + \mathcal{O}(\gamma^2)\Bigr],
  \label{eq:SIT3-master}
  \end{equation}
  where $\mathcal{C}_{\mathrm{state}}[\rho] \ge 0$ is explicitly a \emph{nonlinear} functional of the state $\rho$---for instance, the $\ell_1$-norm basis coherence $2|\rho_{01}|$ for a single qubit in a fixed basis, or a genuine multipartite entanglement monotone that vanishes identically on all separable states (such as concurrence/tangle, negativity, or relative entropy of entanglement) while remaining strictly positive on GHZ or $W$ states. \emph{Correction Regarding Quantum Fisher Information:} Earlier Version~81 wording incorrectly listed ordinary quantum Fisher information $F_Q(\rho, \hat{J})$ as an entanglement monotone vanishing on separable states; in fact, for the separable product state $|+\rangle^{\otimes N}$ and collective generator $\hat{J}_z = \frac{1}{2}\sum_{i=1}^N \hat{Z}_i$, direct evaluation gives $F_Q(|+\rangle^{\otimes N}, \hat{J}_z) = 4\,\mathrm{Var}(\hat{J}_z) = N \ne 0$. Ordinary quantum Fisher information can serve only as an entanglement criterion relative to a declared separability ceiling (for example, the shifted excess $\max(0, F_Q(\rho, \hat{J}_z) - N)$ for $\hat{J}_z$), not as a zero-on-separable functional itself.
  \item \textbf{Option B (State-Specific Linear Entanglement Witness $\hat{W}_{\Psi}$):} For a \emph{fixed} target entangled state (such as a specific GHZ state), a standard linear Hermitian witness $\hat{W}_{\mathrm{GHZ}}$ can satisfy $\mathrm{Tr}(\rho_{\mathrm{sep}}\hat{W}_{\mathrm{GHZ}}) \ge 0$ on separable states while taking a negative value on the target GHZ state, without annihilating all product states or reproducing $2|\rho_{01}|$.
\end{itemize}

\paragraph{Open Status of Downstream Clock, Sorkin ($I_{123}$), and No-Signaling Claims.}
Because the earlier operator $\hat{\mathcal{R}}_{\mathrm{coh}}$ does not exist, all downstream formulas previously attributed to it---specifically GHZ vs.\ $W$ amplification scaling, three-arm clock-beat offsets, and the Sorkin third-order interference coefficient $\kappa_3$---are explicitly marked \textbf{open phenomenological targets} requiring rederivation once a concrete functional $\mathcal{C}_{\mathrm{state}}[\rho]$ (or witness) and microscopic coupling $\gamma$ are specified (\texttt{OPEN-6C}). Furthermore, Version~81 retracts the assertion that a ``phase-only'' state-dependent modification at $\mathcal{O}(\gamma)$ automatically guarantees no-signaling: in nonlinear modifications of quantum mechanics, preventing superluminal signaling requires specifying the full composite-system state-update law, tensor-product extension, and measurement/collapse prescription explicitly (cf.\ Gisin--Polchinski nonlinear quantum mechanics constraints). In the $\gamma \to 0$ limit, standard quantum mechanics and General Relativity are recovered identically.

\bigskip

\subsection{Boundary Between Gauge Connections and Gravitational Curvature}
\label{sec:dual_geometry}

Earlier SIT drafts attempted to derive general electromagnetic curvature directly from a pure phase gradient $A_\mu = (\hbar/e)\partial_\mu\theta$ and to identify spacetime curvature with an ``informational torque'' functional. Version~81 aligns this subsection strictly with \texttt{M\_AUTH\_SIT\_EFT\_v1} and Section~\ref{sec:magnetism-holonomy}:
\begin{itemize}
  \item \textbf{Electromagnetic Sector (Independent Connection, \texttt{OPEN-6B}):} On any smooth, simply connected region, a pure phase gradient $A_\mu^{\mathrm{flat}} = (\hbar/e)\partial_\mu\theta$ has identically vanishing curvature $F_{\mu\nu}^{\mathrm{flat}} = \partial_\mu A_\nu^{\mathrm{flat}} - \partial_\nu A_\mu^{\mathrm{flat}} = 0$ and cannot generate ordinary nonzero magnetic fields ($\mathbf{B} \ne 0$). In canonical \texttt{M\_AUTH\_SIT\_EFT\_v1}, local $U(1)$ gauge invariance is carried by an \emph{independent} dynamical gauge connection $A_\mu$ with arbitrary curvature $F_{\mu\nu} = \partial_\mu A_\nu - \partial_\nu A_\mu$ (massive $m_A^2 = \kappa_c\hat{q}^2 v_0^2/f_{\mathrm{EM}0}$ whenever $\hat{q}v_0 \ne 0$, leaving the visible massless photon dictionary open under \texttt{OPEN-6B}).
  \item \textbf{Gravitational Sector (Variational Einstein Equations):} Spacetime curvature in \texttt{M\_AUTH\_SIT\_EFT\_v1} is governed exclusively by the variational Einstein field equation $G_{\mu\nu} = \kappa_G T_{\mu\nu}^{\mathrm{tot}}$ of Eq.~\eqref{eq:covariant_einstein_sec5}, where scalar gradients of $(\rho_t, \psi)$ enter through their canonical stress--energy tensors $T_{\mu\nu}^{(\rho_t)} + T_{\mu\nu}^{(\psi)} + T_{\mu\nu}^{(\mathrm{pot})}$. Heuristic Berry-phase or ``informational torque'' analogies are kept strictly separate from metric variation (Section~\ref{sec:informational-torque}).
\end{itemize}

\section{Exploratory Multi-Modal Energy-Density Ansatz (Separate from Canonical Action)}
\label{sec:integrated_functional}

This section preserves an exploratory phenomenological ansatz from earlier SIT drafts---\emph{not} derived from unchanged \texttt{M\_AUTH\_SIT\_EFT\_v1} (whose canonical vacuum stress--energy vanishes at $(\rho_0, v_0)$)---exploring how an effective multi-modal energy density might be parameterized in terms of coarse-grained coherence proxies and dynamic rates (see Section~\ref{sec:informational_energy_equivalence} for dimensional and vacuum-subtraction requirements).

\subsection{Phenomenological Local Energy-Density Ansatz}
In this separately postulated heuristic model, an auxiliary energy density is written as
\begin{equation}
    \varepsilon_{\mathrm{SIT}}^{\mathrm{phen}}(x) = \zeta \cdot R_{\mathrm{coh}}(x) \cdot [\rho_t(x)]^2,
    \label{eq:local_energy_density}
\end{equation}
where $\zeta$ is a phenomenological scale parameter (requiring explicit background subtraction relative to vacuum $(\rho_0, v_0)$ and an explicit covariant stress--energy derivation before any comparison with $T_{00}^{\mathrm{tot}}$ of \texttt{M\_AUTH\_SIT\_EFT\_v1}).

This single-mode ansatz is insufficient to describe general quantum systems, motivating the exploratory multi-modal sketch below (distinct from the proved canonical Hamiltonian of Section~\ref{sec:quantum-operator-formalism}).

\subsection{Exploratory Multi-Modal Sketch}
To parameterize multiple coarse-grained oscillatory modes ($i$) heuristically, earlier drafts wrote a schematic mode-sum expression:
\begin{equation}
    \varepsilon_{\mathrm{total}}^{\mathrm{phen}}(x,t) = \sum_i \left\{ \zeta_i R_{\mathrm{coh},i} [\rho_{t,i}]^2 + \frac{1}{2}(\partial_\mu R_{\mathrm{coh},i})^2 + \frac{\kappa}{2}(\partial_\mu \rho_{t,i})^2 \right\}
    \label{eq:full_functional}
\end{equation}
(a classical heuristic mode sum, distinct from the canonical Hamiltonian of \texttt{M\_AUTH\_SIT\_EFT\_v1} in Section~\ref{sec:quantum-operator-formalism} and explicitly omitting the impossible linear coherence operator $\hat{\mathcal{R}}_{\mathrm{coh}}$ retracted in Section~\ref{sec:SIT3}).

\subsection{Boundary Between This Exploratory Sketch and Canonical \texttt{M\_AUTH\_SIT\_EFT\_v1}}
Version~81 clarifies three points where earlier drafts overstated the scope of Eq.~\eqref{eq:full_functional}:
\begin{enumerate}[leftmargin=2em]
  \item \textbf{Gravitational and Cosmological Sector:} In unchanged \texttt{M\_AUTH\_SIT\_EFT\_v1}, metric variation produces $T_{\mu\nu}^{\mathrm{tot}}$ with potential $V_{\mathrm{tot}}(\rho_t, |\psi|) \ge 0$ satisfying $V_{\mathrm{tot}}(\rho_0, v_0) = 0$ (not $V = \zeta R\rho_t^2$, which is nonzero at the constant vacuum, Section~\ref{sec:informational_energy_equivalence}). Dark-matter and dark-energy parameterizations remain open Level-1 cosmological hypotheses (Section~\ref{sec:black-holes-voids}, \texttt{OPEN-6D}).
  \item \textbf{Quantum State Dependence:} Because no fixed linear Hermitian operator $\hat{\mathcal{R}}_{\mathrm{coh}}$ can annihilate all product states while giving positive GHZ eigenvalues (Section~\ref{sec:SIT3}), any state-dependent energy or BEC-versus-thermal gravitational response requires either an action-derived microscopic state-preparation source $J_{\mathrm{prep}}$ (\texttt{OPEN-6C}) or a separately specified nonlinear/phenomenological condensate model (such as Paper~003's Penrose--Onsager contract, Section~\ref{sec:smoking-gun-bec}).
  \item \textbf{Thermodynamic Dissipation:} Closed-system solutions of \texttt{M\_AUTH\_SIT\_EFT\_v1} satisfy $\nabla_\mu J^\mu_{\mathrm{coh}} = 0$ identically; irreversible dissipation is governed in finite-dimensional open quantum systems by unital CPTP purity non-increase and general CPTP relative-entropy contraction (Result~\ref{thm:CoherenceMonotonicity}), rather than by claiming $\partial_\mu J^\mu_{\mathrm{coh}} > 0$ derives the Second Law of Thermodynamics directly from the canonical action.
\end{enumerate}



\section{Exploratory Relational Network Analogy and Resolution of \texttt{C08}}
\label{sec:network-formulation}

Super Information Theory also motivates discrete relational graph toy models for mesoscopic oscillator simulations, provided naive row-normalized stochastic averaging (\texttt{C08}) is avoided and the mathematical boundary between heuristic graph dynamics and the continuum field theory \texttt{M\_AUTH\_SIT\_EFT\_v1} is maintained strictly.

\paragraph{Relational Events and Hermitian Adjacency Amplitudes.}
Let $V = \{v_1, v_2, \dots, v_N\}$ denote a finite set of elemental relational nodes equipped with a Hermitian complex adjacency matrix $A(\tau) = A^\dagger(\tau)$ ($A_{ij}(\tau) = A_{ji}^*(\tau)$) parameterized by internal evolution parameter $\tau$, carrying link phase $\arg A_{ij}(\tau)$ and relational weight $|A_{ij}(\tau)|^2$. Each edge retains a Hermitian exponentially weighted memory kernel
\begin{equation}
M_{ij}(\tau) = \int_0^\tau e^{-(\tau-\sigma)/\tau_m}\,A_{ij}(\sigma)\,d\sigma
\qquad \bigl(M(\tau) = M^\dagger(\tau)\bigr)
\end{equation}
with memory horizon $\tau_m > 0$, while node-level scalar proxies track accumulated activity and local phase alignment:
\begin{equation}
\rho_t^{\mathrm{net}}(i,\tau) = \frac{1}{\tau_0}\int_0^\tau \sum_j |A_{ij}(\sigma)|^2\,d\sigma,
\qquad
\Rcoh^{\mathrm{net}}(i,\tau) = \frac{1}{d_i}\left|\sum_j A_{ij}(\tau)\right|.
\end{equation}

\paragraph{Hermiticity Repair and Explicit Withdrawal of Earlier Continuum-Action Claim.}
Earlier drafts proposed the evolution equation $i\hbar\,\partial_\tau A_{ij} = -\kappa_{\mathrm{net}}(LA)_{ij} - i\lambda_{\mathrm{diss}}A_{ij} + \mu_{\mathrm{mem}}M_{ij} + \nu_{\mathrm{net}}[\rho_t(i)-\rho_t(j)]A_{ij}$ (with symmetric graph Laplacian $L = D - W = L^\dagger$) while declaring $A(\tau) = A^\dagger(\tau)$ Hermitian. That equation does \emph{not} preserve Hermiticity: dividing by $i\hbar$ maps Hermitian contributions into anti-Hermitian derivatives, yielding the nonzero Hermiticity defect
\begin{equation}
\partial_\tau\bigl(A - A^\dagger\bigr) = \frac{i}{\hbar}\Bigl[\kappa_{\mathrm{net}}(LA + AL) - 2\mu_{\mathrm{mem}} M\Bigr] \ne 0
\label{eq:network_hermiticity_defect}
\end{equation}
even on a two-node graph ($L = \bigl(\begin{smallmatrix}1 & -1\\ -1 & 1\end{smallmatrix}\bigr)$, real symmetric $A(0) = \bigl(\begin{smallmatrix}1 & 1\\ 1 & 2\end{smallmatrix}\bigr)$, positive damping $\lambda_{\mathrm{diss}} > 0$), so the trajectory immediately leaves the declared Hermitian state space! To preserve $A(\tau) = A^\dagger(\tau)$ identically along the flow, a consistent exploratory graph toy model must place Hermitian operators inside a commutator $-i[H_{\mathrm{net}}, A]/\hbar$ and real symmetric damping/diffusion/memory terms outside the imaginary unit:
\begin{equation}
\partial_\tau A(\tau)
= -\frac{i}{\hbar}\bigl[H_{\mathrm{net}}(\tau),\, A(\tau)\bigr]
  - \frac{\lambda_{\mathrm{diss}}}{\hbar}\,A(\tau)
  - \frac{\kappa_{\mathrm{diff}}}{2}\bigl\{L,\, A(\tau)\bigr\}
  + \frac{\mu_{\mathrm{mem}}}{\hbar}\,M(\tau),
\label{eq:discrete_super_coherence}
\end{equation}
where $H_{\mathrm{net}}(\tau) := -\kappa_{\mathrm{net}} L + \nu_{\mathrm{net}}\,\mathrm{diag}(\rho_t^{\mathrm{net}}(\tau)) = H_{\mathrm{net}}^\dagger(\tau)$ and $\{L, A\} = LA + AL$, ensuring $\partial_\tau(A - A^\dagger) \equiv 0$ identically whenever $A = A^\dagger$ and $M = M^\dagger$.

\paragraph{Explicit Scope Boundary (Not a Derivation of $S_{\mathrm{AUTH}}$).}
Version~81 explicitly withdraws the earlier claim that taking $a \to 0$ in a first-order discrete graph toy model derives the full relativistic gravitational-gauge continuum action $S_{\mathrm{AUTH}}$ (Eq.~\eqref{eq:canonical_action_sec5}). Equation~\eqref{eq:discrete_super_coherence} is retained strictly as a Hermiticity-preserving phenomenological graph toy model; by contrast, the proved classical $C^4$ continuum limit of SIT applies strictly to the Stage-6 periodic 4D lattice kinetic and plaquette functionals of Section~\ref{sec:stage4-lattice} (with full quantum continuum QFT construction explicitly open under \texttt{OPEN-1}).

\subsection{Mathematical Resolution of Row-Normalized Clock Collapse (\texttt{C08})}
\label{sec:c08-network-resolution}

An important structural obstruction uncovered during formal verification (\texttt{ODQ-6}, formalized in Lean~4.30 module \texttt{HsitTail.SitUv.SITF41F42ProjectionCoherenceTower}, theorem \texttt{unbounded\_fiber\_collapses\_projected\_support}) concerns how discrete network amplitudes are normalized and projected:
\begin{enumerate}[leftmargin=2em]
  \item \textbf{Why Naive Row-Normalized Stochastic Clock Readout Collapses (\texttt{C08} Retired):} If one imposes strict row-stochastic normalization $\sum_{j=1}^N |A_{ij}(\tau)|^2 = 1$ for every vertex $i$ \emph{and} simultaneously defines local coherence as $\Rcoh^{\mathrm{naive}}(i,\tau) := \sum_j |A_{ij}(\tau)|^2$, then $\Rcoh^{\mathrm{naive}}(i,\tau) \equiv 1$ and $\rho_t^{\mathrm{naive}}(i,\tau) = (\tau/\tau_0)\cdot 1$ become spatially uniform across all vertices $i$, collapsing $\nabla\Rcoh \equiv 0$ and $\nabla\rho_t \equiv 0$ identically! Similarly, iterating a row-normalized random-walk transition matrix $P = D^{-1}W$ ($P\mathbf{1} = \mathbf{1}$) without spectral scale separation washes out spatial clock gradients, while unweighted projection across unbounded microscopic fibers collapses projected norm support. Accordingly, naive row-normalized clock readout (\texttt{C08}) is permanently retired in Tier~4.
  \item \textbf{Continuum Spectral Sobolev and 4D Euclidean Wilson Replacements:} Two rigorous mathematical structures replace \texttt{C08} in Version~81:
  \begin{itemize}
    \item \emph{Weighted Sobolev Eigenmode Towers (Section~\ref{sec:stage3-sobolev}):} Expanding fields in the orthonormal eigenbasis $\{\phi_n\}_{n\in\mathbb{N}}$ of the unnormalized Laplace--Beltrami operator $-\Delta_q\phi_n = \mu_n\phi_n$ ($\mu_n \ge 0$, $\mu_n \to \infty$) yields the strict infinite-dimensional Sobolev hierarchy $H^2(\mu) \subsetneq H^1(\mu) \subsetneq \ell^2(\mathbb{N})$ with controlled spectral Galerkin defect bounds $\|u - P_N u\|_{H^1(\mu)}^2 \le (1 + \Lambda_N^2)^{-1}\|u\|_{H^2(\mu)}^2$.
    \item \emph{Compact-Group 4D Hypercubic Wilson Lattices (Section~\ref{sec:stage4-lattice}):} For non-perturbative gauge-coherence physics, edges $(x,\mu)$ on a 4D Euclidean hypercubic torus $(\mathbb{Z}/l\mathbb{Z})^4$ carry unitary group transporters $U(x,\mu) \in G \subset U(n+1)$ integrated with normalized bi-invariant Haar measure $dh(U_e)$, while sites carry unpinned matter/coherence amplitudes $\psi(x) \in \mathbb{C}^{n+1}$ and temporal scalars $\rho_t(x) \in \mathbb{R}$ coupled through $\sum_{x,\mu}\|\psi(x) - U(x,\mu)\psi(x+\hat{\mu})\|^2$, preserving both spatial gradients and exact local gauge invariance at every finite lattice spacing.
  \end{itemize}
\end{enumerate}

\section{Time Density and Phase--Rate Dynamics}\label{sec:timedensity}

Super Information Theory identifies a scalar \emph{time--density field} $\rho_t(x)$, with physical dimension $[\mathrm{T}^{-1}]$, as one of its two primitive variables.  Its operational meaning is set by the phase accumulation of a reference clock carried along any world-line.  Let $\varphi(x)$ denote the phase of the locally maximally coherent mode.  We postulate
\begin{equation}\label{eq:phi-rate}
    \dot{\varphi}(x)\;=\;\frac{S_{\text{coh}}(x)}{\hbar_{\rm eff}},
    \qquad
    \rho_t(x)\;=\;\rho_0\,\Bigl(\dot{\varphi}_0/\dot{\varphi}(x)\Bigr),
\end{equation}
where $S_{\text{coh}}$ is the density of the coherence action and $\hbar_{\rm eff}$ is a fixed microscopic constant that calibrates informational action in ordinary SI units.  Equation~\eqref{eq:phi-rate} replaces the older ``quantum stopwatch'' metaphor with a dimensionally explicit relation: slower phase advance (smaller~$\dot{\varphi}$) corresponds to larger~$\rho_t$ and hence stronger scalar--tensor back-reaction in the action of Sec.~\ref{sec:framework}.  The converse holds for decoherence-dominated regions.

Linearising the canonical field equations of \texttt{M\_AUTH\_SIT\_EFT\_v1} around the uniform vacuum $(\rho_0, v_0)$ produces the coupled symmetric two-scalar linear response operator $(K\nabla^2 - H)$ derived in Sections~\ref{sec:framework} and~\ref{sec:screening-mechanism} (rather than a one-way uncoupled Poisson modification $\nabla^2\Phi = 4\pi G\rho_m + \beta\nabla^2\delta\rho_t$, which is a pedagogical approximation, Section~\ref{sec:experimental-predictions}). Earlier drafts asserted an unqualified laboratory torsion-balance bound $|\beta| < 10^{-5}$; Version~81 retires that standalone bound because short-range fifth-force constraints depend jointly on the mass eigenvalues $(m_-, m_+)$, the source coupling $Q_{\mathrm{eff}}$, and finite-geometry screening integrals (\texttt{OPEN-6A}, Section~\ref{sec:experimental-predictions}). High-accuracy transportable optical clocks probe the disambiguated local field-response coefficient $\alpha_\rho$ via $\Delta\nu/\nu = \alpha_\rho\,\delta\rho_t/\rho_0$ and gravitational-potential redshift slope $\alpha_\Phi = \alpha_\rho\mathcal{T}$ (Appendix~\ref{appendix:ClockConstraints}).

\subsection{Self-Organising Feedback and Coherence Flow}

The global $U(1)$ phase symmetry of the complex field $\psi$ (or gauge-covariant matter current when $\hat{q} \neq 0$) leads, via Noether's theorem, to a conserved $U(1)$ charge current:
\[
    J^\mu_{\mathrm{coh}} = \frac{i\kappa_c}{2} \left( \psi (D^\mu \psi)^* - \psi^* D^\mu \psi \right) = \kappa_c |\psi|^2 (\nabla^\mu \theta - \hat{q} A^\mu),
    \qquad
    \nabla_\mu J^\mu_{\mathrm{coh}} = 0.
\]
This current is the exact classical Noether current of \texttt{M\_AUTH\_SIT\_EFT\_v1} (distinct from operational quantum-state coherence measures, Section~\ref{sec:coherence-conservation}). In unchanged \texttt{M\_AUTH\_SIT\_EFT\_v1}, spacetime curvature is sourced strictly by the symmetric metric stress--energy tensor $T_{\mu\nu}^{\mathrm{tot}}$ (Section~\ref{sec:stress-energy}); earlier drafts introduced an antisymmetric heuristic tensor $\tau_{\mu\nu\rho} = \beta\nabla_{[\mu}\Rcoh\nabla_{\nu]}\rho_t u_\rho$ and claimed a $10\,\mathrm{pm}$ atom-interferometer displacement over a one-meter baseline. Consistent with Section~\ref{sec:informational_torque}, Version~81 retracts both the unproved antisymmetric Einstein source term and the unsupported $10\,\mathrm{pm}$ interferometer forecast as open phenomenological items requiring a complete action-to-apparatus derivation (\texttt{OPEN-6}).

\subsection{Relation to Quantum Interference}

In standard gauge theory and in \texttt{M\_AUTH\_SIT\_EFT\_v1} with an independent gauge connection $A_\mu$ (Section~\ref{sec:magnetism-holonomy}), the electromagnetic Aharonov--Bohm phase around a closed loop enclosing magnetic flux $\Phi_B = \int F$ is $\theta_{\mathrm{AB}} = (e/\hbar)\oint A_\mu\,dx^\mu$ (nonzero because $F_{\mu\nu} = \partial_\mu A_\nu - \partial_\nu A_\mu \neq 0$ inside the solenoid). Earlier drafts wrote $\theta_{\mathrm{AB}} = (e/\hbar)\oint (\nabla\theta)\cdot d\boldsymbol{\ell}$ and claimed electron interferometers already constrain $\lambda \Rcoh$; Version~81 retracts that expression because a single-valued gradient $\nabla\theta$ has vanishing curl and zero line integral on any simply connected patch, while the multiplicative constitutive dressing $f_{\mathrm{EM}}(\Rcoh) = 1 + c_{\mathrm{EM}}\Rcoh^2$ modifies electromagnetic propagation only through explicit $\Rcoh$-dependent constitutive susceptibility rather than replacing $A_\mu$ with $\nabla\theta$ (Section~\ref{sec:magnetism-holonomy}).

\section{Exploratory Quantum-Domain Analogies and Experimental Proposals}
\label{sec:applications-predictions}

The subsections below collect exploratory quantum-domain analogies from earlier drafts while explicitly separating conceptual sketches from proved consequences of \texttt{M\_AUTH\_SIT\_EFT\_v1} and repairing earlier algebraic sign and dimensional inconsistencies.

\subsection{Exploratory Tunneling-Rate Analogy}
\label{sec:quantum-tunneling}

In standard quantum mechanics, one-dimensional barrier transmission is governed by the stationary Schr\"odinger equation $-(\hbar^2/2m)\psi'' + [V(x) - E]\psi = 0$. Earlier SIT drafts explored a phenomenological analogy in which an effective barrier stiffness or inverse diffusion parameter depends on local time density, $-D(\rho_t)\,u'' + [V(x)-E]u = 0$ with WKB exponent $2\int_{x_1}^{x_2}\sqrt{(V(x)-E)/D(\rho_t(x))}\,dx$. Note that this stationary equation is a heuristic phenomenological model (not a direct consequence of \texttt{M\_AUTH\_SIT\_EFT\_v1}). Moreover, Version~81 corrects a sign error in earlier drafts: if $D(\rho_t) \propto 1/\rho_t$ so that the WKB action scales as $\int\sqrt{(V-E)\rho_t}\,dx$, then \emph{increasing} $\rho_t$ \emph{increases} the WKB barrier exponent and \emph{suppresses} (rather than enhances) transmission $T \sim \exp(-2S_{\mathrm{eff}}/\hbar)$; conversely, enhancing transmission at larger $\rho_t$ would require $D(\rho_t)$ to increase with $\rho_t$. Any physical SIT modification to tunneling rates remains an open Level-1 phenomenology question (\texttt{OPEN-6C}).

\subsection{Time--Energy Uncertainty Bounds and Consistency Check}
\label{sec:quantum-time-energy}

The quantum time--energy uncertainty relation $\Delta E\,\Delta t \ge \hbar/2$ (in Mandelstam--Tamm operational form, relating energy spread to the minimum timescale for orthogonal state evolution) is a standard kinematic theorem of Hilbert-space quantum mechanics. Earlier SIT drafts heuristically substituted $\Delta t \approx \rho_t^{-1}$ together with $\Delta E \approx E_{\mathrm{coh}}(1 - \Rcoh)$ while simultaneously associating high coherence ($\Rcoh \to 1$) with large $\rho_t$ and ``greater temporal uncertainty.'' Version~81 explicitly retracts that heuristic combination as internally inconsistent:
\begin{enumerate}[leftmargin=2em]
  \item If $\Delta t = \rho_t^{-1}$, then larger $\rho_t$ gives \emph{smaller} (not greater) $\Delta t$;
  \item Taking $\Rcoh \to 1$ together with large $\rho_t$ drives the product $E_{\mathrm{coh}}(1 - \Rcoh)\rho_t^{-1} \to 0$, violating $\ge \hbar/2$ (and conversely rearranging the inequality as $\rho_t \le 2E_{\mathrm{coh}}(1 - \Rcoh)/\hbar$ would force $\rho_t \to 0$ as $\Rcoh \to 1$, contradicting the manuscript's own high-$\rho_t$ / high-$\Rcoh$ correlation).
\end{enumerate}
Standard quantum uncertainty relations therefore stand on their own Hilbert-space derivation and are neither replaced by nor derived from the classical EFT scalar fields $(\rho_t, \psi)$.

\subsection{SuperTimePosition Coordinate (Exploratory Concept)}
\label{sec:stp}

In the exploratory SuperTimePosition (STP) conceptualization (examined further in Paper~034), one defines a world-line integrated temporal phase count $\Theta(x,t) := \int^t \rho_t(x,t')\,dt'$. In \texttt{M\_AUTH\_SIT\_EFT\_v1}, the dynamical field equations are the covariant wave equations for $(\rho_t, \psi, A_\mu, g_{\mu\nu})$ derived in Section~\ref{sec:framework} (rather than independent wave equations for $\Theta$); whether an STP parameterization admits a non-trivial relativistic observable beyond standard proper time $\tau = \int ds$ remains an open research question.

\subsection{Measurement-Induced Perturbations (Dimensional Repair and Open Status)}
\label{sec:measurement}

Earlier drafts sketched a transient gravitational-potential perturbation during quantum measurement via $\delta V_{\mathrm{grav}} \sim G\,\delta\rho_g/r$ and $\delta\nu/\nu \approx \delta V_{\mathrm{grav}}$. Version~81 notes two dimensional omissions in that sketch---a gravitational potential requires integrating mass density over the active source volume $\delta M = \int_V \delta\rho_g\,d^3x$ (giving $\delta\Phi \sim -G\,\delta M/r$, Section~\ref{sec:informational_energy_equivalence}), and the dimensionless fractional redshift divides by $c^2$ ($\delta\nu/\nu \sim \delta\Phi/c^2$)---and marks any measurement-induced scalar or gravitational transient as an unproved Level-1 hypothesis contingent on deriving an explicit microscopic state-preparation source $J_{\mathrm{prep}}$ (\texttt{OPEN-6C}).

\subsection{Distinction from Position--Momentum Uncertainty}
\label{sec:uncertainty-distinction}

Standard canonical position--momentum uncertainty $\Delta x\,\Delta p \ge \hbar/2$ and Mandelstam--Tamm time--energy bounds belong to quantum kinematics, whereas the classical SIT EFT fields $(\rho_t, \psi)$ govern covariant field dynamics at scales below $\Lambda_{\mathrm{UV}}$; Version~81 keeps these two domains cleanly separated rather than equating $\rho_t^{-1}$ with a quantum measurement duration.

\subsection{Singularity Avoidance Remains an Open Ultraviolet Question}
\label{sec:collapse-prevention}

Because the earlier heuristic inequality $\rho_t \le 2E_{\mathrm{coh}}(1 - \Rcoh)/\hbar$ is internally inconsistent with $\Rcoh \to 1$ at large $\rho_t$ (Section~\ref{sec:quantum-time-energy}), Version~81 retracts the claim that this algebraic rearrangement prevents classical gravitational singularities. Within the effective field theory \texttt{M\_AUTH\_SIT\_EFT\_v1}, curvature behavior above the cutoff scale $\Lambda_{\mathrm{UV}}$ lies outside the EFT regime of validity, and singularity resolution remains an open ultraviolet quantum-gravity problem (\texttt{OPEN-5}).

\subsection{Light, Lasers, and the Gravitational Field}

Coherent light, such as that produced by a laser, represents a macroscopic quantum state in which photons are phase-aligned and occupy a well-defined mode. In conventional physics, the gravitational field generated by a laser is expected to be minuscule, scaling only with the energy of the beam and independent of its coherence properties. This expectation has been formalized in calculations showing that a 1~J laser pulse, regardless of coherence, generates the same infinitesimal spacetime curvature as an equivalent incoherent light pulse with the same energy distribution. Classical general relativity therefore treats coherence as organizational, not causal: it simply concentrates energy, but does not amplify or uniquely source gravity.

In contrast, SIT proposes that coherence---when maximized among a system of photons---induces localized variations in the time-density field, $\rho_t$, which couple directly to spacetime curvature. The hypothesis is that the more perfectly phase-aligned a quantum system, the more efficiently it modulates or ``focuses'' the informational content of its stress-energy, leading to enhanced or anomalous gravitational effects not captured by classical GR.

\subsection{Bose--Einstein Condensates and Coherence-Induced Gravity}

Matter-wave coherence is realized most dramatically in Bose--Einstein condensates (BECs), where a macroscopic population of atoms occupies the same quantum state. In standard physics, the gravitational field produced by a BEC is again expected to reflect only its total mass-energy; coherence is not considered as a causal modifier. SIT, however, extends its conjecture to matter systems: the perfect coherence of a BEC is posited to alter local time-density gradients, potentially generating subtle deviations in gravitational behavior compared to an incoherent ensemble of the same mass.

If SIT is correct, one would expect ultra-coherent systems---whether composed of photons or massive particles---to generate or modulate gravitational fields in ways that cannot be reduced to energy density alone. This prediction sets the stage for experimental challenge.


\subsection{Experimental Evidence and Proposals}

To date, no experiment has definitively observed gravitational anomalies attributable to coherence per se. The gravitational field of laboratory light beams, including intense lasers, remains below current detection thresholds, and all observed phenomena are consistent with classical predictions. Likewise, tests with BECs have not revealed deviations from the equivalence principle: coherent matter falls at the same rate as incoherent matter, within experimental error.

Nevertheless, advancing technology is narrowing the gap between SIT predictions and empirical accessibility. Proposed experiments include:
\begin{itemize}
    \item High-power laser interferometry to search for self-induced gravitational lensing or beam--beam interactions exceeding classical expectations.
    \item Drop tests comparing the free-fall acceleration of coherent BECs versus thermal clouds under ultra-sensitive gravimetry.
    \item Quantum tomography and interferometric protocols measuring declared operational coherence statistics (purity, visibility, or condensate fraction) under matched-preparation controls (Eq.~\eqref{eq:bec_matched_residual}) once an explicit state-to-field map (\texttt{OPEN-6C}) is specified.
\end{itemize}

Furthermore, proposed gravity-mediated entanglement experiments (such as BMV-type setups) probe whether gravitational fields act quantum-mechanically on spatial superpositions; testing any SIT-specific contribution beyond standard linearized quantum gravity requires deriving the corresponding observable shift from a specified action and parameter point (\texttt{OPEN-6}).

\subsection{Scoped Falsifiability}

Consistent with the Scoped Falsification Principle (Section~\ref{sec:experimental-predictions}), a null result excludes only a specified action, parameter point, matter closure, and preparation model whose derived prediction exceeds experimental sensitivity, rather than falsifying the entire research program in the abstract while \texttt{OPEN-6} remains open.
% --- END INPUT sections/03_part1_sec09_to_15.tex ---
% --- BEGIN INPUT sections/04_part1_sec16_to_25.tex ---
\section{The Arrow of Time: Unital Purity Monotonicity, Relative-Entropy Contraction, and Informational Hysteresis}
\label{sec:arrow-of-time}
\label{sec:thermodynamics}
\label{sec:HTheoremProof}

At the microscopic level, the classical field equations of \texttt{M\_AUTH\_SIT\_EFT\_v1} are strictly time--reversal symmetric: replacing $t \mapsto -t$, $\theta \mapsto -\theta$, and $A_i \mapsto -A_i$ leaves the covariant action and its Euler--Lagrange equations invariant. The macroscopic arrow of time therefore cannot be a fundamental asymmetry of the classical field equations; it arises from initial/boundary conditions and coarse-grained open-system dynamics.

Following the operational separation contract of Section~\ref{sec:PrimitiveFieldDefinitions} and the July~2026 CPTP audit of Paper~002, Version~81 strictly distinguishes three different mathematical statements that earlier drafts conflated: (i)~on-shell $U(1)$ Noether charge conservation $\nabla_\mu J^\mu_{\mathrm{coh}} = 0$ for the classical EFT field $\psi$; (ii)~monotonicity of quantum state purity $P(\rho) = \mathrm{Tr}(\rho^2)$, which holds universally \emph{only} for unital CPTP maps (and fails for non-unital channels such as amplitude damping); and (iii)~monotonic contraction of quantum relative entropy $D(\Phi(\rho)\|\Phi(\omega)) \le D(\rho\|\omega)$, which holds for arbitrary CPTP maps (Result~\ref{thm:CoherenceMonotonicity}).

\paragraph{Integrative Context and Narrowed Kinetic Attribution.}
SIT draws conceptual motivation from Shannon entropy, open-system quantum thermodynamics, and rigorous kinetic theory. In particular, the Deng--Hani--Ma (2025) propagation-of-chaos theorem rigorously derives the classical Boltzmann kinetic limit for a specified microscopic hard-sphere gas; Version~81 cites it strictly as an illustrative mathematical benchmark showing how macroscopic irreversibility can be proved from a defined microscopic reversible model, \emph{not} as a proof that trajectories increasing the SIT scalar modulus $\Rcoh$ have vanishing measure (applying kinetic-limit reasoning to SIT fields would require its own separate microscopic model and scaling proof).

A conceptual research theme of SIT is modeling equilibration via local iterative ``signal-dissipation'' updates that reduce local differences across coupled degrees of freedom (Subsection~16.7 below). Such discrete dissipative update models are \emph{separate reduced phenomenological models}, distinct from the time-reversible classical Euler--Lagrange equations of \texttt{M\_AUTH\_SIT\_EFT\_v1}.

In the canonical action, closed-system $U(1)$ symmetry furnishes the exact conserved gauge-invariant current $J^\mu_{\mathrm{coh}} = \kappa_c |\psi|^2 (\nabla^\mu\theta - \hat{q}A^\mu)$ ($\nabla_\mu J^\mu_{\mathrm{coh}} = 0$), which governs Noether charge transport of $\psi$ without by itself deriving open-system entropy production or purity decay.

\begin{result}[Corrected Open-System Monotonicity: Unital CPTP Purity Non-Increase and General CPTP Relative-Entropy Contraction]
\label{thm:CoherenceMonotonicity}
Let $\rho_{\mathrm{sys}}$ be a finite-dimensional quantum state ($d > 1$) evolving under a completely positive, trace-preserving (CPTP) map $\Phi$ or Markovian Lindblad generator, and let $R_{\mathrm{coh}}^{\mathrm{proxy}}(\rho_{\mathrm{sys}}) := (\mathrm{Tr}[\rho_{\mathrm{sys}}^2] - 1/d)/(1 - 1/d)$ denote the operational normalized purity proxy (kept distinct from the classical EFT field amplitude $\Rcoh = |\psi|/\psi_0$ pending \texttt{OPEN-6C}):
\begin{enumerate}[leftmargin=2em]
  \item \textbf{Retraction of Version~79 Universal CPTP Purity Claim:} Contrary to Version~79, density-matrix purity $P(\rho) = \mathrm{Tr}(\rho^2)$ does \emph{not} decrease under every CPTP map, nor does unchanged purity imply unitary evolution. Two exact finite-dimensional counterexamples (imported from July~2026 Paper~002, Section~8) refute both statements:
  \begin{itemize}
    \item \emph{Non-unital amplitude damping increases purity:} For $0 < p \le 1$ with Kraus operators $K_0 = |0\rangle\langle 0| + \sqrt{1-p}\,|1\rangle\langle 1|$ and $K_1 = \sqrt{p}\,|0\rangle\langle 1|$, acting on the maximally mixed qubit $\rho = I/2$ produces $\Phi(\rho) = \mathrm{diag}\bigl((1+p)/2,\,(1-p)/2\bigr)$, whose purity $(1+p^2)/2 > 1/2$ strictly increases.
    \item \emph{Non-unitary complete dephasing preserves purity on diagonal states:} Complete dephasing in the computational basis is non-unitary, yet fixes every already-diagonal density matrix with identically unchanged purity.
  \end{itemize}
  \item \textbf{Exact Unital CPTP Purity Non-Increase:} In fixed finite dimension $d$, purity $P(\rho) = \mathrm{Tr}(\rho^2)$ (and hence $R_{\mathrm{coh}}^{\mathrm{proxy}}$) is non-increasing for \emph{every} state $\rho$ under $\Phi$ if and only if $\Phi$ is \emph{unital} ($\Phi(I) = I$). For continuous-time Lindblad dynamics~\eqref{eq:Lindblad_equation}, unitality holds whenever $\sum_k \gamma_k [L_k, L_k^\dagger] = 0$ (for example, Hermitian jump operators $L_k = L_k^\dagger$), yielding $\frac{d}{dt}\mathrm{Tr}(\rho_{\mathrm{sys}}^2) \le 0$ and $\frac{d}{dt}R_{\mathrm{coh}}^{\mathrm{proxy}} \le 0$ for every state.
  \item \textbf{General CPTP Contraction via Quantum Relative Entropy:} For \emph{any} (including non-unital) CPTP map $\Phi$ and any two states $\rho, \omega$ satisfying $\mathrm{supp}(\rho) \subseteq \mathrm{supp}(\omega)$, the quantum relative entropy obeys the data-processing inequality
  \begin{equation}
  D\bigl(\Phi(\rho)\,\|\,\Phi(\omega)\bigr) \le D(\rho\,\|\,\omega),
  \qquad
  D(\rho\,\|\,\omega) := \mathrm{Tr}\bigl[\rho(\ln\rho - \ln\omega)\bigr],
  \label{eq:relative_entropy_dpi}
  \end{equation}
  providing the mathematically sound general monotonicity law for distinguishability and mismatch reduction (reducing to purity/entropy monotonicity relative to $\omega = I/d$ specifically when $\Phi$ is unital).
\end{enumerate}
\end{result}

\subsection{Open-System Lindblad Dynamics and Scope of Monotonicity}
\label{sec:CoherenceConservationArrow}
To state the thermodynamic dissipation bridge without repeating Version~79's non-unital error, we separate the operational proxy $R_{\mathrm{coh}}^{\mathrm{proxy}}$ from the classical EFT modulus $\Rcoh = |\psi|/\psi_0$ and track the exact hypotheses of Lindblad evolution.

\paragraph{Step 1: Operational Purity Proxy vs.\ EFT Field Amplitude.}
For a local finite-dimensional subsystem of dimension $d > 1$ described by density matrix $\rho_{\mathrm{sys}}(x)$, the normalized purity proxy is
\begin{equation}
    R_{\mathrm{coh}}^{\mathrm{proxy}}(x) = \frac{\mathrm{Tr}[\rho_{\mathrm{sys}}^2(x)] - \frac{1}{d}}{1 - \frac{1}{d}}.
\label{eq:Rcoh_purity_relation}
\end{equation}
As established in Section~\ref{sec:PrimitiveFieldDefinitions} and Paper~002 (Section~4.2), $R_{\mathrm{coh}}^{\mathrm{proxy}}$ is a well-defined quantum-information statistic; connecting it to the classical EFT field amplitude $\Rcoh = |\psi|/\psi_0$ requires an explicit state-to-field calibration kernel (\texttt{OPEN-6C}) rather than an automatic identity.

\paragraph{Step 2: Markovian Lindblad Master Equation.}
Open-system Markovian evolution is governed by the Gorini--Kossakowski--Sudarshan--Lindblad equation:
\begin{equation}
    \frac{d\rho_{\mathrm{sys}}}{dt} = -\frac{i}{\hbar}[H_{\mathrm{sys}}, \rho_{\mathrm{sys}}] + \sum_k \gamma_k \left( L_k \rho_{\mathrm{sys}} L_k^\dagger - \frac{1}{2}\{L_k^\dagger L_k, \rho_{\mathrm{sys}}\} \right),
    \label{eq:Lindblad_equation}
\end{equation}
with decay rates $\gamma_k \ge 0$.

\paragraph{Step 3: Unital Purity Monotonicity vs.\ General Relative-Entropy Contraction.}
Differentiating $P(t) = \mathrm{Tr}(\rho_{\mathrm{sys}}^2(t))$ along~\eqref{eq:Lindblad_equation} yields
\begin{equation}
    \frac{d}{dt} \mathrm{Tr}(\rho_{\mathrm{sys}}^2)
    = -\sum_k \gamma_k\,\mathrm{Tr}\!\left([L_k, \rho_{\mathrm{sys}}]^\dagger [L_k, \rho_{\mathrm{sys}}]\right)
      + \sum_k \gamma_k\,\mathrm{Tr}\!\left(\rho_{\mathrm{sys}}^2 [L_k, L_k^\dagger]\right).
    \label{eq:purity_decay}
\end{equation}
The first sum is manifestly non-positive; the second term vanishes identically for all states $\rho_{\mathrm{sys}}$ if and only if the generator is unital ($\sum_k \gamma_k [L_k, L_k^\dagger] = 0$). Consequently: (i)~under \emph{unital} Lindblad evolution, $\frac{d}{dt}\mathrm{Tr}(\rho_{\mathrm{sys}}^2) \le 0$ and $\frac{d}{dt}R_{\mathrm{coh}}^{\mathrm{proxy}} \le 0$ hold for every state, with equality whenever $[L_k, \rho_{\mathrm{sys}}] = 0$ for all active jump operators (which includes non-unitary dephasing of already-diagonal states); and (ii)~under \emph{non-unital} dissipative cooling or spontaneous emission (such as amplitude damping), purity can strictly increase, while the general monotonic law is relative-entropy contraction $D(\rho_{\mathrm{sys}}(t)\|\omega_{\mathrm{ss}}) \le D(\rho_{\mathrm{sys}}(0)\|\omega_{\mathrm{ss}})$ toward any stationary reference state $\omega_{\mathrm{ss}}$ (Eq.~\eqref{eq:relative_entropy_dpi}).

\subsection{The Geometric Arrow: Irreversibility from Informational Hysteresis}

Within the unital and relative-entropy domains specified above, open-system distinguishability contracts monotonically, while at the conceptual level SIT investigates informational hysteresis (Section~\ref{sec:memory_substrate}) as a complementary geometric picture of macroscopic irreversibility.

While the fundamental field equations of \texttt{M\_AUTH\_SIT\_EFT\_v1} are time-reversal symmetric, macroscopic field configurations and coarse-grained open-system records are generally not. Accumulated phase gradients and dissipated radiation into environmental degrees of freedom make exact time-reversed initial conditions practically inaccessible in macroscopic regimes.

This separates microscopic time-reversal symmetry of the classical action from effective macroscopic irreversibility. The statistical contraction of accessible distinguishability in open systems (Eq.~\eqref{eq:relative_entropy_dpi}) and the coarse-grained spreading of field excitations provide complementary descriptions of irreversible behavior without invoking fundamental microscopic time asymmetry.

\blockquote{
It's not the particle interactions themselves that are irreversible, but how they affect the larger informational field.
}

\subsection{Conceptual Entropy-Bookkeeping Analogy}

At the heuristic level (distinct from a proved theorem of \texttt{M\_AUTH\_SIT\_EFT\_v1}), one may compare a coarse-grained spatial functional of the coherence proxy
\[
S_{\mathrm{coh}} = -\int R_{\mathrm{coh}}^{\mathrm{proxy}}\ln R_{\mathrm{coh}}^{\mathrm{proxy}}\sqrt{-g}\,d^{3}x
\]
against the usual subsystem von~Neumann entropy
\[
S_{\mathrm{dec}} = -\mathrm{Tr}\,\rho\ln\rho.
\]
Version~81 explicitly notes that $S_{\mathrm{coh}} + S_{\mathrm{dec}}$ is \emph{not} a conserved quantity for arbitrary solutions of the field equations or general CPTP channels; rather, the pairing serves strictly as a conceptual bookkeeping analogy, while rigorous classical irreversibility in kinetic theory (such as the Deng--Hani--Ma hard-sphere Boltzmann limit) applies to its own specified microscopic model rather than proving monotonicity of $\Rcoh$.

\subsection{Classical Causality without Retrocausality}

Although the microscopic action is time-symmetric, linearized perturbations of \texttt{M\_AUTH\_SIT\_EFT\_v1} propagate causally inside the light cone ($c_{\mathrm{char}} = 1$). No signal propagates backward in time; apparent non-local quantum correlations reflect entanglement and measurement context rather than superluminal or retrocausal signaling.

\subsection{Long--Range Temporal Coherence (Conceptual Regime)}

Because the global $U(1)$ symmetry of \texttt{M\_AUTH\_SIT\_EFT\_v1} furnishes the conserved gauge-invariant coherence current $J^\mu_{\mathrm{coh}} = \kappa_c |\psi|^2 (\nabla^\mu\theta - \hat{q}A^\mu)$ (Eq.~\eqref{eq:abstract_coherence_current}), isolated coherent domains maintain phase order until external couplings exchange coherence charge or induce decoherence. Earlier Version~79 prose speculated on specific millihertz phononic-superlattice clock slow-downs and attosecond near-horizon frame-dragging amplitudes; Version~81 removes those uncalculated numerical scales because no action-to-apparatus derivation (\texttt{OPEN-6}) currently fixes such magnitudes from \texttt{M\_AUTH\_SIT\_EFT\_v1}.

\subsection{Thermodynamic Dissipation as a Computational Process}

\subsection{Mathematical Formulation: Local Signal-Dissipation Dynamics}

Consider a system of $N$ interacting agents, nodes, or oscillators, each with a local property $Q_i(t)$ at time $t$.
In physical applications, $Q_i(t)$ may represent:

\begin{itemize}
    \item Phase of an oscillator or wave,
    \item Energy or momentum of a particle,
    \item Local amplitude of quantum coherence.
\end{itemize}

Define the difference between two components:
\[
\Delta Q_{ij}(t) = Q_i(t) - Q_j(t)
\]

When two components interact, they partially dissipate their difference according to a local update rule:
\begin{align}
Q_i(t+\delta t) &= Q_i(t) - \frac{\gamma}{2} \Delta Q_{ij}(t) \\
Q_j(t+\delta t) &= Q_j(t) + \frac{\gamma}{2} \Delta Q_{ij}(t)
\label{eq:local-update}
\end{align}
where $0 < \gamma \leq 1$ is a dissipation (or coupling) parameter.
Over many iterations, such updates drive all $Q_i$ toward a common value---realizing equilibrium as a computational process.
\subsection{Mapping Coherence/Decoherence to Signal Dissipation}

The local dissipation rules serve as the microscopic engine for the macroscopic evolution of the theory's central fields:
\begin{itemize}
    \item The \textbf{coherence ratio} $R_{\mathrm{coh}}(x,t)$ evolves as local subsystems interact and their informational states dissipate differences.
\item The \textbf{time-density field} $\rho_t(x,t)$ is driven by the cumulative effect of these microscopic dissipation events, encoding the local density of temporally resolved events.
\end{itemize}

In SIT, $Q_i(t)$ can be interpreted as the local quantum coherence amplitude or phase at site $i$.
Coherence between nodes $i$ and $j$ is maximized when $Q_i = Q_j$; decoherence corresponds to persistent differences or fluctuations.
Thus, each local dissipation event can be viewed as a micro-level act of coherence reinforcement:
\begin{align}
\text{High coherence:}\quad & Q_i \approx Q_j \implies \Delta Q_{ij} \approx 0,\\
\text{Active dissipation:}\quad & |\Delta Q_{ij}| \text{ large} \implies \text{system far from equilibrium (low coherence)}.
\end{align}

The global, macroscopic field evolution---such as the time-density field $\rho_t(x)$---emerges from the statistical aggregation of many such micro-level updates.
When interpreted as a field, the signal dissipation process maps onto a diffusion-like equation:
\begin{equation}
\frac{\partial Q(x,t)}{\partial t} = D \nabla^2 Q(x,t),
\label{eq:dissipation}
\end{equation}
where $Q(x,t)$ is the spatially varying property (e.g., coherence amplitude), and $D$ is an effective diffusion constant derived from $\gamma$ and the interaction network.
\subsection{Macro-Level Evolution of the Time-Density Field}

The collective action of countless signal-dissipation steps at the micro level leads to the smooth evolution of the macro-level time-density field, $\rho_t(x)$, governed by its field equation (see Section~\ref{sec:core-action}):
\[
\Box \rho_t - V'(\rho_t) + \cdots = 0
\]
Here, the ``$\cdots$'' term absorbs all source and coupling terms from matter, coherence, and interactions.
The approach to equilibrium (or attractor states) for $\rho_t$ can be understood as the integrated result of local coherence-dissipation events at every scale.
\subsection{Information-Theoretic Interpretation: Entropy and Coherence}

The signal-dissipation process can be interpreted as local reduction of informational (Shannon or quantum) entropy.
As local differences in $Q_i$ vanish, the system moves toward maximal mutual coherence, corresponding to minimal entropy production.
Conversely, persistent incoherence (large $\Delta Q_{ij}$) maintains or increases entropy.
This micro-to-macro bridge justifies the treatment of coherence conservation and time-density evolution as physical processes with direct thermodynamic and computational meaning.
\subsection{Example: Coherence Relaxation in a Network}

Consider a ring of $N$ oscillators with initial random phases $Q_i(0)$.
At each time step, each oscillator exchanges phase information with its nearest neighbors according to the local update rule (Eq.~\ref{eq:local-update}).
Over many iterations, all $Q_i$ converge to a common phase, maximizing global coherence:
\[
\lim_{t \to \infty} Q_i(t) = Q_{\mathrm{eq}} \quad \forall i
\]
The rate of convergence (relaxation time) depends on $\gamma$ and network topology.
In physical terms, this process models synchronization in neural circuits, quantum decoherence in open systems, or even heat conduction.
This mechanistic foundation ensures that the evolution of $R_{\mathrm{coh}}$ and $\rho_t$ is not left to abstract principle, but arises from a well-defined, physically computable process:
\begin{itemize}
    \item \emph{Coherence increases} as local differences dissipate, corresponding to the growth of order and synchrony in the system.
\item \emph{Decoherence or disorder} results from the persistence or amplification of local differences, inhibiting equilibration.
\item The \emph{time-density field} $\rho_t(x,t)$ evolves as the statistical outcome of these local interactions, linking the micro-dynamics of dissipation to macroscopic gravitational and temporal effects.
\end{itemize}

\subsection{Summary}

By formalizing thermodynamic dissipation as a wave-based, local computational process, SIT unifies the dynamics of coherence, entropy, and time-density into a single theoretical framework.
The macroscopic field evolution described in Section~\ref{sec:core-action} is thus revealed as the statistical outcome of myriad microscopic signal-dissipation events---a mechanistic bridge between information, computation, and fundamental physics.

\subsection{The Coherence Field as a Conceptual Memory Substrate and Informational Hysteresis}
\label{sec:hysteresis_memory}
\label{sec:memory_substrate}

While Result~\ref{thm:CoherenceMonotonicity} establishes strictly \emph{bounded} open-system monotonicity (purity non-increase for unital CPTP maps and quantum relative-entropy contraction $D(\Phi(\rho)\|\Phi(\omega)) \le D(\rho\|\omega)$ for general CPTP maps), SIT additionally explores a conceptual geometric picture termed \textbf{informational hysteresis}. In this phenomenological interpretation, after a localized subsystem exchanges energy and phase correlations with an unobserved environment, the coarse-grained reduced state does not retrace its initial configuration identically.

During a cyclic control trajectory $p$ in a reduced dissipative state space, an effective path-dependent geometric phase offset may be parameterized schematically as
\begin{equation}
\Phi_{\mathrm{geom}}[p] \approx \oint_p C_{\mathrm{rec}}^+[p]\,dt
\label{eq:informational_hysteresis_loop}
\end{equation}
(Appendices~\ref{app:recursion-influence}--\ref{app:prior_work}). Version~81 emphasizes that this hysteresis picture is a \emph{conceptual open-system model}, distinct from both a universal decay law for global $\Rcoh$ and any claimed conservation of total informational entropy (retracted above).

\subsection{Scope of Open-System Monotonicity and Holonomy}

In summary, Result~\ref{thm:CoherenceMonotonicity} replaces earlier universal global-coherence decay and global entropy-conservation assertions with established, bounded open-system quantum theorems (unital CPTP purity non-increase and general CPTP relative-entropy contraction). Parallel transport of a charged field around closed loops yields standard $U(1)$ gauge holonomy $\exp\bigl(i\hat{q}\oint_\gamma A_\mu dx^\mu\bigr)$ under the independent connection $A_\mu$ of \texttt{M\_AUTH\_SIT\_EFT\_v1}, while deriving any explicit thermodynamic arrow for the classical scalar fields $(\rho_t, \psi)$ from microscopic coarse-graining remains an open problem distinct from the Deng--Hani--Ma hard-sphere Boltzmann theorem.

\section{Wave--Particle Duality and Coherence--Decoherence Phenomenology}
\label{sec:duality}

At the conceptual level, wave--particle complementarity is closely related to environment-induced decoherence: high interferometric fringe visibility produces interference patterns, whereas environmental entanglement suppresses off-diagonal density-matrix elements and recovers classical trajectory statistics.

In semiclassical quantization, the Bohr--Sommerfeld / Einstein--Brillouin--Keller condition
\begin{equation}\label{eq:EBK}
   \oint p_i\,dq^i \;=\; 2\pi n\hbar,
   \qquad n\in\mathbb{Z},
\end{equation}
selects constructive single-valued phase quantization. Following the operational separation contract (Section~\ref{sec:PrimitiveFieldDefinitions}), Version~81 distinguishes standard quantum interferometric visibility (an operational density-matrix statistic) from the classical $U(1)$ Noether current $J^\mu_{\mathrm{coh}} = \kappa_c|\psi|^2(\nabla^\mu\theta - \hat{q}A^\mu)$ of \texttt{M\_AUTH\_SIT\_EFT\_v1}: the Noether current does \emph{not} automatically equal fringe visibility, nor does unchanged \texttt{M\_AUTH\_SIT\_EFT\_v1} derive quantitative crossover curves for Ramsey phase-echo, Mach--Zehnder, or atom-fountain experiments prior to specifying the microscopic detector and preparation couplings (\texttt{OPEN-6C}). Those interferometric platforms remain prospective conceptual analogies rather than completed SIT predictions.

\section{Measurement-Induced Coherence Gradients in the Two-Slit Geometry}
\label{sec:two-slit}

In a standard two-slit experiment interference arises when the
coherence ratio \(R_{\text{coh}}\) is uniform across both paths. As a
prospective conceptual hypothesis, Super Information Theory suggests
that a monitoring device might perturb this balance by inducing a
local coherence perturbation \(\Delta R_{\text{coh}}\) and, through
the canonical scalar coupling of \texttt{M\_AUTH\_SIT\_EFT\_v1}, a
corresponding temporal-density shift \(\delta\rho_t\). In the
canonical action, phase evolution couples to the independent gauge
connection \(A_\mu\) via the covariant derivative
\(D_\mu\psi = (\nabla_\mu - i\hat{q}A_\mu)\psi\) rather than through a
flat gradient ansatz \(A_i = (\hbar/e)\partial_i\arg\psi\) (whose
local curl vanishes identically on smooth patches). Whether any
monitoring apparatus produces a measurable fringe displacement remains
an open microphysical problem contingent on deriving the detector
source functional \(J_{\mathrm{prep}}\) and the electromagnetic
coupling dictionary (\texttt{OPEN-6B}/\texttt{OPEN-6C}); any such
effect would be strongly suppressed by existing post-Newtonian and
short-range bounds on SIT's scalar couplings and is not yet a
quantitatively derived prediction of the canonical action.

A prospective experimental test would replace a projective
``which-way'' detector with a tunable low-noise weak phase probe once
the detector-to-field transfer function is derived under
\texttt{OPEN-6C}, testing for any continuous shift of the interference
envelope above ordinary environmental decoherence.

\subsection{Broader Consequences}

Conceptually, because the covariant coherence-field dynamics links
microscopic phase evolution and macroscopic spacetime response in the
same effective action, weak-measurement interferometry offers a
prospective tabletop testbed alongside torsion-balance and clock
experiments. Once an explicit detector source and fringe-shift formula
are derived from \texttt{M\_AUTH\_SIT\_EFT\_v1} (\texttt{OPEN-6C}),
a non-observation would constrain the specific coupling parameter
region and preparation model tested, in accordance with the Scoped
Falsification Principle of Section~\ref{sec:experimental-predictions}.

\subsection{Neural and Cosmological Echoes}

The regulated horizon that emerges around an over-coupled detector
is mathematically identical to the coherence boundaries that
organise cortical phase waves and to the curvature cut-offs that
stabilise collapsing stars.  In each case a conserved Noether
current shunts excess coherence into decoherent radiation,
enforcing a finite amplitude--frequency product and preventing
singularity formation.  The laboratory experiment therefore
connects three scales---quantum, neural, and cosmic---through a single
gauge-holonomy principle already embedded in the master action.



\section{Stability, Causality, Ghost Freedom, and Mathematical Consistency}
\label{sec:stability-causality}

We now consolidate the classical and linearized quantum stability proofs for the authoritative $(-,+,+,+)$ effective field theory \texttt{M\_AUTH\_SIT\_EFT\_v1}, backed by machine-checked Lean~4.30 theorems (\texttt{SitQg.Stage2.QG11To14}, \texttt{SitQg.Stage2.QG15To19}, \texttt{SITQG09\allowbreak{}Linearized\allowbreak{}Fluctuation\allowbreak{}And\allowbreak{}Obstruction}).

\subsection{Stability, Absence of Ghosts, and Positive-Definite Radial Hessian}

\subsubsection{Quadratic Expansion in Mostly-Plus Signature $(-,+,+,+)$}

Expanding the bosonic fields of $S_{\mathrm{AUTH}}$ around the zero-cosmological-constant broken-symmetry vacuum
\[
g_{\mu\nu} = \eta_{\mu\nu} + h_{\mu\nu}, \qquad
\rho_t(x) = \rho_0 + \delta\rho_t(x), \qquad
|\psi(x)| = v_0 + \delta|\psi|(x), \qquad
A_\mu(x) = \delta A_\mu(x),
\]
with $\eta_{\mu\nu} = \mathrm{diag}(-1,+1,+1,+1)$, the scalar kinetic terms $-\frac{\kappa_t}{2}\eta^{\mu\nu}\partial_\mu\delta\rho_t\partial_\nu\delta\rho_t - \frac{\kappa_c}{2}\eta^{\mu\nu}\partial_\mu\delta|\psi|\partial_\nu\delta|\psi|$ evaluate explicitly to
\begin{equation}
\mathcal{L}_{\mathrm{scal}}^{(2)}
= +\frac{\kappa_t}{2}\Bigl[(\partial_t\delta\rho_t)^2 - (\nabla\delta\rho_t)^2\Bigr]
  + \frac{\kappa_c}{2}\Bigl[(\partial_t\delta|\psi|)^2 - (\nabla\delta|\psi|)^2\Bigr]
  - \frac{1}{2}
  \begin{pmatrix}\delta\rho_t & \delta|\psi|\end{pmatrix}
  H
  \begin{pmatrix}\delta\rho_t \\ \delta|\psi|\end{pmatrix},
\label{eq:quadratic_scalar_lagrangian_sec19}
\end{equation}
where $H$ is the $2\times 2$ radial Hessian matrix~\eqref{eq:radial_Hessian_sec5}. In unitary gauge ($\theta = 0$ absorbed into longitudinal $\delta A_\mu$ for $v_0 > 0$), the gauge-phase sector yields the canonical Proca Lagrangian
\begin{equation}
\mathcal{L}_{\mathrm{vec}}^{(2)}
= -\frac{f_{\mathrm{EM}}(R_0)}{4}\,\delta F_{\mu\nu}\delta F^{\mu\nu}
  - \frac{1}{2}\kappa_c\hat{q}^2 v_0^2\,\eta^{\mu\nu}\delta A_\mu\delta A_\nu,
\qquad
m_A^2 = \frac{\hat{q}^2\kappa_c v_0^2}{f_{\mathrm{EM}}(R_0)} > 0,
\label{eq:proca_mass_sec19}
\end{equation}
reducing smoothly at $v_0 = 0$ ($f_{\mathrm{EM}}(0)=1$) to two massless Maxwell polarizations plus one massless Goldstone phase mode. Simultaneously, ADM scalar--vector--tensor decomposition of the linearized Einstein--Hilbert sector yields the transverse-traceless graviton Lagrangian
\begin{equation}
\mathcal{L}_{\mathrm{TT}}^{(2)}
= \frac{1}{8\kappa_G}\Bigl[(\partial_t h_{ij}^{\mathrm{TT}})^2 - (\partial_k h_{ij}^{\mathrm{TT}})^2\Bigr],
\qquad
c_T^2 = 1,
\label{eq:graviton_TT_sec19}
\end{equation}
while the elliptic ADM Hamiltonian constraint $\Delta\Psi = +\frac{\kappa_G}{2}T_{00}^{(2)}$ eliminates the non-propagating conformal trace $\Psi$ and the Gauss constraint eliminates $\delta A_0$.

\subsubsection{Ghost Freedom, Tachyonic Stability, and Exact Polarizations Count}

Under the canonical positivity domain ($\kappa_G, \kappa_t, \kappa_c, \lambda_\rho, \lambda_\psi, m_{\mathrm{link}}, \psi_0, \rho_0, M_{\mathrm{sym}} > 0$, $c_{\mathrm{EM}}, \lambda_{\mathrm{portal}}, \hat{q}, v_0 \ge 0$):
\begin{itemize}
  \item \textbf{Exact Physical Mode Count ($7$ Pure-SIT / $8$ Matter-Coupled):} In pure-SIT bosonic vacuum ($\mathcal{L}_{\mathrm{matter}} = 0$), exactly \textbf{$7$ physical bosonic polarizations} ($14$ phase-space degrees of freedom) propagate: $2$ TT gravitons $(h_+, h_\times)$, $2$ transverse vector modes, $1$ longitudinal Proca mode ($v_0 > 0$) or phase scalar ($v_0 = 0$), and $2$ massive real scalar modes from $(\delta\rho_t, \delta|\psi|)$. When coupled to a representative neutral screened scalar matter field $\chi$ in curved $3+1$D ADM canonical gravity (\texttt{Stage3ADM}, Section~\ref{sec:stage3-curved-adm}), Dirac--Bergmann constraint reduction yields $36 - 2\times 10 = 16$ physical phase-space dimensions, i.e.\ \textbf{$8$ propagating bosonic modes} ($7 + 1$).
  \item \textbf{Strict Absence of Ghosts:} Every propagating kinetic coefficient ($\frac{1}{8\kappa_G} > 0$, $\frac{\kappa_t}{2} > 0$, $\frac{\kappa_c}{2} > 0$, $\frac{f_{\mathrm{EM}}(R_0)}{2} \ge \frac{1}{2} > 0$) is strictly positive.
  \item \textbf{Strict Absence of Tachyons:} For $v_0 > 0$ and $K = \mathrm{diag}(\kappa_t, \kappa_c) \succ 0$, $\mathrm{tr}(H) > 0$ and $\det(H) = 2v_0^2(2\lambda_\rho\lambda_\psi\rho_0^2 + \lambda_\psi m_{\mathrm{link}}^2 + 4\lambda_\rho\rho_0^2\lambda_{\mathrm{portal}}) > 0$ guarantee strictly positive physical squared eigenmasses $0 < \frac{\det(H)}{\max(\kappa_t,\kappa_c)\,\mathrm{tr}(H)} \le m_1^2 \le m_2^2$ of $M_{\mathrm{phys}}^2 = K^{-1/2} H K^{-1/2}$ ($H u_a = m_a^2 K u_a$; at unit kinetic normalization $K = I_2$ in \texttt{BM\_\allowbreak{}QG16\_\allowbreak{}UNIT\_\allowbreak{}MASS\_\allowbreak{}EIGENVALUES}, $H = \bigl(\begin{smallmatrix}2 & -2\\ -2 & 6\end{smallmatrix}\bigr)$ gives $m_\mp^2 = 4 \mp 2\sqrt{2} > 0$).
\end{itemize}

\subsection{Perturbative EFT Renormalizability and Weinberg Derivative Power Counting}

\subsubsection{Renormalizable Scalar-Gauge Core}

In flat spacetime ($G \to 0$) with dimension-6 dressings switched off ($c_{\mathrm{EM}} \to 0$, $M_{\mathrm{sym}} \to \infty$), $\mathcal{L}_{\mathrm{coh}} - \frac{1}{4}F_{\mu\nu}F^{\mu\nu}$ is a standard dimension-$\le 4$ power-counting renormalizable Abelian Higgs + portal singlet QFT in $3+1$ dimensions.

\subsubsection{Weinberg Loop Suppression Below $\Lambda_{\mathrm{UV}}$}

Including Einstein--Hilbert gravity ($\kappa_G = M_{\mathrm{Pl}}^{-2}$) and the dimension-6 operators suppressed by $\Lambda_{\mathrm{EM}}^{-2}$ and $\Lambda_{\mathrm{sym}}^{-2}$, any connected $L$-loop Feynman diagram with $V$ vertices, $I$ internal lines, and $E_{\mathrm{ext}}$ external legs satisfies Euler's topological identity $I = L + V - 1$, yielding Weinberg's derivative power-counting exponent
\begin{equation}
\sum_{v=1}^V (n_v - 2) = E_{\mathrm{ext}} + 2L - 2 \quad (\text{grav.}),
\qquad
\mathcal{M}_L(E) \;\lesssim\; \left(\frac{E}{\Lambda_{\mathrm{UV}}}\right)^{2L + 2}
\quad (2 \to 2,\ E \ll \Lambda_{\mathrm{UV}}),
\label{eq:weinberg_power_counting_sec19}
\end{equation}
proved monotonically decreasing in $L$ for $E < \Lambda_{\mathrm{UV}}$ in Lean~4.30 (\texttt{weinberg\allowbreak{}Loop\allowbreak{}Power\_\allowbreak{}monotone}; e.g.\ at $E/\Lambda_{\mathrm{UV}} = 10^{-2}$, suppression factors are $10^{-4}, 10^{-8}, 10^{-12}$ at $L = 0, 1, 2$).

\subsection{Energy Conditions in Signature $(-,+,+,+)$}

For a canonical real scalar sector with $T_{\mu\nu}^{(\phi)} = \kappa_\phi\bigl(\partial_\mu\phi\partial_\nu\phi - \frac{1}{2}g_{\mu\nu}g^{\alpha\beta}\partial_\alpha\phi\partial_\beta\phi\bigr) - g_{\mu\nu}V(\phi)$ ($\kappa_\phi > 0$, $V(\phi) \ge 0$):
\begin{itemize}
  \item \textbf{Null Energy Condition (NEC):} For every null vector $n^\mu$ ($g_{\mu\nu}n^\mu n^\nu = 0$), $T_{\mu\nu}^{(\phi)}n^\mu n^\nu = \kappa_\phi(n^\mu\partial_\mu\phi)^2 \ge 0$, and $T_{\mu\nu}^{(\mathrm{EM})}n^\mu n^\nu = f_{\mathrm{EM}}(\Rcoh)(n^\mu F_{\mu\lambda})(n_\nu F^{\nu\lambda}) \ge 0$ since $f_{\mathrm{EM}} \ge 1$. Thus the NEC holds identically across the entire pure-SIT bosonic core.
  \item \textbf{Weak Energy Condition (WEC):} For every unit timelike vector $u^\mu$ ($g_{\mu\nu}u^\mu u^\nu = -1$), choosing a local Lorentz rest frame gives $T_{\mu\nu}^{(\mathrm{SIT})}u^\mu u^\nu = \frac{\kappa_t}{2}[(\partial_0\rho_t)^2 + (\nabla\rho_t)^2] + \frac{\kappa_c}{2}[|D_0\psi|^2 + |\mathbf{D}\psi|^2] + \frac{f_{\mathrm{EM}}}{2}(\mathbf{E}^2 + \mathbf{B}^2) + V_{\mathrm{tot}}(\rho_t,|\psi|) \ge 0$, so the WEC holds globally by $V_{\mathrm{tot}} \ge 0$.
  \item \textbf{Strong Energy Condition (SEC):} As in standard inflationary and dark-energy scalar theories, $(T_{\mu\nu} - \frac{1}{2}g_{\mu\nu}T)u^\mu u^\nu$ can become negative in potential-dominated regimes ($V_{\mathrm{tot}} > \text{kinetic}$), permitting late-time cosmic acceleration without ghost or NEC pathology.
\end{itemize}

\subsection{Causality and Strict Hyperbolicity}

All Euler--Lagrange equations~\eqref{eq:summary_einstein_sec8}--\eqref{eq:summary_maxwell_sec8} (in harmonic/Lorenz gauge) share the second-order principal symbol $\sigma_2(\xi) = -g^{\mu\nu}\xi_\mu\xi_\nu\,I$. Because $f_{\mathrm{EM}}(\Rcoh)$ multiplies $F_{\mu\nu}F^{\mu\nu}$ isotropically without derivative couplings of the form $\partial_\mu\rho_t\partial_\nu\rho_t F^{\mu\alpha}F^\nu_{\ \alpha}$, all characteristic cones coincide exactly with the physical metric null cone $g^{\mu\nu}\xi_\mu\xi_\nu = 0$ ($c_{\mathrm{char}} = c_T = 1$), strictly forbidding superluminal signal propagation.

\subsection{Local Cauchy Problem and Well-Posedness}

On any globally hyperbolic spacetime $(\mathcal{M}^4, g_{\mu\nu})$ with smooth initial data on a spacelike Cauchy slice $\Sigma_0$ satisfying the elliptic ADM Hamiltonian, momentum, and $U(1)$ Gauss constraints, Leray--Choquet-Bruhat--Yvonne hyperbolic PDE theory guarantees local-in-time existence, uniqueness (up to diffeomorphism and $U(1)$ gauge), and causal domain-of-dependence stability; full nonlinear curved $3+1$D constraint preservation under evolution is established at the local formal 3-jet level modulo explicit spatial divergences $\partial_k B^k$ on the physical Sylvester Riemannian positive-definite cone \texttt{IsSPD3} via Dirac--Bergmann Poisson closure in Section~\ref{sec:stage3-curved-adm} (while infinite-dimensional Sobolev/Banach constraint domains with Regge--Teitelboim asymptotic boundary conditions and global infinite-time multi-chart Cauchy PDE regularity on generic curved data remain noted under \texttt{OPEN-2B}/\texttt{OPEN-2}).

\subsection{Summary of Mathematical Rigor}

\texttt{M\_AUTH\_SIT\_EFT\_v1} is free of ghosts and tachyons, satisfies the NEC and WEC globally, propagates all bosonic modes causally on the physical light cone, closes its nonlinear $3+1$D ADM constraints first-class, and admits controlled Weinberg loop power counting below $\Lambda_{\mathrm{UV}}$.

\section{Reductions and Boundaries with Established Theories}
\label{app:reductions}

This section summarizes the exact algebraic reductions, unbroken-phase limits, heavy-scalar approximations, and explicit scope boundaries relating \texttt{M\_AUTH\_SIT\_EFT\_v1} to General Relativity, standard Quantum Mechanics, and classical kinetic theory.

\subsection{Metric Variation}

Starting from the SIT action (Equation~\eqref{eq:SIT_action}), the gravitational field equations arise from variation with respect to the metric \(g_{\mu\nu}\):
\[
\delta S_{\mathrm{SIT}}/\delta g^{\mu\nu} = 0.
\]

\subsection{Scope of Weak-Field Post-Newtonian Bounds}

Earlier drafts quoted informal estimates such as $|\alpha\,\delta\rho_t| \lesssim 10^{-5}$ and $|\alpha| \lesssim 10^{-2}$ from Lunar Laser Ranging, Cassini Doppler tracking, and binary-pulsar timing. Consistent with Stage~6 (Section~\ref{sec:stage6-action-diagnostics}) and the blocked six-record primary provenance ledger (\texttt{OPEN-6D}), Version~81 retires those informal estimates: constraining the parameters of \texttt{M\_AUTH\_SIT\_EFT\_v1} against Solar System, torsion-balance, and clock observations requires deriving the full coupled gravitational and material-charge responses from the stated action first.

\subsection{Algebraic Metric Reduction vs.\ Coupled Solution Conditions}

In canonical \texttt{M\_AUTH\_SIT\_EFT\_v1} (which couples $\rho_t$ minimally to $g_{\mu\nu}$ rather than as a non-minimal Brans--Dicke scalar and does not rescale $\hbar$), substituting constant vacuum fields $(\rho_t, |\psi|) = (\rho_0, v_0)$ with pure-gauge phase $D_\mu\psi = 0$ ($V_{\mathrm{tot}} = 0$, $\nabla V_{\mathrm{tot}} = 0$) reduces the \emph{metric} equation algebraically to
\[
G_{\mu\nu} = 8\pi G_{\mathrm{eff}}\,T_{\mu\nu}^{\mathrm{SM}} + 8\pi G\,f_{\mathrm{EM}}(R_0)\,T_{\mu\nu}^{\mathrm{EM},0},
\qquad
G_{\mathrm{eff}} = G\,f_{\mathrm{SM}}(R_0).
\]
This algebraic metric reduction constitutes an \emph{exact solution of the full coupled theory} only when the scalar, gauge, and matter equations are simultaneously satisfied: at $v_0 > 0$, neutral dust ($\mathcal{L}_{\mathrm{matter}} = -\rho_m < 0$) leaves a nonzero scalar-equation residual $-v_0\rho_m/(\psi_0^2 M_{\mathrm{sym}}^2) \ne 0$, and $[D_\mu, D_\nu]\psi = -i\hat{q}F_{\mu\nu}\psi = 0$ forces $F_{\mu\nu} = 0$ whenever $\hat{q}v_0 \ne 0$ (Appendix~D.2). Source-free vacuum ($\mathcal{L}_{\mathrm{matter}} = 0, F_{\mu\nu} = 0$) is an exact solution at $(\rho_0, v_0)$; broader exact sourced GR-plus-Maxwell reduction ($G_{\mathrm{eff}} = G$, $f_{\mathrm{EM}} = 1$, $m_A = 0$) holds in the separate unbroken parameter limit $v_0 = 0$; and heavy-scalar decoupling ($m_1, m_2 \gg E_{\mathrm{obs}}$, where $m_1^2, m_2^2$ solve $H u_a = m_a^2 K u_a$) provides a controlled asymptotic approximation.

\subsection{Standard Quantum Evolution Boundary (Not Derived from Classical EFT)}

Consistent with the operational separation contract (Section~\ref{sec:PrimitiveFieldDefinitions}), \texttt{M\_AUTH\_SIT\_EFT\_v1} is a classical effective field theory for $(\rho_t, \psi, A_\mu, g_{\mu\nu})$ and does \emph{not} derive the von Neumann equation $i\hbar\,\partial_t\rho = [\hat{H}, \rho]$ or the single-particle Schr\"odinger equation $i\hbar\,\partial_t\psi_{\mathrm{QM}} = \hat{H}\psi_{\mathrm{QM}}$ from classical scalar fields. Standard Hilbert-space quantum mechanics applies untouched in the tested weak-field constant-background limit, while any quantum-matter back-reaction onto $(\rho_t, \psi)$ remains open under \texttt{OPEN-6C}.

\subsection{External Kinetic-Theory Correspondence (Not Derived from $S_{\mathrm{AUTH}}$)}

In the decoupled constant-background limit where $(\rho_t, \psi)$ perturbations are negligible, dilute classical gases obey the standard Boltzmann kinetic equation
\[
\frac{\partial f}{\partial t} + \mathbf{v} \cdot \nabla_x f + \mathbf{F} \cdot \nabla_p f = C[f]
\]
under the usual microscopic collisional assumptions of statistical mechanics (for instance, the hard-sphere Boltzmann--Grad theorem of Deng--Hani--Ma~2025). Consistent with Appendix~\ref{app:KineticLimit}, \texttt{M\_AUTH\_SIT\_EFT\_v1} does \emph{not} derive the collision operator $C[f]$ or Navier--Stokes closure from the classical SIT scalars $(\rho_t, \psi)$; any SIT-specific transport correction remains an open research problem.

\section{Informational Symmetry and CPT Balance}
\label{sec:CPT-balance}

Super Information Theory seeks CPT balance without appealing to an external ``mirror universe.''  
The master action \texttt{M\_AUTH\_SIT\_EFT\_v1} already furnishes the conserved gauge-invariant Noether current
$J^\mu_{\mathrm{coh}} = \kappa_c|\psi|^2(\nabla^\mu\theta - \hat{q}A^\mu)$ (Eq.~\eqref{eq:abstract_coherence_current}),
whose covariant divergence vanishes on shell and enforces global phase-charge conservation.  
Coupled to the independent gauge connection $A_\mu$ through $D_\mu\psi = (\nabla_\mu - i\hat{q}A_\mu)\psi$ (rather than a pure-gradient flat connection, whose local field strength vanishes identically), local changes in coherence charge are balanced by current flux within a single spacetime manifold.  
Apparent macroscopic time-reversal asymmetries are therefore attributed to coarse graining and boundary conditions,
not to a second time-reversed cosmos.  

For completeness, the earlier "anti-universe" speculation is now treated as a historical note in Appendix~\ref{app:speculative}.

On a cosmological scale, the exploratory ``Halfway Universe'' concept is discussed as an open Level-1 phenomenological hypothesis (distinct from a proved symmetry identity) in Section~\ref{sec:blackholesvoids}.

\subsection{Coherence Saturation and Exploratory Horizon Analogies}
\label{sec:info-horizons}

Earlier drafts introduced the heuristic term ``informational horizon'' to imagine a boundary where $\Rcoh$ approaches a saturation value and blocks classical curvature singularities or accounts for quantum-eraser interference. Version~81 explicitly scopes those statements as speculative conceptual analogies rather than proved consequences of \texttt{M\_AUTH\_SIT\_EFT\_v1}. In the canonical action, $\psi$ is an unbounded complex scalar regulated dynamically by its quartic potential $U(|\psi|) = \frac{\lambda_\psi}{4}(|\psi|^2 - v_0^2)^2$ (rather than a hard kinematic ceiling $\Rcoh \le 1$), no black-hole singularity-resolution theorem has been proved inside the low-energy EFT regime below $\Lambda_{\mathrm{UV}}$ (\texttt{OPEN-5}), and quantum-eraser statistics are already fully described by standard Hilbert-space quantum mechanics without invoking a spacetime horizon.

\subsection{Gravitational Stability and Ultraviolet Cutoff Scope}
\label{sec:grav-stability}

Within the classical low-energy regime of \texttt{M\_AUTH\_SIT\_EFT\_v1}, local well-posedness and energy positivity follow from the ghost-free kinetic signature ($\kappa_t > 0, \kappa_c > 0$) and the nonnegative $Z_2$-lifted scalar potential $V_{\mathrm{tot}} \ge 0$ (Sections~\ref{sec:framework}--\ref{sec:screening-mechanism}). Whether an ultraviolet completion above $\Lambda_{\mathrm{UV}} = \min(M_{\mathrm{Pl}}, \Lambda_{\mathrm{EM}}, \Lambda_{\mathrm{sym}})$ resolves gravitational collapse singularities in supernovae or early-universe cosmology remains an open quantum-gravity problem (\texttt{OPEN-5}) rather than a derived theorem of the classical EFT action.

\subsection{Summary of Bounded Empirical Benchmarks}

Torsion-balance short-range tests (Lee~2020), the kinetic Boltzmann--Grad theorem of Deng--Hani--Ma~(2025) in its hard-sphere regime, and standard electromagnetic Aharonov--Bohm holonomy occupy distinct scientific roles: Deng--Hani--Ma provides an external kinetic-theory benchmark (not a constraint on SIT scalar couplings, Section~\ref{sec:coherence-conservation}), Aharonov--Bohm phase measures gauge curvature of an independent $U(1)$ connection $A_\mu$ (Section~\ref{sec:magnetism-holonomy}), and short-range gravitational bounds constrain \texttt{M\_AUTH\_SIT\_EFT\_v1} only through the full coupled screening and finite-geometry likelihood program (\texttt{OPEN-6A}, Section~\ref{sec:experimental-predictions}). None of these benchmarks renders singularity resolution or quantum measurement foundations solved within the current monograph.



\section{Renormalisation-Group Flow, Weinberg Power Counting, and Seeley--DeWitt Heat-Kernel Structure}
\label{sec:rg-heat-kernel}

For SIT's laboratory and astrophysical phenomenology to remain radiatively consistent across scales $E \ll \Lambda_{\mathrm{UV}}$, its couplings must evolve benignly under renormalization-group (RG) flow while its one-loop gravitational and dimension-6 divergences follow standard Wilsonian effective field theory.

\subsection{Sub-Planckian Matter-Scalar Toy Setup and One-Loop Beta Functions}

In a minimal sub-Planckian renormalizable subsector ($\mu \ll \Lambda_{\mathrm{UV}}$) coupling a real time-density fluctuation $\delta\rho_t$ with quartic $\frac{\lambda_\rho}{4!}(\delta\rho_t)^4$, complex coherence scalar $\psi$ with $\frac{\lambda_c}{4!}|\psi|^4$ and portal $\frac{\lambda_{c\rho}}{4}(\delta\rho_t)^2|\psi|^2$, a benchmark Dirac fermion with Yukawa $\alpha_Y\delta\rho_t\bar{\Psi}\Psi$, and an Abelian gauge field $A_\mu$ with coupling $e$, dimensional regularization in $d = 4 - \epsilon$ ($\overline{\mathrm{MS}}$) gives the standard one-loop beta functions $\mu\,dX/d\mu = \beta_X$:
\begin{align}
\beta_{\lambda_\rho} &= \frac{3\lambda_\rho^2 - 24\alpha_Y^4 + \lambda_{c\rho}^2}{16\pi^2}, \label{eq:beta_lambda_rho_sec22}\\
\beta_{\alpha_Y} &= \frac{5\alpha_Y^3 - 6\alpha_Y e^2}{16\pi^2}, \label{eq:beta_alpha_sec22}
\end{align}
remaining perturbative with zero Landau pole across terrestrial, collider, and cosmological scales whenever initial couplings satisfy $|\alpha_Y|, \lambda_\rho, \lambda_{c\rho} \lesssim 10^{-1}$. Meanwhile, in canonical \texttt{M\_AUTH\_SIT\_EFT\_v1}, direct linear couplings to $F_{\mu\nu}F^{\mu\nu}$ and $\mathcal{L}_{\mathrm{matter}}$ are forbidden by $\mathbb{Z}_2$ invariance of $|\psi|^2$, entering only at dimension-6 via $\Lambda_{\mathrm{EM}}^{-2}|\psi|^2 F^2$ and $\Lambda_{\mathrm{sym}}^{-2}|\psi|^2\mathcal{L}_{\mathrm{matter}}$ and enjoying automatic Planck/GUT-scale suppression~\eqref{eq:cutoff_ratio_identity}.

\subsection{Covariant One-Loop Seeley--DeWitt Heat-Kernel Expansion}

Integrating out quadratic bosonic fluctuations around a smooth background $(g_{\mu\nu}, \rho_t, \psi, A_\mu)$ produces a Laplace-type fluctuation operator $\Delta = -g^{\mu\nu}\nabla_\mu\nabla_\nu - E$ with bundle curvature $\Omega_{\mu\nu} = [\nabla_\mu, \nabla_\nu]$. In proper-time heat-kernel regularization with UV cutoff $\Lambda_{\mathrm{reg}} \le \Lambda_{\mathrm{UV}}$,
\begin{equation}
\mathrm{Tr}\,e^{-s\Delta}
\;\sim\;
\frac{1}{(4\pi s)^2}\sum_{k=0}^\infty s^k \int d^4x\,\sqrt{-g}\,\mathrm{tr}\,a_k(x)
\qquad (s \to 0^+),
\label{eq:heat_kernel_asymptotic}
\end{equation}
the first three Gilkey--Seeley--DeWitt invariants govern all one-loop UV divergences in four dimensions:
\begin{enumerate}[leftmargin=2em]
  \item \textbf{Quartic Vacuum Divergence ($a_0$):}
  \begin{equation}
  a_0(x) = I,
  \end{equation}
  generating an $O(\Lambda_{\mathrm{reg}}^4 / (16\pi^2))$ shift absorbed into the bare vacuum energy counterterm $\delta\Lambda_{\mathrm{vac}}$, preserving the physical renormalization condition $V_{\mathrm{tot}}(\rho_0, v_0) = 0$.
  \item \textbf{Quadratic Mass and Einstein--Hilbert Divergence ($a_1$, \texttt{ODQ-5}):}
  \begin{equation}
  a_1(x) = \frac{1}{6}\,R\,I + E,
  \end{equation}
  generating $O(\Lambda_{\mathrm{reg}}^2/(16\pi^2))$ renormalizations of the scalar mass parameters and shifting the bare inverse Newton constant $1/(16\pi G_{\mathrm{bare}})$. Crucially (\texttt{ODQ-5}), \texttt{M\_AUTH\_SIT\_EFT\_v1} retains $\frac{R}{16\pi G} = \frac{M_{\mathrm{Pl}}^2}{2}R$ as a \emph{primitive} tree-level term in the effective action rather than adopting naive one-loop Sakharov induced gravity (retired branch \texttt{C12}, discussed in Section~\ref{sec:emergence_sm_g_moved}), because quadratic cutoff shifts depend on regularization scheme and unknown UV physics above $\Lambda_{\mathrm{UV}}$.
  \item \textbf{Logarithmic Curvature-Squared and Dimension-6 Counterterms ($a_2$):}
  \begin{equation}
  a_2(x)
  = \frac{1}{180}\bigl(R_{\mu\nu\rho\sigma}R^{\mu\nu\rho\sigma} - R_{\mu\nu}R^{\mu\nu}\bigr)I
    + \frac{1}{6}\,\Box_g\!\left(E + \frac{1}{5}R\right)I
    + \frac{1}{2}\left(E + \frac{1}{6}R\right)^2
    + \frac{1}{12}\,\Omega_{\mu\nu}\Omega^{\mu\nu},
  \label{eq:seeley_dewitt_a2}
  \end{equation}
  generating logarithmic $\ln(\Lambda_{\mathrm{reg}}/\mu)$ running of dimension-4 couplings together with one-loop higher-curvature counterterms $c_1(\mu)R^2 + c_2(\mu)R_{\mu\nu}R^{\mu\nu}$ (as well as Euler--Heisenberg-like $O(\Lambda_{\mathrm{EM}}^{-4})(F_{\mu\nu}F^{\mu\nu})^2$ operators).
\end{enumerate}

\subsection{Two-Loop Goroff--Sagnotti Proliferation and the Explicit \texttt{OPEN-1} Boundary}

At two loops in pure and scalar-coupled Einstein gravity, the non-vanishing on-shell Goroff--Sagnotti cubic Riemann counterterm $\frac{209}{2880(16\pi^2)^2\epsilon}\,\frac{1}{M_{\mathrm{Pl}}^2}R_{\mu\nu}^{\ \ \rho\sigma}R_{\rho\sigma}^{\ \ \alpha\beta}R_{\alpha\beta}^{\ \ \mu\nu}$ appears, followed by an infinite tower of higher-dimension operators as $L \to \infty$. Below $\Lambda_{\mathrm{UV}}$, every such operator is suppressed by $(E/\Lambda_{\mathrm{UV}})^{2L+2} \ll 1$ and leaves all laboratory, Solar System, and sub-Planckian cosmological predictions completely predictive and finite; at or above $E \ge \Lambda_{\mathrm{UV}}$, perturbative EFT breaks down, and constructing a regulator-independent non-perturbative 4D quantum-gravity UV completion remains explicitly open (\texttt{OPEN-1}, Tier~5).

\section{Accelerated Worldlines: Dimensional Audit and Open Self-Force Status}
\label{sec:electron-accel}

Earlier SIT drafts proposed an informal accelerated-particle proper-time correction written as $\delta\rho_t \simeq \beta\,G m_0 / (a\,r_e)$ (with dimensionless $\beta$ asserted to satisfy a torsion-balance bound $|\beta| < 10^{-5}$ and fractional proper-time shift $< 10^{-15}$). Version~81 explicitly \emph{retires} both that formula and those numerical claims for two concrete reasons:
\begin{enumerate}[leftmargin=2em]
  \item \textbf{Dimensional Inconsistency:} In SI units with dimensionless $\beta$, rest mass $m_0$ ($[\mathrm{M}]$), proper acceleration $a$ ($[\mathrm{L\,T^{-2}}]$), and classical radius $r_e$ ($[\mathrm{L}]$), the ratio $G m_0 / (a\,r_e)$ has physical dimension $[\mathrm{L^3\,M^{-1}\,T^{-2}}][\mathrm{M}] / ([\mathrm{L\,T^{-2}}][\mathrm{L}]) = [\mathrm{L}]$ (\emph{length}), contradicting the assigned inverse-time dimension $[\rho_t] = [\mathrm{T^{-1}}]$ of the temporal-density scalar (Section~\ref{sec:atomic-clocks}).
  \item \textbf{Open Self-Force and Torsion-Balance Bridge (\texttt{OPEN-6}):} Consistent with Sections~\ref{sec:timedensity}, \ref{sec:experimental-predictions}, and~\ref{sec:stage6-action-diagnostics}, standalone $|\beta| < 10^{-5}$ bounds are retired, and any scalar radiation-reaction or self-force shift along an accelerated worldline in \texttt{M\_AUTH\_SIT\_EFT\_v1} requires an explicit green-function regularization and microscopic worldline coupling derivation (\texttt{OPEN-6A}/\texttt{OPEN-6B}).
\end{enumerate}

\section{Operational Coherence Proxies vs.\ Open Quantitative Response Maps}
\label{sec:quantitative}

Earlier drafts heuristically identified $\Rcoh$ with a discrete state-count fraction $N_{\mathrm{coherent}}/N_{\mathrm{total}}$, substituted $\rho_t = \rho_0\exp[\alpha\Rcoh]$, and quoted illustrative figures including a $\Delta\nu/\nu \approx 5\times 10^{-11}$ cavity-modulated optical-clock shift (at $\Delta\Rcoh \sim 10^{-8}, \alpha \simeq 10^{-2}$) and a predicted $1\%$ galaxy-cluster lensing mass-profile correction. Consistent with the operational separation contract (Section~\ref{sec:PrimitiveFieldDefinitions}) and the Scoped Falsification Principle (Section~\ref{sec:experimental-predictions}), Version~81 retires both the $5\times 10^{-11}$ clock forecast and the $1\%$ cluster mass-profile figure as unsupported quantitative predictions: in \texttt{M\_AUTH\_SIT\_EFT\_v1}, $\Rcoh = |\psi|/\psi_0$ is a classical field modulus ratio rather than a discrete state count, deriving a cavity-modulated clock shift requires resolving both the visible $U(1)$ photon dictionary (\texttt{OPEN-6B}) and the laboratory preparation source $J_{\mathrm{prep}}$ (\texttt{OPEN-6C}), and deriving any cluster-scale weak-lensing residual requires a complete cosmological perturbation and mass-calibration calculation (\texttt{OPEN-6D}). Likewise, heuristic phase-dissipation expressions such as $dI/dt = -\gamma\int (\Delta\phi)^2\,d^3x$ serve solely as conceptual open-system motifs rather than canonical Hamiltonian equations of $S_{\mathrm{AUTH}}$.

\section{Conditional Parameterization of Atomic Clock Frequency Shifts in SIT}
\label{sec:atomic-clock-prediction}
\label{sec:atomic-clocks}

Section~\ref{sec:atomic-clock-prediction} formulates the disambiguated conditional parameterization for atomic clock frequency shifts under postulated $\rho_t$ variations (contingent on \texttt{OPEN-6B}/\texttt{OPEN-6C}), separating gravitational-potential redshift tests from hypothetical laboratory-prepared coherence benchmarks and distinguishing published primary measurements from quarantined sensitivity targets.

\subsection{Physical Context, Disambiguated Response Symbols, and Assumptions} % anchor: second occurrence

Super Information Theory introduces a scalar field \(\rho_t(x)\), the \emph{time-density field}, with physical dimension inverse time \([T^{-1}]\). This field is hypothesized to modulate the local structure of proper time or effective electromagnetic couplings, potentially affecting quantum transition frequencies measurable by atomic clocks (contingent on the visible \(U(1)\) photon dictionary, \texttt{OPEN-6B}). To repair a notation collision present in Versions~79--80.2---which reused the single symbol \(\alpha_{\mathrm{eff}}\) for both the local field response and the gravitational-potential redshift slope, thereby double-counting the transfer factor \(\mathcal{T}\) whenever \(\mathcal{T} \ne 1\)---Version~81 distinguishes the two coefficients strictly:
\begin{itemize}
    \item \(\rho_t(x)\) varies slowly on scales relevant to laboratory clocks and gravitational potentials.
    \item Local coupling of \(\rho_t\) to the electromagnetic sector is parameterized by a function \(f_2(\rho_t)\) expanded to first order using the \emph{local field-response coefficient} \(\alpha_\rho\):
    \[
    f_2(\rho_t) = 1 + \alpha_\rho \frac{\delta \rho_t}{\rho_0} + \mathcal{O}(\delta \rho_t^2),
    \]
    where \(\rho_0\) is the vacuum expectation value of \(\rho_t\).
    \item For gravitational-potential variations \(\Delta\Phi/c^2\), defining the dimensionless laboratory transfer factor \(\mathcal{T} := (\rho_0^{-1}\partial\rho_t/\partial(\Phi/c^2))_{\mathrm{lab}}\) yields \(\delta\rho_t/\rho_0 = \mathcal{T}\,(\Delta\Phi/c^2)\) and defines the \emph{gravitational-potential redshift slope}
    \[
    \alpha_\Phi := \frac{d\ln\nu}{d(\Phi/c^2)} = \alpha_\rho\,\mathcal{T}.
    \]
    \item Empirically, the cited Tokyo Skytree transportable optical-lattice clock test (Takamoto~\emph{et~al.}, 2020) reported \(\alpha_{\mathrm{RS}} = (+1.4 \pm 9.1)\times 10^{-5}\) across a \(450\,\mathrm{m}\) height difference, while NIST (2026) achieved \(3.2\times 10^{-18}\) fractional optical-clock ratio uncertainty. The legacy \(3\times 10^{-8}\) value cited in Versions~79--80.2 is formally quarantined in Version~81 as a \emph{hypothetical future sensitivity benchmark}, not a published \(95\%\)~CL experimental bound.
\end{itemize}

\subsection{Derivation of the Gravitational Potential-Induced Shift}

When atomic transition frequencies depend on electromagnetic couplings modulated by \(f_2(\rho_t)\), the local fractional frequency shift at spacetime point \(x\) is governed strictly by the local field-response coefficient \(\alpha_\rho\):
\[
\frac{\Delta \nu}{\nu}(x) = \alpha_\rho \frac{\delta \rho_t(x)}{\rho_0}.
\]
For a gravitational-potential-driven variation \(\delta\rho_t/\rho_0 = \mathcal{T}\,(\Delta\Phi/c^2)\), substituting into the local response gives without double-counting \(\mathcal{T}\):
\[
\left(\frac{\Delta \nu}{\nu}\right)_{\mathrm{SIT}}
= \alpha_\rho\,\frac{\delta \rho_t}{\rho_0}
= \alpha_\rho\,\mathcal{T}\,\frac{\Delta\Phi}{c^2}
= \alpha_\Phi\,\frac{\Delta\Phi}{c^2}.
\]
This additive correction is distinct from the pure GR redshift \(\Delta\Phi/c^2\) and requires the visible \(U(1)\) photon dictionary (\texttt{OPEN-6B}) plus an action-derived terrestrial scalar profile \(\mathcal{T}\) before any primary clock dataset can constrain \texttt{M\_AUTH\_SIT\_EFT\_v1} unconditionally.

\subsection{Empirical Provenance vs.\ Hypothetical Sensitivity Benchmarks}

\textbf{Primary experimental provenance vs.\ quarantined hypothetical target.} The published Tokyo Skytree optical-lattice clock test (Takamoto~\emph{et~al.}, 2020) constrains gravitational-redshift deviations at \(\alpha_{\mathrm{RS}} = (+1.4 \pm 9.1)\times 10^{-5}\) (\(1\sigma\) combined), whereas the historical figure
\[
|\alpha_\Phi|_{\mathrm{hyp}} \lesssim 3 \times 10^{-8} \quad \text{(hypothetical future sensitivity target; NOT an empirical 95\% CL bound)}
\]
would require \(1.47\times 10^{-21}\) fractional clock accuracy over the \(450\,\mathrm{m}\) Skytree baseline (\(\Delta\Phi/c^2 = 4.91\times 10^{-14}\)). Version~81 therefore quarantines \(3\times 10^{-8}\) strictly as a hypothetical future sensitivity target (see Section~\ref{sec:stage6-action-diagnostics} and Appendix~\ref{appendix:ClockConstraints}).

Vacuum \(\rho_0\) corresponds in Planck-scale exploratory estimates to
\[
\rho_0 \sim \frac{1}{t_{\mathrm{Planck}}} \approx 5.4 \times 10^{43} \ \mathrm{s}^{-1},
\]
while in canonical EFT benchmark \texttt{BM\_PHYS\_EFT\_1} we set \(\rho_0 = 2.0\times 10^{16}\,\mathrm{GeV}\). Forecasts are expressed in terms of \(\delta\rho_t/\rho_0\) with \(\alpha_\rho\), or in terms of \(\Delta\Phi/c^2\) with \(\alpha_\Phi = \alpha_\rho\mathcal{T}\).

\textbf{Conditional metrological status.} Until the visible electromagnetic \(U(1)\) dictionary (\texttt{OPEN-6B}) and terrestrial transfer factor \(\mathcal{T}\) are derived from \texttt{M\_AUTH\_SIT\_EFT\_v1}, clock comparisons bound the phenomenological product \(\alpha_\Phi = \alpha_\rho\mathcal{T}\) rather than imposing an unconditional constraint on the unchanged action.

\subsection{Test 2: Conditional Template for Engineered Coherence-Induced Shifts (\texttt{OPEN-6B}/\texttt{OPEN-6C})}
A second prospective experimental modality explores whether directly manipulating a laboratory system's quantum coherence can induce a differential clock shift at fixed gravitational potential. Crucially, under Stage-6 \texttt{OPEN-6} discipline, this channel is explicitly a \emph{conditional future observable template} rather than a fixed or non-tunable quantitative prediction of unchanged \texttt{M\_AUTH\_SIT\_EFT\_v1}:
\begin{enumerate}[leftmargin=2em]
  \item \textbf{Distinct Preparation-Response Requirement:} Whereas a potential-driven clock shift uses the unified dictionary $y_{\mathrm{clock}} = \alpha_\rho r$ ($r := \delta\rho_t/\rho_0$) paired with a derived gravitational transfer relation $r = \mathcal{T} p$ ($p := \Delta\Phi/c^2$) to give $y_{\mathrm{clock}} = \alpha_\Phi p$ ($\alpha_\Phi = \alpha_\rho\mathcal{T}$), a laboratory-prepared coherent state requires its own action-derived state-preparation source $J_{\mathrm{prep}}$ (\texttt{OPEN-6C}) together with an explicit visible electromagnetic $U(1)$ photon dictionary (\texttt{OPEN-6B}). The gravitational-potential transfer factor $\mathcal{T}$ cannot simply be carried over into a distinct state-preparation geometry.
  \item \textbf{Upper Bounds Do Not Imply Nonzero Lower Bounds:} Any historical exploratory scale (such as $\Delta\nu/\nu \sim 10^{-15}$ for an illustrative parameter ansatz) is strictly conditional and cannot be inferred from an experimental upper limit on $\alpha_\Phi$. An experimental upper bound on a coupling never establishes a nonzero lower bound on a predicted signal.
  \item \textbf{Precise Falsification Scope:} Comparing two identical optical clocks held at the same gravitational potential (one in an engineered coherent state, one decohered) can exclude a \emph{fully specified parameter benchmark} once its preparation source $J_{\mathrm{prep}}$, visible-photon response coefficients, and experimental nuisance parameters are derived and fixed; a null result cannot exclude the entire SIT framework while \texttt{OPEN-6B} and \texttt{OPEN-6C} remain unresolved.
\end{enumerate}

\subsection{Experimental Feasibility and Conditional Falsifiability Scope}

Modern optical-lattice and quantum-logic atomic clocks achieve optical frequency-ratio measurement uncertainties at or below $3.2\times 10^{-18}$ across phase-stabilized optical fiber links (NIST/BACON, \emph{Phys.\ Rev.\ Lett.}\ \textbf{137}, 033201, 2026; DOI \texttt{10.1103/g865-9mk1}, \texttt{arXiv:2512.21428}) and test terrestrial gravitational redshifts at $\alpha_{\mathrm{RS}} = (+1.4 \pm 9.1)\times 10^{-5}$ ($1\sigma$ combined; Tokyo Skytree, Takamoto~\emph{et~al.}, 2020). Converting these primary measurements into quantitative exclusion contours on \texttt{M\_AUTH\_SIT\_EFT\_v1} requires (i)~deriving the visible $U(1)$ response coefficients and terrestrial transfer factor $\mathcal{T}$ for potential-driven tests ($y_{\mathrm{clock}} = \alpha_\Phi p = \alpha_\rho\mathcal{T} p$), and (ii)~deriving $J_{\mathrm{prep}}$ and apparatus nuisance budgets for engineered-coherence tests. Until those mappings (\texttt{OPEN-6B}/\texttt{OPEN-6C}) are closed, both channels remain explicitly blocked in the Stage-6 provenance manifest (\texttt{primary\_measurements\_manifest\_v2.json}).
% --- END INPUT sections/04_part1_sec16_to_25.tex ---
% --- BEGIN INPUT sections/05_part2_sec26_to_31.tex ---
\part{Experimental Program and Falsifiability}

\subsubsection{Cold-Atom Interferometry Phase Shifts (Exploratory, Conditional on \texttt{OPEN-6A}/\texttt{OPEN-6C})}

Cold-atom interferometry is tracked as a prospective exploratory channel for probing whether engineered gradients in macroscopic quantum coherence or local time density can source observable phase shifts. Unchanged \texttt{M\_AUTH\_SIT\_EFT\_v1} does \emph{not} currently derive a fixed numerical phase-shift forecast for cold-atom fountains or optical-lattice interferometers (earlier draft estimates such as $\Delta\phi \sim 10^{-3}\,\mathrm{rad}$ are removed here as unsupported).

\textbf{Derivation Prerequisites and Measurement Strategy:}
\begin{itemize}
    \item First derive the microscopic atomic scalar charge (\texttt{OPEN-6A}) and the gauge-invariant state-preparation source $J_{\mathrm{prep}}$ (\texttt{OPEN-6C}) for a specified cold-atom preparation protocol, controlling ordinary changes in density, temperature, and energy.
    \item Only once a concrete parameter point and preparation model yield a derived phase-shift prediction exceeding interferometric phase noise can a cold-atom experiment test or exclude that specific model.
\end{itemize}

\subsubsection{Quantum Entanglement and Space-Based Bell Tests (Exploratory)}

As a Tier-4 conceptual research question outside the currently derived predictions of \texttt{M\_AUTH\_SIT\_EFT\_v1}, one may ask whether coherence-modulated gravitational backgrounds alter entanglement-correlation statistics across differing gravitational potentials (ground-based versus orbital platforms). Any quantitative test requires deriving an explicit coupling between the scalar background and photonic/matter correlation observables (\texttt{OPEN-6B}/\texttt{OPEN-6C}) and bounding all standard relativistic and optical systematics.

\subsubsection{Probing Coherence Phase Holonomy via Electromagnetism (Conditional on \texttt{OPEN-6B})}

In the canonical action \texttt{M\_AUTH\_SIT\_EFT\_v1}, the internal gauge potential $A_\mu$ is a principal-bundle connection (reducing locally to a pure gradient only on flat patches where $F_{\mu\nu} = 0$) and acquires a large broken-phase Higgs/Proca mass $m_A^2 = \kappa_c \hat{q}^2 v_0^2 / f_0$ whenever $\hat{q} v_0 \ne 0$, remaining distinct from the massless visible Standard Model photon unless an explicit visible-sector identification or mixing portal is derived (\texttt{OPEN-6B}). Prospective Aharonov--Bohm and SQUID magnetometry probes of beyond-Maxwell holonomy effects therefore remain strictly conditional on resolving \texttt{OPEN-6B} (Section~\ref{sec:magnetism-holonomy}).

\subsubsection{Noether's Theorem and Informational Symmetry Checks (Exploratory)}

Within \texttt{M\_AUTH\_SIT\_EFT\_v1}, global $U(1)$ invariance yields the exact gauge-invariant conserved current $J^\mu_{\mathrm{coh}} = \kappa_c|\psi|^2(\nabla^\mu\theta - \hat{q}A^\mu)$ (Eq.~\eqref{eq:abstract_coherence_current}). Translating that exact field-theoretic conservation law into laboratory atomic-clock or interferometric observables remains contingent on the microphysical charge and preparation dictionaries (\texttt{OPEN-6A}--\texttt{OPEN-6C}).

\subsection{Condensed Matter Exploration: Oscillator Network Synchrony (Conditional on \texttt{OPEN-6C})}

Earlier exploratory drafts conjectured that macroscopic arrays of coupled oscillators (such as superconducting Josephson junctions or optomechanical resonators) might source a distinctive gravitational perturbation upon entering a phase-locked state. Under \texttt{QG49} discipline, unchanged \texttt{M\_AUTH\_SIT\_EFT\_v1} does \emph{not} derive an anomalous or non-tunable gravitational perturbation from classical or quantum phase locking. Any serious condensed-matter gravimetry proposal must first define a gauge-invariant observable for the synchronized preparation, derive from the action whether driving phase locking changes the scalar source vector $J_{\mathrm{prep}}$ after ordinary changes in temperature, mass density, and electromagnetic energy are held fixed (\texttt{OPEN-6C}), and verify whether the resulting response exceeds zero (a preparation-specific null theorem remains an open scientific outcome).

\subsection{Astrophysical and Cosmological Explorations (Conditional on \texttt{OPEN-6D})}
\label{sec:part2-astro-probes}
\subsubsection{Gravitational Lensing Residuals (Tier-4 Exploratory)}

Earlier draft revisions listed speculative percent-level gravitational lensing deviations; those numerical percentages are removed here because unchanged \texttt{M\_AUTH\_SIT\_EFT\_v1} has not yet been evolved through a cosmological Friedmann--Robertson--Walker background and linear/nonlinear weak- and strong-lensing perturbation analysis (\texttt{OPEN-6D}). High-resolution cluster and galaxy-scale lensing surveys (JWST, Euclid, Rubin Observatory) provide empirical benchmarks only once a concrete cosmological parameter branch and its derived metric-slip prediction are computed from the action.

\subsubsection{Cosmological Questions: Dark Matter, Dark Energy, and Expansion History (Tier-4 Exploratory)}

Whether extensions or cosmological solutions of SIT can account for galactic rotation curves, late-time acceleration, or expansion-rate tensions without standard cold dark matter and $\Lambda$ remains an unproved Tier-4 research hypothesis (no 2--5\% cosmological deviation forecast is derived from \texttt{M\_AUTH\_SIT\_EFT\_v1}). Testing any cosmological SIT proposal against supernovae, baryon acoustic oscillations, Cosmic Microwave Background anisotropies, and weak-lensing shear requires first solving the full coupled background and perturbation equations (\texttt{OPEN-6D}).

\subsection*{Experimental Signatures and Required Sensitivities}

The following experimental domains offer promising avenues for testing the theory, each with a distinct observable signature and required sensitivity:

\begin{description}
    \item[Atomic Clock Tests (Conditional on \texttt{OPEN-6B}/\texttt{OPEN-6C})] Two prospective channels are tracked: (1)~gravitational-potential redshift tests ($y_{\mathrm{clock}} = \alpha_\Phi p = \alpha_\rho\mathcal{T} p$), anchored to primary data such as Tokyo Skytree's $(+1.4\pm 9.1)\times 10^{-5}$ and NIST/BACON frequency-ratio uncertainties $\le 3.2\times 10^{-18}$; and (2)~engineered-coherence differential clock comparisons contingent on deriving the state-preparation source $J_{\mathrm{prep}}$ and visible $U(1)$ dictionary. Neither channel currently implies an unconditional fixed signal floor.

    \item[Cold-Atom Interferometry (Conditional on \texttt{OPEN-6A}/\texttt{OPEN-6C})] Exploratory phase-shift and differential-acceleration protocols contingent on deriving $J_{\mathrm{prep}}$ and microscopic atomic scalar charges.

    \item[Quantum Bell Tests] Exploratory study of whether coherence-modulated environments alter correlation statistics under controlled systematics.

    \item[Condensed Matter Gravimetry (Conditional on \texttt{OPEN-6A}/\texttt{OPEN-6C})] Exploratory gravimetric tests of synchronized oscillator arrays once the microscopic source coupling $J_{\mathrm{prep}}$ and composite material charges are derived from \texttt{M\_AUTH\_SIT\_EFT\_v1}.

    \item[Gravitational Lensing (Level-1 Exploratory)] Cluster-scale weak and strong lensing templates tested against Euclid, JWST, and Rubin Observatory surveys.

    \item[Cosmological Surveys (Level-1 Exploratory)] Background expansion and structure-growth comparisons at multi-percent survey precision.

    \item[Electromagnetism--Coherence Link] Probing principal-bundle phase holonomy in Aharonov--Bohm-type experiments (Section~\ref{sec:magnetism-holonomy}), contingent on the visible $U(1)$ dictionary (\texttt{OPEN-6B}).
\end{description}

\subsection{Experimental Timeline and Roadmap}

We organize empirical development around a staged verification sequence governed by \texttt{OPEN-6}:

\textbf{Near-Term (Action-to-Apparatus Derivation and Data Audit):}
\begin{itemize}
    \item Complete the Lee~(2020) dissertation data-availability audit, patterned 18-fold/120-fold Pt--epoxy torque reproduction, and composite material charge analysis (\texttt{OPEN-6A}/\texttt{OPEN-6D}).
    \item Resolve the visible $U(1)$ photon dictionary (\texttt{OPEN-6B}) and laboratory preparation source $J_{\mathrm{prep}}$ (\texttt{OPEN-6C}) through the canonical $(K, H)$ propagator identity.
\end{itemize}

\textbf{Medium-Term (Benchmark-Specific Confrontation):}
\begin{itemize}
    \item Joint two-mode torsion-balance, optical-clock, and atom-interferometer likelihood exclusions against locked primary provenance records (\texttt{primary\_measurements\_manifest\_v2.json}).
    \item Precision gravitational lensing observational comparisons.
\end{itemize}

\textbf{Long-Term (Next-Generation Metrology):}
\begin{itemize}
    \item Space-based optical/nuclear clock networks and long-baseline atom interferometers once benchmark predictions and nuisance budgets are fully specified.
\end{itemize}

This staged roadmap separates completed mathematical analysis from open action-to-apparatus derivations.

\section{Empirical Content and Parameter Structure of \texttt{M\_AUTH\_SIT\_EFT\_v1}}
\label{sec:empirical-content}
See Section~\ref{sec:atomic-clock-prediction} and Appendix~\ref{appendix:ClockConstraints} for the unified optical-clock response dictionary ($y_{\mathrm{clock}} = \alpha_\rho r$, with $y_{\mathrm{clock}} = \alpha_\Phi p$ and $\alpha_\Phi = \alpha_\rho\mathcal{T}$ for potential-driven configurations where $r = \mathcal{T} p$ is derived).

A defining requirement for a viable effective field theory is transparent accounting of its actual coupling functions and parameter dependencies. Here we connect empirical parameter fixing directly to the canonical action \texttt{M\_AUTH\_SIT\_EFT\_v1}.

\subsection{Canonical Coupling Functions vs.\ Historical Exponential Ans\"atze}

In the authoritative Version~81 action \texttt{M\_AUTH\_SIT\_EFT\_v1} (Eq.~\eqref{eq:abstract_master_action}), the constitutive dressings of the gauge and matter sectors are the dimensionless, strictly vacuum-positive functions
\begin{equation}
f_{\mathrm{EM}}(\Rcoh) = 1 + c_{\mathrm{EM}}\Rcoh^2 \ge 1,
\qquad
f_{\mathrm{SM}}(\Rcoh) = 1 + \frac{\Rcoh^2}{2 M_{\mathrm{sym}}^2} \ge 1,
\end{equation}
coupled to the real scalar $\rho_t$ (mass dimension $[M]^1$) through the portal potential $U_{\mathrm{link}}(\rho_t, |\psi|)$ and linearized kinetic/Hessian matrices $(K, H)$. By contrast, earlier exploratory drafts (Versions~78--79) wrote a separate heuristic exponential ansatz of the form $f_{\mathrm{alt}} = \exp[\tilde{\beta}(\rho_t/\rho_0 - 1)] \approx 1 + \tilde{\beta}(\delta\rho_t/\rho_0)$ (using a dimensionless ratio $\rho_t/\rho_0$, since multiplying dimension-$1$ $\rho_0$ or $\delta\rho_t$ by a dimensionless parameter $\beta$ would be dimensionally inconsistent). Any such exponential coupling represents a \emph{distinct} alternative model requiring its own complete derivation and is not part of canonical \texttt{M\_AUTH\_SIT\_EFT\_v1}.

\subsection{Parameter Dependencies and Experimental Constraints}

Canonical \texttt{M\_AUTH\_SIT\_EFT\_v1} carries multiple independent physical parameters—the scales $(\psi_0, \rho_0, v_0, m_{\mathrm{link}})$, quartic/portal couplings $(\lambda_\rho, \lambda_\psi, \lambda_{\mathrm{portal}})$, gauge charge $\hat{q}$, and dimension-$6$ portal parameters $(c_{\mathrm{EM}}, M_{\mathrm{sym}})$—which enter linearized exchange forces jointly through the machine-checked propagator identity $(p^2 K + H)^{-1} = V\,\mathrm{diag}(1/(p^2+m_a^2))\,V^\top$ (\texttt{Stage6Addendum.\allowbreak{}ActionResponsePropagatorIdentity}). Consequently, fixing a single experimental bound constrains one combination of parameters rather than rendering all remaining predictions parameter-independent; full predictive confrontation across multiple channels requires specifying the shared canonical parameter tuple together with the microphysical mappings in \texttt{OPEN-6}.

\subsection{Concrete Example: Unified Clock-Response Dictionary}

Consider a conditional electromagnetic-sector coupling (pending \texttt{OPEN-6B}) probed by optical-lattice clocks. Using the unified notation of Section~\ref{sec:atomic-clock-prediction} and Appendix~\ref{appendix:ClockConstraints} with $r := \delta\rho_t/\rho_0$ and $p := \Delta\Phi/c^2$, the local fractional frequency shift is governed strictly by $\alpha_\rho$:
\[
y_{\mathrm{clock}} = \alpha_\rho\,r = \alpha_\rho\,\frac{\delta\rho_t}{\rho_0}.
\]
For a potential-driven configuration where $r = \mathcal{T} p$ has been derived from the field equations, this becomes without double-counting $\mathcal{T}$:
\[
y_{\mathrm{clock}} = \alpha_\Phi\,p = \alpha_\rho\,\mathcal{T}\,\frac{\Delta\Phi}{c^2},
\qquad
\alpha_\Phi = \alpha_\rho\,\mathcal{T}.
\]
A laboratory-prepared coherence state requires its own action-derived source-response relation (\texttt{OPEN-6C}) rather than reusing $\mathcal{T}$. The published Tokyo Skytree campaign (Takamoto~\emph{et~al.}, 2020) reported $\alpha_{\mathrm{RS}} = (+1.4 \pm 9.1)\times 10^{-5}$, and the NIST/BACON report (\texttt{arXiv:2512.21428}) achieved optical frequency-ratio measurement uncertainties $\le 3.2\times 10^{-18}$; the legacy $|\alpha_\Phi|_{\mathrm{hyp}} \lesssim 3 \times 10^{-8}$ benchmark is strictly a quarantined hypothetical future sensitivity target.

\section{Prospective Coherence-Gravity Test: Bose--Einstein Condensate vs.\ Thermal Cloud (\texttt{OPEN-6A}/\texttt{OPEN-6C})}
\label{sec:smoking-gun-bec}
\label{sec:bec_smoking_gun}

One long-horizon conceptual test asks whether a macroscopic quantum condensate ($^{87}\mathrm{Rb}$ BEC) and an equal-mass thermal cloud of the same isotope exhibit any differential gravitational acceleration beyond General Relativity. Under Version~81 epistemic governance---integrating the Penrose--Onsager condensate-source boundary of Paper~003 and the experimental-control framework of Paper~014---this experiment is framed strictly as a \emph{prospective future test contingent on resolving \texttt{OPEN-6A} and \texttt{OPEN-6C}}, not as a currently derived nonzero signal prediction of unchanged \texttt{M\_AUTH\_SIT\_EFT\_v1}.

\subsection{Standard Physics vs.\ Conditional SIT Condensate-Source Hypothesis}

Under the Equivalence Principle of General Relativity, gravitational coupling depends solely on total stress--energy and not on internal condensate coherence. Equally, unchanged \texttt{M\_AUTH\_SIT\_EFT\_v1} (like the baseline July~2026 conformal action of Paper~002/003) does \emph{not} by itself distinguish a Bose--Einstein condensate from a matched thermal cloud of equal mass-energy. Producing a condensate-dependent gravitational differential requires either (i)~deriving the microscopic composite atomic scalar charge $q_{\mathrm{eff}}(^{87}\mathrm{Rb})$ (\texttt{OPEN-6A}) and the laboratory state-preparation source $J_{\mathrm{prep}}$ (\texttt{OPEN-6C}) directly from a specified many-body matter action (showing explicitly how a Penrose--Onsager condensate eigenmode sources local fields consistently with stress--energy conservation), or (ii)~adding an explicit phenomenological condensate-source operator with its own declared coupling and conservation law (as formulated for the July~2026 neutral-scalar benchmark in Paper~003, without mislabeling that added hypothesis as an automatic consequence of \texttt{M\_AUTH\_SIT\_EFT\_v1}). At canonical benchmark \texttt{BM\_PHYS\_EFT\_1} ($m_1 \approx 2.45\times 10^{14}\,\mathrm{GeV}$, range $\lambda_1 \approx 8.06\times 10^{-31}\,\mathrm{m}$), tree-level scalar exchange is exponentially suppressed at macroscopic scales.

\subsection{Matched-Preparation Residual Protocol and Systematic Controls (Paper~014)}

In principle, comparing the same atomic species in two distinct coherence preparations eliminates isotopic composition differences (in contrast to dual-isotope $^{85}\mathrm{Rb}/^{87}\mathrm{Rb}$ tests such as Asenbaum~\emph{et~al.}, 2020), but introduces strong causal confounding because condensate fraction covaries with density profile, mean-field interaction energy, temperature, expansion velocity, and trap history (Paper~014). Following Paper~002 (Section~9.1) and Paper~014, a defensible protocol therefore targets the matched-preparation residual
\begin{equation}
\Delta a_{\mathrm{res}} := (a_{\mathrm{high\text{-}C}} - a_{\mathrm{low\text{-}C}}) - \Delta a_{\mathrm{standard}},
\label{eq:bec_matched_residual}
\end{equation}
requiring explicit controls across three steps:
\begin{enumerate}
    \item \textbf{Characterize Matched Preparations and Standard Systematics:} Prepare equal-atom-number $^{87}\mathrm{Rb}$ ensembles above and below condensation (measured by explicit condensate fraction $f_{\mathrm{BEC}}$ rather than equating $f_{\mathrm{BEC}}$ with the classical field $\Rcoh$ by fiat) and model $\Delta a_{\mathrm{standard}}$ across density profile, mean-field energy, magnetic sensitivity, vibration, wave-front aberration, thermal expansion, and center-of-mass trajectory.
    \item \textbf{Gradient Reversal and Preparation-Error Identifiability:} Verify that any candidate residual tracks a preregistered reversal or translation of the modeled source gradient while total preparation energy and geometry remain matched, and check whether ordinary preparation uncertainties absorb the signal.
    \item \textbf{Compare Differential Acceleration Against a Derived Benchmark:} Evaluate $\Delta a_{\mathrm{res}}/g_E$ only against a fully specified theoretical benchmark (including action, $J_{\mathrm{prep}}$, composite charges, apparatus transfer function, and covariance model) that fixes a concrete quantitative prediction.
\end{enumerate}

\subsection{Exact Falsification Scope}

\begin{itemize}
    \item \textbf{Standard Physics Predicts:} Identical gravitational acceleration ($\delta a = 0$) within experimental systematics.
    \item \textbf{Conditional SIT Benchmark Scope:} An experimental upper limit on a coupling (such as a clock redshift bound) never establishes a nonzero lower limit on a predicted BEC signal. Any legacy target near $3\times 10^{-8}$ ($29.4\,\mu\mathrm{Gal}$, Appendix~\ref{app:BEC-benchmark}) is strictly a quarantined hypothetical sensitivity scale. A null measurement can exclude a \emph{fully specified parameter benchmark} once its preparation source $J_{\mathrm{prep}}$, material response, and nuisance parameters are derived and fixed; it cannot exclude the entire SIT framework while those mappings remain open.
\end{itemize}

\section{Experimental Channels, Operational Definitions, and Open Action-to-Apparatus Bridges}
\label{sec:experimental-predictions}
\label{sec:experimental-predictions-falsifiability}

Building upon the unified clock dictionary and Stage-6 provenance ledger (\texttt{primary\_measurements\_manifest\_v2.json}), this section summarizes the operational definitions and prospective laboratory channels together with their explicit \texttt{OPEN-6} prerequisites:

\begin{itemize}
    \item \textbf{Sub-Millimeter Torsion-Balance Channel (Lee~\emph{et~al.}, 2020):} Currently the single weak-field channel directly evaluated from unchanged \texttt{M\_AUTH\_SIT\_EFT\_v1} under a declared uniform-slab dust closure ($\lambda_1 \approx 8.06\times 10^{-31}\,\mathrm{m}$ at \texttt{BM\_PHYS\_EFT\_1}, Section~\ref{sec:stage6-action-diagnostics}); advancing to a full patterned rotating-mass torque likelihood for the 18-fold/120-fold Pt--epoxy bodies on BK7 supports with source and detector contrasts $\rho_{\mathrm{Pt}} q_{\mathrm{Pt}} - \rho_{\mathrm{epoxy}} q_{\mathrm{epoxy}}$ is \texttt{OPEN-6A}/\texttt{OPEN-6D}.

    \item \textbf{Optical-Clock and BEC Channels (Blocked Pending \texttt{OPEN-6A}--\texttt{OPEN-6C}):} Prospective differential-acceleration ($^{87}\mathrm{Rb}$ BEC vs.\ thermal cloud) and synchronized optical-clock comparisons remain conditional templates requiring action-derived $J_{\mathrm{prep}}$, composite material charges, and the visible $U(1)$ photon dictionary before specific nonzero signal magnitudes can be predicted.
\end{itemize}

\subsection{Operational Definitions and Observable Quantities}

To connect formal fields to experiments, we rely on the operational definitions established in Section~\ref{sec:operational-definitions} \label{sec:def}, while separating operational proxies from action-derived source terms (\texttt{OPEN-6}):

\begin{itemize}
    \item \textbf{Time-Density Field $\rho_t(x)$:} Real scalar field with natural mass dimension $[M]^1$ entering \texttt{M\_AUTH\_SIT\_EFT\_v1}; operational clock readouts map through $y_{\mathrm{clock}} = \alpha_\rho r$ ($r = \delta\rho_t/\rho_0$) once \texttt{OPEN-6B} fixes $\alpha_\rho$.
    \item \textbf{Coherence Ratio $R_{\mathrm{coh}}(x)$:} Dimensionless modulus ratio $|\psi(x)|/\psi_0 \ge 0$; linking laboratory purity or visibility to classical field displacements requires deriving $J_{\mathrm{prep}}$ (\texttt{OPEN-6C}).
    \item \textbf{Derived Quantities:} Astrophysical lensing, torsion-balance torques, and clock comparisons test concrete parameter tuples once the corresponding forward transfer functions are derived.
\end{itemize}

\subsection{Prospective Laboratory and Astrophysical Signatures (Contingent on \texttt{OPEN-6})}

Subject to the Stage-6 epistemic boundaries above, five prospective observational channels are tracked for future benchmark-specific testing:

\begin{enumerate}
    \item \textbf{Coherence-Dependent Gravitational Tests:} Comparing equal-mass coherent versus incoherent ensembles once $J_{\mathrm{prep}}$ and composite material charges (\texttt{OPEN-6A}/\texttt{OPEN-6C}) fix a concrete benchmark prediction.

    \item \textbf{Laboratory-Scale Frequency Shifts (Unified Clock Dictionary):} Using the consistent clock dictionary (zero double counting of $\mathcal{T}$),
    \[
    y_{\mathrm{clock}} = \alpha_\rho\,r \quad \left(r := \frac{\delta\rho_t}{\rho_0}\right),
    \qquad
    y_{\mathrm{clock}} = \alpha_\Phi\,p \quad \left(p := \frac{\Delta\Phi}{c^2},\;\alpha_\Phi = \alpha_\rho\mathcal{T}\right),
    \]
    for potential-driven configurations once $r = \mathcal{T} p$ and $\alpha_\rho$ are derived (\texttt{OPEN-6B}), while laboratory-prepared coherence states require their own separate preparation response (\texttt{OPEN-6C}).

    \item \textbf{Gravitational Lensing Deviations (Level-1 Exploratory):} Analyzing cluster-scale weak and strong lensing maps for systematic residuals relative to $\Lambda$CDM.

    \item \textbf{Decoherence-Environment Tests:} Searching for systematic shifts under controlled decoherence environments once $J_{\mathrm{prep}}$ is derived.

    \item \textbf{Bundle-Holonomy Phase Tests:} Precision Aharonov--Bohm and geometric-phase probes contingent on the visible $U(1)$ dictionary (\texttt{OPEN-6B}).
\end{enumerate}

In the uniform-background and heavy-scalar decoupling limits (such as \texttt{BM\_PHYS\_EFT\_1}, where scalar eigenmasses lie near $10^{14}\text{--}10^{16}\,\mathrm{GeV}$), the classical continuum action \texttt{M\_AUTH\_SIT\_EFT\_v1} reduces to General Relativity and standard Quantum Field Theory with exponentially suppressed short-range scalar exchanges. However, consistency of any specific parameter branch with all established experiments is \emph{not} guaranteed automatically by taking a decoupling limit; it must be verified channel by channel once microscopic material charges, the visible electromagnetic dictionary, environmental screening, and apparatus transfer functions are derived (\texttt{OPEN-6}).

\subsection{Exploratory Sensitivity Targets and Missing-Derivation Checklist (Not Derived Predictions)}

The canonical continuum action \texttt{M\_AUTH\_SIT\_EFT\_v1} does \emph{not} currently derive unconditional laboratory signal amplitudes for optical clocks, Bose--Einstein condensates, weak gravitational lensing, engineered decoherence-gravity couplings, or anomalous electromagnetic responses. Earlier draft revisions displayed speculative order-of-magnitude forecasts alongside universal ``falsifiability conditions''; those speculative amplitudes are removed here under \texttt{QG49} epistemic discipline. Table~\ref{tab:exploratory-missing-derivations} instead catalogs each prospective experimental channel alongside its current reference experimental precision, the exact missing microphysical derivation required before any quantitative prediction exists, and the strictly scoped falsification rule that governs hypothesis testing once a derivation is supplied.

\begin{table}[htbp]
\centering
\caption{Exploratory laboratory and astrophysical test channels for SIT, current reference experimental precision, missing action-to-apparatus derivations, and scoped exclusion rules. No entry in this table is a derived prediction of unchanged \texttt{M\_AUTH\_SIT\_EFT\_v1}.}
\label{tab:exploratory-missing-derivations}
\footnotesize
\begin{tabular}{|p{3.3cm}|p{2.9cm}|p{4.2cm}|p{4.1cm}|}
\hline
\textbf{Prospective Channel (Exploratory)} & \textbf{Reference Experimental Precision} & \textbf{Missing Derivation Required (\texttt{OPEN-6})} & \textbf{Scoped Exclusion Rule (Once Derived)} \\
\hline
\textbf{Optical atomic clock frequency shifts} (potential-driven or prepared states) & $\le 3.2\times 10^{-18}$ optical frequency-ratio precision (NIST/BACON, 2026) & Derive matter-to-clock coupling $\alpha_\rho$, potential transfer $\mathcal{T}$ (\texttt{OPEN-6B}), and laboratory preparation source $J_{\mathrm{prep}}$ (\texttt{OPEN-6C}) & A null result excludes only a specified action, parameter point, and preparation model whose derived shift exceeds experimental uncertainty \\
\hline
\textbf{Short-range gravity \& coherent quantum matter} (torsion balances, BECs, interferometers) & Sub-millimeter ISL bounds (E\"ot-Wash, Lee~2020) and atom interferometry ($\sim 10^{-12}$) & Derive macroscopic source charge $J_{\mathrm{src}}$ beyond pressureless-dust closure (\texttt{OPEN-6A}) and coherent-state source $J_{\mathrm{prep}}$ (\texttt{OPEN-6C}) & For a declared closure, null short-range results bound the $(\alpha_a, \lambda_a)$ parameter point (e.g.\ heavy \texttt{BM\_PHYS\_EFT\_1} is unconstrained; only explicit light-scalar parameter points are excluded) \\
\hline
\textbf{Weak and strong gravitational lensing residuals} (clusters, large-scale structure) & Percent-to-sub-percent cosmic shear residuals (current and Stage-IV surveys) & Derive cosmological background evolution, scalar perturbation growth, and non-linear lensing slip from \texttt{M\_AUTH\_SIT\_EFT\_v1} (\texttt{OPEN-6D}) & Excludes only a specified cosmological action extension and parameter point whose derived lensing residual contradicts observations \\
\hline
\textbf{Decoherence-dependent gravity tests} (optomechanics, macroscopic superpositions) & Precision differential gravimetry and optomechanical force sensing & Derive an explicit microphysical coupling between laboratory decoherence rates and $T_{\mu\nu}$ from the canonical action (\texttt{OPEN-6C}) & No quantitative exclusion exists until a concrete action-to-apparatus derivation maps decoherence protocols to $\delta g_{\mu\nu}$ \\
\hline
\textbf{Bundle-holonomy electromagnetic probes} (precision magnetometry, SQUIDs) & Precision magnetometry and Aharonov--Bohm phase interferometry & Establish the visible $U(1)_{\mathrm{EM}}$ identification and prove any non-standard correction beyond Maxwell holonomy (\texttt{OPEN-6B}) & Excludes only an explicit beyond-Maxwell extension once its derived field deviation exceeds magnetometer sensitivity \\
\hline
\end{tabular}
\end{table}

\noindent
\textbf{Scoped Falsification Principle:} A null experimental outcome never falsifies an entire conceptual framework in the abstract; it excludes only a \emph{specified action, parameter point, matter closure, and laboratory preparation model} after the corresponding observable prediction has been derived from that action and shown to exceed the apparatus's verified sensitivity. Conversely, where \texttt{BM\_PHYS\_EFT\_1} places the scalar eigenmasses near $10^{14}\text{--}10^{18}\,\mathrm{GeV}$ ($\lambda_a \sim 10^{-31}\,\mathrm{m}$), unchanged \texttt{M\_AUTH\_SIT\_EFT\_v1} predicts exponentially suppressed Yukawa deviations well below every current laboratory threshold.

\subsection{Observational Identifiability Cautions and Standing Implication Guardrails (Papers~002, 007, 014, and~055)}
\label{subsec:identifiability-and-qgtcd-guardrails}

Four methodological safeguards established in the July~2026 portfolio (Papers~002, 007, 014, and~055) govern all empirical interpretations and future extensions of Version~81:
\begin{enumerate}[leftmargin=2em]
  \item \textbf{Mass Calibration and Finite-Range Metric-Potential Ratios (Paper~002 \& Paper~007):} For finite-range Yukawa scalar corrections, a local ratio of metric potentials is \emph{not} automatically identical to the effective parameter fitted by a specific Solar System or laboratory experiment; light-propagation Shapiro observables and the Keplerian gravitational mass inferred from massive-body orbital dynamics must be calibrated jointly.
  \item \textbf{Source Form-Factor Absorption Bias (Paper~002 \& Paper~014):} Synthetic paired-fit benchmarks in Paper~002 demonstrate that inside a single fixed apparatus geometry, an incorrect finite-source form factor can be absorbed silently into a freely fitted coupling even when residuals look small. Multi-geometry or gradient-reversal controls and exact apparatus geometry modeling (\texttt{OPEN-6D}) are therefore mandatory before claiming a parameter fit.
  \item \textbf{Clock-Transition vs.\ Event-Count Observable Separation (Paper~007):} Atomic transition-frequency shifts ($y_{\mathrm{clock}} = \alpha_\rho r$, with gravitational potential transfer $\mathcal{T}$ entering strictly once via $\alpha_\Phi = \alpha_\rho\mathcal{T}$) and point-process event counts per unit proper time are distinct operational observables until an explicit calibrated detector model (accounting for exposure, back-action, correlations, and dead time) relates them.
  \item \textbf{Standing QGTCD Counterexample Guardrails (Paper~055):} Four bounded non-implication certificates from the \emph{QGTCD Constraints Note} (Paper~055) stand as permanent review filters against overclaiming microscopic derivations: (a)~a macroscopic drift-diffusion continuum limit does not uniquely identify a microscopic transition kernel; (b)~a probe-dependent overdamped drift law does not automatically establish universal gravitational free fall; (c)~additive velocity noise does not by itself establish thermal equilibrium; and (d)~matching a weak-field clock-rate relation does not determine a unique relativistic or quantum completion. Any proposed derivation crossing one of these boundaries must state explicitly which additional premise excludes the countermodel and whether that premise is derived or postulated.
\end{enumerate}

\subsection{Distinguishing SIT from Established Physics}

\begin{itemize}
    \item \textbf{Matching Established Physics:}  
    For uniform, maximally decohered systems, or in the heavy-scalar decoupling regime of \texttt{BM\_PHYS\_EFT\_1}, unchanged \texttt{M\_AUTH\_SIT\_EFT\_v1} reduces to General Relativity and standard Quantum Field Theory with exponentially suppressed scalar corrections.
    \item \textbf{Distinctive Prospective Regimes (Conditional on \texttt{OPEN-6}):}  
    Any measurable departure from standard physics would require either a light-mass scalar parameter branch coupled to laboratory matter (subject to existing torsion-balance, clock, and EP bounds) or an explicit, action-derived coupling in macroscopic quantum coherence, high-precision timekeeping, or cosmological backgrounds (\texttt{OPEN-6A}--\texttt{OPEN-6D}).
    \item \textbf{Scoped Testability and Parameter-Space Exclusion:}  
    Once a specific action extension, matter closure, and preparation protocol yield a quantitative observable prediction, precision experiments that fail to observe that derived effect exclude that specific model and parameter point.
\end{itemize}

\subsection{Summary}

Testing SIT empirically requires separating the machine-checked mathematical core of \texttt{M\_AUTH\_SIT\_EFT\_v1} from exploratory phenomenology. The Stage-6 action-to-propagator pipeline (\texttt{Stage6Addendum.\allowbreak{}ActionResponsePropagatorIdentity}) establishes the exact two-pole response structure for arbitrary canonical parameter points and source--detector charges, while \texttt{OPEN-6A}--\texttt{OPEN-6D} identify the exact microphysical derivations needed before optical clocks, quantum interferometers, gravitational lensing surveys, or precision magnetometers can test specific beyond-standard predictions.

\subsection{Exploratory Program: Testing Coherence-Conservation Hypotheses}
\label{sec:falsifiability_coherence_conservation}

At the conceptual and exploratory level (Tier-4 research hypotheses, not yet derived from unchanged \texttt{M\_AUTH\_SIT\_EFT\_v1}), Super Information Theory investigates whether a generalized coherence-conservation balance governs coupled quantum, gravitational, and biological/neural systems. These exploratory hypotheses motivate future theoretical derivation work rather than representing completed quantitative predictions.

\paragraph{Prospective Exploratory Research Channels}

Once explicit microscopic source and detector couplings (\texttt{OPEN-6B}--\texttt{OPEN-6D}) are derived, three exploratory avenues could be evaluated against experiment:

\begin{itemize}
    \item \textbf{Coherence-Dependent Gravitational Source Tests (\texttt{OPEN-6C}):} Investigating whether an action-derived coupling between macroscopic quantum coherence (e.g.\ phase-locked optical fields or Bose--Einstein condensates) and stress-energy produces a calculable gravitational shift beyond standard energy-momentum sourcing. Any quantitative test requires deriving $J_{\mathrm{prep}}$ and comparing the predicted amplitude for a specified parameter point against torsion-balance or atom-interferometry sensitivity.
    \item \textbf{Neural Coherence Budget Models (Phenomenological, Outside \texttt{M\_AUTH\_SIT\_EFT\_v1}):} Exploring whether information-theoretic coherence trade-offs across neural frequency bands or populations in EEG/MEG recordings obey constrained budget dynamics beyond conventional biophysical models. This domain is an independent phenomenological hypothesis and is not derived from the continuum gravitational EFT action.
    \item \textbf{Hybrid Quantum-Informational Protocols (\texttt{OPEN-6C}):} Examining whether engineered quantum-coherence redistribution exhibits lawful informational conservation constraints beyond standard open-quantum-system master equations once explicit interaction Hamiltonians are specified.
\end{itemize}

\paragraph{Scoped Empirical Evaluation Criteria}

Each exploratory hypothesis above can only be tested or excluded under strictly scoped conditions:
\begin{itemize}
    \item A concrete action, parameter point, matter closure, and preparation model must first yield a derived quantitative prediction exceeding experimental uncertainty.
    \item A null result at that derived sensitivity then excludes that specific model and parameter choice, without implying an unsupported claim about untested parameter regimes or unrelated sectors.
    \item In biological and neural applications, any putative coherence-budget signature must first be defined operationally and shown to differ quantitatively from conventional neurophysiological and statistical-mechanics baselines.
\end{itemize}

\paragraph{Summary}

Framed with explicit \texttt{OPEN-6} prerequisites and scoped exclusion rules, the coherence-conservation program serves as an honest roadmap connecting the machine-checked mathematical core of SIT to future action-derived laboratory and cross-domain tests.

\subsection{Core Experimental Tests in Physics}

\section{Conceptual Perspectives on Quantum Uncertainty and Decoherence}
\label{sec:quantum_reinterpret}

Consistent with Section~\ref{subsec:uncertainty-vs-coherence}, Version~81 separates established Hilbert-space quantum mechanics from the classical $U(1)$ Noether current of \texttt{M\_AUTH\_SIT\_EFT\_v1}: neither the Heisenberg--Robertson uncertainty relation ($\Delta x\,\Delta p \ge \hbar/2$) nor wave--particle complementarity is derived from a universal conserved coherence budget (which is refuted by Hadamard basis rotations and the maximally mixed state $I/2$).

\subsection{Variance Uncertainty vs.\ Operational Basis Coherence}

In standard quantum theory, non-commutativity of conjugate operators ($[\hat{x}, \hat{p}] = i\hbar$) enforces variance and entropic uncertainty relations across incompatible measurement bases. SIT adopts these established quantum-information trade-offs as conceptual inspiration when comparing localized versus delocalized states, without claiming that a single fixed-basis coherence functional is conserved across complementary bases or across neural frequency bands.

\subsection{Environment-Induced Decoherence and Fringe Visibility}

Similarly, the transition from high-visibility interference patterns to classical trajectory statistics is described by standard environment-induced decoherence (exporting phase correlations into unobserved environmental degrees of freedom, governed by unital purity non-increase or general relative-entropy contraction in Result~\ref{thm:CoherenceMonotonicity}). Connecting interferometric fringe visibility to classical scalar perturbations $(\delta\psi, \delta\rho_t)$ remains an open state-to-field mapping problem (\texttt{OPEN-6C}).

\section{Illustrative Radial Green--Function Visualization (Decoupled One-Way Toy Model)}
\label{sec:sphericalresonancesrevisited}

\textbf{Model-Relation Qualification:} In the canonical quadratic action of \texttt{M\_AUTH\_SIT\_EFT\_v1} (Appendix~E and Stage~6), the static two-scalar response $(K\nabla^2 - H)\,\delta\Phi = J_{\mathrm{src}}$ carries a \emph{symmetric, reciprocal} off-diagonal portal Hessian ($H_{12} = H_{21} = -2 m_{\mathrm{link}}\sqrt{\lambda_{\mathrm{portal}}}\,v_0$), yielding a two-pole Yukawa superposition governed by $(p^2 K + H)^{-1} = V\,\mathrm{diag}((p^2+m_a^2)^{-1})\,V^\top$ rather than a triangular one-way system. By contrast, to visualize qualitatively how a static Yukawa screening profile behaves in three spatial dimensions, earlier exploratory notes considered a simplified \emph{illustrative one-way forced system} (valid only as a toy pedagogical approximation when back-reaction from $\delta\Rcoh$ onto $\delta\rho_t$ is artificially neglected):
\begin{align}
   \bigl(\nabla^2-\mu_t^2\bigr)\,\delta\rho_t(r) &= -\,4\pi G\,\delta
   m\,\delta^{(3)}(\boldsymbol{r}),\label{eq:greens-rhot}\\[4pt]
   \bigl(\nabla^2-\mu_R^2\bigr)\,\delta R_{\text{coh}}(r) &= -\,\gamma\,
   \delta\rho_t(r),\label{eq:greens-R}
\end{align}
with toy inverse ranges $\mu_t, \mu_R > 0$. Its static, spherically symmetric Green functions are
\[
   \delta\rho_t(r)=\frac{G\,\delta m}{r}\,e^{-\mu_t r},
   \qquad
   \delta R_{\text{coh}}(r)=
   \gamma\,\frac{G\,\delta m}{r}\,
   \frac{e^{-\mu_t r}-e^{-\mu_R r}}{\mu_R^2-\mu_t^2}
   \quad (\mu_R \ne \mu_t),
\]
illustrating ordinary three-dimensional $e^{-\mu r}/r$ Yukawa decay (reducing to $1/r$ as $\mu_t \to 0$).

\paragraph{Relation to Canonical Action and Experimental Bounds (Retiring Legacy $\beta < 10^{-5}$ Claim).}
Earlier drafts used Eq.~\eqref{eq:greens-rhot} to assert a blanket torsion-balance bound $\beta < 10^{-5}$ for $\mu_t^{-1} \gtrsim 0.1\,\mathrm{m}$ without a primary experimental provenance record or action-to-detector map. Version~81 retires that ungrounded bound in favor of the canonical two-mode propagator $(p^2 K + H)^{-1}$ and conditional weak-field E\"ot-Wash analysis of Section~\ref{sec:stage6-action-diagnostics} (where \texttt{BM\_PHYS\_EFT\_1} has $\lambda_1 \approx 8.056\times 10^{-31}\,\mathrm{m}$ and full patterned rotating-mass likelihood reproduction from Lee's dissertation Appendix~A remains \texttt{OPEN-6A}/\texttt{OPEN-6D}).

\paragraph{Yukawa Screening Length (Not a Causal Horizon or Tunneling Derivation).}
At radii $r \gg \mu_t^{-1}$, the static Yukawa factor $e^{-\mu_t r} \ll 1$ simply reflects ordinary exponential Yukawa screening of a massive scalar propagator. Earlier drafts referred to $r \approx \mu_t^{-1}$ as an ``informational horizon'' governing quantum tunneling amplitudes; Version~81 retracts that identification because a static screened Green function defines a classical spatial screening length (Compton wavelength), \emph{not} a relativistic causal horizon and \emph{not} a derivation of quantum tunneling probabilities.

\section{Open Questions and Future Research Directions}
\label{sec:open-questions-future-directions}

Super Information Theory presents several profound open questions, spanning theoretical, empirical, and philosophical domains. Here we clearly outline these challenges and propose specific future research directions. Verifying SIT demands multiscale simulations that cover twenty-five orders of magnitude in length and fifteen in time.  Adaptive-mesh relativistic codes accelerated on GPU clusters are therefore mandatory, as are quantum-inspired algorithms that compress phase-space evolution into tractable tensor networks.  On the experimental side the key limitation is precision: optical-lattice clocks must push below the $10^{-19}$ stability frontier, interferometers must resolve sub-milliradian phase drifts and lensing surveys must reduce systematics to parts per thousand.  All three goals are technologically plausible within the coming decade but require coordinated international effort.  Finally, because SIT is irreducibly interdisciplinary, sustained collaboration among metrologists, neuroscientists, cosmologists and computer scientists is essential; open-source tool-chains and shared data standards are already being drafted to make that collaboration routine.

\subsection{Empirical Validation and Experimental Challenges}

Direct empirical verification is crucial for establishing SIT's validity:
\begin{itemize}
    \item Conducting high-precision laboratory tests of coherence-induced gravitational phenomena, using advanced atomic clock arrays and cold-atom interferometry.
    \item Utilizing astrophysical observations---such as precise gravitational lensing, galaxy rotation curves, and cosmic microwave background anisotropies---to differentiate SIT from standard cosmological models and alternative theories (MOND, dark energy).
    \item Developing precise quantum computational simulations to systematically test SIT's coherence--time-density predictions under controlled experimental conditions.
\end{itemize}

\subsection{Philosophical and Interdisciplinary Investigations}

SIT's informational ontology raises significant interdisciplinary questions:
\begin{itemize}
    \item Examining how an informational foundation for reality impacts philosophical debates on causality, determinism, free will, and agency.
    \item Developing cognitive and neuroscientific experiments---leveraging advanced neuroimaging and brain--computer interfaces---to empirically test SIT's coherence-based model of consciousness.
    \item Exploring practical implications of SIT principles for future artificial intelligence systems, particularly their coherent informational synchronization capabilities and ethical considerations.
\end{itemize}

\subsection{Quantum Computational Modelling and Simulations}

A promising avenue for theoretical validation involves quantum computational modelling:
\begin{itemize}
    \item Leveraging quantum computational simulations to rigorously test SIT's predictions about coherence--decoherence dynamics and gravitational interactions at scales inaccessible to classical computational methods.
    \item Exploring hybrid quantum--classical computational frameworks to better characterize SIT's unique quantum-gravitational predictions.
    \item Utilizing AI-driven simulation platforms to refine coherence field models, improving theoretical accuracy and guiding future experimental setups.
\end{itemize}

\subsection{Statistical and Methodological Rigor}

Future research must emphasize rigorous statistical methodologies:
\begin{itemize}
    \item Employing Bayesian inference, adaptive filtering, and rigorous hypothesis testing to clearly differentiate SIT-predicted phenomena from experimental noise or systematic error.
    \item Cross-validation across multiple experimental modalities---atomic clocks, interferometry, lensing, and cosmological surveys---to robustly confirm SIT predictions.
\end{itemize}

These explicit directions ensure Super Information Theory remains empirically rigorous, theoretically precise, and methodologically transparent, positioning it clearly for interdisciplinary engagement and robust scientific validation.

\subsection{Interdisciplinary Bridges}

SIT explicitly promotes interdisciplinary collaborations, integrating quantum physics, neuroscience, and AI research through:

\begin{itemize}
\item Detailed explorations of quantum coherence's role in macroscopic gravitational phenomena.
\item Neuroscientific parallels explicitly established between neural synchronization and quantum informational coherence.
\item AI frameworks explicitly modeled upon natural coherence-driven adaptive informational systems.
\end{itemize}

\subsection{Future Directions and Open Challenges}

Future directions for SIT are explicitly defined through mathematical formalization, empirical validation, and interdisciplinary research:

\begin{itemize}
\item Rigorous PDE/Lagrangian formalisms explicitly linking quantum coherence and time-density fields.
\item Precision empirical tests explicitly designed, including atomic clocks, gravitational lensing, and quantum interference protocols.
\item Cross-disciplinary collaborations explicitly aimed at coherent informational modeling in neural, astrophysical, and cognitive contexts.
\end{itemize}

\subsection{Philosophical Outlook}

Because coherence and decoherence are linked by the Noether-conserved information current \(J_{\mathrm{coh}}^{\mu}\), neither branch is metaphysically privileged; reality is the standing wave between them.  Consciousness, in this reading, is an emergent watermark where biological free-energy minimization intersects the cosmic least-mismatch drive.
% --- END INPUT sections/05_part2_sec26_to_31.tex ---
% --- BEGIN INPUT sections/06_part3_sec32_to_42.tex ---
\part{Interdisciplinary Synthesis, Neuroscience, and Cosmological Frontiers}

\section{Conceptual and Interdisciplinary Implications}
\label{sec:conceptual-implications}

\subsection{Cognitive and Neural Consequences}

By framing information as physical coherence, SIT directly connects fundamental physics to the mechanisms of cognition and consciousness:

\begin{itemize}
    \item \textbf{Oscillatory Binding and Perceptual Integration:}  
    The neural process of binding disparate signals into unified perception is realized as a local increase in $R_{\mathrm{coh}}$, with wave-based signal dissipation (Section~\ref{sec:thermodynamics}) driving the system toward coherent, conscious states.
    \item \textbf{Predictive Coding and Active Inference:}  
    SIT's thermodynamic/computational framework aligns with the Free Energy Principle in neuroscience, suggesting that brains minimize prediction error by iteratively reducing local incoherence---mirroring the minimization of informational entropy.
    \item \textbf{Memory, Learning, and Agency:}  
    Network-level plasticity and learning correspond to the creation and maintenance of stable high-coherence attractor states---encoding information not just in static synaptic weights, but in dynamical, oscillatory patterns.
\end{itemize}

\subsection{Information-Theoretic and Computational Implications}

SIT provides a unified physical foundation for concepts from quantum information, classical information theory, and computation:

\begin{itemize}
    \item \textbf{Information as an Attractor:}  
    The tendency of both quantum and classical systems to evolve toward stable, low-entropy (high-coherence) states explains the emergence of structure, memory, and agency in natural and artificial systems.
    \item \textbf{Computational Thermodynamics:}  
    The signal-dissipation laws (Section~\ref{sec:thermodynamics}) provide a stepwise, local, and computable model for the approach to equilibrium, reconciling the ``arrow of time'' in thermodynamics with quantum reversibility.
    \item \textbf{Measurement as Coherence Redistribution:}  
    Both quantum measurement and neural information extraction become special cases of coherence redistribution---quantitatively described by changes in $R_{\mathrm{coh}}$ and associated entropy flow.
\end{itemize}

\subsection{Broader Scientific and Philosophical Consequences}

SIT's mathematically anchored unification of coherence, information, and dynamics opens new avenues for interdisciplinary science:

\begin{itemize}
    \item \textbf{Physics:}  
    SIT suggests that gravitational, electromagnetic, and kinetic phenomena are limiting cases of informational dynamics---implying new routes to quantum gravity and a ``computation-first'' ontology.
    \item \textbf{Neuroscience and Cognitive Science:}  
    By rendering subjective phenomena (such as awareness or agency) in terms of physically measurable coherence, SIT enables rigorous empirical studies of consciousness and information integration.
    \item \textbf{Artificial Intelligence and Technology:}  
    Technologies inspired by SIT may use engineered coherence (in quantum computers, neuromorphic chips, or network architectures) to realize new forms of efficient, adaptive computation.
    \item \textbf{Philosophy of Science:}  
    The theory reframes ``emergence'' as the statistical consequence of mechanistic, computable micro-dynamics---bridging the gap between reductionist physics and holistic, emergent phenomena.
\end{itemize}

\subsection{Summary}

By grounding interdisciplinary implications in rigorous mathematics and operational definitions, SIT provides a robust framework for uniting physical, biological, computational, and philosophical perspectives. The theory clarifies the nature of information as a physical, measurable property, and offers a new lens on the emergence of structure, cognition, and agency in the universe.

\section{Information as an Organizing Attractor (Conceptual Program and Lyapunov Counterexamples)}

At the conceptual level, SIT asks when dissipative or open-system informational dynamics drive states toward coherent attractors subject to energy and environmental constraints. However, Version~81 explicitly corrects two mathematical errors in earlier drafts of this section: the false claim of a universal Lyapunov functional for every solution of the reversible canonical field equations, and the conflation of amplitude $R$ with squared density $R^2$ in the continuity equation.

\subsection{Why $\mathcal{H}[R] = -\int R\ln R$ Is Not a Universal Lyapunov Functional}

In \texttt{M\_AUTH\_SIT\_EFT\_v1}, the conserved global $U(1)$ Noether current for $\psi = |\psi|e^{i\theta}$ is $J^\mu_{\mathrm{coh}} = \kappa_c|\psi|^2(\nabla^\mu\theta - \hat{q}A^\mu)$ (Section~\ref{sec:coherence-conservation}), whose classical field equations are Hamiltonian and time-reversal symmetric. Earlier drafts defined
\begin{equation}
\mathcal{H}[R_{\mathrm{coh}}]
   \;=\;
   -\!\int R_{\mathrm{coh}}\ln R_{\mathrm{coh}}\,
          \sqrt{-g}\,d^{3}x
\tag{A2}\label{A2}
\end{equation}
and claimed $(d\mathcal{H}/dt) \le 0$ for every solution of the field equations. Two explicit counterexamples refute that universal Lyapunov claim:
\begin{enumerate}[leftmargin=2em]
  \item \textbf{Opposite Sign Under Diffusion ($\partial_t R = D\Delta R$):} Even for the manuscript's own phenomenological diffusion model $\partial_t R = D\Delta R$ ($D > 0$) on a periodic spatial domain with $R > 0$, differentiating~\eqref{A2} and integrating by parts yields
  \[
  \frac{d\mathcal{H}}{dt}
  = -\int (\ln R + 1)\,D\Delta R\,d^3x
  = +D\int \frac{|\nabla R|^2}{R}\,d^3x \ge 0,
  \]
  which is strictly \emph{increasing} (not non-increasing) for every nonconstant positive solution! (By contrast, the opposite-sign functional $+\int R\ln R\,d^3x$ or the Dirichlet energy $\frac{1}{2}\int|\nabla R|^2\,d^3x$ decreases monotonically under pure diffusion with appropriate boundary conditions.)
  \item \textbf{Counterexample in Reversible Canonical Scalar--Gravity Dynamics:} For admissible spatially homogeneous initial data of the reversible canonical scalar--gravity system (with strictly positive scalar mass matrix) initialized at $R_{\mathrm{coh}} = 1$ with $\dot{R}_{\mathrm{coh}} = -\varepsilon$ ($\varepsilon > 0$), $\frac{d}{dt}(-R_{\mathrm{coh}}\ln R_{\mathrm{coh}}) = -(\ln R_{\mathrm{coh}} + 1)\dot{R}_{\mathrm{coh}} = +\varepsilon > 0$ per unit spatial volume, directly refuting $d\mathcal{H}/dt \le 0$ across all solutions.
\end{enumerate}
Accordingly, the organizing-attractor proposal is retained strictly as a research direction requiring a \emph{specified dissipative reduced model} (declaring the dissipative dynamics, functional sign, normalization, boundary conditions, and environment explicitly before proving any Lyapunov decay theorem).

\subsection{Quantum Scale: Squared-Density Continuity vs.\ Amplitude}

For a nonrelativistic complex field $\psi = |\psi|e^{i\theta}$ evolving under the standard Schr\"odinger equation, the exact Madelung/Noether continuity equation conserves the \emph{squared modulus} (density) $|\psi|^2 = R^2$ rather than the bare amplitude $R$:
\begin{equation}
\partial_{t}\bigl(R^2\bigr)
   +\frac{\hbar}{m}
     \nabla\!\cdot\!\bigl(R^2\nabla\theta\bigr)
   = 0
\tag{A3}\label{A3}
\end{equation}
(earlier drafts wrote $R$ linearly in place of $R^2$). Laser cooling, phonon condensation, and superconducting interiors serve as qualitative physical analogies for low-dissipation attractors in open quantum systems (Tier-4 conceptual heuristics, not derived predictions of \texttt{M\_AUTH\_SIT\_EFT\_v1}). Whether a concrete action-derived phase-diffusion bound emerges in narrow-line optical clocks remains open pending \texttt{OPEN-6B}/\texttt{OPEN-6C}.

\subsection{Gravitational Scale: Heuristic Curvature Analogy (Tier-4 Conceptual Sketch)}

In an exploratory heuristic model writing $\rho_{t}= \rho_{0}\exp(\alpha R_{\mathrm{coh}})$ (distinct from the canonical quadratic potential of \texttt{M\_AUTH\_SIT\_EFT\_v1}), non-minimal curvature-gradient terms take the schematic form
\[
G_{\mu\nu}=8\pi G\,T_{\mu\nu}
+\alpha\,\bigl(\nabla_{\!\mu}\nabla_{\!\nu}-g_{\mu\nu}\Box\bigr)R_{\mathrm{coh}} .
\tag{A4}\label{A4}
\]
When $\nabla_{\mu}R_{\mathrm{coh}}=0$, Eq.~\eqref{A4} reduces to Einstein gravity with an effective constant shift. No quantitative weak-lensing signal amplitude across galaxy cluster outskirts is derived from unchanged \texttt{M\_AUTH\_SIT\_EFT\_v1} (earlier exploratory $\delta R_{\mathrm{coh}}\sim 10^{-3}$ sketches are retired pending a full cosmological derivation under \texttt{OPEN-6D}).

\subsection{Neural Scale: Predictive Synchronisation (Phenomenological Analogy)}

For a finite network of $N$ oscillatory degrees of freedom, phenomenological phase dynamics may be modeled by a Kuramoto coupling analog,
\[
\dot{\theta}_{i}
   = \omega_{i}
   - \frac{\kappa}{N}\sum_{j} 
     R_{ij}\,\sin(\theta_{i}-\theta_{j}),
\tag{A5}\label{A5}
\]
with symmetric coupling weights $R_{ij} = R_{ji}$ (distinct from the classical EFT scalar $\Rcoh = |\psi|/\psi_0$). Earlier drafts asserted that the global order parameter $r e^{i\Psi} = N^{-1}\sum_j e^{i\theta_j}$ satisfies $\dot{r} = -\partial\mathcal{H}_{\mathrm{net}}/\partial\Psi$; Version~81 retracts that identity because any phase-shift-invariant interaction potential $\mathcal{V}(\boldsymbol{\theta}) = -\frac{\kappa}{2N}\sum_{i,j} R_{ij}\cos(\theta_i - \theta_j)$ is invariant under uniform shifts $\theta_j \mapsto \theta_j + \alpha$ ($\Psi \mapsto \Psi + \alpha$), so $\partial\mathcal{V}/\partial\Psi \equiv 0$ identically! Instead, for identical frequencies ($\omega_i = \omega$) the exact gradient flow acts in the $N$-angle coordinates via $\dot{\theta}_i = \omega - \partial_{\theta_i}\mathcal{V}$, and in the uniform-coupling limit ($R_{ij} \equiv 1$, where $\mathcal{V} = -\frac{\kappa N}{2} r^2$) direct differentiation of $r e^{i\Psi} = N^{-1}\sum_j e^{i\theta_j}$ gives
\begin{equation}
\dot{r} = +\frac{\kappa r}{N}\sum_{j=1}^N \sin^2(\theta_j - \Psi) \ge 0
\label{eq:kuramoto_order_param_corrected}
\end{equation}
(whereas for spatially varying $R_{ij}$, the dynamics does not close in $(r, \Psi)$ alone). This phenomenological neuroscience model is explicitly distinct from the continuum gravitational EFT \texttt{M\_AUTH\_SIT\_EFT\_v1}.

\subsection{Technological and Cosmological Cascades (Exploratory Scaling Hypothesis)}

In an exploratory scale-free ansatz (not derived from the massive portal spectrum of \texttt{M\_AUTH\_SIT\_EFT\_v1}), a putative scaling regime with anomalous dimension $\eta$ would yield a power-law two-point spectrum
\[
\langle R_{\mathrm{coh}}(k)R_{\mathrm{coh}}(-k)\rangle 
   \;\sim\; k^{\eta-3}.
\tag{A6}\label{A6}
\]
Determining whether any concrete renormalization-group fixed point exists for a specified SIT action remains an open theoretical problem rather than an established prediction.

\subsection{Exploratory Program Checklist (Not Derived Predictions)}

\begin{enumerate}[label=\textbullet,leftmargin=1.8em]
\item \emph{Atomic clocks}: derive the visible $U(1)$ dictionary and driving response (\texttt{OPEN-6B}/\texttt{OPEN-6C}) before comparing against optical-clock stability data.  
\item \emph{Weak lensing}: derive cosmological background and shear-perturbation equations from the action (\texttt{OPEN-6D}) before quoting cluster-outskirt residuals.  
\item \emph{MEG coherence}: evaluate phenomenological Kuramoto-type budget models against conventional neurophysiological baselines independently of gravitational EFT.  
\item \emph{Complex networks}: treat graph-scaling analogies separately from physical spacetime field theory.
\end{enumerate}

These exploratory cross-domain analogies illustrate the conceptual origin of SIT while maintaining a strict \texttt{QG49} boundary separating them from the derived theorems of \texttt{M\_AUTH\_SIT\_EFT\_v1}.

\section{Conceptual Cross-Scale Metaphor: From Planetary Accretion to Morphogenetic Repair}
\label{sec:evolving-configuration}

At the level of conceptual motivation (distinct from the literal coupled field equations of \texttt{M\_AUTH\_SIT\_EFT\_v1}), earlier SIT narratives sketched a damped phenomenological relaxation metaphor,
\[
\Box R_{\text{coh}}^{\mathrm{eff}}+\gamma\,\partial_{t}R_{\text{coh}}^{\mathrm{eff}}+\mu_{R}^{2}R_{\text{coh}}^{\mathrm{eff}}=-\frac{\delta\mathcal{H}_{\mathrm{eff}}}{\delta R_{\text{coh}}^{\mathrm{eff}}},
\qquad
\Box\rho_{t}^{\mathrm{eff}}+\mu_{t}^{2}\rho_{t}^{\mathrm{eff}}=-\alpha\,\mu_{R}^{2}R_{\text{coh}}^{\mathrm{eff}},
\]
to compare dissipative self-organization motifs across astrophysics, morphogenesis, and computation. Version~81 explicitly marks this section as a heuristic cross-domain metaphor rather than claiming that planets, biological organisms, and software execute a single canonical PDE derived from $S_{\mathrm{AUTH}}$.

\subsection*{Cosmic accretion as large--scale phase locking}

In the protoplanetary nebula collisional damping reduces random velocities and allows neighbouring mass shells to share orbital phase.  The coherence ratio therefore climbs, $\partial_{t}R_{\text{coh}}>0$, raising the local time density.  The augmented $\rho_{t}$ slows orbital clock rate in that annulus, trapping additional dust and reinforcing the phase lock.  The feedback terminates only when the gradient energy in $R_{\text{coh}}$ balances radiative losses, at which point a planetesimal has formed.  Subsequent migration of $R_{\text{coh}}$ fronts explains why resonant orbits cluster in harmonic chains: adjacent wells fall into phase to minimise interfacial tension in $\partial_{r}R_{\text{coh}}$.

\subsection*{Morphogenetic repair as mesoscopic error correction}

Michael Levin's experiments reveal that amputated planarian fragments restore missing heads by guiding voltage, calcium and transcriptional waves along precise pathways.  SIT interprets those waves as biochemical carriers of $R_{\text{coh}}$.  Injury lowers coherence at the cut site; the surrounding tissue responds by diffusing $R_{\text{coh}}$ inward until the gradient vanishes.  Cells decode the local value of $\rho_{t}=\rho_{0}\exp[\alpha R_{\text{coh}}]$; zones of high time density divide more slowly, giving stem cells the chronological window required to differentiate properly.  

A natural question is whether this heuristic analogy suggests ``differential aging'' effects in ordinary organisms. Even within speculative scalar-coupling analogies, dense macroscopic media would lie deep inside the screened regime (Section~\ref{sec:screening-mechanism}), and consistent with Section~\ref{sec:neuroscience-implications} below, Version~81 makes \emph{no} claim that planarian bioelectric signaling or developmental timing is governed by the relativistic spacetime scalars $(\psi, \rho_t)$ of \texttt{M\_AUTH\_SIT\_EFT\_v1}.

In Michael Levin's bioelectric morphogenesis framework itself, target morphology acts as a biochemical and bioelectric attractor across tissue scales; SIT cites that phenomenon purely as an inspiring cross-domain self-organization analogy.

\subsection*{Thought and technology as iterative phase refinement}

At the level of conceptual analogy, both cortical assemblies and collaborative engineering workflows refine noisy candidates into coherent representations through iterative feedback and pruning, without implying a shared physical field equation.

\subsection*{Conceptual summary of cross-scale motifs}

Dust-cloud damping, bioelectric tissue patterning, and neural synchronization illustrate recurring dissipative self-organization motifs across independent scientific disciplines; in Version~81, those motifs serve strictly as conceptual inspiration rather than literal solutions of $S_{\mathrm{AUTH}}$.

\section{Integrative Feedback Loops and Self-Referential Organisation (Prospective Analogies)}
\label{sec:integrative-feedback}

At the level of cybernetic and systems-neuroscience analogy (kept strictly separate from the canonical gravitational-gauge EFT \texttt{M\_AUTH\_SIT\_EFT\_v1}, Section~\ref{sec:neuroscience-implications}), integrative feedback loops provide a useful conceptual motif for modeling adaptive open systems. As a \emph{proposed prospective experimental protocol} (rather than an already-completed validation of SIT), one can design closed-loop human--AI virtual-reality neurofeedback studies in which high-density EEG from immersed participants is streamed in real time to a phase-tracking model that adapts visual and auditory cues to test whether closed-loop sensory pacing modulates operational neural synchrony metrics ($R_{\mathrm{coh}}^{\mathrm{neural}}$, Section~\ref{sec:SITneuroPrediction}). Similarly, earlier drafts wrote a schematic drift--diffusion equation
\[
\partial_t u_{\mathrm{eff}}
   \;=\;
   D\nabla^{2} u_{\mathrm{eff}}
   - \nabla\!\cdot\!\bigl(u_{\mathrm{eff}}\nabla\Phi_{\mathrm{eff}}\bigr)
\]
and asserted that ``every physical system obeys the same schematic evolution law'' from trapped ions to neural synchrony and galaxy clusters. Version~81 explicitly \emph{retracts} that universal-law assertion: closed quantum systems evolve unitarily, \texttt{M\_AUTH\_SIT\_EFT\_v1} is a second-order hyperbolic Lorentzian field theory rather than a parabolic drift--diffusion equation, and cortical neural dynamics are modeled on their own independent neurophysiological track without identifying EEG/LFP synchrony with the canonical spacetime fields $(\psi, \rho_t)$.

\section{Implications for Neuroscience, Cognition, and Consciousness (Independent Empirical Track)}
\label{sec:neuroscience-implications}
\label{sec:neuroscience-cognition-implications}

Super Information Theory shares historical motivation with computational neuroscience, where phase-synchronized oscillations and coincidence detection play central functional roles. Consistent with the operational separation contract (Section~\ref{sec:PrimitiveFieldDefinitions}) and the independent neuroscience track of \emph{Learning to Feel} (Paper~114, v1.2), Version~81 explicitly separates operational neurophysiological variables from the spacetime scalar fields $(\psi, \rho_t)$ of \texttt{M\_AUTH\_SIT\_EFT\_v1}: no derivation currently identifies cortical EEG/LFP phase synchrony with the canonical gravitational/gauge fields or their matter-source terms.

\subsection{Neuroscience-Inspired Predictive Synchronization (Analogical Framing)}

Predictive coding, active inference, and the Free Energy Principle motivate structural analogies with phase-synchronized oscillatory networks. Rather than claiming that the brain is a mesoscopic limit of the gravitational EFT \texttt{M\_AUTH\_SIT\_EFT\_v1}, we treat cortical phase synchronization strictly as an independent phenomenological modeling domain using operational neural variables.

\subsection{Neural Dynamics and Operational Analytic-Signal Variables}

Drawing on Self Aware Networks, neural oscillations, phase locking, and network synchrony can be parameterized operationally by a complex analytic-signal envelope (e.g.\ via Hilbert transform of band-limited LFP/EEG activity),
\begin{equation}
Z_{\mathrm{neural}}(\mathbf{x}, t) := A_{\mathrm{neural}}(\mathbf{x}, t)\,e^{i\phi_{\mathrm{neural}}(\mathbf{x}, t)},
\label{eq:neural_analytic_signal}
\end{equation}
with real nonnegative amplitude $A_{\mathrm{neural}} \ge 0$ and oscillation phase $\phi_{\mathrm{neural}}$, kept explicitly distinct from the spacetime field $\psi(x) = |\psi(x)|e^{i\theta(x)}$ and its dimensionless ratio $\Rcoh = |\psi|/\psi_0$.

\subsection{Neural Coincidence Detection and Correction of Earlier Field-Current Conflation}
\label{sec:neural_coherence}

Pyramidal neurons exhibit coincidence sensitivity when excitatory postsynaptic potentials arrive within a short temporal window ($1\text{--}5\,\mathrm{ms}$) relative to local beta--gamma cycles, producing graded membrane depolarizations, $\mathrm{Ca}^{2+}$ influx, and continuous spike-timing/burst modulation rather than idealized noiseless Boolean bits. That neurophysiological observation motivates the conceptual slogan \emph{Bit from Coincidence, Wave from Bit} for modeling cortical signaling and synaptic plasticity.

\paragraph{Explicit Mathematical Corrections to Earlier Drafts.}
Version~81 corrects three inconsistencies present in earlier revisions of this subsection:
\begin{enumerate}[leftmargin=2em]
  \item \textbf{Real Nonnegative $\Rcoh$ Has No Nonconstant Argument:} Earlier drafts wrote a current $J_\mu^{\mathrm{coh}} = \kappa_c \Rcoh\,\partial_\mu\arg\Rcoh$. Because $\Rcoh = |\psi|/\psi_0 \ge 0$ is real and nonnegative, $\arg\Rcoh = 0$ wherever $\Rcoh > 0$ and is undefined at $\Rcoh = 0$; that expression does not generate a nontrivial current. In \texttt{M\_AUTH\_SIT\_EFT\_v1}, the canonical $U(1)$ Noether current depends on the phase $\theta = \arg\psi$ of the complex field $\psi$ and the independent gauge connection $A_\mu$: $J^\mu_{\mathrm{coh}} = \kappa_c|\psi|^2(\nabla^\mu\theta - \hat{q}A^\mu)$ (Eq.~\eqref{eq:abstract_coherence_current}).
  \item \textbf{Amplitude $|\psi|$ Enters Multiple Action Terms:} Contrary to earlier wording stating that $\Rcoh$ enters the master action only through $J^\mu_{\mathrm{coh}}$, the modulus $|\psi| = \psi_0\Rcoh$ enters \texttt{M\_AUTH\_SIT\_EFT\_v1} directly through its radial gradient energy $-\frac{\kappa_c}{2}(\partial|\psi|)^2$, the potential $V_{\mathrm{tot}}(\rho_t, |\psi|)$, and the dressings $f_{\mathrm{EM}}(\Rcoh)$ and $f_{\mathrm{SM}}(\Rcoh)$.
  \item \textbf{Separation of Neural Activity from Spacetime Dipoles and Gravitational RG Flow:} Version~81 retracts the claims that cortical beta-burst coincidence currents act as gravitational ``time-density dipoles'' warping spacetime phase space or that synaptic long-term potentiation/depression (LTP/LTD) obeys the same renormalization-group flow as gravitational couplings. Simply substituting $\arg\psi$ for $\arg\Rcoh$ would not bridge cortical neurophysiology to \texttt{M\_AUTH\_SIT\_EFT\_v1}; any future physical connection would require defining a concrete observable, units, coupling dynamics, and independently testable consequences beyond established bioelectromagnetism.
\end{enumerate}

\subsection*{Consequences for Learning and Prediction (Empirical Neuroscience Track)}

Within empirical neuroscience alone (without invoking gravitational EFT), coincidence-window timing, suprathreshold depolarization duration, and phase-dependent plasticity provide testable neurophysiological hypotheses (developed independently in \emph{Learning to Feel}, Paper~114 v1.2, using calibrated comparators, held-out interventions, and receiving-influence measurements).

As an exploratory cross-scale analogy (not derived from \texttt{M\_AUTH\_SIT\_EFT\_v1}), one may compare synchronization clustering in neural networks with structure formation in physics while keeping their governing equations and empirical tests strictly decoupled.

\subsection{Emergence of Consciousness and Self-Awareness}

Within SIT, consciousness and agency are modeled heuristically as high-order informational phenomena arising from globally synchronized coherence among neural networks. Subjective experience and cognition are modeled as time-dependent informational wave patterns stabilized by neural phase synchrony. This extends the Self Aware Networks theory of consciousness (ToC), situating the brain's informational cycles within the larger framework of universal coherence dynamics. Notably, SIT hypothesizes that transient increases in neural phase coherence correlate with spikes in informational and temporal density---potentially manifesting as moments of heightened self-awareness, insight, or conscious unity.

Moreover, SIT suggests that the brain's oscillatory organization may form complex, quasi-crystalline patterns, encoding informational complexity and adaptability. This analogy extends naturally to artificial intelligence, where coherence-driven network synchrony could underpin computational forms of cognition, bridging biological and artificial consciousness within a unified informational logic.

\subsection{Prospective Experimental Connections: EEG/MEG and VR/AR/BCI}

At the phenomenological neural-modeling level, SIT proposes that experimentally induced changes in neural synchrony, measured via EEG or MEG, may exhibit quantifiable correlates in cognitive performance and task-related awareness. Immersive AR/VR and brain--computer interface (BCI) setups offer motivating observations and prospective experimental testbeds for probing these proposed coherence-synchronization analogies rather than completed empirical validations of SIT:
\begin{itemize}
\item Motivating neurofeedback paradigms exploring whether closed-loop phase synchronization correlates with cognitive task performance.
\item Prospective real-time neural phase-tracking experiments within immersive VR/AR and BCI platforms designed to test proposed phenomenological coherence models.
\end{itemize}

\subsection{Summary: Informational Coherence Across Scales}

Super Information Theory outlines a conceptual framework linking quantum, gravitational, and neural synchronization analogies under a common informational vocabulary. We emphasize that biological and cognitive applications in Part~III are motivating analogies and proposed phenomenological models, kept strictly separate from the machine-checked spacetime EFT derivations of Parts~I, II, and~IV.

\subsection{Neural Oscillations and Predictive Coding}

SIT coherence dynamics directly analogize neural phase synchronization mechanisms observed in predictive coding frameworks (Friston's Free Energy Principle). Neural synchronization corresponds to informational coherence, and predictive uncertainty maps onto informational decoherence gradients.

\subsection{Proposed Phenomenological Neural-Synchrony Model}
\label{sec:SITneuroPrediction}

To make the neural-synchrony analogy operationally testable as a standalone phenomenological model (separate from the microscopic spacetime action \texttt{M\_AUTH\_SIT\_EFT\_v1}), we formulate a prospective test relating neural covariance purity and external oscillatory perturbations:

\begin{prediction}[Prospective Phenomenological Neural-Synchrony Test]
Let $R_{\mathrm{coh}}^{\mathrm{neural}}(t)$ denote an operational local synchrony metric defined as the normalized purity of the EEG/LFP covariance matrix in a given cortical area. Under the proposed dissipative neural-synchrony model, externally modulating oscillatory synchrony (e.g., via transcranial alternating current stimulation, tACS) produces a measurable change in the integrated synchrony functional
\[
\mathcal{C}_{\mathrm{neural}}(t) := \int_{\text{region}} R_{\mathrm{coh}}^{\mathrm{neural}}(x, t)\, d^3x,
\]
with $\frac{d \mathcal{C}_{\mathrm{neural}}}{dt} \leq 0$ under purely dissipative perturbations (such as general anesthesia) in the absence of external driving.

\textbf{Scoped Falsifiability:} Observing $\mathcal{C}_{\mathrm{neural}}(t)$ to increase systematically under purely decohering pharmacological perturbations without compensating input would falsify this specific phenomenological neural-synchrony model (without by itself testing the microscopic spacetime EFT \texttt{M\_AUTH\_SIT\_EFT\_v1}).
\end{prediction}

\paragraph{Prospective Experimental Protocol:}
\begin{itemize}
    \item Record high-density EEG or LFP from cortical tissue in vivo or in vitro.
    \item Apply tACS, optogenetic, or pharmacological manipulations to increase or decrease oscillatory synchrony.
    \item Compute $R_{\mathrm{coh}}^{\mathrm{neural}}$ at each location and time, and monitor $\mathcal{C}_{\mathrm{neural}}(t)$.
    \item Test monotonicity under known decohering interventions (e.g., general anesthesia).
\end{itemize}

\paragraph{Current Sensitivity:}
State-of-the-art signal processing allows $<1\%$ changes in neural synchrony and coherence to be resolved across hundreds of channels (see \textit{Nature Neuroscience} 2019, 22:807--819).

\paragraph{Implication:}
This phenomenological formulation operationalizes the proposed neural-synchrony analogy for prospective testing while keeping it clearly distinct from action-derived spacetime physics.

\section{Neural Vector Embeddings: Dendritic Configurations as Informational Attractors}

Within Super Information Theory (SIT), dendritic structures of neurons serve as biological instantiations of vector embeddings, dynamically encoding learned statistical distributions of incoming temporal and spatial coincidence patterns. This process reflects iterative informational synchronization, capturing spatiotemporal relationships through adaptive dendritic morphology and synaptic connectivity. No conceptual ideas from the original exposition have been lost; instead, this section now explicitly aligns neural phenomena with SIT's informational formalism.

Dendrites function as units of mixed selectivity, responding selectively to complex combinations of inputs rather than individual stimuli. This selectivity is structurally embedded within dendritic configurations, manifesting as high-dimensional embeddings of phase wave differentials representing coincident input patterns. Thus, dendritic morphology physically realizes stable informational attractor states, concretely encoding experienced patterns into lasting biological forms.

The embedding analogy indexes neuronal phase-wave differentials defined relative to baseline operational neural synchrony ($R_{\mathrm{coh}}^{\mathrm{neural}}$, kept distinct from the spacetime EFT modulus $\Rcoh = |\psi|/\psi_0$). Selective dendritic growth and synaptic pruning implement biological signal filtering that parallels, at a purely conceptual level, synchronization and dispersion motifs studied in physics.

Memory formation within dendritic architectures can be modeled phenomenologically as stabilization of low-prediction-error attractors across neuronal populations.

Prospective empirical evaluation of this neurocomputational analogy uses calcium imaging and high-density electrophysiology to measure dendritic activity and connectivity patterns alongside computational vector-embedding benchmarks (independently of gravitational EFT).

\subsection{Dendritic Architectures as Stored Matrices of Learned Relationships}

Viewing dendritic structures simultaneously as vector embeddings and biological lookup tables clarifies how neurons encode relational memories as matrices of learned statistical regularities. This perspective integrates structural plasticity mechanisms, including long-term potentiation (LTP) and depression (LTD), with functional retrieval processes mediated by nonlinear dendritic spikes. Such spikes function analogously to computational lookup operations, retrieving learned activation patterns corresponding to specific synaptic inputs.

In operational neural variables (Section~\ref{sec:neuroscience-implications}), dendritic input--output mappings may be parameterized phenomenologically as
\begin{equation}
    D_{ij}(t) = f_{\mathrm{neural}}\!\left(r_{\mathrm{osc}}(t),\, R_{\mathrm{coh}}^{\mathrm{neural}}(t),\, \mathbf{S}(t)\right),
\end{equation}
where $D_{ij}(t)$ denotes an effective dendritic response gain at time $t$ as a function of local oscillatory rate $r_{\mathrm{osc}}(t)$, operational synchrony $R_{\mathrm{coh}}^{\mathrm{neural}}(t)$, and synaptic input vector $\mathbf{S}(t)$ (distinct from the canonical spacetime fields $\rho_t, \Rcoh$).

Memory thus becomes a distributed, tensorial phenomenon wherein network-wide dendritic matrices collectively form higher-dimensional embeddings of complex associative relationships. Such embeddings encode not merely isolated memories, but comprehensive relational mappings across neuronal populations, adhering to the following structured analogy:

\begin{itemize}
    \item \textbf{Synaptic input patterns}: represented as vectors encoding incoming signals.
    \item \textbf{Dendritic embedding and retrieval}: operationalized through learned matrices encoding relational mappings from inputs to dendritic outputs.
    \item \textbf{Neuronal outputs}: resulting vectors composed of graded dendritic signals and spike-induced action potentials.
    \item \textbf{Network-wide memory encoding}: represented as higher-order tensors, encompassing distributed dendritic matrices across multiple neurons.
\end{itemize}

This rigorous framing clarifies relationships between structural dendritic modifications and functional synaptic plasticity, uniting neuroscience with formal computational analogies such as vector embeddings, associative memory, and attention mechanisms. The integration strengthens SIT's explanatory power, aligning biological neural structures explicitly with its formal informational dialect.

\subsection{Dendritic Vector Embeddings}

Biological dendritic structures encode relational memories analogous to high-dimensional tensor embeddings within artificial neural networks. Viewing each dendrite as a learned multidimensional associative embedding supplies a motivating neurobiological analogy and proposed representational correspondence for SIT's coherence-based modeling, rather than completed biological validation.

\section{Phase Wave Differential Tokens and Traveling Waves in Neural Assemblies}

We define the neural \emph{phase wave differential token} as a deviation ($\Delta$) from the coherent oscillatory pattern characteristic of a neural array, cluster, cortical column, or other defined neuronal assembly. Mathematically, this deviation can be represented explicitly as:
\[
\Phi_{\text{token}} = \Delta (\phi_{\text{neuron}} - \phi_{\text{group}})
\]
where $\phi_{\text{neuron}}$ is the instantaneous phase of an individual neuron's oscillation and $\phi_{\text{group}}$ is the mean phase of its synchronized ensemble.

A spike differing in timing or magnitude from the synchronized ensemble introduces a local \emph{phase wave differential}. This initiates traveling waves governed by a diffusion-attenuation equation, now explicitly defined as:
\[
\frac{\partial \Phi}{\partial t} = D \nabla^2 \Phi - \gamma \Phi
\]
Here, $D$ is the diffusion constant capturing wave propagation across the neural network, and $\gamma$ characterizes attenuation through inhibitory interactions and synaptic constraints.

Traveling waves propagate through neural tissue analogously to sequential domino cascades. Each wavefront activates dendritically embedded vector memories, triggering cascades of neuronal excitation and inhibitory interactions. Thus, memory retrieval and cognitive processes arise dynamically from these traveling waves.

Within this phenomenological cortical model (separate from \texttt{M\_AUTH\_SIT\_EFT\_v1}), operational ensemble synchrony $R_{\mathrm{coh}}^{\mathrm{neural}}$ can be coupled to phase-wave potentials via a coarse-grained neural drift--diffusion ansatz,
\[
\frac{\partial R_{\mathrm{coh}}^{\mathrm{neural}}}{\partial t} = D \nabla^2 R_{\mathrm{coh}}^{\mathrm{neural}} - \nabla \cdot \bigl(R_{\mathrm{coh}}^{\mathrm{neural}} \nabla \Phi_{\mathrm{neural}}\bigr),
\]
providing a standalone neurocomputational hypothesis for how traveling phase differentials modulate cortical ensemble synchrony.

\subsection{Quantum-Inspired Neural Information Processing}

Neural information processing aligns explicitly with SIT's quantum coherence dialect, incorporating quantum-inspired computational analogies:

\begin{itemize}
\item Neural coherence-decoherence dynamics as analogs to quantum computational processes.
\item Emergence of neural interference patterns, entanglement analogs (long-range correlations), and competitive neural state superpositions.
\end{itemize}

\subsection{Prospective Neuroscientific Models and Experiments}

Prospective neuroscientific research directions motivated by these phenomenological analogies include:

\begin{itemize}
\item Advanced EEG/MEG studies designed to test proposed coherence-based phenomenological models.
\item Empirical investigations into phase-synchronization mechanisms underlying attention, memory retrieval, and perceptual binding.
\item Computational neural models incorporating proposed coherence-synchronization metrics for prospective experimental comparison.
\end{itemize}

\subsection{Technological Applications and Brain-Computer Interfaces}

SIT's explicit coherence framework opens novel technological possibilities for brain-computer interfaces (BCIs):

\begin{itemize}
\item Development of coherence-based neural interfaces optimized for neural synchronization and enhanced user-device interaction.
\item Neuromorphic computational architectures explicitly inspired by coherence-based principles, offering significant advances in computational efficiency and adaptability.
\end{itemize}

\subsection{Philosophical and Ethical Considerations}

Understanding neural processes as coherence-driven phenomena introduces significant philosophical and ethical dimensions:

\begin{itemize}
\item Reevaluation of the mind-body problem through a rigorous informational ontology, reframing debates surrounding consciousness and subjectivity.
\item Ethical analyses addressing potential societal impacts and responsible governance of coherence-based neural enhancements.
\end{itemize}

\subsection{Summary of Neuroscientific Impact}

SIT rigorously grounds neuroscience within a coherent informational framework, explicitly reshaping our understanding of neural cognition and consciousness. By systematically connecting neural dynamics to quantum informational processes, SIT promotes interdisciplinary advancements across cognitive science, neurotechnology, and philosophy.

\section{Implications for Artificial Intelligence and Computation}
\label{sec:ai-computation-implications}

SIT explicitly introduces informational coherence principles, profoundly transforming artificial intelligence (AI) research and computational paradigms.

\subsection{AI and Computational Architectures}

Artificial Intelligence architectures naturally reflect coherence dynamics:
\begin{itemize}
\item Transformer self-attention corresponds to coherence-based informational focusing.
\item Sparse routing architectures in neural networks parallel SIT's phase-wave differential computations.
\item Neuromorphic computing architectures explicitly model coherence-decoherence transitions, aligning AI computation with biological neural dynamics.
\end{itemize}

\subsection{Quantum-Inspired Computational Paradigms}

AI models explicitly incorporating quantum-inspired coherence-decoherence mechanisms promise significant advances in computational efficiency and optimization:

\begin{itemize}
\item Explicit coherence-based AI algorithms enhancing computational efficiency through quantum-like transitions.
\item Quantum-inspired neural networks explicitly leveraging coherence dynamics for superior adaptive performance.
\end{itemize}

\subsection{Neural Networks and Informational Coherence}

Explicitly reframing artificial neural networks (ANNs) as coherence detectors elucidates their fundamental computational role, improving architectural and algorithmic design through rigorous coherence-based principles.

\subsection{Adaptive, Self-Organizing AI Systems}

Explicit SIT principles facilitate self-organizing adaptive AI systems, leveraging coherence-driven informational dynamics for robust autonomous behavior.

\subsection{Quantum Computing and Quantum Algorithms}

SIT explicitly informs novel quantum algorithm designs, introducing coherence mechanisms applicable to quantum error correction, optimization, and machine learning.

\subsection{Ethical and Societal Implications of Coherence-Based AI}

SIT explicitly addresses ethical and societal considerations:

\begin{itemize}
\item Societal impact analyses explicitly addressing coherence-based AI systems' autonomy, transparency, and ethical governance.
\item Development of explicit ethical frameworks ensuring responsible deployment and societal integration.
\end{itemize}

Furthermore, explicit coherence-driven AI education and democratization promise profound societal transformations toward inclusivity, cognitive diversity, and enhanced collective innovation.

\subsection{Empirical Validation and Future Research}

Future explicit computational modeling and empirical validation efforts include detailed simulations and practical coherence-based implementations within existing AI frameworks.

\subsection{Summary of Impact on AI and Computation}

SIT explicitly grounds AI research within rigorous informational coherence dynamics, promising substantial improvements in computational efficiency, adaptability, and societal integration. AI systems actively evolve informational coherence, dynamically shaping novel emergent complexity through explicit coherence-based computational paradigms.


\section{Integrative Insights from Related Frameworks}

Super Information Theory (SIT) is not an isolated conceptual structure; rather, it synthesizes key insights from several complementary theoretical frameworks across quantum physics, gravitational theory, thermodynamics, and computational neuroscience. Here, we explicitly illustrate how SIT draws upon related lines of work including \textit{Super Dark Time}, \textit{SuperTimePosition}, \textit{Micah's New Law of Thermodynamics}, \textit{Self Aware Networks (SAN)}, and \textit{Quantum Gradient Time Crystal Dilation} within an exploratory conceptual synthesis spanning quantum, gravitational, and neural analogies.

SIT rigorously synthesizes key prior insights without conceptual loss, including:

\begin{itemize}
\item Time-density variations explicitly correlated with quantum coherence-decoherence processes.
\item Emergence of gravitational phenomena from coherence-driven thermodynamics, explicitly linked through local time-density fields.
\item Quantum Coherence Coordinates (QCC) introduced as explicit constructs bridging classical informational theories with contemporary quantum-gravitational concepts.
\end{itemize}

\subsection{Quantum-Gravitational Computational Cycles: \textit{Super Dark Time} and \textit{SuperTimePosition}}

The frameworks of \textit{Super Dark Time} and \textit{SuperTimePosition} posit deterministic quantum-gravitational phenomena arising from ultrafast phase oscillations occurring below observational timescales. In \textit{SuperTimePosition}, quantum entities oscillate deterministically between particle-like localizations and wave-like delocalizations at extremely high frequencies. Quantum randomness arises observationally due to synchronization limitations between these rapid internal cycles and slower external measurements.

SIT incorporates these insights explicitly by framing quantum coherence and decoherence as synchronized and unsynchronized states within informational oscillatory processes. Coherence emerges naturally as stable synchronization of quantum-gravitational oscillations, whereas decoherence manifests through partial or incomplete sampling of these deterministic cycles. Hence, quantum randomness and gravitational effects share a unified deterministic informational foundation.

\subsection{Wave-Based Dissipation and Coherence: \textit{Micah's New Law of Thermodynamics}}

According to \textit{Micah's New Law of Thermodynamics}, physical and informational systems universally tend toward equilibrium through iterative dissipation of differences, occurring via wave-based phase interactions. Differences in energy, phase, momentum, or informational content dissipate through iterative exchanges, driving systems toward synchronized equilibria.

Within SIT, informational coherence directly corresponds to minimal phase differences among quantum-gravitational oscillators, while decoherence arises from persistent unsynchronized states. Mathematically, this dissipation-driven synchronization is explicitly described through the coherence--decoherence dynamics:
\[
\frac{d\rho_t}{dt} = \sum_{i,j}\alpha \sin(\Delta\phi_{ij}),
\]
where \(\Delta\phi_{ij}\) quantifies phase differences dissipating towards coherence. Thus, SIT incorporates Micah's Law rigorously, positioning coherence as a natural thermodynamic equilibrium state.

\subsection{Quantum Origins of the Time-Density Field: \textit{Super Dark Time}}

The \textit{Super Dark Time} framework articulates gravity as emerging directly from local quantum interference patterns shaping a fundamental quantum-mechanical time-density field. Quantum coherence patterns---stable interference effects---yield elevated local time densities (gravitational wells). Conversely, decoherence disrupts these patterns, decreasing local time densities.

This concept is explicitly captured in SIT through the coherence-dependent formulation of time-density:
\[
\rho_t(R_{\mathrm{coh}}) = \rho_0 + \gamma \cdot \mathrm{Re}\left[\sum_{n,m} C_n C_m^* e^{i(\phi_n - \phi_m)}\right],
\]
where quantum amplitudes \(C_n, C_m\), phases \(\phi_n, \phi_m\), and gravitational coupling constant \(\gamma\) precisely describe how coherence dynamically governs gravitational phenomena. Thus, SIT rigorously links quantum mechanical foundations with gravitational dynamics.

\subsection{Predictive Synchrony and Oscillatory Dynamics: \textit{Self Aware Networks}}

The \textit{Self Aware Networks (SAN)} framework proposes consciousness and cognitive functions as emergent phenomena arising from multi-scale oscillatory synchrony in neural networks. Neural oscillations encode predictive coding mechanisms, active inference, and perceptual synchronization, providing robust biological models for information processing.

Translating these biological insights rigorously into SIT, coherence is defined explicitly as a multi-scale synchronization metric. Informational coherence corresponds to globally synchronized informational states across quantum-gravitational scales. Decoherence reflects multi-scale informational desynchronization and predictive uncertainty. Thus, SAN's biological oscillatory models inform the quantum-gravitational synchronization logic central to SIT, bridging cognitive neuroscience and fundamental physics.

\subsection{Significance and Interdisciplinary Synthesis}

Integrating these insights from \textit{Super Dark Time}, \textit{SuperTimePosition}, \textit{Micah's New Law}, and \textit{Self Aware Networks}, SIT emerges as a rigorously unified theory that significantly strengthens its empirical and theoretical foundations by:

\begin{enumerate}
    \item Clarifying theoretical relationships between coherence and decoherence processes.
    \item Explicitly linking quantum-gravitational phenomena to deterministic synchronization mechanisms.
    \item Grounding thermodynamic dissipation rigorously within informational coherence dynamics.
    \item Providing biologically informed analogies from predictive neural synchronization models, enhancing explanatory and empirical relevance.
\end{enumerate}

This synthesis positions SIT as an integrative framework at the nexus of quantum mechanics, gravitational physics, thermodynamics, and cognitive neuroscience, demonstrating broad interdisciplinary applicability and robust empirical testability.

Tracing the intellectual trajectory leading to Super Information Theory highlights critical conceptual threads---particularly the characterization of information as a "dynamic substrate" underlying self-organizing processes across scales. The foundational synthesis of molecular signaling and neural oscillatory dynamics, first explicitly described in \textit{Bridging Molecular Mechanisms and Neural Oscillatory Dynamics}, introduced scale-invariant informational processing logic. Additionally, the collaborative, open-source approach pioneered in early Self Aware Networks discussions (GitHub \& YouTube, 2022) emphasized methodological transparency, interdisciplinarity, and democratization of scientific inquiry. These roots underscore SIT's commitment to integrative scientific rigor and methodological openness.

\section{Theoretical and Mathematical Directions}

To further refine and empirically validate Super Information Theory, this section proposes explicit theoretical directions synthesizing Predictive Coding, Karl Friston's Free Energy Principle and Active Inference, Feynman's Path Integral formulation, and recent quantum gravity advancements. These directions are structured around a key theoretical challenge:

\subsection{Proposed Connections to Quantum Symmetry Principles}

Recent parallel literature on time-symmetric open quantum dynamics, such as the symmetrical quantum Langevin equations studied by Guff et al.~(2025), provides motivating conceptual context and prospective experimental analogies for SIT's coherence--decoherence framework. Exploring whether such symmetrical open-system formulations connect to SIT's action and whether they yield observable signatures in atomic-clock synchronization or cold-atom interferometry remains an open prospective research direction (pending \texttt{OPEN-6}), rather than an existing empirical validation of SIT's foundational claims.

\section{Conclusion and Broader Impact}
\label{sec:conclusion-impact}

Super Information Theory (SIT) outlines an interdisciplinary research program exploring informational coherence across physical effective field theory and biological/computational analogies. This section summarizes the motivating goals and prospective directions of Part~III while maintaining the strict \texttt{QG49} boundary separating formal EFT and lattice theorems (Parts~I, II, and~IV) from open empirical derivations (\texttt{OPEN-6}).

Super Information Theory (SIT) frames foundational questions in physics and cosmology by investigating how dynamical coherence and time-density fields interact with spacetime geometry and matter within a controlled sub-Planckian effective field theory.

\subsection{Synthesis of Core Contributions}

At the conceptual level, SIT organizes its informational variables through Quantum Coherence Coordinates (QCC), the balanced ``halfway universe'' heuristic, and dissipation-driven synchronization analogies motivated by Micah's New Law, while anchoring its formal mathematical claims strictly in the canonical effective action \texttt{M\_AUTH\_SIT\_EFT\_v1} and its 4D lattice extension (Part~IV).

\subsection{Proposed Interface Between Quantum Coherence and Gravity}

By coupling dynamical coherence and time-density fields to Einstein--Hilbert gravity in \texttt{M\_AUTH\_SIT\_EFT\_v1}, SIT provides a controlled sub-Planckian EFT framework for studying coherence-coupled gravitational dynamics, while keeping ultraviolet quantum gravity (\texttt{OPEN-1}), quantum Wheeler--DeWitt closure (\texttt{OPEN-2}), and microphysical matter couplings (\texttt{OPEN-6}) explicitly open.

\subsection{Prospective Cosmological Hypotheses}

At the exploratory cosmological level, SIT motivates investigating whether spatial or epochal variations in coherence-coupled fields could contribute to gravitational lensing, structure formation, or expansion-rate discrepancies such as the Hubble tension. Turning these qualitative hypotheses into quantitative cosmological predictions requires deriving the coupled cosmological background and perturbation equations directly from the canonical action (\texttt{OPEN-6D}), which remains open work.

\subsection{Informational Cosmology and Structure Formation}

SIT also invites exploratory modeling of large-scale cosmic structures---galaxy clusters, filaments, and voids---alongside coherence-gradient analogies, distinct from established $\Lambda$CDM cosmological perturbation theory until derived directly from \texttt{M\_AUTH\_SIT\_EFT\_v1}.

\subsection{Prospective Experimental Pathways}

Prospective experimental venues---including precision atomic clock comparisons, torsion-balance inverse-square-law and equivalence-principle tests, quantum interferometry, and astrophysical surveys---define the arenas in which SIT can be rigorously tested once the open action-to-observable derivations (\texttt{OPEN-6}) are completed.

\subsection{Prospective Laboratory and Cosmological Targets}

Rather than claiming already-derived quantitative deviations in gravitational lensing profiles, time dilation, or cosmic microwave background anisotropies, Version~81 explicitly classifies those observables as \emph{prospective laboratory and cosmological targets} contingent on resolving \texttt{OPEN-6A}--\texttt{OPEN-6D} (microscopic material charges, the visible $U(1)_{\mathrm{EM}}$ dictionary, preparation-specific sources $J_{\mathrm{prep}}$, and full apparatus/cosmological transfer functions). Once those derivations are completed for a specified parameter state, atomic-clock comparisons, torsion balances, gravitational-wave detectors, and astrophysical surveys will supply scoped, quantitative tests capable of bounding or excluding that parameter state.

\subsection{Interdisciplinary and Philosophical Implications}

Philosophically, SIT explores conceptual bridges across physical, biological, and cognitive modeling through coherence-based informational language, clearly distinguishing heuristic analogies from established mathematical derivations. Likewise, the companion neuroscience paper \emph{Learning to Feel} (Paper~114: initial edition October~17, 2026; current main paper v1.2 dated October~20, 2026, with technical supplement retaining v1.1) pursues its own empirical neuroscience track without identifying neural synchrony or learned receiving influence with the canonical SIT spacetime scalar fields $(\psi, \rho_t)$ or their matter-source terms.

\subsection{Quantum Foundations and Determinism}

At the interpretational level, SIT explores connections with deterministic phase-oscillation analogies such as the SuperTimePosition (STP) and Gravitational Phase-Lock (Paper~034) framework, framed explicitly as a conditional measurement-dependent model program rather than an established derivation replacing standard quantum mechanics.

\subsection{Broader Societal and Technological Impact}

Across artificial intelligence, neuroscience, and foundational physics, SIT's coherence-based modeling vocabulary motivates exploratory algorithmic architectures and interdisciplinary collaboration.

\subsection{Future Research Directions}

Future SIT research focuses on closing the explicitly demarcated mathematical and observational obligations (\texttt{OPEN-1} through \texttt{OPEN-6}), including full action-to-apparatus derivations, independent experimental control reproductions, and rigorous variational and quantum continuum-limit analysis.

\subsection{Summary of Scope}

In summary, Super Information Theory Version~81 pairs a machine-checked sub-Planckian EFT and 4D lattice core with an explicitly scoped interdisciplinary research program, separating proved individual mathematical theorems from motivating analogies, prospective experimental targets, and unresolved open problems.
% --- END INPUT sections/06_part3_sec32_to_42.tex ---
% --- BEGIN INPUT sections/07_part3_sec43_to_50.tex ---
\section{Core Predictions and Resolutions to Open Problems}
\label{sec:core-predictions-open}

\subsection{Physical Grounding for the Free Energy Principle}

Super Information Theory provides a candidate field-theoretic underpinning for macroscopic principles of biological and adaptive self-organization, most notably Karl Friston's Free Energy Principle (FEP). The FEP posits that persistent open systems act so as to bound and minimize a variational free-energy functional serving as an upper bound on sensory surprise. Because FEP is formulated as a substrate-independent statistical constraint, SIT complements it by specifying a concrete physical mechanism through which free-energy-like Lyapunov functionals arise in open dissipative media without postulating teleological forces.

\paragraph{Ontological Status of Effective Functionals.}
No new fundamental ``purpose'' primitive is introduced into $S_{\mathrm{AUTH}}$. Any candidate control or memory functional $\Phi_{\mathrm{eff}}$ is strictly a \emph{system-specific reduced effective functional} obtained after coarse-graining open-system environmental couplings in the Feynman--Vernon influence-functional formalism (Appendices~\ref{app:recursion-influence}--\ref{app:clean-recursion}):
\begin{itemize}
  \item \textbf{Prediction Error as Phase-Wave Differential:} Sensory mismatch corresponds physically to localized spatial or temporal gradients in $(\Rcoh, \theta)$ between internal oscillatory assemblies and external boundary driving.
  \item \textbf{Active Inference as Dissipative Coherence Reconfiguration:} Open subsystems export entropy into their environment while relaxing internal phase-wave differentials toward phase-locked attractors, realizing perception and motor control as damped gradient flows on the reduced informational landscape.
  \item \textbf{Stationary Action as the Underlying Law:} Apparent free-energy minimization is an emergent reduced description of trajectories governed at the microscopic level by $\delta S_{\mathrm{AUTH}} = 0$ coupled to a thermal/dissipative bath.
\end{itemize}

\section{Exploratory Phenomenological Energy-Density Ansatz (Separate from Canonical Action)}
\label{sec:informational_energy_equivalence}

This section preserves the exploratory ``SIT~4.0'' phenomenological energy-density proposal as a \emph{separately specified heuristic idea}, clearly demarcating it from the authoritative continuum action \texttt{M\_AUTH\_SIT\_EFT\_v1} and repairing two dimensional and consistency errors present in earlier drafts.

\subsection{The Phenomenological Energy-Density Ansatz and Vacuum-Mismatch Boundary}

Earlier SIT~4.0 drafts postulated an auxiliary local energy-density ansatz
\begin{equation}
    \varepsilon_{\text{SIT}}^{\mathrm{phen}}(x) = \zeta \cdot |\psi(x)| \cdot [\rho_t(x)]^2,
    \label{eq:sit_energy_equivalence}
\end{equation}
where $\zeta$ is a phenomenological dimensional scale parameter (such that $[\zeta]\cdot[|\psi|]\cdot[\rho_t]^2 = [\text{energy density}] = [M L^{-1} T^{-2}]$ in SI units, or $[M]^1$ when $[|\psi|] = [\rho_t] = [M]^1$ in natural units). Two strict mathematical boundaries separate Eq.~\eqref{eq:sit_energy_equivalence} from unchanged \texttt{M\_AUTH\_SIT\_EFT\_v1}:
\begin{enumerate}[leftmargin=2em]
  \item \textbf{Not Derived from Unchanged \texttt{M\_AUTH\_SIT\_EFT\_v1} (Vacuum Consistency Check):} At the canonical constant broken-phase vacuum $(\rho_0, v_0 > 0, A_\mu = 0)$ of \texttt{M\_AUTH\_SIT\_EFT\_v1}, all field gradients, $F_{\mu\nu}$, and the total scalar potential $V_{\mathrm{tot}}(\rho_0, v_0) = 0$ vanish identically, so the canonical scalar--gauge stress--energy tensor is strictly $T_{\mu\nu}^{\mathrm{scal+EM}} = 0$. By contrast, evaluating Eq.~\eqref{eq:sit_energy_equivalence} on the same constant positive vacuum yields $\zeta v_0 \rho_0^2 > 0$. Therefore Eq.~\eqref{eq:sit_energy_equivalence} cannot serve simultaneously as a new postulate and as a consequence of the unchanged action; connecting it to canonical stress--energy would require an explicit background subtraction (e.g.\ relative to $(\rho_0, v_0)$) or a separate covariant action operator with proved $\nabla^\mu T_{\mu\nu} = 0$.
  \item \textbf{Mass Density vs.\ Integrated Mass (Restoring Volume Integration):} Converting an energy density $\varepsilon$ via $\varepsilon = m c^2 / V$ produces a \emph{mass density} $\rho_{\mathrm{mass}}(x) := \varepsilon_{\mathrm{SIT}}^{\mathrm{phen}}(x)/c^2$, not a total particle mass. For a spatially uniform body of volume $V$, total mass is $M = V \varepsilon_{\mathrm{SIT}}^{\mathrm{phen}} / c^2$; for a general spatial configuration, an integrated mass requires an explicit energy integral over a spatial hypersurface:
  \begin{equation}
      \rho_{\mathrm{mass}}^{\mathrm{phen}}(x) = \frac{\zeta}{c^2}\,|\psi(x)|\,[\rho_t(x)]^2,
      \qquad
      M_{\mathrm{info}}^{\mathrm{phen}} = \int_{\mathcal{V}} \rho_{\mathrm{mass}}^{\mathrm{phen}}(x)\,d^3x.
      \label{eq:mass_from_sit}
  \end{equation}
\end{enumerate}

\subsection{Conceptual Structure-vs.-Activity Interpretation}

Qualitatively, Eq.~\eqref{eq:sit_energy_equivalence} expresses an exploratory conceptual contrast between localized structural order (large $|\psi|$) and rapid dynamical activity (large $\rho_t$). Version~81 explicitly retracts any claim that Eq.~\eqref{eq:mass_from_sit} derives elementary particle rest masses from first principles, explains galactic dark matter, or accounts for human consciousness; those remain open speculative questions outside the proved derivations of \texttt{M\_AUTH\_SIT\_EFT\_v1}.



\section{Coherence-Regulated Hydrodynamics and the Clay Navier--Stokes Boundary}
\label{sec:navier-stokes-boundary}

In the classical kinetic and hydrodynamic limit (Section~\ref{sec:stability-causality} and Appendix~\ref{app:KineticLimit}), velocity moments of the coarse-grained Boltzmann distribution recover the compressible and incompressible Navier--Stokes equations. Within Super Information Theory, macroscopic fluid fields---velocity $\mathbf{v}(\mathbf{x},t)$, pressure $p(\mathbf{x},t)$, and vorticity $\boldsymbol{\omega} = \nabla \times \mathbf{v}$---represent hydrodynamic moments of an underlying physical substrate whose microscopic gradients are coupled to $\Rcoh(x)$ and $\rho_t(x)$.

\subsection{Physical Regularization via Coherence Saturation and Dispersive Hyperviscosity}

When a physical fluid approaches extreme small-scale shear or vortex stretching, two structural mechanisms intrinsic to \texttt{M\_AUTH\_SIT\_EFT\_v1} activate at the mesoscopic cutoff scale:
\begin{enumerate}[leftmargin=2em]
  \item \textbf{Bounded Normalized Coherence and Positive Gradient Energy:} Normalized interferometric purity satisfies $0 \le \Rcoh^{\mathrm{proxy}}(x) \le 1$, while spatial gradients of $(\rho_t, \psi)$ carry positive-definite $H^1$ energy density $\frac{\kappa_t}{2}|\nabla\rho_t|^2 + \frac{\kappa_c}{2}|\mathbf{D}\psi|^2 + V_{\mathrm{tot}}(\rho_t,|\psi|) \ge 0$.
  \item \textbf{State-Dependent Dispersive and Hyperviscous Regularization:} Coarse-graining the coupled $(\rho_t, \psi)$ stress tensor and open-system Lindblad dissipation into the hydrodynamic momentum balance supplements constant Newtonian viscosity $\nu\Delta\mathbf{v}$ with Bohm/Korteweg-type dispersive capillary stresses and high-wavenumber coherence-saturation damping,
  \begin{equation}
  \partial_t\mathbf{v} + (\mathbf{v}\cdot\nabla)\mathbf{v} + \frac{1}{\rho}\nabla p
  = \nu\Delta\mathbf{v} - \nu_{\mathrm{coh}}\,(-\Delta)^\alpha\mathbf{v} + \frac{\kappa_c}{2\rho}\nabla\bigl(\Delta|\psi|/|\psi|\bigr),
  \qquad \alpha \ge \frac{5}{4},
  \label{eq:sit_modified_navier_stokes}
  \end{equation}
  for which Lions' classical hyperviscosity theorem guarantees global smooth solutions whenever $\alpha \ge 5/4$.
\end{enumerate}

\subsection{Explicit Epistemic Boundary: Unmodified Clay Navier--Stokes Remains \texttt{NOT\_SOLVED}}

In strict compliance with the \texttt{QG49} five-tier epistemic ledger (Appendix~\ref{app:epistemic-ledger-lean}), it is essential to state with complete mathematical precision what this mechanism does and does \emph{not} establish:
\begin{itemize}
  \item \textbf{Physical Regularization vs.\ Pure Mathematical Millennium Problem:} Equation~\eqref{eq:sit_modified_navier_stokes} provides a \emph{modified physical hydrodynamic model} in which short-distance coherence gradient penalties prevent physical energy density from diverging. It does \textbf{not} constitute a proof of global smooth solutions (nor a finite-time blow-up counterexample) for the \emph{unmodified} pure incompressible 3D Eulerian Clay Millennium Navier--Stokes equations ($\partial_t\mathbf{v} + (\mathbf{v}\cdot\nabla)\mathbf{v} + \nabla p = \nu\Delta\mathbf{v}$, $\nabla\cdot\mathbf{v} = 0$ with constant $\nu > 0$ and zero higher-order dispersion), which remains strictly \textbf{\texttt{NOT\_SOLVED}} (Tier~5).
  \item \textbf{Formal Ladyzhenskaya Embedding Interface (\texttt{hLady}, \texttt{OPEN-3}):} Within the Lean~4.30 Sobolev formalization (\texttt{SITStage3\allowbreak{}Infinite\allowbreak{}Sobolev\allowbreak{}Coherence\allowbreak{}Tower}), sequence-space $H^2(\mu) \subsetneq H^1(\mu) \subsetneq \ell^2$ embeddings and $O(\Lambda_N^{-1})$ Galerkin bounds are kernel-proved unconditionally, whereas 3D continuum spatial $L^4$ Ladyzhenskaya interpolation ($\|u\|_{L^4}^4 \le C\|u\|_{L^2}\|\nabla u\|_{L^2}^3$) is isolated transparently as an explicit conditional hypothesis \texttt{hLady} (\texttt{W\_ADV\_TD\_02}, \texttt{OPEN-3}).
\end{itemize}

\section{Informational Torque: Conceptual Proposal and Mathematical Corrections}
\label{sec:informational-torque}
\label{sec:informational_torque}

Earlier SIT drafts introduced an exploratory geometric concept termed ``informational torque'' to probe whether non-parallel amplitude and phase gradients of $\psi(x)$ could play a role analogous to Berry curvature in sourcing spacetime curvature. Version~81 preserves the conceptual question while explicitly correcting three mathematical errors in the earlier formulas and grounding gravitational curvature strictly in the variational Einstein equations of \texttt{M\_AUTH\_SIT\_EFT\_v1}.

\subsection{Retraction of Earlier Flat-Connection Torque and Singular Curvature Identities}

Three concrete mathematical checks invalidate the earlier equations of this section:
\begin{enumerate}[leftmargin=2em]
  \item \textbf{Flat Phase-Gradient Connection Cannot Generate Nonzero Torque:} Earlier drafts wrote a Berry-like connection $\mathbf{A}_{\mathrm{Berry}} \propto \nabla\theta$ and identified its curvature with $\boldsymbol{\tau}_{\mathrm{old}} = (\hbar/e)|\psi|\nabla|\psi|\times\nabla\theta$. On any smooth patch, however, the curl of a gradient vanishes identically ($\nabla\times\nabla\theta = \mathbf{0}$). As an explicit counterexample, setting $|\psi|(x,y) = x + 2$ and $\theta(x,y) = y$ gives $\nabla\times(\nabla\theta) = \mathbf{0}$ everywhere while $|\psi|\nabla|\psi|\times\nabla\theta = (0, 0, x+2) \ne \mathbf{0}$, proving that $\boldsymbol{\tau}_{\mathrm{old}}$ is \emph{not} the curvature of a connection proportional to $\nabla\theta$.
  \item \textbf{Local Gauge Non-Invariance of $\nabla\theta$:} Under a local $U(1)$ gauge transformation $\psi \mapsto e^{i\hat{q}\chi}\psi$, $A_\mu \mapsto A_\mu + \partial_\mu\chi$, the bare phase gradient $\nabla\theta$ shifts by $\hat{q}\nabla\chi$. Any gauge-invariant spatial cross product built from $\psi$ must replace $\nabla\theta$ with the covariant superflow vector $\nabla\theta - \hat{q}\mathbf{A}$ (proportional to the canonical Noether current $\mathbf{J}_{\mathrm{coh}} = \kappa_c|\psi|^2(\nabla\theta - \hat{q}\mathbf{A})$):
  \begin{equation}
  \boldsymbol{\tau}_{\mathrm{cov}} := \kappa_c\,|\psi|\,\nabla|\psi| \times \bigl(\nabla\theta - \hat{q}\mathbf{A}\bigr)
  = \frac{1}{2}\,\nabla\bigl(|\psi|^2\bigr) \times \bigl(\kappa_c(\nabla\theta - \hat{q}\mathbf{A})\bigr).
  \label{eq:gauge_invariant_torque_def}
  \end{equation}
  While Eq.~\eqref{eq:gauge_invariant_torque_def} repairs local $U(1)$ gauge invariance, it is simply an algebraic field composite and does \emph{not} derive spacetime curvature.
  \item \textbf{Retraction of the Singular $R_{\mu\nu\rho}{}^\sigma \propto \tau_{\mu\nu\rho}{}^\sigma / \Box\rho_t$ Expression:} Earlier drafts displayed a formula $R_{\mu\nu\rho}{}^\sigma = \frac{16\pi G}{c^4}\tau_{\mu\nu\rho}{}^\sigma / \Box\rho_t$ attributed to varying the master action. Version~81 explicitly retracts that expression: metric variation of \texttt{M\_AUTH\_SIT\_EFT\_v1} produces the standard Einstein equation $G_{\mu\nu} = \kappa_G T_{\mu\nu}^{\mathrm{tot}}$ (Eq.~\eqref{eq:covariant_einstein_sec5}), no derivation yields the displayed rank-4 tensor ratio, its denominator $\Box\rho_t$ vanishes identically on canonical constant-field backgrounds, and vacuum General Relativity admits nonzero Riemann curvature ($R_{\mu\nu\rho}{}^\sigma \ne 0$) even when $(\rho_t, \psi)$ are strictly constant.
\end{enumerate}

\subsection{Exploratory Experimental Template (Conditional on \texttt{OPEN-6C})}

The conceptual question above motivates asking whether optical-lattice atom interferometers with controlled phase and contrast gradients could probe transverse coherence-gradient responses once a concrete coupling is derived. Under \texttt{QG49} discipline, unchanged \texttt{M\_AUTH\_SIT\_EFT\_v1} does \emph{not} derive a quantitative transverse displacement forecast (earlier picometer estimates remain retired pending derivation of $J_{\mathrm{prep}}$ and material couplings under \texttt{OPEN-6A}/\texttt{OPEN-6C}).

\subsection{Scope of Any Future Geometric-Phase / Berry-Curvature Construction}

To construct a mathematically valid Berry curvature in a future extension, one must specify an explicit family of quantum states $|\Psi(\boldsymbol{\lambda})\rangle$ over a declared parameter manifold $\boldsymbol{\lambda}$ (with nontrivial connection $\mathcal{A}_a = i\langle\Psi|\partial_{\lambda^a}\Psi\rangle$ and two-form $\mathcal{F}_{ab} = \partial_a\mathcal{A}_b - \partial_b\mathcal{A}_a$, distinct from the spacetime gradient $\nabla\theta$ of a single scalar phase). In the present canonical EFT \texttt{M\_AUTH\_SIT\_EFT\_v1}, gravitational curvature is sourced strictly by $T_{\mu\nu}^{\mathrm{tot}}$ and gauge curvature is carried by the independent $U(1)$ connection $A_\mu$ (Section~\ref{sec:magnetism-holonomy}, \texttt{OPEN-6B}).

\vspace{1cm}

\section{Magnetism, Gauge Holonomy, and the Flat-Connection Boundary}
\label{sec:magnetism-holonomy}
\label{sec:magnetism_as_holonomy}

Earlier drafts of Super Information Theory explored whether ordinary magnetic fields could be derived purely from gradients of the coherence phase $\theta(x) = \arg\psi(x)$. Version~81 replaces that heuristic identification with a strict geometric separation matching Part~II, Paper~002, and Companion Paper~III: the canonical action \texttt{M\_AUTH\_SIT\_EFT\_v1} carries an \emph{independent} internal $U(1)$ gauge connection $A_\mu$ whose relationship to visible electromagnetism remains open (\texttt{OPEN-6B}), whereas setting $A_\mu$ equal to a local phase gradient yields a flat (pure-gauge) connection.

\subsection{Historical Flat-Connection Ansatz, Its $F_{\mu\nu} = 0$ Limitation, and the Canonical Independent Connection}

In the canonical action \texttt{M\_AUTH\_SIT\_EFT\_v1} (Part~II), local $U(1)$ gauge symmetry of the complex coherence field $\psi(x)$ requires introducing an \emph{independent} principal-bundle connection 1-form $A_\mu$ inside the covariant derivative $D_\mu\psi = (\nabla_\mu - i\hat{q}A_\mu)\psi$, with gauge-invariant curvature $F_{\mu\nu} = \partial_\mu A_\nu - \partial_\nu A_\mu$ and kinetic term $-\frac{1}{4}f_{\mathrm{EM}}(R_{\mathrm{coh}})F_{\mu\nu}F^{\mu\nu}$. By contrast, earlier SIT drafts examined a historical/conjectural flat-connection ansatz writing
\begin{equation}
   A_\mu^{\mathrm{flat}} \; := \; \frac{\hbar}{e}\,\partial_\mu\theta(x).
\end{equation}
This flat-connection expression is \emph{not} the general connection required by local $U(1)$ gauge symmetry. On any simply connected region where $\theta(x)$ is smooth and locally single-valued, the curl of a gradient vanishes identically:
\begin{equation}
   F_{\mu\nu}^{\mathrm{flat}} \;=\; \partial_\mu A_\nu^{\mathrm{flat}} - \partial_\nu A_\mu^{\mathrm{flat}}
   \;=\; \frac{\hbar}{e}\bigl(\partial_\mu\partial_\nu - \partial_\nu\partial_\mu\bigr)\theta(x) \;=\; 0,
   \qquad
   \mathbf{B}^{\mathrm{flat}} \;=\; \nabla \times \mathbf{A}^{\mathrm{flat}} \;=\; 0.
\end{equation}
Nontrivial global holonomy can exist on multiply connected domains or around singular vortex/defect cores where $\theta(x)$ is multivalued, but a pure phase gradient does \emph{not} establish a general construction of ordinary non-flat magnetic fields ($\mathbf{B} \ne 0$, $F_{\mu\nu} \ne 0$). Throughout Version~81, general magnetic fields require an independent gauge connection $A_\mu$ with nonvanishing curvature $F_{\mu\nu} \ne 0$, and matching SIT's internal $U(1)$ sector to visible electromagnetism remains an explicitly open microphysical task (\texttt{OPEN-6B}).

As standard electromagnetic background, the Aharonov--Bohm (AB) effect demonstrates the physical observability of closed-loop $U(1)$ holonomy even where $\mathbf{B} = 0$ along the particle paths: an electron beam split around an enclosed magnetic flux $\Phi_B$ acquires the standard quantum-mechanical phase shift
\[
  \theta_{\text{AB}}
  =
  \frac{e}{\hbar}\oint A_i\,dx^i = \frac{e}{\hbar}\Phi_B.
\]
Any additional SIT interpretation linking the visible electromagnetic AB connection to the coherence bundle $\psi$ is explicitly conditional on establishing the visible-sector gauge dictionary (\texttt{OPEN-6B}); the established AB interference pattern is standard quantum electrodynamics, not an independent empirical validation of SIT.

\subsection{Electron--Beam Deflection Background and Withdrawal of the Legacy CRT Bound}
\label{sec:magnet}

\paragraph{Standard electromagnetic baseline.}
In a cathode-ray tube (CRT) or an electron microscope, the ordinary Lorentz force
$\boldsymbol{F}=e\,\boldsymbol{v}\!\times\!\boldsymbol{B}$ bends an electron
trajectory into an arc of radius
$R_{\!B}=mv/(eB)$, with the same curvature arising in standard quantum mechanics from minimal substitution
$\boldsymbol{p}\to\boldsymbol{p}-e\boldsymbol{A}$ for a genuine non-flat electromagnetic connection ($\nabla\times\boldsymbol{A} = \boldsymbol{B} \ne 0$).

\paragraph{Conditional SIT discussion and withdrawal of legacy $\lambda R_{\mathrm{coh}} \lesssim 10^{-8}$ figure.}
Earlier exploratory drafts reinterpreted classical CRT deflection ($v \simeq 2.0\times 10^{7}\,\mathrm{m\,s^{-1}}$, $B \simeq 10^{-2}\,\mathrm{T}$, $90^\circ$ bend over $L = 5\,\mathrm{cm}$) using the flat-connection phase-gradient ansatz to quote a heuristic bound $\lambda R_{\mathrm{coh}} \lesssim 10^{-8}$ and describe electron deflection as a calibrated probe of the coherence field. Version~81 explicitly \emph{withdraws} that CRT-derived $10^{-8}$ figure as an empirical SIT bound and withdraws the ``calibrated probe'' characterization. Because ordinary Lorentz deflection occurs in a region with $\boldsymbol{B} \ne 0$ (which cannot be generated by a smooth pure gradient $\partial_\mu\theta$) and because neither the visible $U(1)_{\mathrm{EM}}$ dictionary (\texttt{OPEN-6B}) nor the apparatus transfer function (\texttt{OPEN-6D}) has been derived from \texttt{M\_AUTH\_SIT\_EFT\_v1}, CRT deflection currently stands solely as established electromagnetic background rather than an empirical SIT constraint.

\paragraph{Scoped conclusion.}
Electron deflection remains a standard electromagnetic phenomenon. Any future SIT-specific magnetic or interferometric bound requires first resolving how visible electromagnetism embeds alongside or mixes with SIT's internal $U(1)$ connection (\texttt{OPEN-6B}) and deriving the corresponding observable shift from the canonical action.

\subsection{Conditional Decoupling Hypothesis and Consistency Check}

Under the separate conjectural hypothesis that static magnetic fields couple predominantly through gauge/phase degrees of freedom rather than sourcing large modulus ($|\psi|$) or temporal-density ($\rho_t$) gradients, magnetic setups and scalar-gravity setups would probe distinct sectors of the theory. However, in the canonical action \texttt{M\_AUTH\_SIT\_EFT\_v1}, $-\frac{1}{4}f_{\mathrm{EM}}(R_{\mathrm{coh}})F_{\mu\nu}F^{\mu\nu}$ couples $F_{\mu\nu}F^{\mu\nu}$ directly to $R_{\mathrm{coh}}^2$ whenever $c_{\mathrm{EM}} \ne 0$ and the visible electromagnetic field is identified with $F_{\mu\nu}$ (which must be evaluated consistently with \texttt{OPEN-6B}).

\textbf{Prospective Setup:} Place a Cavendish-type torsion balance or precision gravimeter next to a stable high-field superconducting magnet (e.g., 10~Tesla), controlling for ordinary magnetic susceptibility and classical electromagnetic stress-energy.

\textbf{The Consistency Check (Conditional on the Gauge Dictionary, \texttt{OPEN-6B}):} Under a conjectural pure-phase decoupling model in which static magnetic fields do not source $|\psi|$ or $\rho_t$ gradients beyond standard electromagnetic stress-energy $T_{\mu\nu}^{\mathrm{EM}}$, no anomalous differential gravitational attraction is induced. Conversely, deriving a quantitative prediction from \texttt{M\_AUTH\_SIT\_EFT\_v1} requires resolving \texttt{OPEN-6B} (including the status of $c_{\mathrm{EM}} R_{\mathrm{coh}}^2 F^2$ for visible magnetic fields).


\section{Deriving the Standard Model and Gravity from First Principles}
\label{sec:emergence_sm_g_moved}
\label{sec:vortical_defects}

Having established the mathematical consistency of Super Information Theory (SIT), this section moves the theory from phenomenological unification to fundamental explanation.
We now examine how Super Information Theory interfaces with the Standard Model particle spectrum, flavor mixing, and gravitational coupling, separating rigorously machine-checked Stage-4 chiral-gauge theorems from open topological research conjectures under \texttt{QG49} epistemic discipline.

\subsection{A Conjectured Topological Particle Spectrum from the Coherence Field}

A central ambition of fundamental physics is to derive the observed spectrum of elementary particles from a minimal set of first principles. Super Information Theory (SIT) proposes a path toward this goal by conjecturing that elementary particles are stable, quantized topological excitations of the fundamental complex coherence field, $\psi(x)$. In this framework, the diverse properties of particles---both fermions and bosons---emerge as calculable invariants of an underlying informational geometry.

\subsubsection{Stress-Energy and the Emergence of Gravity}

This topological classification is not merely a labeling scheme; it has direct physical consequences for the spacetime metric. A topological defect in the coherence field carries energy and momentum, and therefore sources gravitational curvature. Its stress-energy tensor follows directly from varying the SIT master action with respect to the metric $g_{\mu\nu}$:
\begin{equation}
    T_{\mu\nu}^{\text{vortex}} = \kappa_{c}\left(\nabla_{\mu} R_{\text{coh}}\nabla_{\nu} R_{\text{coh}} - \frac{1}{2} g_{\mu\nu}\nabla_{\alpha} R_{\text{coh}}\nabla^{\alpha} R_{\text{coh}}\right) + T_{\mu\nu}^{\text{core}},
\end{equation}
where $T_{\mu\nu}^{\text{core}}$ encodes the short-distance structure within the defect's core radius where the phase becomes ill-defined. Because this $T_{\mu\nu}^{\text{vortex}}$ appears on the right side of the Einstein field equations, $G_{\mu\nu} = 8\pi G T_{\mu\nu}$, each stable winding necessarily sources curvature. This leads to a profound conclusion: \textbf{gravity is not an \emph{extra} field but the inevitable geometric response to the energy stored in a coherence defect.}

\subsubsection{The Stability Condition and Informational Structure of Particles}

Not all topological configurations are physically realized as stable particles. A vortex must be energetically stable against decay into radiation. Whether a given vortex preserves macroscopic coherence or dissipates depends on the balance between its internal gradient energy (informational tension) and the decohering couplings that radiate energy away. We define the coarse-grained energy per unit length of a defect as:
\begin{equation}
    \mathcal{E} = \pi\kappa_{c}\ln(r_{\infty}/r_{0}) - \Gamma,
\end{equation}
where $r_{0}$ is the core radius, $r_{\infty}$ is the system size, and $\Gamma$ is the dissipation rate determined by environment-induced decoherence. This leads to two distinct physical regimes:
\begin{itemize}
    \item $\mathcal{E} > 0$: A \textbf{coherence-sustaining, long-lived defect}. The informational tension dominates, creating a stable, massive, gravitating particle.
    \item $\mathcal{E} < 0$: A \textbf{turbulent, rapidly decaying defect}. Dissipation dominates, leading to a transient, radiation-dominated excitation.
\end{itemize}
This stability condition is governed by an excitability ratio. Configurations with a low ratio are stable particles, whereas those with a high ratio are fleeting resonances. Furthermore, the internal structure of the defect affects the local rate of time via the fundamental SIT constraint $\rho_{t} = \rho_{0}\exp[\alpha R_{\text{coh}}]$. At the core of a strongly wound vortex, the coherence field must vanish ($R_{\text{coh}} \rightarrow 0$), causing the time-density to approach its vacuum value, $\rho_{t} \rightarrow \rho_{0}$, which dilates proper time. This provides the concrete mechanism by which coherence defects seed both particle identity and spacetime geometry.

\subsubsection{A Unified View of Fermions and Bosons: Sources and Interactions}

SIT makes a fundamental distinction between the two classes of elementary particles, grounding it in their distinct informational roles.
\begin{itemize}
    \item \textbf{Fermions as Stable Topological Sources:} The fundamental particles of matter (quarks and leptons) are identified as stable, localized, and topologically protected defects (solitons, vortices, or knots) in the coherence field. Their persistence is guaranteed by topological invariants. They are the stable \emph{sources} of informational structure in the universe. The Pauli Exclusion Principle arises as a topological constraint preventing two identical defects from coexisting at the same location.

    \item \textbf{Bosons as Quantized Propagators of Interaction:} The fundamental force carriers are identified not as stable defects, but as the quantized, propagating excitations of the fields that mediate interactions \emph{between} the fermionic sources. They are the carriers of informational change, phase communication, and topological reconfiguration.
\end{itemize}

\subsubsection{Topological Classification of Fermion Families and Generations}

The classification of fermionic matter rests on the topological properties of the coherence manifold, $M_{\text{coh}}$, which represents the space of all possible stable configurations of the $\psi$ field.

\begin{itemize}
    \item \textbf{Particle Families as Homotopy Classes:} The primary division of fermions into distinct families is determined by the topology of $M_{\text{coh}}$. We postulate that different particle families correspond to elements of different homotopy groups, $\pi_n(M_{\text{coh}})$. For instance, the fundamental group $\pi_1(M_{\text{coh}})$ might classify leptonic states, while a higher-order group like $\pi_3(M_{\text{coh}})$ could correspond to the more complex internal structure of quarks.

    \item \textbf{Generations as Nested Topological Structures:} The existence of three generations of matter is proposed to arise from higher-order, nested topological structures. A second-generation particle would be a stable topological excitation of the field configuration that constitutes a first-generation particle. The increasing mass of higher generations is a natural consequence of the increased gradient energy and informational complexity required to sustain these nested configurations.
\end{itemize}

\subsubsection{Quantum Numbers and Bosons from Geometric Principles}

Within SIT, conserved quantum numbers and their associated force-carrying bosons are not abstract labels but emerge from the geometry of the coherence field and its phase.

\begin{itemize}
    \item \textbf{Electromagnetism (Canonical Gauge Boundary, \texttt{OPEN-6B}):} In the canonical action \texttt{M\_AUTH\_SIT\_EFT\_v1}, local $U(1)$ gauge invariance is carried by an \emph{independent} internal connection $A_\mu$ (distinct from a flat phase gradient $\partial_\mu\theta$), which acquires a broken-phase mass $m_A^2 = \kappa_c \hat{q}^2 v_0^2 / f_{\mathrm{EM}0}$ whenever $\hat{q}v_0 \ne 0$. Identifying the massless visible \textbf{photon} therefore cannot rely on taking $A_\mu \propto \partial_\mu\theta$ or equating the broken internal connection directly with visible electromagnetism; instead, it requires either a separate unbroken visible $U(1)_{\mathrm{EM}}$ gauge field decoupled from the broken coherence vacuum expectation value $v_0$ or an explicit visible-versus-hidden gauge dictionary (\texttt{OPEN-6B}, with tree-level kinetic mixing analyzed in Stage~6).

    \item \textbf{Strong Force:} The $SU(3)$ color charge of quarks emerges from a more complex, non-Abelian topological structure. The \textbf{gluons} are the eight quantized excitations of the connections on this non-Abelian structure, mediating transitions between the three stable color states.

    \item \textbf{Weak Force:} The weak force changes a particle's topological identity. The massive \textbf{W and Z bosons} are high-energy, propagating excitations of the coherence field capable of inducing topological reconnections, thereby transforming one type of fermionic knot into another. Their large mass reflects the high energy barrier for this fundamental reconfiguration.
\end{itemize}

\subsubsection{Spin and Topological Statistics}
The distinction between fermions and bosons is encoded in the behavior of their corresponding topological excitations under a $2\pi$ rotation. Fermionic modes are identified with field configurations that acquire a non-trivial phase factor (a minus sign) upon such a rotation (e.g., a skyrmion or a twisted loop), naturally embedding the spin-statistics theorem into the geometry of the coherence field.

\subsubsection{The Origin of Mass Scale and the Higgs Boson from the SIT Vacuum}
The parameters that govern mass are determined by the self-consistent properties of the SIT vacuum. The informational potential $V(|\psi|, \rho_t)$ possesses a non-trivial minimum that defines the vacuum state, and the curvature of this potential sets the fundamental energy scale for all excitations. Therefore, mass does not arise from coupling to a separate Higgs field, but from the energy required to "lift" a topological excitation out of the stable SIT vacuum.

In this picture, the experimentally observed \textbf{Higgs boson} is reinterpreted. It is not a quantum of a separate, fundamental field. Instead, we propose it is a massive, scalar, non-topological excitation of the \textbf{SIT vacuum field itself}. It is a propagating ripple in the background "informational stiffness" of spacetime. This excitation naturally couples to the topological defects (fermions), affecting their stability and energy, thereby reproducing the observed Higgs mechanism without requiring an additional fundamental field.

\subsubsection{Roadmap to a Unique and Complete Spectrum}

This framework moves beyond analogy to a concrete program of proof. The ultimate goal is to demonstrate that the complete set of stable, low-energy topological solutions to the non-linear SIT field equations is precisely the Standard Model particle spectrum, with no extraneous states. This program contrasts with purely algebraic approaches by being fundamentally geometric and dynamic. The roadmap to this proof involves four key steps:
\begin{enumerate}
    \item \textbf{Define the full non-linear field equations} for the complex field $\psi$ and $\rho_t$, including the self-interaction terms specified by the informational potential $V(|\psi|, \rho_t)$.
    \item \textbf{Analyze the stability of solitonic solutions} to these equations, identifying which configurations are energetically stable and long-lived.
    \item \textbf{Classify all stable solutions} using the tools of algebraic topology (homotopy, homology, and knot theory) to generate a complete set of topological invariants.
    \item \textbf{Prove a one-to-one correspondence} between this set of classified topological excitations and the observed particles of the Standard Model, and demonstrate that no other stable solutions are permitted in the low-energy regime.
\end{enumerate}
If successful, this research program would reframe the Standard Model of particle physics not as a collection of fundamental postulates, but as the potential emergent consequence of a deeper, unified informational geometry.

\subsection{A Proposed Mechanism for Mass and Particle Mixing from Informational Stability}

Following the conjecture that particles are stable, topological excitations of the coherence field, we now outline a proposed mechanism for explaining their fundamental properties: mass and mixing. This section conjectures that the observed mass hierarchy and mixing matrices of the Standard Model are not arbitrary parameters but could be emergent consequences of the stability and geometry of these underlying informational modes.

\subsubsection{The Mass-Resilience Relation and the Informational Resilience Index}

In SIT, mass is not a fundamental property but is instead a measure of a coherence mode's resistance to decoherence---its \textit{informational resilience}. A more resilient, stable, or topologically complex coherence pattern requires more energy to create or disrupt, manifests as a more significant local thickening of the time-density field $\rho_t$, and is thus perceived as having greater mass.

We formalize this concept by defining the \textbf{Informational Resilience Index}, $I_R(i)$, for the $i$-th particle mode. This dimensionless index is a calculable functional of the particle's corresponding coherence field configuration, $R_{\text{coh}}^{(i)}(x)$, representing the total action cost of the static field configuration:
\begin{equation}
    I_R(i) = \int_{\mathbb{R}^3} \mathcal{L}_{\text{static}}\left( R_{\text{coh}}^{(i)}, \nabla R_{\text{coh}}^{(i)} \right) d^3x = \int_{\mathbb{R}^3} \left( \frac{\kappa_c}{2} |\nabla R_{\text{coh}}^{(i)}|^2 + V(R_{\text{coh}}^{(i)}) \right) d^3x,
    \label{eq:resilience_index_consolidated}
\end{equation}
where $\mathcal{L}_{\text{static}}$ is the static Lagrangian density for the $R_{\text{coh}}$ field.

The mass of the $i$-th particle is then determined by the \textbf{Mass-Resilience Relation}:
\begin{equation}
    m_i = m_0 \exp(k \cdot I_R(i)).
    \label{eq:mass_resilience_relation_consolidated}
\end{equation}
Here, $m_0$ is a fundamental mass scale set by the vacuum energy of the informational fields, and $k$ is a dimensionless constant parameterizing the informational medium's "stiffness." This exponential relationship naturally generates a large hierarchy of masses from O(1) differences in the Informational Resilience Index.

\subsubsection{Particle Mixing as Quantum Tunneling Between Coherence Modes}

SIT reinterprets particle interactions and flavor changes as quantum tunneling events between different stable topological minima in the configuration space of the $R_{\text{coh}}$ field. The CKM and PMNS mixing matrices, which parameterize these transitions, are therefore derivable from the geometry of this configuration space.

The matrix element $V_{ij}$ describing the transition amplitude between particle state $i$ and state $j$ is given by the path integral over all field configurations that interpolate between the two stable modes. In the semi-classical (WKB) approximation, this tunneling amplitude is dominated by the overlap of the wave-functionals corresponding to the two particle modes:
\begin{equation}
    V_{ij} \approx \mathcal{N} \int \Psi_j^*[R_{\text{coh}}] \Psi_i[R_{\text{coh}}] \mathcal{D}R_{\text{coh}},
    \label{eq:mixing_overlap_integral_consolidated}
\end{equation}
where $\Psi_i[R_{\text{coh}}]$ is the wave-functional for the $i$-th particle's coherence mode. This implies that the mixing angles are determined by the geometric "shape" and proximity of the different stable coherence modes.


\paragraph{Explicit \texttt{QG49} Scope Note on Three-Generation Topological Homotopy (\texttt{NOT\_SOLVED}).}
While Sections~48.1--48.2 formulate a concrete geometric research program modeling fermion families and mass hierarchies via solitonic resilience indices $I_R(i)$, deriving the exact number of three chiral fermion generations from homotopy groups $\pi_3(\mathcal{M}_{\mathrm{coh}})$ (or nested skyrmion moduli) has \textbf{not} been proved and remains an explicit open conjecture (\textbf{\texttt{3-generation homotopy: NOT\_\allowbreak{}SOLVED}}, Tier~5). What \emph{is} rigorously proved in Lean~4.30 across Stage~4 Phase~4 and Stage~5.1 (detailed below in Subsection~48.4 and Section~\ref{sec:stage4-matter-holonomy}) is the finite-matrix Neuberger--Ginsparg--Wilson chiral overlap algebra (with explicit non-zero operator witness $D_{\mathrm{ov}} \ne 0$) and the exact injective unitary group embedding $SU(3)_c \times SU(2)_L \times U(1)_Y \hookrightarrow U(6)$ coupled to a positive massive fermion-determinant SIT+SM probability measure.

\subsection{Fundamental Einstein--Hilbert Gravity vs.\ Retirement of Naive One-Loop Induced Gravity (\texttt{C12})}

\subsubsection{The Informational Stress--Energy Tensor}

Within \texttt{M\_AUTH\_SIT\_EFT\_v1}, the scalar sector carries the canonical $(-,+,+,+)$ stress--energy tensor
\begin{equation}
T_{\mu\nu}^{(\mathrm{info})}
= \kappa_t\,\partial_\mu\rho_t\partial_\nu\rho_t
  + \kappa_c\,\mathrm{Re}\bigl[(D_\mu\psi)^*(D_\nu\psi)\bigr]
  - g_{\mu\nu}\left[\frac{\kappa_t}{2}(\partial\rho_t)^2 + \frac{\kappa_c}{2}|D\psi|^2 + V_{\mathrm{tot}}(\rho_t,|\psi|)\right],
\label{eq:info_tension_tensor_consolidated}
\end{equation}
whose vacuum expectation value at $(\rho_0, v_0)$ vanishes classically ($V_{\mathrm{tot}}(\rho_0, v_0) = 0$).

\subsubsection{Retirement of Sakharov One-Loop Induced Gravity (\texttt{C12}, \texttt{ODQ-5})}

In Version~79 (former Subsection~48.3.2), an exploratory Sakharov induced-gravity ansatz proposed setting the bare Einstein--Hilbert coefficient to zero and computing $1/(16\pi G)$ entirely from one-loop matter/coherence fluctuations $\sim \Lambda_{\mathrm{UV}}^2/(16\pi^2)$. Formal audit in Stages~2--3 (\texttt{ODQ-5}) retired that branch (\texttt{C12}, Tier~4) for three decisive mathematical reasons:
\begin{enumerate}[leftmargin=2em]
  \item Without a primitive tree-level $\frac{M_{\mathrm{Pl}}^2}{2}R$ kinetic term, the tree-level metric propagator does not exist and the $13\times 13$ curved ADM velocity Hessian loses its $(5+,1-)$ DeWitt gravitational block (Section~\ref{sec:stage3-curved-adm}).
  \item As established by the Seeley--DeWitt coefficient $a_1 = \frac{1}{6}R I + E$ in Section~\ref{sec:rg-heat-kernel}, one-loop quadratic $\Lambda_{\mathrm{reg}}^2$ shifts depend on the arbitrary regulator scheme and receive uncontrolled contributions from unknown trans-Planckian degrees of freedom at $E \ge \Lambda_{\mathrm{UV}}$.
  \item Retaining $\mathcal{L}_{\mathrm{EH}} = \frac{R}{16\pi G} = \frac{M_{\mathrm{Pl}}^2}{2}R$ as a primitive low-energy parameter absorbed against $a_1$ counterterms yields a completely consistent, predictive Wilsonian EFT for all $E \ll \Lambda_{\mathrm{UV}}$, while honestly leaving the microscopic origin of $M_{\mathrm{Pl}}$ above $\Lambda_{\mathrm{UV}}$ to \texttt{OPEN-1} (\texttt{NOT\_SOLVED}).
\end{enumerate}

\subsection{Proved Finite-Matrix Ginsparg--Wilson Chiral Fermions, Non-Zero Neuberger Overlap Construction, and $SU(3)_c\times SU(2)_L\times U(1)_Y \hookrightarrow U(6)$ Interface}

To couple chiral Standard Model multiplets to SIT's $(\rho_t, \psi)$ fields on a 4D lattice with explicit epistemic demarcation between machine-checked finite-matrix theorems and analytical Brillouin-zone dispersion, Stage~4 Phase~4 (\texttt{Stage4\allowbreak{}Matter.\allowbreak{}Chiral\allowbreak{}Overlap\allowbreak{}SMInterface}, theorem \texttt{SITStage4\allowbreak{}Proofs.\allowbreak{}chiral\_\allowbreak{}overlap\_\allowbreak{}sm\_\allowbreak{}interface\_\allowbreak{}holds}) and Stage~5.1 (\texttt{Stage5\allowbreak{}Extensions.\allowbreak{}Matter\allowbreak{}And\allowbreak{}ADMDomain\allowbreak{}Refinements}, $51$ declarations) establish in Lean~4.30:
\begin{enumerate}[leftmargin=2em]
  \item \textbf{Finite-Matrix Ginsparg--Wilson Algebra and Non-Zero Neuberger Overlap Witness:} Frozen Stage~4 proves that for any matrix $D \in M_N(\mathbb{C})$ with Hermitian $\gamma_5$ ($\gamma_5^2 = I$, $D^\dagger = \gamma_5 D \gamma_5$) obeying $D\gamma_5 + \gamma_5 D = a D\gamma_5 D$ ($a > 0$), $\hat{\gamma}_5 := \gamma_5(I - a D)$ satisfies $\hat{\gamma}_5^2 = I$, $D\hat{\gamma}_5 + \gamma_5 D = 0$, $D + D^\dagger = a D^\dagger D$, and $D D^\dagger = D^\dagger D$ (noting candidly that $D = 0$ satisfies that abstract premise alone). Stage~5.1 eliminates triviality constructively by building $D_{\mathrm{ov}} = \frac{1}{a}(I - \gamma_5 H_W)$ from a $\gamma_5$-Hermitian unitary Wilson sign factor $H_W$ and proving both Ginsparg--Wilson compliance (\texttt{neuberger\allowbreak{}Overlap\_\allowbreak{}satisfies\_\allowbreak{}ginsparg\_\allowbreak{}wilson}) and $D_{\mathrm{ov}} \ne 0$ (\texttt{neuberger\allowbreak{}Overlap\_\allowbreak{}witness\_\allowbreak{}nontrivial}).
  \item \textbf{Injective Standard Model Gauge Multiplet Embedding:} Using block-diagonal unitary matrices, $\mathrm{embedSM} : U(3) \times U(2) \times U(1) \hookrightarrow U(6)$ embeds $SU(3)_c \times SU(2)_L \times U(1)_Y$ injectively as a multiplicative group homomorphism preserving both trace additivity $\mathrm{Tr}(\mathrm{embedSM}(g_3,g_2,g_1)) = \mathrm{Tr}\,g_3 + \mathrm{Tr}\,g_2 + \mathrm{Tr}\,g_1$ and determinant multiplicativity $\det(\mathrm{embedSM}(g_3,g_2,g_1)) = \det g_3 \cdot \det g_2 \cdot \det g_1$.
  \item \textbf{Constructive Massive Fermion-Determinant SIT+SM Probability Measure:} Frozen Stage~4 proves gauge-invariant measure normalization $W_{\min} \le Z_{\mathrm{SIT+SM}} \le W_{\max}$ conditional on any continuous gauge-invariant functional $0 < W_{\min} \le W_{\mathrm{ferm}} \le W_{\max}$, while Stage~5.1 constructs the concrete massive determinant weight $W_{\mathrm{ferm}}(U) := \det(D_{\mathrm{ov}}[U]^\dagger D_{\mathrm{ov}}[U] + m_f^2 I_N)$ ($m_f > 0$) and proves the Hermitian eigenvalue bounds $0 < m_f^{2N} \le W_{\mathrm{ferm}}(U) \le ((2/a)^2 + m_f^2)^N$ and exact unitary gauge invariance (\texttt{det\_\allowbreak{}conj\allowbreak{}Transpose\_\allowbreak{}mul\_\allowbreak{}self\_\allowbreak{}add\_\allowbreak{}mass\allowbreak{}Sq\_\allowbreak{}ge}, \texttt{massive\allowbreak{}Fermion\allowbreak{}Det\allowbreak{}Weight\_\allowbreak{}constructs\_\allowbreak{}sm\_\allowbreak{}interface}). Meanwhile, momentum-space single-pole dispersion at $p = 0$ excluding the $15$ Brillouin-zone corner doublers ($0 < m_0 < 2r$), exponential locality on L\"uscher-admissible plaquette fields, and massless single-Weyl chiral determinant Berry phases (\texttt{OPEN-5A}) remain analytical/open outside Lean scope (Section~\ref{sec:stage4-matter-holonomy}).
\end{enumerate}

\subsection{Conditional Isolation of the Level-2 Arithmetic Spectral Bridge (\texttt{HELD\_PER\_ODQ7})}

Earlier exploratory drafts investigated an operator-theoretic Level-2 arithmetic correspondence comparing logarithmic prime sources and Riemann zeta zeros to resonant modes of an auxiliary operator $H_{\mathrm{SIT}}$. Under \texttt{QG49} epistemic hygiene (\texttt{ODQ-7}, formalized in Lean~4.30 module \texttt{SitQg.Stage2.QG31To35} and \texttt{SITF43\allowbreak{}To\allowbreak{}F46\allowbreak{}Riemann\allowbreak{}Bridge}), this arithmetic sector is \textbf{logically decoupled in both directions} from the physical \texttt{M\_AUTH\_SIT\_EFT\_v1} gravitational, gauge, and experimental core:
\begin{itemize}
  \item Every physical theorem in Parts~I--IV holds unconditionally without any number-theoretic hypothesis.
  \item Conversely, SIT makes \textbf{zero claim} to an unconditional proof of the Riemann Hypothesis (\textbf{\texttt{OPEN-4: NOT\_SOLVED}}, \texttt{HELD\_PER\_ODQ7}, \texttt{rh\_solved = false}) or of any computational complexity separation such as $P \ne NP$ (\textbf{\texttt{P = NP mixing gap: NOT\_\allowbreak{}SOLVED}}). Any spectral correspondence between self-adjoint operators and zeta zeros remains strictly conditional on unproved spectral-realization hypotheses isolated in Tier~3/Tier~5.
\end{itemize}


\section{Non-Perturbative 4D Haar Wilson-Coherence Resolution of Gauge Gating and the Mass Gap Boundary}
\label{sec:mass-gap}

A central question in non-Abelian gauge physics is how short-range massive/confined gauge excitations arise and how SIT's coherence-gated gauge dynamics behave as coherence approaches zero ($\Rcoh \to 0$) versus a nonzero condensate ($v_0 > 0$). Version~81 replaces the informal 2025 continuum $f \to 0$ canonical-Hamiltonian heuristic (\texttt{C09}) with the \textbf{exact Stage-4 two-phase non-perturbative 4D hypercubic Euclidean Haar Wilson-coherence lattice resolution}, proved in Lean~4.30 (\texttt{Stage4Lattice/*}, Section~\ref{sec:stage4-lattice}), while maintaining a rock-solid boundary separating SIT's coupled lattice/EFT results from the unsolved Clay Millennium pure Yang--Mills continuum problem.

\subsection{Why Continuum $f \to 0$ (\texttt{C09}) Fails Canonically Yet Resolves on a Compact 4D Haar Lattice}

In late-2025 exploratory notes (\texttt{SIT-60} and \texttt{Coherence\_\allowbreak{}Gated\_\allowbreak{}Yang\_\allowbreak{}Mills\_\allowbreak{}Gap 3.\allowbreak{}txt}, explicitly footnoted there as non-Wightman/non-OS), setting the gauge kinetic prefactor $f(\rho_t, |\psi|)$ to vanish at $|\psi| = 0$ ($f(\rho_t, 0) = 0$) was proposed to penalize zero-coherence chromoelectric fluctuations. Within \textbf{continuum canonical Hamiltonian EFT}, however, Legendre transformation of $-\frac{1}{4}f(\rho_t,|\psi|)F_{\mu\nu}^a F^{a\mu\nu}$ produces the electric kinetic energy density
\begin{equation}
\mathcal{H}_{\mathrm{elec}}^{\mathrm{cont}} = \frac{1}{2 f(\rho_t, |\psi|)}\,\Pi_i^a \Pi^{ai},
\label{eq:continuum_pi_over_f_singularity}
\end{equation}
which diverges catastrophically as $f \to 0$, destroying uniform ellipticity/hyperbolicity and preventing weak-field Maxwell recovery. For this reason, $f(\rho_t,0) = 0$ (\texttt{C09}) is permanently Tier-4 excluded from continuum \texttt{M\_AUTH\_SIT\_EFT\_v1}, where $f_{\mathrm{EM}}(\Rcoh) = 1 + c_{\mathrm{EM}}\Rcoh^2 \ge 1$ holds strictly (\texttt{ODQ-3}).

On a \textbf{compact-group 4D Euclidean hypercubic Wilson lattice} $\Lambda_l = (\mathbb{Z}/l\mathbb{Z})^4$ with link variables $U(x,\mu) \in G \subset U(n+1)$ integrated against normalized bi-invariant Haar probability measure $dh(U_e)$, no division by $f$ ever occurs! Each oriented plaquette $(x,\mu,\nu)$ carries non-negative Wilson energy density $E_\square(U) = 1 - \frac{1}{n+1}\mathrm{Re}\,\mathrm{Tr}\,U_\square(x,\mu,\nu) \in [0, 2]$ multiplied by the symmetric 4-corner geometric-mean coherence gate $f_\square(\rho,\psi) \ge 0$ (\texttt{plaquetteGate}) in the unshifted Stage-4 coupled Euclidean action (\textbf{Model~L1}, \texttt{sit\allowbreak{}Wilson\allowbreak{}Coherence\allowbreak{}Action}):
\begin{equation}
S_{\mathrm{SIT}}^{\mathrm{lat}}(U,\psi,\rho)
= \beta\sum_{\square} f_\square(\rho,\psi)\,E_\square(U)
  + \kappa_c\sum_{x,\mu}\bigl\|\psi(x) - U(x,\mu)\psi(x+\hat{\mu})\bigr\|^2
  + \sum_x V_{\mathrm{site}}\bigl(\rho(x), \|\psi(x)\|^2\bigr).
\label{eq:sec49_coupled_lattice_action}
\end{equation}
Because $S_{\mathrm{SIT}}^{\mathrm{lat}} \ge 0$ for all $\beta, \kappa_c \ge 0$ and $f_{\mathrm{gate}}, V_{\mathrm{site}} \ge 0$, the Boltzmann weight satisfies $0 < \exp(-S_{\mathrm{SIT}}^{\mathrm{lat}}) \le 1$ pointwise, and compact Haar integration proves (\texttt{SITStage4\allowbreak{}Proofs.\allowbreak{}c09\_\allowbreak{}nonperturbative\_\allowbreak{}finiteness\_\allowbreak{}holds})
\begin{equation}
0 < Z_{\mathrm{SIT}}(G, l, \beta, \kappa_c) \le 1,
\label{eq:sec49_Z_bounds}
\end{equation}
yielding a bona fide, locally gauge-invariant probability measure $d\mu_{\mathrm{SIT}}$ even when $f_{\mathrm{gate}}(\rho, 0) = 0$. Furthermore, while unshifted Model~L1~\eqref{eq:sec49_coupled_lattice_action} governs fixed-background and conditional gauge expectations yet exhibits a reduced two-state gate-kernel sign obstruction in its isolated cross-slab factor $\exp(-\beta g_0 g_1)$ (\texttt{unshifted\_\allowbreak{}gate\_\allowbreak{}two\_\allowbreak{}by\_\allowbreak{}two\_\allowbreak{}det\_\allowbreak{}neg}, motivating an arithmetic baseline rather than constituting a full gauge-integrated coupled L1 no-go theorem), replacing the constant baseline $1$ inside $f_\square E_\square(U)$ by the 4-corner arithmetic average $\bar{f}_\square = \frac{1}{4}\sum_{c\in\square} f_c$ defines unified \textbf{Model~L3} (\texttt{sitAMGMGatedOSAction} in \texttt{Stage5\allowbreak{}Extensions.\allowbreak{}Unified\allowbreak{}AMGMGated\allowbreak{}OSPositivity}, Section~\ref{sec:stage4-lattice}), which differs from Model~L1 by the gauge-link-independent sum $\beta\,\mathrm{amgmCoherenceGap} + \kappa_t\sum_{x,\mu}(\rho(x+e_\mu)-\rho(x))^2 \ge 0$ (isolating the pure AM--GM site-gradient counterterm at $\kappa_t = 0$ or against a reference action carrying the same clock-hopping term) and simultaneously achieves within a single action exact zero-matter ($\delta_0 \otimes \mu_{\mathrm{temp}}$ and conditional $\psi = 0$) Haar disorder (with zero single-plaquette expectation for $l \ge 2$), identical conditional link law $P_{\mathrm{L3}}(U\mid\rho,\psi) = P_{\mathrm{L1}}(U\mid\rho,\psi)$, exact action equality on spatially constant backgrounds paired with explicit continuum matching of the $O(a^{-2})$ / marginal / $O(a^2)$ AM--GM radial scalar kinetic shift $\frac{3\beta}{8a^2}\int \frac{(\partial_\mu f)^2}{f}\,d^4x$ (Section~\ref{sec:stage4-lattice}), and coupled $(U,\psi,\rho)$ $2L_t$ reflection positivity across compact-ball and unbounded/Gaussian site priors.

\subsection{Two-Phase Unification: Unbroken Haar Disorder, Activation Floor, Broken-Phase EFT, and First-Order Polyakov Barrier}

The single non-perturbative lattice measure~\eqref{eq:sec49_coupled_lattice_action} unifies four rigorous regimes proved in Lean~4.30:
\begin{enumerate}[leftmargin=2em]
  \item \textbf{Exact $\psi \equiv 0$ Ultra-Local Haar Disorder (\texttt{haar\_\allowbreak{}coherence\_\allowbreak{}gating\_\allowbreak{}floor\_\allowbreak{}holds}):} In the unbroken zero-coherence sector $\psi \equiv 0$ with $f_{\mathrm{gate}}(\rho,0) = 0$, the conditional link law reduces to pure product Haar measure $\prod_e dh(U_e)$. By $SU(2)$ antithetic Haar invariance ($u \mapsto -u$), single-plaquette traces have zero mean and connected two-plaquette correlators across distinct Euclidean time slices ($x_{0,4} \ne x_{1,4}$) vanish identically ($\langle \mathrm{Tr}\,U_{\square_0}\,\mathrm{Tr}\,U_{\square_1}^*\rangle_c = 0$), corresponding to infinite lattice mass gap / zero correlation length on the zero-coherence branch.
  \item \textbf{Quantitative Coherence-Activation Bound:} Turning on any finite site coherence profile $\psi(x)$ with $0 \le f_{\mathrm{gate}}(r,s) \le L_f s$ shifts bounded gauge expectations ($|O(U)| \le M$) from their $\beta = 0$ Haar baseline by at most
  \begin{equation}
  \Bigl|\langle O\rangle_{\beta,\psi} - \langle O\rangle_{0,\psi}\Bigr|
  \le 48\,M\,\beta\,L_f\sum_{x\in\Lambda_l}\|\psi(x)\|^2,
  \label{eq:sec49_48_bound}
  \end{equation}
  while any coercive site potential $V_{\mathrm{site}}(r,s) \ge \frac{1}{2}\Delta_{\mathrm{coh}} s$ ($\Delta_{\mathrm{coh}} > 0$) imposes a strictly positive energetic cost $\sum_x V_{\mathrm{site}}(\rho(x),\|\psi(x)\|^2) \ge \frac{1}{2}\Delta_{\mathrm{coh}}\sum_x\|\psi(x)\|^2 > 0$ on every non-zero coherence excitation $\psi \ne 0$.
  \item \textbf{Broken-Phase $O(a^2)$ Recovery of \texttt{M\_AUTH\_SIT\_EFT\_v1} (\texttt{broken\_\allowbreak{}phase\_\allowbreak{}eft\_\allowbreak{}unification\_\allowbreak{}holds}):} In a constant broken-magnitude background $\|\psi(x)\|^2 = v_0^2 > 0$, $\rho(x) = \rho_0$ with $f_{\mathrm{gate}}(\rho_0, v_0^2\Rcoh^2) = f_0 + c_2\Rcoh^2$ ($f_0 > 0$), dividing by $f_0$ reproduces the exact Stage-2 constitutive dressing $f_{\mathrm{EM}}(\Rcoh) = 1 + (c_2/f_0)\Rcoh^2 \ge 1$, generates a gauge-invariant Proca screening scale $m_A^2 = \hat{q}^2\kappa_c v_0^2 / f_0 > 0$ (consistent with Elitzur's theorem $\langle U_{e,00}\rangle = 0$ without gauge fixing), and recovers the continuum quadratic Yang--Mills trace $\lim_{\varepsilon\to 0}\varepsilon^{-2}[1 - \frac{1}{n+1}\mathrm{Re}\,\mathrm{Tr}\exp(\varepsilon F)] = \frac{1}{2(n+1)}\|F\|_F^2$ at $O(a^2)$.
  \item \textbf{Heavy-Coherence First-Order Polyakov Barrier (\texttt{compact\_\allowbreak{}polyakov\_\allowbreak{}first\_\allowbreak{}order\_\allowbreak{}transition\_\allowbreak{}holds}):} Coupling the thermal Polyakov loop character $u_1$ to a heavy coherence mode $\Rcoh$ of mass $m_{\mathrm{coh}} > 0$ via $V_{\mathrm{joint}}(a_T, u_1, \Rcoh) = (1-a_T)u_1 + b_0 u_1^2 + c u_1^3 + \frac{1}{2}m_{\mathrm{coh}}^2\Rcoh^2 - \kappa u_1\Rcoh$ ($c > 0$) and completing the square yields $b_{\mathrm{eff}} = b_0 - \frac{\kappa^2}{2 m_{\mathrm{coh}}^2} < 0$ whenever $\kappa^2 > 2 b_0 m_{\mathrm{coh}}^2$. At critical temperature $a_{\mathrm{crit}} = 1 - b_{\mathrm{eff}}^2/(4c) < 1$, the confined ($u_1 = 0$) and deconfined ($u_* = -b_{\mathrm{eff}}/(2c) > 0$) phases coexist degenerately across the exact positive cubic free-energy barrier
  \begin{equation}
  V_{\mathrm{eff}}\!\left(a_{\mathrm{crit}},\,\frac{u_*}{3}\right) = \frac{4}{27}\,c\,u_*^3 > 0.
  \label{eq:sec49_polyakov_barrier}
  \end{equation}
\end{enumerate}

\subsection{Explicit Clay Millennium Boundary: Continuum Pure Yang--Mills Mass Gap Remains \texttt{NOT\_SOLVED}}

In strict adherence to \texttt{QG49} Tier~5, neither the coupled SIT coherence-lattice gap mechanisms above nor the October 2026 pure-gauge scaling-limit construction (\texttt{OpenAI-YM-20261028}, evaluated in \texttt{yang-mills-group/YM} \texttt{assessment-1.pdf} Section~8) resolves the Clay Millennium Prize problem for \textbf{unmodified four-dimensional continuum pure Yang--Mills theory} ($f \equiv 1$ with zero scalar fields): proving infinite-volume uniqueness and a strictly positive Osterwalder--Schrader Hamiltonian spectral gap $\sigma(H_{\mathrm{YM}}) \subset \{0\} \cup [\Delta, \infty)$ ($\Delta > 0$) in the $a_N \to 0$, $L_p \to \infty$ continuum limit remains \textbf{\texttt{NOT\_SOLVED}}, and is isolated conditionally inside Stage-4 Target~\texttt{(7)} (\texttt{Open\allowbreak{}AIUVFixed\allowbreak{}Point\allowbreak{}Assumptions -> Yang\allowbreak{}Mills.\allowbreak{}Joint\allowbreak{}Existence}, Section~\ref{sec:2025-vs-2026-ym}).

\section{Black Holes, Voids, and the Halfway Universe (Exploratory Level-1 Cosmology Program)}
\label{sec:black-holes-voids}
\label{sec:blackholesvoids}

\textbf{Level-1 Exploratory Status:} The cosmological scenarios in this section are \emph{exploratory Level-1 phenomenological hypotheses} separate from unchanged \texttt{M\_AUTH\_SIT\_EFT\_v1} (whose canonical vacuum satisfies $V_{\mathrm{tot}}(\rho_0, v_0) = 0$ and whose cosmological perturbation theory remains open under \texttt{OPEN-6D}). Version~81 preserves these cosmological research questions while correcting four algebraic and logical inconsistencies from earlier drafts.

\subsection{Halfway Universe Heuristic and Sign Correction of the Contrast Integral}

Conceptually, the ``Halfway Universe'' proposal investigates whether cosmic structure formation exhibits a dynamical balance between compact high-density gravitational regions (black holes and clusters) and expanding under-dense cosmic voids, motivating open questions about Bekenstein--Hawking entropy bookkeeping and late-time cosmic acceleration (without claiming to resolve the black-hole information paradox or derive a conserved cosmic coherence budget from \texttt{M\_AUTH\_SIT\_EFT\_v1}).

\paragraph{Sign Correction of the Earlier Zero-Sum Equation.}
Earlier drafts wrote $\int_{\mathcal{V}}[\rho_t(\mathbf{x}) - \rho_{t,0}]\,d^3x = \alpha\int_{\mathcal{V}}\Rcoh\,d^3x \approx 0$ and argued that positive black-hole contributions cancel against cosmic voids. Because $\Rcoh \ge 0$ everywhere and the void limit sets $\Rcoh \to 0$ and $\rho_t \to \rho_{t,0}$, voids contribute $0$ (not negative values) to the integral, so cancellation cannot occur as written. Formulating a genuine zero-mean cosmic contrast requires defining a \emph{signed} fluctuation relative to a specified spatial mean background $\bar{R}_{\mathrm{bg}}(t) := \langle\Rcoh\rangle_{\mathcal{V}}(t) > 0$,
\begin{equation}
\delta\Rcoh(\mathbf{x}, t) := \Rcoh(\mathbf{x}, t) - \bar{R}_{\mathrm{bg}}(t),
\qquad
\int_{\mathcal{V}} \delta\Rcoh(\mathbf{x}, t)\,d^3x = 0,
\label{eq:signed_cosmic_contrast}
\end{equation}
where overdense regions ($\delta\Rcoh > 0$) and underdense voids ($\delta\Rcoh < 0$) carry opposite signs by construction; whether any nonlinear solution of a cosmological SIT action links $\delta\rho_t$ linearly to $\delta\Rcoh$ remains an open Level-1 problem.

\subsection{Exploratory Dark-Sector Parameterizations and Corrected Friedmann Linearization}
\label{sec:darkmatter}

Rather than claiming that unchanged \texttt{M\_AUTH\_SIT\_EFT\_v1} already replaces particle dark matter or derives the baryonic Tully--Fisher relation, we state explicitly what each exploratory Level-1 phenomenological ansatz assumes and correct the Friedmann linearization factor:
\begin{enumerate}[leftmargin=2em]
  \item \textbf{Single-Galaxy $1/r^2$ Density Profile vs.\ Baryonic Tully--Fisher Relation:} In a phenomenological modified-Poisson ansatz $\nabla^2 V_{\mathrm{grav}}(\mathbf{x}) = 4\pi G\bigl[\rho_b(\mathbf{x}) + \rho_{\mathrm{eff,coh}}(\mathbf{x})\bigr]$ with effective density $\rho_{\mathrm{eff,coh}} := \alpha_{\mathrm{dm}}\Rcoh$ (where $\alpha_{\mathrm{dm}}$ carries mass-density dimensions), assuming a flat circular velocity curve $v(r) \approx v_{\mathrm{flat}}$ merely infers an isothermal halo profile $\rho_{\mathrm{eff,coh}}(r) \approx v_{\mathrm{flat}}^2/(4\pi G r^2)$ for that galaxy. It does \emph{not} derive the cross-galaxy baryonic Tully--Fisher scaling relation ($M_b \propto v_{\mathrm{flat}}^4$, McGaugh~2012), which links total baryonic mass across distinct galaxies to asymptotic rotation speed and remains an open target for any cosmological extension.
  \item \textbf{Corrected Friedmann Linearization Factor ($2\rho_{\mathrm{tot}}$, Not $6\rho_{\mathrm{tot}}$):} In an exploratory flat ($k = 0$) Friedmann parameterization adding an effective coherence energy-density contribution $\alpha_{\mathrm{cos}}\langle\Rcoh\rangle(z)$ (with $\alpha_{\mathrm{cos}}$ carrying energy-density dimensions),
  \begin{equation}
  H^2(z) = \frac{8\pi G}{3}\,\rho_{\mathrm{tot}}(z),
  \qquad
  \rho_{\mathrm{tot}}(z) := \rho_m(z) + \alpha_{\mathrm{cos}}\langle\Rcoh\rangle(z) + \rho_\Lambda,
  \label{eq:friedmann_exploratory}
  \end{equation}
  differentiating $H^2 = \frac{8\pi G}{3}\rho_{\mathrm{tot}}$ at fixed $\rho_m, \rho_\Lambda$ gives $2H\,\delta H = \frac{8\pi G}{3}\alpha_{\mathrm{cos}}\,\delta\langle\Rcoh\rangle$, yielding the corrected leading fractional shift
  \begin{equation}
  \frac{\delta H_0}{H_0} \approx \frac{\alpha_{\mathrm{cos}}}{2\,\rho_{\mathrm{tot}}}\,\delta\langle\Rcoh\rangle
  \label{eq:hubble_shift_corrected}
  \end{equation}
  (correcting the erroneous factor $1/(6\rho_{\mathrm{tot}})$ printed in earlier drafts). Whether $\langle\Rcoh\rangle(z)$ evolves cosmologically in a way consistent with CMB acoustic peaks, BAO, and supernova data remains an open Level-1 calculation (\texttt{OPEN-6D}).
  \item \textbf{Dimensional and Transfer-Function Requirements for Lensing and Flow Templates:} Any prospective weak-lensing deflection or galaxy--void outflow parameterization of the schematic form $\delta\theta(\mathbf{x}) \sim \mathcal{K}_{\mathrm{lens}}\,\nabla_\perp\Rcoh(\mathbf{x})$ or $v_{\mathrm{flow}} \sim \mathcal{K}_{\mathrm{flow}}\,\Delta\Rcoh$ requires explicit dimensional coupling kernels $(\mathcal{K}_{\mathrm{lens}}, \mathcal{K}_{\mathrm{flow}})$ and cosmological transfer functions derived from a specified gravitational action (subject to the mass-calibration and form-factor identifiability cautions of Section~\ref{subsec:identifiability-and-qgtcd-guardrails}). In uniform backgrounds, unchanged \texttt{M\_AUTH\_SIT\_EFT\_v1} reduces to standard General Relativity.
\end{enumerate}
% --- END INPUT sections/07_part3_sec43_to_50.tex ---
% --- BEGIN INPUT sections/08_part4_sec51_to_52.tex ---
\part{Machine-Verified Formalization, Curved 3+1D ADM Gravity, and 4D Non-Perturbative Lattice Gauge Theory}

\section{Infinite-Dimensional Sobolev Coherence Tower, Galerkin Convergence, and Sharp Spatial Cutoff Loop Anomalies}
\label{sec:stage3-sobolev}

Early machine-checked tranches of the SIT verification program (\texttt{SITTranche\allowbreak{}DCoherence\allowbreak{}Tower\allowbreak{}And\allowbreak{}Transport}) modeled ultraviolet mode projections over finite index sets \texttt{Finset N}. Stage~3 Phase~3 removed that finite-carrier restriction by constructing an infinite-dimensional Hilbert sequence carrier and weighted Sobolev tower in Lean~4.30 (\texttt{SITStage3\allowbreak{}Infinite\allowbreak{}Sobolev\allowbreak{}Coherence\allowbreak{}Tower}, $66$ strict standard-axiom theorems audited in \texttt{SITStage3Phase3Audit}), while Stage~4 Phase~0 paired that theorem stack with the November~2026 perturbative diagnosis of sharp spatial mode cutoffs (\emph{The Hawking--Page transition at weak coupling}, \texttt{hexagon:2611.00175v1}) to establish why non-perturbative 4D gauge-coherence physics must transition from spatial Galerkin projectors $P_N$ to a 4D hypercubic Euclidean Wilson lattice.

\subsection{Infinite-Dimensional $\ell^2(\mathbb{N})$ Carrier and Strict Weighted Sobolev Nesting $H^2(\mu) \subsetneq H^1(\mu) \subsetneq \ell^2$}

On a compact Riemannian three-manifold $(\Sigma^3, q)$ (such as the periodic three-torus $\mathbb{T}^3$ or round three-sphere $S^3$), the non-negative Laplace--Beltrami operator $-\Delta_q \ge 0$ possesses a countable orthonormal basis $\{\phi_n\}_{n\in\mathbb{N}}$ in $L^2(\Sigma^3, d\mu_q)$ with eigenvalues $0 \le \mu_0 \le \mu_1 \le \dots \to +\infty$. Expanding real scalar field components in this eigenbasis identifies $L^2(\Sigma^3, d\mu_q)$ isometrically with the infinite-dimensional real sequence space
\begin{equation}
\mathrm{L2Seq} := \ell^2(\mathbb{N}; \mathbb{R})
= \left\{ u : \mathbb{N} \to \mathbb{R} \;\middle|\; \|u\|_{\ell^2}^2 := \sum_{n=0}^\infty |u_n|^2 < \infty \right\},
\label{eq:L2Seq_def}
\end{equation}
equipped with canonical basis vectors $e_n := \mathrm{lp.single}\;2\;n\;1$ satisfying $\langle e_n, e_m\rangle = \delta_{nm}$ and $\|e_n\|_{\ell^2} = 1$. Because $\{e_n\}_{n\in\mathbb{N}}$ is an infinite linearly independent family, Lean theorem \texttt{l2\allowbreak{}Seq\_\allowbreak{}not\_\allowbreak{}finite\allowbreak{}Dimensional} proves unconditionally that $\neg\,\mathrm{FiniteDimensional}\,\mathbb{R}\,\mathrm{L2Seq}$.

Given any non-negative spectral weight sequence $\mu : \mathbb{N} \to \mathbb{R}_{\ge 0}$, the kinetic-energy ($H^1$) and elliptic-operator graph ($H^2$) domains are formalized as $\mathbb{R}$-submodules of $\mathrm{L2Seq}$:
\begin{align}
H^1(\mu) &:= \left\{ u \in \mathrm{L2Seq} \;\middle|\; \|u\|_{H^1(\mu)}^2 := \sum_{n=0}^\infty \bigl(1 + \mu_n\bigr) u_n^2 < \infty \right\}, \label{eq:H1_mu_def}\\
H^2(\mu) &:= \left\{ u \in \mathrm{L2Seq} \;\middle|\; \|u\|_{H^2(\mu)}^2 := \sum_{n=0}^\infty \bigl(1 + \mu_n\bigr)^2 u_n^2 < \infty \right\}. \label{eq:H2_mu_def}
\end{align}
From $1 \le 1 + \mu_n \le (1 + \mu_n)^2$, monotonicity of non-negative infinite series yields $\|u\|_{\ell^2}^2 \le \|u\|_{H^1(\mu)}^2 \le \|u\|_{H^2(\mu)}^2$ and continuous inclusions $H^2(\mu) \subseteq H^1(\mu) \subseteq \mathrm{L2Seq}$.

\begin{theorem}[Strict Infinite-Dimensional Sobolev Subspace Separation]
\label{thm:sobolev_strict_nesting}
For the concrete unbounded spectral sequence $\mu_n^{\mathrm{std}} := (n + 1)^2$ ($n \in \mathbb{N}$), the Sobolev submodules satisfy strict proper inclusions
\begin{equation}
H^2(\mu^{\mathrm{std}}) \;\subsetneq\; H^1(\mu^{\mathrm{std}}) \;\subsetneq\; \ell^2(\mathbb{N}; \mathbb{R}).
\label{eq:strict_sobolev_chain}
\end{equation}
Specifically, $w^{(1)}_n := (n + 1)^{-2}$ satisfies $w^{(1)} \in H^1(\mu^{\mathrm{std}}) \setminus H^2(\mu^{\mathrm{std}})$, and $w^{(0)}_n := (n + 1)^{-1}$ satisfies $w^{(0)} \in \ell^2(\mathbb{N}) \setminus H^1(\mu^{\mathrm{std}})$ (\texttt{h2\allowbreak{}Submodule\_\allowbreak{}strict\_\allowbreak{}lt\_\allowbreak{}h1\allowbreak{}Submodule}, \texttt{h1\allowbreak{}Submodule\_\allowbreak{}strict\_\allowbreak{}lt\_\allowbreak{}top}).
\end{theorem}

\begin{proof}
For $w^{(1)}_n = (n+1)^{-2}$, $\sum_{n=0}^\infty (w^{(1)}_n)^2 = \sum_{m=1}^\infty m^{-4} < \infty$ and $\sum_{n=0}^\infty (1 + (n+1)^2)(n+1)^{-4} = \sum_{m=1}^\infty (m^{-4} + m^{-2}) < \infty$ by $p$-series convergence ($p > 1$), hence $w^{(1)} \in H^1(\mu^{\mathrm{std}})$. However, $(1 + (n+1)^2)^2 (w^{(1)}_n)^2 = \bigl((n+1)^{-2} + 1\bigr)^2 \ge 1$ fails to converge to zero as $n \to \infty$, so $w^{(1)} \notin H^2(\mu^{\mathrm{std}})$. Analogously, $(w^{(0)}_n)^2 = (n+1)^{-2}$ is summable, whereas $(1 + (n+1)^2)(w^{(0)}_n)^2 \ge 1$ diverges, proving $w^{(0)} \in \ell^2 \setminus H^1(\mu^{\mathrm{std}})$.
\end{proof}

\paragraph{Weyl Asymptotics Remark (\texttt{Finding (e)}).}
Note that $\mu_n^{\mathrm{std}} = (n+1)^2 = O(n^2)$ models 1D Sturm--Liouville / circle eigenvalue growth, whereas on a compact 3-manifold $(\Sigma^3, q)$ Weyl's law gives $\mu_n \sim c_W n^{2/3}$. Both satisfy $\mu_n \ge 0$, $\mu_n \to +\infty$, and strict nesting $H^2(\mu) \subsetneq H^1(\mu) \subsetneq \ell^2$ (for $\mu_n \sim n^{2/3}$, $u_n = (n+1)^{-7/10} \in \ell^2 \setminus H^1$ and $u_n = (n+1)^{-11/10} \in H^1 \setminus H^2$), while the denser 3D spectrum $O(n^{2/3})$ explains why pointwise 3D $L^4(\Sigma^3)$ Ladyzhenskaya control requires geometric 3D Sobolev interpolation (\texttt{hLady}, \texttt{OPEN-3}) rather than 1D sequence majorization.

\subsection{Spectral Galerkin Projectors $P_N$, Strong Convergence, and Quantitative Tail Rates}

For each mode cutoff $N \in \mathbb{N}$, define the continuous linear spectral projector $P_N : \mathrm{L2Seq} \to \mathrm{L2Seq}$ (\texttt{galerkinProj N}) by coordinate truncation $(P_N u)_n := u_n$ for $n < N$ and $0$ for $n \ge N$. The audited Lean~4.30 suite establishes:
\begin{enumerate}[leftmargin=2em]
  \item \textbf{Orthogonal Projector Algebra:} $\|P_N\|_{\mathrm{op}} \le 1$, $P_N^2 = P_N$, $P_N P_M = P_{\min(N,M)}$, $\langle P_N u, v\rangle = \langle u, P_N v\rangle$ ($P_N^* = P_N$), $P_N(\mathrm{L2Seq}) \subset H^2(\mu)$, and exact Pythagorean decomposition $\|u\|_{\ell^2}^2 = \|P_N u\|_{\ell^2}^2 + \|u - P_N u\|_{\ell^2}^2$.
  \item \textbf{Strong Operator Convergence vs.\ Uniform-Norm Obstruction:} For every $u \in \mathrm{L2Seq}$, $\|u - P_N u\|_{\ell^2} \to 0$ as $N \to \infty$ (\texttt{galerkin\allowbreak{}Proj\_\allowbreak{}tendsto\_\allowbreak{}nhds}). Conversely, for every $N \in \mathbb{N}$, $\|e_N - P_N e_N\|_{\ell^2} = \|e_N\|_{\ell^2} = 1$ (\texttt{galerkin\allowbreak{}Proj\_\allowbreak{}tail\_\allowbreak{}basis\allowbreak{}Vec\_\allowbreak{}norm}), proving that $P_N \not\to I$ in the uniform operator norm $\|\cdot\|_{\mathrm{op}}$ precisely because $\mathrm{L2Seq}$ is infinite-dimensional.
  \item \textbf{Quantitative Sobolev Tail Decay Rates:} Whenever $u \in H^2(\mu)$ and $\mu_n \ge \Lambda_N^2 > 0$ for all $n \ge N$, coordinatewise majorization $(1 + \Lambda_N^2) \le (1 + \mu_n)$ on $n \ge N$ yields exact ultraviolet error bounds (\texttt{galerkin\allowbreak{}Proj\_\allowbreak{}l2\_\allowbreak{}defect\_\allowbreak{}bound}, \texttt{galerkin\allowbreak{}Proj\_\allowbreak{}h1\_\allowbreak{}defect\_\allowbreak{}bound}):
  \begin{equation}
  \|u - P_N u\|_{\ell^2}^2 \le \frac{1}{\bigl(1 + \Lambda_N^2\bigr)^2}\,\|u\|_{H^2(\mu)}^2,
  \qquad
  \|u - P_N u\|_{H^1(\mu)}^2 \le \frac{1}{1 + \Lambda_N^2}\,\|u\|_{H^2(\mu)}^2.
  \label{eq:sobolev_tail_bounds}
  \end{equation}
  \item \textbf{Isometric Gauge Conjugation:} For any linear isometry $U : \mathrm{L2Seq} \simeq_{\ell_i} \mathrm{L2Seq}$, the conjugated projector $P_N^{(U)} := U P_N U^{-1}$ is self-adjoint, idempotent, intertwines $U$ ($P_N^{(U)}(Uu) = U(P_N u)$), converges strongly to $I$, and transfers both quantitative tail bounds~\eqref{eq:sobolev_tail_bounds} onto $U(H^2(\mu))$. (Per \texttt{Finding (e)}, this formalizes abstract Hilbert-space unitary conjugation on real $\mathrm{L2Seq}$, whereas continuum local multiplication $\psi(x) \mapsto e^{i\alpha(x)}\psi(x)$ intertwining $-\mathbf{D}^2$ on $L^2(\Sigma^3,\mathbb{C})$ belongs to residual \texttt{OPEN-3}.)
\end{enumerate}

\subsection{Why Sharp Spatial Mode Cutoffs $P_N$ Break Boost Covariance in 4D Gauge Loops (\texttt{hexagon:2611.00175v1})}
\label{sec:sharp-cutoff-anomaly}

While spectral Galerkin projectors $P_N$ cleanly control fixed-slice Sobolev approximations, a crucial question is whether sharp spatial or spherical-harmonic mode truncation ($n \le J$) can define a Lorentz/Euclidean-covariant non-perturbative 4D quantum gauge theory. Sections~6.1 and~9 of the November~2026 perturbative analysis of 4D gauge theory on $\mathbb{R}\times S^3$ (\emph{The Hawking--Page transition at weak coupling}, \texttt{hexagon:2611.00175v1}, \url{https://hexagonmath.org/pdf/2611.00175v1}) prove quantitatively that it cannot, due to three distinct loop-level obstructions:
\begin{enumerate}[leftmargin=2em]
  \item \textbf{Christ--Lee Operator-Ordering Potentials ($V_{\mathrm{CL}}$):} Eliminating $A_0$ via Gauss's law in Coulomb gauge ($\nabla\cdot\mathbf{A} = 0$) generates the nonlocal Coulomb interaction $H_{\mathrm{C}} = \frac{g_{\mathrm{YM}}^2}{2}\int\!\!\int \rho^a(x) K_{ab}[\mathbf{A}](x,y)\rho^b(y)$ with Faddeev--Popov kernel $K = M^{-1}(-\nabla^2)M^{-1}$ ($M = -\nabla\cdot\mathbf{D}[\mathbf{A}]$). Passing from Weyl-ordered canonical variables to normal-ordered Fock operators induces explicit $O(\hbar^2)$ Christ--Lee local/quasi-local counter-potentials $V_{\mathrm{CL}}$ required to preserve Ward identities.
  \item \textbf{Quadratic Scalar Mass Counterterms from Boost Boundary Violations:} Because Lorentz and conformal boost generators couple adjacent harmonic shells $n \leftrightarrow n+1$, truncating sharply at $n \le J$ severs boost cancellations across the cutoff shell $n = J$. Explicit summation of scalar self-energy diagrams at harmonic cutoffs $J = 0, 1, \dots, 6$ yields
  \[
  \Sigma_{\mathrm{scal}}(J) = \bigl[-1.594,\;-29.477,\;-84.509,\;-170.044,\;-281.395,\;-422.012,\;-588.717\bigr]
  \]
  (with $J=0$ matching the exact plane-wave matrix model value $-51/32 = -1.59375$), fitted to high precision by $-13.9(3)\,J^2 - 15\,J$ (ordinary least-squares coefficient $-13.9722$), corresponding across modes $E = 1, 2, 3$ ($\delta E \simeq \delta m^2/(2E)$) to a severe cutoff-dependent quadratic mass shift $\delta m^2 \approx -28\,J^2$ (versus the covariant supersymmetric target $(2\zeta_3 - 1)/4 \approx 0.351$).
  \item \textbf{Finite $-\frac{4}{3}g^2$ Coulomb Propagator Shift and Multi-Trace Anomalies:} Incomplete cancellation between transverse-gluon and instantaneous-Coulomb loops under sharp truncation induces both non-covariant double-trace terms proportional to $(n_B - n_F)$ ($1/60$ in pure Yang--Mills) and a finite non-zero renormalization of the Coulomb gauge-kinetic kernel:
  \begin{equation}
  \delta h_4 = -\frac{4}{3}\,h_2^{\mathrm{Coul}},
  \qquad
  (-\nabla^2)^{-1} \;\longrightarrow\; \left(1 - \frac{4}{3}g^2\right)(-\nabla^2)^{-1}
  \label{eq:coulomb_4_3_anomaly}
  \end{equation}
  (confirmed by independent numerical sector fits $-1.3332$ and $-1.3296$).
\end{enumerate}

\paragraph{Architectural Consequence for SIT.}
Equation~\eqref{eq:coulomb_4_3_anomaly} establishes rigorously why sharp spatial mode cutoffs $P_N$ must not be used to construct non-perturbative 4D gauge-coherence measures. Replacing $P_N$ with a \textbf{4D hypercubic Euclidean Wilson lattice} (Section~\ref{sec:stage4-lattice}) preserves exact local compact-group gauge invariance at every finite spacing $a$, satisfies hypercubic axis-permutation $S_4$ symmetry and Osterwalder--Schrader reflection positivity (for both pure Wilson gauge fields in Target~(4) and AM--GM-gated coupled $(U,\psi,\rho_t)$ configurations in Model~L3, Section~\ref{sec:stage4-lattice}) without boost-boundary counterterms, and supports compact Haar integration directly.

\SITConditionalFigure{fig6_eft_hierarchy_scales}{Multi-decade energy hierarchy of the canonical effective field theory \texttt{M\_AUTH\_SIT\_EFT\_v1} (\texttt{BM\_PHYS\_EFT\_1}) from laboratory $^{87}\mathrm{Rb}$ BEC ($8.6\times 10^{-22}\,\mathrm{GeV}$) and optical/nuclear clocks ($\sim 10^{-9}\,\mathrm{GeV}$) through Standard Model electroweak/LHC scales and SIT portal/Proca masses ($m_A = 6.06\times 10^{13}\,\mathrm{GeV}$, $v_0 = 2.0\times 10^{14}\,\mathrm{GeV}$, $m_1 = 2.45\times 10^{14}\,\mathrm{GeV}$, $m_2 = 2.83\times 10^{16}\,\mathrm{GeV}$) up to the dimension-6 and gravitational cutoffs $\Lambda_{\mathrm{UV}} = \Lambda_{\mathrm{sym}} = 1.414\times 10^{18}\,\mathrm{GeV}$, $\Lambda_{\mathrm{EM}} = 2.0\times 10^{18}\,\mathrm{GeV}$, and $M_{\mathrm{Pl}} = 2.435\times 10^{18}\,\mathrm{GeV}$, alongside Weinberg loop suppression $(E/\Lambda_{\mathrm{UV}})^{2L+2}$ for $L = 0,1,2,3$.}{fig:eft-hierarchy-scales}{Displays the 40-order-of-magnitude hierarchy of \texttt{BM\_PHYS\_EFT\_1} energy scales and verifies exact identity $(\Lambda_{\mathrm{sym}}/\Lambda_{\mathrm{EM}})^2 = (f_{\mathrm{EM}0}-1)/(f_{\mathrm{SM}0}-1) = 0.5$ ($\mathrm{rtol} \le 8.1\times 10^{-16}$) together with monotone Weinberg loop suppression $(E/\Lambda_{\mathrm{UV}})^{2L}$ ($9.8\times 10^{-29}$ per loop at $\sqrt{s}=14\,\mathrm{TeV}$).}

\section{Curved $3+1$D ADM Canonical Gravity, $13\times 13$ Cholesky--DeWitt Legendre Inversion, and Dirac--Bergmann Closure}
\label{sec:stage3-curved-adm}

To verify that the non-minimal constitutive dressings $f_{\mathrm{EM}}(\Rcoh)$ and $f_{\mathrm{SM}}(\Rcoh)$ and portal potential $U_{\mathrm{link}}(\rho_t, \psi)$ do not spoil general covariance or introduce anomalous second-class/ternary constraints at the full nonlinear curved level, Stage~3 Phase~4, Stage~5, and Stage~5.1 carried out a complete $3+1$D Arnowitt--Deser--Misner (ADM) Hamiltonian analysis of \texttt{M\_AUTH\_SIT\_EFT\_v1}, verified across $74/74$ symbolic and exact rational reverse-AD checks in SymPy~1.14.0 (\texttt{verify\_\allowbreak{}adm\_\allowbreak{}hypersurface\_\allowbreak{}closure.\allowbreak{}py}), $181$ strict standard-axiom kernel declarations in Stage~3 (\texttt{Stage3ADM/*}, \texttt{SITStage3\allowbreak{}Curved\allowbreak{}ADMHypersurface\allowbreak{}Closure}), $227$ unconstrained 6-component off-diagonal pointwise 3-jet kernel declarations in Stage~5 (\texttt{Stage5\allowbreak{}Extensions.\allowbreak{}Off\allowbreak{}Diagonal\allowbreak{}Curved\allowbreak{}ADMClosure}), and Sylvester positive-definite domain specialization \texttt{IsSPD3} in Stage~5.1 (\texttt{Stage5\allowbreak{}Extensions.\allowbreak{}Matter\allowbreak{}And\allowbreak{}ADMDomain\allowbreak{}Refinements}, $51$ declarations).

\subsection{$3+1$D ADM Foliation, $36$-Dimensional Phase Space, and $13\times 13$ DeWitt--Cholesky Legendre Inversion}

In signature $(-,+,+,+)$, foliating $(\mathcal{M}^4, g_{\mu\nu})$ into spatial hypersurfaces $(\Sigma_t, q_{ij})$ with lapse $N > 0$ and shift $N^i$ gives
\begin{equation}
ds^2 = -\bigl(N^2 - N_k N^k\bigr)dt^2 + 2 N_i\,dt\,dx^i + q_{ij}\,dx^i dx^j,
\qquad
\sqrt{-g} = N\sqrt{q},
\label{eq:adm_metric_foliation}
\end{equation}
with future-directed unit normal $n_\mu = (-N,0,0,0)$, $n^\mu = (N^{-1}, -N^i/N)$ ($n_\mu n^\mu = -1$), and extrinsic curvature $K_{ij} = \frac{1}{2N}(\dot{q}_{ij} - D_i N_j - D_j N_i)$. Writing $\psi = u + iv$ in real components and modeling $\mathcal{L}_{\mathrm{matter}} = -\frac{1}{2}g^{\mu\nu}\partial_\mu\chi\partial_\nu\chi - V_{\mathrm{mat}}(\chi)$ by a representative neutral screened scalar field $\chi$, the full canonical system spans \textbf{$18$ configuration coordinates},
\[
\mathcal{Q} = \bigl(q_{ij}\;[6],\; A_i\;[3],\; \rho_t\;[1],\; u\;[1],\; v\;[1],\; \chi\;[1],\; N\;[1],\; N^i\;[3],\; A_0\;[1]\bigr),
\]
and $18$ conjugate momenta $\mathcal{P}$, forming a \textbf{$36$-dimensional unconstrained canonical phase space}.

\paragraph{Primary Constraints and $13\times 13$ Dynamical Velocity Hessian.}
Because $\mathcal{L}_{\mathrm{AUTH}}$ contains no time derivatives of $(N, N^i, A_0)$, five primary first-class constraints hold identically:
\begin{equation}
\pi_N \approx 0, \qquad \pi_{N^i} \approx 0 \quad (i=1,2,3), \qquad \pi_{A_0} \approx 0.
\label{eq:adm_primary_constraints}
\end{equation}
Across the remaining $13$ dynamical velocities $(\dot{q}_{ij}[6], \dot{\rho}_t[1], \dot{u}[1], \dot{v}[1], \dot{A}_i[3], \dot{\chi}[1])$, the ultralocal Hessian matrix $\mathcal{M}_{AB} = \partial^2\mathcal{L}_{\mathrm{ADM}}/\partial\dot{Q}^A\partial\dot{Q}^B$ block-diagonalizes cleanly into five sectors:
\begin{enumerate}[leftmargin=2em]
  \item \textbf{Gravitational $6\times 6$ DeWitt Block:} $\pi^{ij} = -\frac{\sqrt{q}}{2\kappa_G}(K^{ij} - q^{ij}K)$ with inverse DeWitt supermetric $G_{ijkl} = \frac{1}{2}(q_{ik}q_{jl} + q_{il}q_{jk} - q_{ij}q_{kl})$ of signature $(5+, 1-)$ (five positive shear modes plus one negative trace mode constrained by $\mathcal{H}_\perp \approx 0$).
  \item \textbf{Positive Matter-Gauge Blocks ($1 + 2 + 3 + 1 = 7$):} $\pi_\rho = \frac{\sqrt{q}\kappa_t}{N}(\dot{\rho}_t - N^k\partial_k\rho_t)$, $\pi_u = \frac{\sqrt{q}\kappa_c}{N}(\dot{u} + \hat{q}A_0 v - N^k D_k u)$, $\pi_v = \frac{\sqrt{q}\kappa_c}{N}(\dot{v} - \hat{q}A_0 u - N^k D_k v)$, $\Pi^i = \frac{\sqrt{q}f_{\mathrm{EM}}(\Rcoh)}{N}q^{ij}(\dot{A}_j - \partial_j A_0 - N^k F_{kj})$, and $\pi_\chi = \frac{\sqrt{q}f_{\mathrm{SM}}(\Rcoh)}{N}(\dot{\chi} - N^k\partial_k\chi)$.
\end{enumerate}
Parameterizing single-chart symmetric positive-definite spatial metrics via an upper-triangular Cholesky factor $q = L^{\mathsf{T}} L$ (\texttt{Chol3} with positive diagonal entries $L_{11}, L_{22}, L_{33} > 0$, so $\det q = (L_{11}L_{22}L_{33})^2 > 0$ automatically), \texttt{Stage3\allowbreak{}ADM.\allowbreak{}Pillar2\allowbreak{}Legendre} proves in Lean~4.30 that the $13\times 13$ velocity Hessian determinant is strictly non-vanishing ($\det\mathcal{M}_{13\times 13} \ne 0$) and the Legendre transformation $\dot{Q}^A \leftrightarrow P_A$ is globally bijective on every SPD chart.

\subsection{Secondary Constraints and Exact Gauss Sign Crosswalk (\texttt{Finding (c)})}

Legendre inversion casts the canonical Lagrangian into exact Dirac--ADM form:
\begin{equation}
\mathcal{L}_{\mathrm{ADM}}
= \pi^{ij}\dot{q}_{ij} + \pi_\rho\dot{\rho}_t + \pi_u\dot{u} + \pi_v\dot{v} + \Pi^i\dot{A}_i + \pi_\chi\dot{\chi}
  - \Bigl(N\,\mathcal{H}_\perp + N^i\,\mathcal{H}_i^{\mathrm{cov}} - A_0\,G_{\mathrm{Gauss}}\Bigr),
\label{eq:L_ADM_canonical}
\end{equation}
where preservation of Eq.~\eqref{eq:adm_primary_constraints} ($\dot{\pi}_N, \dot{\pi}_{N^i}, \dot{\pi}_{A_0} \approx 0$) enforces five secondary constraints:
\begin{align}
\mathcal{H}_\perp &= \frac{2\kappa_G}{\sqrt{q}}\Bigl(q_{ik}q_{jl} - \frac{1}{2}q_{ij}q_{kl}\Bigr)\pi^{ij}\pi^{kl}
  - \frac{\sqrt{q}}{2\kappa_G}R^{(3)}
  + \frac{1}{2\sqrt{q}}\left[\frac{\pi_\rho^2}{\kappa_t} + \frac{\pi_u^2 + \pi_v^2}{\kappa_c} + \frac{q_{ij}\Pi^i\Pi^j}{f_{\mathrm{EM}}} + \frac{\pi_\chi^2}{f_{\mathrm{SM}}}\right]
  \nonumber\\
&\quad + \sqrt{q}\Bigl[
    \frac{\kappa_t}{2}q^{ij}\partial_i\rho_t\partial_j\rho_t
    + \frac{\kappa_c}{2}q^{ij}(D_i u D_j u + D_i v D_j v)
    \nonumber\\
&\qquad\qquad
    + \frac{f_{\mathrm{EM}}}{4}q^{ik}q^{jl}F_{ij}F_{kl}
    + \frac{f_{\mathrm{SM}}}{2}q^{ij}\partial_i\chi\partial_j\chi
    + V_{\mathrm{tot}} + f_{\mathrm{SM}}V_{\mathrm{mat}}
  \Bigr], \label{eq:H_perp_def}\\
\mathcal{H}_i^{\mathrm{cov}} &= -2 q_{ij} D_k\pi^{jk} + \pi_\rho\partial_i\rho_t + \pi_u D_i u + \pi_v D_i v + F_{ij}\Pi^j + \pi_\chi\partial_i\chi, \label{eq:Hi_cov_def}\\
G_{\mathrm{Gauss}} &:= +\partial_i\Pi^i - \hat{q}\bigl(u\pi_v - v\pi_u\bigr). \label{eq:G_Gauss_def}
\end{align}

\paragraph{Exact Gauss Sign Crosswalk (\texttt{Finding (c)}).}
In the Lean~4.30 formalization (\texttt{Stage3\allowbreak{}ADM.\allowbreak{}Pillar4\allowbreak{}Curved\allowbreak{}Closure\allowbreak{}A}), the field \texttt{ndGauss} is defined with opposite overall sign, $\mathrm{ndGauss}(j) := -\sum_{k=0}^2 \partial_k\Pi_k + \hat{q}(u\pi_v - v\pi_u) = -G_{\mathrm{Gauss}}$. Consequently, the compiled Lean identity \texttt{nd\allowbreak{}Htot\_\allowbreak{}cov\_\allowbreak{}Lie\_\allowbreak{}equiv},
\[
\mathrm{ndHtotCov}(j, i) + A_i\,\mathrm{ndGauss}(j) = \mathrm{ndHtotLie}(j, i),
\]
translates under $\mathrm{ndGauss} = -G_{\mathrm{Gauss}}$ identically into the standard covariant-vs-Lie momentum generator decomposition:
\begin{equation}
\mathcal{H}_i^{\mathrm{cov}} = \mathcal{H}_i^{\mathrm{Lie}} + A_i\,G_{\mathrm{Gauss}}.
\label{eq:Hi_cov_Lie_identity}
\end{equation}

\subsection{First-Class Dirac--Bergmann Hypersurface Deformation Closure and $36 - 2\times 10 = 16$ Mode Count}

On local formal finite-order 3-jet bundles (\texttt{SymMat3Jet} and \texttt{SPDOff\allowbreak{}Diag\allowbreak{}Coupled\allowbreak{}ADMJet}), Lean~4.30 proves exact pointwise differential-polynomial identities with explicit spatial total-divergence remainders $\sum_{k=1}^3\partial_k B^k(x)$ (such as $\mathcal{P}_{\mathrm{Lie}}(\xi, N) = \mathcal{H}_\perp(\mathcal{L}_\xi N) + \partial_k B_{\mathrm{tot}}^k$). On closed compact spatial slices ($\partial\Sigma_t = \emptyset$) or for smooth compactly supported lapse, shift, and gauge test functions $(N_1, N_2, \xi^i, \eta^i, \lambda) \in C_c^\infty(\Sigma_t)$ where $\int_{\Sigma_t}\partial_k B^k\,d^3x = 0$ (while functional-analytic Fr\'echet differentiability on infinite-dimensional Sobolev/Banach phase spaces and asymptotically flat Regge--Teitelboim surface-integral boundary conditions belong to \texttt{OPEN-2B}), integrating these pointwise 3-jet identities yields smeared equal-time first-class Poisson closure:
\begin{align}
\bigl\{\mathcal{H}_\perp[N_1],\,\mathcal{H}_\perp[N_2]\bigr\}
&= \mathcal{H}_i^{\mathrm{cov}}[S^i] + \int_{\Sigma} \partial_k\bigl(2 q_{ij}\pi^{jk} S^i\bigr)\,d^3x
= \mathcal{H}_i^{\mathrm{Lie}}[S^i] + G_{\mathrm{Gauss}}[S^i A_i] + \text{bdy}, \label{eq:pb_H_H}\\
\bigl\{\mathcal{D}_{\mathrm{Lie}}[\xi],\,\mathcal{H}_\perp[N]\bigr\}
&= \mathcal{H}_\perp[\mathcal{L}_\xi N] + \text{bdy}, \label{eq:pb_DLie_H}\\
\bigl\{\mathcal{D}_{\mathrm{cov}}[\xi],\,\mathcal{H}_\perp[N]\bigr\}
&= \mathcal{H}_\perp[\mathcal{L}_\xi N] + G_{\mathrm{Gauss}}\!\left[\frac{N q_{ij}\xi^i\Pi^j}{\sqrt{q}\,f_{\mathrm{EM}}(\Rcoh)}\right] + \text{bdy}, \label{eq:pb_Dcov_H}\\
\bigl\{\mathcal{D}_{\mathrm{cov}}[\xi],\,\mathcal{D}_{\mathrm{cov}}[\eta]\bigr\}
&= \mathcal{D}_{\mathrm{cov}}\bigl[[\xi,\eta]\bigr] - G_{\mathrm{Gauss}}\bigl[\xi^i\eta^j F_{ij}\bigr] + \text{bdy}, \label{eq:pb_Dcov_Dcov}\\
\bigl\{\mathcal{H}_\perp[N],\,G_{\mathrm{Gauss}}[\lambda]\bigr\} &= \text{bdy} \approx 0,
\qquad
\bigl\{G_{\mathrm{Gauss}}[\lambda_1],\,G_{\mathrm{Gauss}}[\lambda_2]\bigr\} = 0, \label{eq:pb_H_G}
\end{align}
where $S^i := q^{ij}(N_1\partial_j N_2 - N_2\partial_j N_1)$ is the universal DeWitt hypersurface deformation vector!

\paragraph{Why $f_{\mathrm{EM}}(\Rcoh)$ and $f_{\mathrm{SM}}(\Rcoh)$ Generate Zero Anomaly in $\{\mathcal{H}_\perp, \mathcal{H}_\perp\}$.}
A decisive non-trivial check (SymPy Suite~6 and \texttt{Stage3\allowbreak{}ADM.\allowbreak{}Pillar4\allowbreak{}Curved\allowbreak{}Closure\allowbreak{}B/C}) explains why the non-minimal coherence dressings preserve Eq.~\eqref{eq:pb_H_H}: because $f_{\mathrm{EM}}((u^2+v^2)/\psi_0^2)$ and $f_{\mathrm{SM}}((u^2+v^2)/\psi_0^2)$ depend algebraically on $(u,v)$ with \emph{zero spatial derivatives} $\partial_i u, \partial_i v$, their Euler--Lagrange variations $\partial\mathcal{H}_\perp/\partial u$ and $\partial\mathcal{H}_\perp/\partial v$ multiply ultralocal $\pi_u, \pi_v$ terms proportionally to $(N_1 N_2 - N_2 N_1) \equiv 0$, contributing \textbf{zero} obstruction to the derivative-antisymmetric structure $(N_1\partial_j N_2 - N_2\partial_j N_1)$.

\paragraph{Dirac--Bergmann Degree-of-Freedom Count (\texttt{Finding (b)}).}
Because the $5$ primary constraints~\eqref{eq:adm_primary_constraints} and $5$ secondary constraints $(\mathcal{H}_\perp, \mathcal{H}_i^{\mathrm{cov}}, G_{\mathrm{Gauss}})$ commute weakly on the constraint surface with \textbf{zero ternary constraints}, all $10$ constraints are first-class, leaving
\begin{equation}
\dim\mathcal{P}_{\mathrm{phys}} = 36 - 2 \times 10 = 16 = 2 \times 8
\label{eq:dirac_bergmann_16_dims}
\end{equation}
physical phase-space dimensions per spatial point: \textbf{$8$ propagating bosonic modes} ($7$ pure-SIT modes $[2\text{ graviton TT} + 2\text{ transverse }U(1) + 1\text{ longitudinal Proca/phase} + 2\text{ radial }(\rho_t,|\psi|)]$ plus $1$ neutral matter scalar $\chi$).

\begin{table}[htbp]
\centering
\caption{Summary of curved $3+1$D ADM verification suites across SymPy~1.14.0 ($74/74$ checks) and Lean~4.30 (\texttt{Stage3ADM/*}, $181$ strict declarations + \texttt{Stage5\allowbreak{}Extensions.\allowbreak{}Off\allowbreak{}Diagonal\allowbreak{}Curved\allowbreak{}ADMClosure}, $227$ strict declarations + \texttt{Stage5\allowbreak{}Extensions.\allowbreak{}Matter\allowbreak{}And\allowbreak{}ADMDomain\allowbreak{}Refinements} Sylvester \texttt{IsSPD3} specialization), with exact \texttt{Finding (d)} and formal 3-jet scope boundaries.}
\label{tab:curved-adm-suites}
\vspace{0.4em}
\footnotesize
\begin{tabular}{|p{3.0cm}|p{7.1cm}|r|p{2.6cm}|}
\hline
\textbf{Suite / Pillar} & \textbf{Mathematical Scope Verified} & \textbf{Count} & \textbf{Engine / Status} \\
\hline
SymPy Suite 1 / \texttt{Pillar1} & $3+1$D ADM kinematics, $5$ primary constraints, shift divergence, Gauss sign lock & $12$ + $43$ & SymPy + Lean (\texttt{PASS}) \\
\hline
SymPy Suite 2 / \texttt{Pillar2} & $13\times 13$ velocity Hessian factorization $(5+,1-) \oplus 7+$, \texttt{Chol3} SPD Legendre bijection & $10$ + $26$ & SymPy + Lean (\texttt{PASS}) \\
\hline
SymPy Suite 3 / \texttt{Pillar3} & $1+1$D matter $a(x)$ and $x$-reduced 3D $\mathrm{diag}(a(x),b(x),c(x))$ with $R^{(3)} \ne 0$ & $24$ + $31$ & SymPy + Lean (\texttt{PASS}) \\
\hline
SymPy Suite 4 / \texttt{Pillar4A--E} & 3D anisotropic $\mathrm{diag}(a,b,c)$ $\{\mathcal{H}_\perp,\mathcal{H}_\perp\}$, $\{\mathcal{D},\mathcal{H}_\perp\}$ (Palatini IBP \#1/\#2 in Lean) + 6-comp.\ $\{\mathcal{D},\mathcal{D}\}$, $\{\mathcal{D},G\}$ & $12$ + $41$ & SymPy + Lean (\texttt{PASS}) \\
\hline
SymPy Suite 5 (Jet AD) & Full 6-component off-diagonal $q_{ij}(x,y,z)$ raw 2nd-order $R^{(3)}$ reverse-AD polynomial jet $\{\mathcal{H}_\perp,\mathcal{H}_\perp\}$ at $\det q = 144, 900$ & $8$ & SymPy exact (\texttt{PASS}) \\
\hline
SymPy Suite 6 / \texttt{Pillar5} & Non-minimal $f'_{\mathrm{EM}}, f'_{\mathrm{SM}}$ cancellation on off-diagonal $q_{ij}$ + Dirac $36 - 2\times 10 = 16$ count & $8$ + $4$ & SymPy + Lean (\texttt{PASS}) \\
\hline
Stage 5 / \texttt{OffDiagonalADM} & Pointwise 3-jet unconstrained 6-component symmetric $3\times 3$ off-diagonal $q_{ij}$ ($\det q > 0$, specialized to Sylvester cone \texttt{IsSPD3}) $\delta_\xi R^{(3)}$, weight-$+1$ transport \& $\{\mathcal{D}_{\mathrm{Lie}},\mathcal{H}_\perp\}$, $\{\mathcal{D}_{\mathrm{cov}},\mathcal{H}_\perp\}$ modulo $\partial_k B^k$ & $227$ & Lean + $\mathbb{Q}(\sqrt{\det q})$ (\texttt{PASS}) \\
\hline
\end{tabular}
\end{table}

\paragraph{Explicit Lean vs.\ SymPy vs.\ \texttt{OPEN-2B}/\texttt{OPEN-2} Scope Boundaries (\texttt{Finding (d)} \& Sylvester \texttt{IsSPD3}).}
Per \texttt{v1\_1} canonical candor: (1) Lean \texttt{Pillar4\allowbreak{}Curved\allowbreak{}Closure\allowbreak{}E/D} initiates gravitational $\{\mathcal{H}_\perp,\mathcal{H}_\perp\}$ cancellation after Palatini integration by parts \#1 and formalizes IBP \#2, while raw second-order symbolic differentiation of $R^{(3)}(q,\partial q,\partial^2 q)$ in $\{\mathcal{H}_\perp,\mathcal{H}_\perp\}$ on full 6-component off-diagonal $q_{ij}(x,y,z)$ is verified in SymPy Suites~4--5; (2) Stage-5 module \texttt{Stage5\allowbreak{}Extensions.\allowbreak{}Off\allowbreak{}Diagonal\allowbreak{}Curved\allowbreak{}ADMClosure} ($227$ declarations: \texttt{offdiag\_\allowbreak{}curved\_\allowbreak{}adm\_\allowbreak{}DLie\_\allowbreak{}Hperp\_\allowbreak{}closure}, \texttt{offdiag\_\allowbreak{}curved\_\allowbreak{}adm\_\allowbreak{}Dcov\_\allowbreak{}Hperp\_\allowbreak{}closure}, \texttt{od\_\allowbreak{}Htot\_\allowbreak{}weight1\_\allowbreak{}transport}, \texttt{odDelR3\_eq}) proves exact local pointwise 3-jet unconstrained 6-component symmetric $3\times 3$ off-diagonal $q_{ij}(x,y,z)$ spatial-diffeomorphism / super-Hamiltonian $\{\mathcal{D},\mathcal{H}_\perp\}$ closure modulo explicit divergences $\partial_k B^k$ directly in Lean~4.30 over the invertible locus $\det q > 0$ (closing \texttt{Finding (d)\#3}); (3) because $\det q > 0$ alone on symmetric $3\times 3$ matrices admits non-Riemannian $(-,-,+)$ counterexamples such as $\mathrm{diag}(-1,-1,1)$ ($\det = +1 > 0$, proved non-SPD in \texttt{s\_\allowbreak{}lorentz\_\allowbreak{}det\_\allowbreak{}pos\_\allowbreak{}not\_\allowbreak{}spd}), Stage-5.1/5.2 (\texttt{Stage5\allowbreak{}Extensions.\allowbreak{}Matter\allowbreak{}And\allowbreak{}ADMDomain\allowbreak{}Refinements}, $66$ declarations) enforces Sylvester's leading-principal-minor criterion \texttt{IsSPD3} ($q_{11} > 0 \wedge q_{11}q_{22}-q_{12}^2 > 0 \wedge \det q > 0$), proves equivalence with quadratic positive-definiteness (\texttt{isSPD3\_iff\_posDef}), and establishes simultaneous off-diagonal $\{\mathcal{D},\mathcal{H}_\perp\}$ closure and matter/electric/magnetic energy positivity on \texttt{SPDOff\allowbreak{}Diag\allowbreak{}Coupled\allowbreak{}ADMJet} (\texttt{spd\_\allowbreak{}offdiag\_\allowbreak{}curved\_\allowbreak{}adm\_\allowbreak{}closure\_\allowbreak{}and\_\allowbreak{}energy\_\allowbreak{}positivity}); and (4) passing from single-chart pointwise formal 3-jet identities modulo $\partial_k B^k$ to integrated constraint algebras on infinite-dimensional Sobolev/Banach phase spaces with Regge--Teitelboim asymptotic boundary conditions (\texttt{OPEN-2B}), together with quantum Wheeler--DeWitt operator factor-ordering, Schwinger commutator anomalies, global multi-chart Cauchy PDE existence, and non-perturbative Dirac observables (\texttt{OPEN-2}), remains explicitly open (Tier~5).

\subsection{Stage-6 Action-Response Diagnostics, Conditional Weak-Field E\"ot-Wash Non-Exclusion, and Primary Experimental Provenance}
\label{sec:stage6-action-diagnostics}

Version~81 replaces the retired legacy Cassini/Yukawa figure (\texttt{fig4}, which relied on an unproved conformal Jordan-frame identification and a hand-coded \texttt{np.piecewise} exclusion sketch) with two Stage-6 Lean~4.30 formal modules (\texttt{Stage6FunctionalResponse.\allowbreak{}ActionIdentityDiagnostics} [$6$ theorems] and \texttt{Stage6FunctionalResponse.\allowbreak{}UnchangedActionWeakFieldObservable} [$5$ theorems]), a conditional weak-field E\"ot-Wash non-exclusion calculation directly from unchanged \texttt{M\_AUTH\_SIT\_EFT\_v1}, and six locked primary experimental provenance records:
\begin{enumerate}[leftmargin=2em]
  \item \textbf{Multiplicative vs.\ Conformal Matter Obstruction and $U(1)$ Kinetic-Mixing Diagnostics (\texttt{ActionIdentityDiagnostics}, $6$ theorems):} For any real scalar matter field $\chi$ in four dimensions, comparing the unchanged multiplicative portal $+\sqrt{-g}\,f_{\mathrm{SM}}\mathcal{L}_{\mathrm{matter}}$ against a conformal Jordan metric $\tilde{g}_{\mu\nu} = A^2 g_{\mu\nu}$ (where $\sqrt{-\tilde{g}}\tilde{g}^{\mu\nu} = A^2\sqrt{-g}g^{\mu\nu}$ while $\sqrt{-\tilde{g}} = A^4\sqrt{-g}$) forces $A^2 = A^4 = f_{\mathrm{SM}} \implies f_{\mathrm{SM}} = 1$ (\texttt{multiplicative\_\allowbreak{}vs\_\allowbreak{}conformal\_\allowbreak{}scalar\_\allowbreak{}coeff\_\allowbreak{}forces\_\allowbreak{}unit}, \texttt{multiplicative\_\allowbreak{}vs\_\allowbreak{}conformal\_\allowbreak{}mass\_\allowbreak{}mismatch\_\allowbreak{}of\_\allowbreak{}ne\_\allowbreak{}one}). Canonical scalar normalization $\chi = e^{-z/2}X$ ($f = e^z$) preserves free scalar mass at constant $f_0$, rescales quartic couplings as $\lambda/f_0$, and satisfies $\nabla_\mu(f T^{\mu\nu}) = \mathcal{L}_m \partial^\nu f$ (\texttt{canonical\_\allowbreak{}scalar\_\allowbreak{}wave\_\allowbreak{}operator\_\allowbreak{}identity}, \texttt{canonical\_\allowbreak{}scalar\_\allowbreak{}lagrangian\_\allowbreak{}divergence\_\allowbreak{}identity}, \texttt{canonical\_\allowbreak{}scalar\_\allowbreak{}constant\_\allowbreak{}prefactor\_\allowbreak{}free\_\allowbreak{}mass\_\allowbreak{}and\_\allowbreak{}quartic}). In the gauge sector, $2\times 2$ visible/hidden $U(1)$ kinetic mixing with a massive hidden mode preserves the massless visible-photon residue at $1/z_v$ independently of the hidden normalization $z_h$ (\texttt{kinetic\_\allowbreak{}mixing\_\allowbreak{}visible\_\allowbreak{}propagator\_\allowbreak{}decomp\_\allowbreak{}and\_\allowbreak{}residue}). Consequently, earlier Jordan-frame Cassini illustrations (such as the nominal $24.24\times$ margin under conformal Model~C2 [$A = f_{\mathrm{SM}}^{\pm 1}$], or the $1/4$-smaller coupling $\alpha_{\mathrm{conf}}^{(\mathrm{C1})} = \frac{1}{4}\alpha_{\mathrm{sym}}$ under conformal Model~C1 [$g^{(J)} = f_{\mathrm{SM}}^{-1}g^{(E)}$, $A = f_{\mathrm{SM}}^{-1/2}$], Section~\ref{sec:ppn-shapiro-exact}) remain strictly \emph{conditional} on explicitly adopting one chosen conformal extension rather than an unconditional prediction of unchanged \texttt{M\_AUTH\_SIT\_EFT\_v1}.
  \item \textbf{Conditional Weak-Field E\"ot-Wash Non-Exclusion Directly from Unchanged \texttt{M\_AUTH\_SIT\_EFT\_v1} (\texttt{UnchangedActionWeakFieldObservable}, $5$ theorems):} Propagating the shared canonical generalized eigenproblem $H v = m^2 K_{\mathrm{eff}}(\beta) v$ (incorporating the Stage-6 lattice AM--GM shift $\Delta K_{\mathrm{rad}} = 3\beta/200{,}000{,}002 = 3.299999967\times 10^{-8}$ at $\beta = 2.2$) gives $m_1 = 2.449459082\times 10^{14}\,\mathrm{GeV}$ ($\lambda_1 = 8.055941\times 10^{-31}\,\mathrm{m}$) and $m_2 = 2.828462483\times 10^{16}\,\mathrm{GeV}$ ($\lambda_2 = 6.976480\times 10^{-33}\,\mathrm{m}$) (\texttt{unchanged\_\allowbreak{}action\_\allowbreak{}two\_\allowbreak{}mode\_\allowbreak{}yukawa\_\allowbreak{}decomposition\_\allowbreak{}and\_\allowbreak{}sum\_\allowbreak{}rule}, \texttt{amgm\_\allowbreak{}kinetic\_\allowbreak{}shift\_\allowbreak{}effect\_\allowbreak{}on\_\allowbreak{}spectrum\_\allowbreak{}and\_\allowbreak{}yukawa}). Evaluating the exact finite-thickness parallel-slab normal-force ratio $R_{\mathrm{slab}}(d)$ (\texttt{eotwash\_\allowbreak{}parallel\_\allowbreak{}slab\_\allowbreak{}yukawa\_\allowbreak{}integral\_\allowbreak{}exact}, \texttt{eotwash\_\allowbreak{}two\_\allowbreak{}mode\_\allowbreak{}lee2020\_\allowbreak{}certificate}) against the Lee~\emph{et~al.}\ (2020) E\"ot-Wash anchor ($\lambda < 38.6\,\mu\mathrm{m}$ at $|\alpha_Y| = 1$, $95\%$~CL, $d = 52\,\mu\mathrm{m}$) yields:
  \begin{itemize}
    \item Under a declared pressureless-dust worldline closure ($q_{\mathrm{eff}} := -\langle\mathcal{L}_m\rangle/\rho_m = 1$), $\alpha_1 = 4.744521\times 10^{-7}$, $\alpha_2 = 1.186308\times 10^{-11}$, $d/\lambda_1 = 6.45486\times 10^{25}$, and $\log_{10} R_{\mathrm{slab}}(52\,\mu\mathrm{m}) = -2.80331\times 10^{25}$: canonical \texttt{BM\_PHYS\_EFT\_1} is strictly not excluded yet exponentially indistinguishable from Newtonian gravity at $52\,\mu\mathrm{m}$.
    \item Under a homogeneous free Klein--Gordon scalar closure, cycle-averaging yields $\langle\mathcal{L}_m\rangle = 0$ (\texttt{free\_\allowbreak{}klein\_\allowbreak{}gordon\_\allowbreak{}cycle\_\allowbreak{}averaged\_\allowbreak{}lagrangian\_\allowbreak{}zero}), so $q_{\mathrm{eff}} = 0$ and $\alpha_1 = \alpha_2 = 0$ identically.
    \item As a deliberately tuned low-mass parameter-space illustration ($25.68$ decades below \texttt{BM\_PHYS\_EFT\_1}), setting $(\alpha_1, \lambda_1) = (1.0, 38.6\,\mu\mathrm{m})$ ($m_1 = 5.112\,\mathrm{meV}$) yields a uniform-slab normal-force ratio $R_{\mathrm{slab}}(52\,\mu\mathrm{m}) = 0.02797$ ($2.80\%$).
  \end{itemize}
  \textbf{Four Mandatory Modeling Caveats:} (i)~neither the $q_{\mathrm{eff}} = 1$ dust closure nor the $q_{\mathrm{eff}} = 0$ free-scalar closure derives the microscopic composite-material scalar charges of real laboratory nuclei/atoms---for Lee~\emph{et~al.}\ (2020), both source and detector scalar contrasts enter jointly via $\rho_{\mathrm{Pt}} q_{\mathrm{Pt}} - \rho_{\mathrm{epoxy}} q_{\mathrm{epoxy}}$ (\texttt{OPEN-6A}); (ii)~$R_{\mathrm{slab}}(d)$ models the normal force between uniform parallel slabs, whereas Lee~\emph{et~al.}\ (2020) measured torsional harmonics from rotating 18-fold and 120-fold patterned platinum attractor and detector bodies with epoxy-filled holes mounted on BK7 glass supports (where the $54$-frequency signal is the third harmonic of the 18-fold pattern)---reproducing the Newtonian fit and published $38.6\,\mu\mathrm{m}$ single-Yukawa exclusion after auditing Lee's dissertation (\emph{A Fourier-Bessel Test of the Gravitational Inverse-Square Law}, 2020: Chapter~11 Table~11.4 single-Yukawa fit parameters and Appendix~A corrected run-by-run harmonic torques, geometries, calibrations, separation errors, and nuisance priors) before performing a joint two-mode confidence analysis remains \texttt{OPEN-6D} (now directly instantiated on canonical \texttt{AuthModelParams} via \texttt{Stage6Addendum.\allowbreak{}ActionResponsePropagatorIdentity}, $5$ theorems, proving $0 < m_0^2 \le m_1^2$, automatic $p^2 + m_a^2 \ne 0$ for $p^2 \ge 0$, and $(p^2 K + H)^{-1} = V\,\mathrm{diag}(1/(p^2+m_a^2))\,V^\top$); (iii)~the $m_1 = 5.112\,\mathrm{meV}$ ($2.80\%$) illustration is strictly a uniform-slab parameter-space witness, not a prediction of \texttt{BM\_PHYS\_EFT\_1}; and (iv)~SymPy/NumPy agreement verifies numerical implementation, not physical modeling assumptions.
  \item \textbf{Corrected Six-Record Primary Provenance (\texttt{primary\_measurements\_manifest\_v2.json}) and Legacy Quarantine:} Multi-channel empirical fitting across Cassini ($\gamma_{\mathrm{PPN}}-1 = (+2.1\pm 2.3)\times 10^{-5}$, Bertotti~\emph{et~al.}\ 2003), Tokyo Skytree ($\alpha_{\mathrm{RS}} = (+1.4\pm 9.1)\times 10^{-5}$, $1\sigma$ combined, Takamoto~\emph{et~al.}\ 2020), MICROSCOPE ($\eta(\mathrm{Ti,Pt}) = [-1.5\pm 2.3(\mathrm{stat})\pm 1.5(\mathrm{syst})]\times 10^{-15}$, Touboul~\emph{et~al.}\ 2022), Asenbaum~\emph{et~al.}\ (2020) ($\eta(^{85}\mathrm{Rb},^{87}\mathrm{Rb}) = [+1.6\pm 1.8(\mathrm{stat})\pm 3.4(\mathrm{syst})]\times 10^{-12}$, dual-isotope rather than BEC vs.\ thermal), Lee~\emph{et~al.}\ (2020), and NIST/BACON (\emph{Phys.\ Rev.\ Lett.}\ \textbf{137}, 033201, published July~14, 2026; DOI \texttt{10.1103/g865-9mk1}, \texttt{arXiv:2512.21428}, optical frequency-ratio measurement uncertainties at or below $3.2\times 10^{-18}$ across phase-stabilized optical fiber links---an experimental uncertainty budget rather than a universal clock floor or direct SIT coupling bound) remains fail-closed \texttt{BLOCKED} pending \texttt{OPEN-6}. Legacy Version~80.2 figures \texttt{fig4} and \texttt{fig5}, the misattributed $3\times 10^{-8}$ Skytree bound and $29.4\,\mu\mathrm{Gal}$ BEC ceiling (reclassified as hypothetical sensitivity targets), and the hand-coded \texttt{np.piecewise} Yukawa sketch are permanently quarantined.
\end{enumerate}

The complete 40-order-of-magnitude energy-scale hierarchy and Weinberg perturbative loop suppression of \texttt{M\_AUTH\_SIT\_EFT\_v1} at canonical benchmark \texttt{BM\_PHYS\_EFT\_1} remain displayed in Figure~\ref{fig:eft-hierarchy-scales}, while legacy Version~80.2 figures plotting the quarantined \texttt{np.piecewise} Yukawa sketch (\texttt{fig4}) and $3\times 10^{-8}$ Skytree benchmark (\texttt{fig5}) remain permanently retired in Version~81.
% --- END INPUT sections/08_part4_sec51_to_52.tex ---
% --- BEGIN INPUT sections/09_part4_sec53_to_55.tex ---
\section{Non-Perturbative 4D Hypercubic Wilson-Coherence Lattice Gauge Theory, Osterwalder--Schrader Reflection Positivity, and Polyakov Phase Structure}
\label{sec:stage4-lattice}

Guided by the sharp-cutoff diagnosis of Section~\ref{sec:sharp-cutoff-anomaly}, Stage~4 formulated and proved in Lean~4.30 (\texttt{Stage4Lattice/*}, compiled against Mathlib \texttt{c5ea0035} with zero \texttt{sorry}, zero \texttt{admit}, zero custom axioms, and zero computational \texttt{native\_decide}) a non-perturbative 4D hypercubic Euclidean Wilson lattice gauge theory coupled gauge-invariantly to SIT's coherence and temporal-density fields. Adapting the public Apache-2.0 Wilson/Wightman specification of \texttt{OpenAI-YM-20261028} (\texttt{Public\allowbreak{}Gauge\allowbreak{}And\allowbreak{}Wightman.\allowbreak{}lean}, \texttt{Public\allowbreak{}Wilson\allowbreak{}Lattice.\allowbreak{}lean}) to Lean~4.30, the formalization locks ten target obligations inside the frozen comparator specification \texttt{SITStage4\allowbreak{}Spec.\allowbreak{}SITStage4\allowbreak{}Joint\allowbreak{}Target} (\texttt{SITStage4\allowbreak{}Challenge\allowbreak{}Spec.\allowbreak{}lean}) and discharges all ten in \texttt{SITStage4\allowbreak{}Solution.\allowbreak{}lean}.

\subsection{Compact Matrix Gauge Groups, Haar Witnesses, and Coupled 4D Lattice Kinematics}

Let $G \subset U(n+1)$ be a closed compact matrix Lie group (\texttt{GaugeGroup}) equipped with its unique normalized bi-invariant Haar probability measure $dh$ (\texttt{Concrete\allowbreak{}Gauge\allowbreak{}Compact\allowbreak{}Haar}: $\int_G dh(u) = 1$, left/right/inversion invariant). Two concrete witnesses anchor Abelian and non-Abelian sectors:
\begin{itemize}
  \item \textbf{Abelian Witness (\texttt{ConcreteU1Gauge}):} $G = U(1) \subset U(1)$ ($n=0$, non-simple: $\neg\mathrm{u1GaugeGroup.IsSimple}$), parameterized by $e^{i\theta}$ with Haar measure $(2\pi)^{-1}d\theta$.
  \item \textbf{Simple Non-Abelian Witness (\texttt{ConcreteSU2Gauge}):} $G = SU(2) \subset U(2)$ ($n=1$, simple: $\mathrm{su2GaugeGroup.IsSimple}$), isomorphic to the unit three-sphere $S^3 \subset \mathbb{H}$ via $U(q) = q_0 I_2 + i\sum_{a=1}^3 q_a\sigma_a$ ($q_0^2 + q_1^2 + q_2^2 + q_3^2 = 1$) with normalized $O(4)$-invariant $S^3$ surface measure ($\int_0^\pi \frac{2}{\pi}\sin^2\chi\,d\chi = 1$).
\end{itemize}

On a periodic four-dimensional hypercubic lattice $\Lambda_l := (\mathbb{Z}/l\mathbb{Z})^4$ ($l \in \mathbb{N}_+$, spacetime axes $\mu \in \mathrm{Fin}\;4 = \{0,1,2,3\}$ with Euclidean time axis $\tau = 3$), a coupled configuration $(U, \psi, \rho)$ consists of:
\begin{enumerate}[leftmargin=2em]
  \item \textbf{Directed Gauge Links:} $U : \Lambda_l \times \mathrm{Fin}\;4 \to G$, defining elementary oriented plaquette holonomies
  \begin{equation}
  U_\square(x,\mu,\nu) := U(x,\mu)\,U(x+\hat{\mu},\nu)\,U(x+\hat{\nu},\mu)^{-1}\,U(x,\nu)^{-1} \in G
  \qquad (\mu < \nu),
  \label{eq:plaquette_def_stage4}
  \end{equation}
  and non-negative plaquette energy density $E_\square(U; x,\mu,\nu) := 1 - \frac{1}{n+1}\mathrm{Re}\,\mathrm{Tr}\,U_\square(x,\mu,\nu) \in [0, 2]$.
  \item \textbf{Fundamental Coherence Multiplet and Temporal Scalar:} Site fields $\psi : \Lambda_l \to \mathbb{C}^{n+1}$ with squared Hermitian norm $\|\psi(x)\|^2 := \sum_{a=0}^n |\psi_a(x)|^2$ and gauge-singlet temporal field $\rho : \Lambda_l \to \mathbb{R}$.
  \item \textbf{Covariant Link Hopping and Symmetric 4-Corner Plaquette Gate:} Each link $(x,\mu)$ carries gauge-covariant hopping norm squared $\mathcal{K}_{\mathrm{hop}}(U,\psi;x,\mu) := \|\psi(x) - U(x,\mu)\psi(x+\hat{\mu})\|^2 \ge 0$, while each elementary plaquette $(x,\mu,\nu)$ is gated symmetrically across all four corners by the geometric mean (\texttt{plaquetteGate})
  \begin{equation}
  \begin{aligned}
  f_\square(\rho,\psi; x,\mu,\nu) &:= \Bigl[
    f_{\mathrm{gate}}\bigl(\rho(x),\|\psi(x)\|^2\bigr)\,
    f_{\mathrm{gate}}\bigl(\rho(x+\hat{\mu}),\|\psi(x+\hat{\mu})\|^2\bigr)\\
  &\qquad\times
    f_{\mathrm{gate}}\bigl(\rho(x+\hat{\nu}),\|\psi(x+\hat{\nu})\|^2\bigr)\,
    f_{\mathrm{gate}}\bigl(\rho(x+\hat{\mu}+\hat{\nu}),\|\psi(x+\hat{\mu}+\hat{\nu})\|^2\bigr)
  \Bigr]^{1/4}.
  \end{aligned}
  \label{eq:four_corner_plaquette_gate}
  \end{equation}
\end{enumerate}

The frozen Stage-4 unshifted coupled Euclidean action (\textbf{Model~L1}, \texttt{sit\allowbreak{}Wilson\allowbreak{}Coherence\allowbreak{}Action} / \texttt{sitWilsonCoherenceLaw}) is
\begin{equation}
\begin{aligned}
S_{\mathrm{SIT}}^{\mathrm{lat}}(U,\psi,\rho)
&= \beta\sum_{x\in\Lambda_l}\sum_{0\le\mu<\nu\le 3} f_\square(\rho,\psi;x,\mu,\nu)\,E_\square(U;x,\mu,\nu)\\
&\quad + \kappa_c\sum_{x\in\Lambda_l}\sum_{\mu=0}^3 \bigl\|\psi(x) - U(x,\mu)\psi(x+\hat{\mu})\bigr\|^2
  + \sum_{x\in\Lambda_l} V_{\mathrm{site}}\bigl(\rho(x), \|\psi(x)\|^2\bigr),
\end{aligned}
\label{eq:full_sit_lattice_action}
\end{equation}
whose unified AM--GM baseline-shifted Osterwalder--Schrader extension (\textbf{Model~L3}, \texttt{sitAMGMGatedOSAction} / \texttt{sitAMGMGatedOSLaw}, sharing the identical conditional gauge law $P(U\mid\rho,\psi)$ and constant-background action while inducing a quantified radial-scalar gradient stiffness on inhomogeneous $(\rho,\psi)$ fluctuations) is established in Subsection~53.2 below.

\subsection{Targets (1)--(3): Non-Perturbative Finiteness, Exact Gauge Invariance, and Broken-Phase EFT Unification}

\begin{theorem}[Targets (1)--(3): Finiteness, Local Gauge Invariance, and Broken-Phase Unification]
\label{thm:stage4_targets_1_to_3}
Let $\beta \ge 0$, $\kappa_c \ge 0$, let $f_{\mathrm{gate}}, V_{\mathrm{site}} : \mathbb{R}^2 \to \mathbb{R}_{\ge 0}$ be continuous non-negative functions (allowing $f_{\mathrm{gate}}(r,0) = 0$), and let $\mu_{\mathrm{site}}$ be any unitary-invariant probability prior on $\mathbb{C}^{n+1}\times\mathbb{R}$ (explicitly witnessed by \texttt{sitSitePriorWitness}, assigning positive mass both to $\{\psi = 0\}$ and to $\{\|\psi\|^2 = 1\}$). Then:
\begin{enumerate}[leftmargin=2em]
  \item \textbf{Non-Perturbative Finiteness (\texttt{c09\_\allowbreak{}nonperturbative\_\allowbreak{}finiteness\_\allowbreak{}holds}):} The pure Wilson energy is $S_4$-axis-permutation invariant, $0 < Z_W(G,l,\beta) \le M$, and the coupled partition function satisfies
  \begin{equation}
  0 < Z_{\mathrm{SIT}}(G,l,\beta,\kappa_c) := \int \exp\bigl(-S_{\mathrm{SIT}}^{\mathrm{lat}}(U,\psi,\rho)\bigr)\,d\mu_0(U,\psi,\rho) \le 1,
  \end{equation}
  defining a normalized probability measure $d\mu_{\mathrm{SIT}}$ under which every bounded continuous observable $O$ is integrable, even when $f_{\mathrm{gate}}(r,0) = 0$.
  \item \textbf{Exact Local Non-Abelian Gauge Invariance (\texttt{local\_\allowbreak{}nonabelian\_\allowbreak{}gauge\_\allowbreak{}invariance\_\allowbreak{}holds}):} Under local lattice gauge transformations $\Omega : \Lambda_l \to G$ acting via
  \[
  U^\Omega(x,\mu) = \Omega(x)\,U(x,\mu)\,\Omega(x+\hat{\mu})^{-1}, \qquad
  \psi^\Omega(x) = \Omega(x)\psi(x), \qquad
  \rho^\Omega(x) = \rho(x),
  \]
  plaquettes transform covariantly ($U_\square^\Omega(x,\mu,\nu) = \Omega(x)U_\square(x,\mu,\nu)\Omega(x)^{-1}$), every closed lattice Wilson loop $\mathrm{Re}\,\mathrm{Tr}\,W_\gamma(U)$, every term of $S_{\mathrm{SIT}}^{\mathrm{lat}}$, and the full pushforward probability measure $(\mathrm{gaugeTransformCoupled}\;\Omega)_* d\mu_{\mathrm{SIT}} = d\mu_{\mathrm{SIT}}$ are strictly invariant, and $g.\mathrm{lieAlgebra}$ satisfies the matrix Jacobi cocycle identity.
  \item \textbf{Elitzur's Theorem and Broken-Phase \texttt{M\_AUTH\_SIT\_EFT\_v1} Limit (\texttt{broken\_\allowbreak{}phase\_\allowbreak{}eft\_\allowbreak{}unification\_\allowbreak{}holds}):} Without gauge fixing, single-link expectations vanish identically ($\int_{SU(2)} U_{e,00}\,d\mu_W = 0$ for all $l \ge 2$). On gauge-invariant constant-modulus configurations $\|\psi(x)\|^2 = v_0^2 > 0$, $\rho(x) = \rho_0$ with $f_{\mathrm{gate}}(\rho_0, v_0^2\Rcoh^2) = f_0 + c_2\Rcoh^2$ ($f_0 > 0$), $f_\square \equiv f_0$, $f_{\mathrm{gate}}/f_0 = 1 + (c_2/f_0)\Rcoh^2 = f_{\mathrm{EM}}(\Rcoh) \ge 1$ matches Stage-2 \texttt{AuthModelParams}, and for every skew-Hermitian $F \in \mathfrak{u}(n+1)$ ($F^\dagger = -F$),
  \begin{equation}
  \lim_{\varepsilon \to 0,\;\varepsilon\ne 0}
  \frac{1 - \frac{1}{n+1}\mathrm{Re}\,\mathrm{Tr}\exp(\varepsilon F)}{\varepsilon^2}
  = \frac{-\mathrm{Re}\,\mathrm{Tr}(F^2)}{2(n+1)}
  = \frac{\sum_{i,j}|F_{ij}|^2}{2(n+1)} \ge 0.
  \label{eq:lie_exp_continuum_trace_limit}
  \end{equation}
\end{enumerate}
\end{theorem}

\SITConditionalFigure{fig1_su2_phase_diagram}{Exploratory finite-volume 4D $SU(2)$ Wilson-coherence Markov-chain Monte Carlo (MCMC) consistency benchmarks on $4^4$ ($240$ thermalization + $200\times 3$ strided = $840$ sweeps/chain) and $6^4$ ($160$ thermalization + $120\times 2$ strided = $400$ sweeps/chain) hypercubic lattices (\texttt{simulate\_\allowbreak{}su2\_\allowbreak{}wilson\_\allowbreak{}coherence\_\allowbreak{}4d.\allowbreak{}py}, hybrid $S^3$/Kennedy--Pendleton heatbath + overrelaxation with independent Cold $U=I$ and Hot $U\sim\mathrm{Haar}$ starts, median Cold/Hot gap $0.84\sigma$, Madras--Sokal $\tau_{\mathrm{int}} \le 5.69$, treated throughout strictly as empirical finite-lattice convergence diagnostics rather than formal Markov-chain mixing certificates): (a) mean plaquette action $\langle 1 - \frac{1}{2}\mathrm{Re}\,\mathrm{Tr}\,U_\square\rangle$ across $\beta \in [0.5, 4.0]$ comparing unshifted Model~L1 (\texttt{sit\allowbreak{}Wilson\allowbreak{}Coherence\allowbreak{}Action}) against unified AM--GM OS-positive Model~L3 (\texttt{sitAMGMGatedOSAction}) and bounded-uniform ball prior ($R_{\max} = 2.5$, bare Gaussian weight $0$ vs.\ $1/8$, prior-shape sensitivity check), showing exact $\psi = 0$ ($f_{\mathrm{gate}}=0$, \texttt{C09} Haar) disorder $\langle E_\square\rangle = 1.000 \pm 0.001$ and Target~(5)-bounded coherent trajectories relative to pure-YM strong/weak asymptotes; (b) Cold/Hot diagnostic check and smooth non-singular activation of $\langle|\psi(x)|\rangle$ and $\langle E_\square\rangle$ at $\beta = 2.5$ on $4^4$ ($840$ sweeps) from the $\psi = 0$ ($Z > 0$) Haar limit across $v_0 \in [0, 1.25]$ and hopping stiffness $\kappa_c \in \{0.0, 0.4, 1.0\}$, confirming numerical concordance across Model~L1, Model~L3, and bounded-uniform prior $R_{\max} = 2.5$; (c) fixed-weight ($\mathrm{wt}=1/8$) \texttt{truncated\_\allowbreak{}gaussian\_\allowbreak{}exhaustion\_\allowbreak{}sequence} on $4^4$ ($\beta = 2.2, \kappa_c = 0.4, v_0 = 1.0$, $1{,}140$ sweeps/chain, $600$ pooled Cold+Hot measurements per $R_{\max}$) scanning $R_{\max} \in \{1.15, 1.40, 1.80, 2.50, 3.00, \infty\}$ to realize \texttt{truncateSitePrior} exhaustion directly; and (d) $4\mathrm{D}$ periodic torus refinement $L \in \{8, 16, 32, 64\}$ ($a = 1/L$) confirming the centered-plaquette $O(a^2)$ AM--GM gradient law $\bar{f}_p - f_\square = \frac{a^2}{8f}[(\partial_\mu f)^2+(\partial_\nu f)^2] + O(a^4)$ and second-order convergence across Regimes~A, B, and~C of \texttt{amgm\_\allowbreak{}three\_\allowbreak{}eft\_\allowbreak{}scaling\_\allowbreak{}regimes}.}{fig:su2-phase-diagram}{Exploratory finite-volume $4^4$ and $6^4$ $SU(2)$ MCMC benchmarks verifying smooth non-singular activation from zero-matter Haar disorder ($\langle E_\square\rangle = 1$), bounded-uniform prior-shape sensitivity, fixed-weight truncated-Gaussian $R_{\max}$ exhaustion, and $L\in\{8,16,32,64\}$ AM--GM continuum refinement.}

\subsection{Target (4) and Stage-5.1 Unified AM--GM Resolution: Osterwalder--Schrader Reflection Positivity Across Models L1, L2, and L3}

To connect Euclidean lattice expectations to a positive-definite physical Hilbert space and self-adjoint Hamiltonian transfer operator $\hat{T} = e^{-a\hat{H}}$, Target~\texttt{(4)} (\texttt{Reflection\allowbreak{}Positivity\allowbreak{}Split.\allowbreak{}.\allowbreak{}Reflection\allowbreak{}Positivity\allowbreak{}OS}, theorem \texttt{os\_\allowbreak{}reflection\_\allowbreak{}positivity\_\allowbreak{}holds}) proves link-reflection positivity on periodic lattices of arbitrary even temporal extent $2 L_t$ ($L_t \in \mathbb{N}_+$).

Define the fixed-point-free half-integer Euclidean time reflection across the dual hyperplane $\{t = -1/2,\, L_t - 1/2\}$:
\begin{align}
\theta_{\mathrm{site}}(x)_i &:= \begin{cases}
  -x_3 - 1 \pmod{2 L_t}, & i = 3,\\
  x_i, & i \in \{0,1,2\},
\end{cases}
\nonumber\\[0.3em]
(\Theta U)(x,\mu) &:= \begin{cases}
  U\bigl(\theta_{\mathrm{site}}(x + \hat{e}_3),\,3\bigr)^{-1}, & \mu = 3,\\
  U\bigl(\theta_{\mathrm{site}}(x),\,\mu\bigr), & \mu \in \{0,1,2\}.
\end{cases}
\label{eq:os_time_reflection_def}
\end{align}
Simultaneously, decomposing the Wilson action across adjacent 3D spatial slices $U_0, U_1 : (\mathbb{Z}/l\mathbb{Z})^3 \times \mathrm{Fin}\;3 \to G$ connected by temporal links $W : (\mathbb{Z}/l\mathbb{Z})^3 \to G$ with symmetric slab action $S_{\mathrm{slab}}(U_0,U_1,W) = \frac{1}{2}S_{\mathrm{sp}}(U_0) + S_{\mathrm{temp}}(U_0,U_1,W) + \frac{1}{2}S_{\mathrm{sp}}(U_1)$ defines the integrated transfer kernel
\begin{equation}
K_\beta(U_0, U_1) := \int_{G^{V_3}} \exp\bigl(-\beta\,S_{\mathrm{slab}}(U_0, U_1, W)\bigr)\,\prod_{\mathbf{x}} dh\bigl(W(\mathbf{x})\bigr).
\label{eq:slab_transfer_kernel}
\end{equation}
Expanding $\exp\bigl(\frac{\beta}{n+1}\mathrm{Re}\,\mathrm{Tr}\,U_\square\bigr)$ in positive-definite matrix entry polynomials, integrating out hinge temporal links across $t = -1$ and $t = L_t - 1$ via Haar inversion invariance, and applying dominated convergence establishes for all $\beta \ge 0$ and every continuous observable $F$ supported strictly in the positive-time half-lattice $[0, L_t - 1]$:
\begin{align}
K_\beta(U_0, U_1) = K_\beta(U_1, U_0) &> 0,
\qquad
\mathrm{Re}\int \overline{f(U_0)}\,f(U_1)\,K_\beta(U_0,U_1)\,d\mu_{\mathrm{sp}}^2 \ge 0,
\nonumber\\
\mathrm{Re}\int \overline{F(\Theta U)}\,F(U)\,d\mu_W(U) &\ge 0.
\label{eq:os_reflection_positivity_ineq}
\end{align}

Lifting Eq.~\eqref{eq:os_reflection_positivity_ineq} to fluctuating coupled $(U, \psi, \rho_t)$ lattice configurations requires distinguishing three explicit lattice models proved in Lean~4.30:
\begin{enumerate}[leftmargin=2em]
  \item \textbf{Model~L1 (\texttt{sit\allowbreak{}Wilson\allowbreak{}Coherence\allowbreak{}Action} / \texttt{sitWilsonCoherenceLaw}, frozen Stage~4):} Uses unshifted $\beta\sum_p \mathrm{plaquetteGate}\cdot E_p(U)$~\eqref{eq:full_sit_lattice_action}, which vanishes identically at $\psi = 0$ (\texttt{C09} Haar disorder) and is Euclidean time-reflection symmetric (\texttt{sit\allowbreak{}Wilson\allowbreak{}Coherence\allowbreak{}Action\_\allowbreak{}time\allowbreak{}Reflect\_\allowbreak{}eq}, \texttt{sit\allowbreak{}Wilson\allowbreak{}Coherence\allowbreak{}Law\_\allowbreak{}time\allowbreak{}Reflect\_\allowbreak{}eq}), yet encounters a \emph{reduced two-state gate-kernel obstruction} across reflection-straddling temporal plaquettes due to the $-\beta g_0 g_1$ zero-point cross-term: isolating the scalar factor $K_{\mathrm{unshifted}}(\beta; g_0, g_1) = \exp(-\beta g_0 g_1)$ on $\{0, 1\}$ yields the $2\times 2$ matrix $K = \bigl(\begin{smallmatrix}1 & 1\\ 1 & m_{11}\end{smallmatrix}\bigr)$ with hypothesis $m_{11} = e^{-\beta} < 1$ ($\beta > 0$) and negative determinant $\det K = m_{11} - 1 < 0$ (\texttt{unshifted\_\allowbreak{}gate\_\allowbreak{}two\_\allowbreak{}by\_\allowbreak{}two\_\allowbreak{}det\_\allowbreak{}neg}, \texttt{unshifted\_\allowbreak{}gate\_\allowbreak{}exp\_\allowbreak{}neg\_\allowbreak{}beta\_\allowbreak{}not\_\allowbreak{}psd}; note that this theorem is a reduced $2\times 2$ gate-kernel obstruction motivating Model~L3's arithmetic baseline rather than a full coupled Model~L1 transfer-measure no-go theorem after Haar link integration and covariant matter hoppings).
  \item \textbf{Model~L2 (\texttt{sitOSCoupledAction} / \texttt{sitOSCoupledLaw} in \texttt{Stage5\allowbreak{}Extensions.\allowbreak{}Coupled\allowbreak{}OSReflection\allowbreak{}Positivity}, $153$ declarations):} Establishes covariant matter and clock hopping polarizations (\texttt{cov\allowbreak{}Hopping\_\allowbreak{}polarization}, \texttt{clock\allowbreak{}Hopping\_\allowbreak{}polarization}), positive-semidefiniteness of the integrated coupled slab transfer kernel (\texttt{coupled\allowbreak{}Slab\allowbreak{}Transfer\allowbreak{}Kernel\_\allowbreak{}psd}), and coupled $2L_t$ OS reflection positivity (\texttt{coupled\_\allowbreak{}os\_\allowbreak{}reflection\_\allowbreak{}positivity\_\allowbreak{}2\allowbreak{}Lt}, \texttt{coupled\_\allowbreak{}os\_\allowbreak{}reflection\_\allowbreak{}positivity\_\allowbreak{}witness\_\allowbreak{}2\allowbreak{}Lt}) on compact-support site priors (\texttt{Is\allowbreak{}Compactly\allowbreak{}Supported\allowbreak{}Site\allowbreak{}Prior}), \emph{but} retains ungated $\beta\,\mathrm{wilsonEnergy}(U)$ at $\psi = 0$ whenever $\beta > 0$ (so it does not reduce to ultra-local Haar disorder at $\psi = 0$ for $\beta > 0$) while its periodic \texttt{fSlabGate} sum collapses by index shift (\texttt{periodic\_\allowbreak{}temporal\_\allowbreak{}gate\_\allowbreak{}sum\_\allowbreak{}eq}) to an ordinary single-site sum.
  \item \textbf{Model~L3 (\texttt{sitAMGMGatedOSAction} / \texttt{sitAMGMGatedOSLaw} in \texttt{Stage5\allowbreak{}Extensions.\allowbreak{}Unified\allowbreak{}AMGMGated\allowbreak{}OSPositivity}, $98$ declarations --- Unified Stage-5.1 Resolution):} Replaces the constant plaquette zero-point baseline $1$ by the 4-corner arithmetic average $\bar{f}_p := \frac{1}{4}\sum_{c\in p} f_c$ (\texttt{plaquette\allowbreak{}Arithmetic\allowbreak{}Gate}),
  \begin{align}
  S_{\mathrm{L3}}(U,\psi,\rho_t)
  &= \beta\sum_p\left(\bar{f}_p(\rho_t,\psi) - \frac{f_\square(\rho_t,\psi)}{n+1}\mathrm{Re}\,\mathrm{Tr}\,U_p\right)
    + \kappa_c\sum_{x,\mu}\|D_\mu\psi(x)\|^2 \nonumber\\
  &\quad + \kappa_t\sum_{x,\mu}(\Delta_\mu\rho_t(x))^2
    + \sum_x V_{\mathrm{site}}\bigl(\rho_t(x),\|\psi(x)\|^2\bigr),
  \label{eq:model_l3_amgm_action}
  \end{align}
  which by the 4-corner AM--GM inequality $f_\square \le \bar{f}_p$ (\texttt{plaquette\allowbreak{}Gate\_\allowbreak{}le\_\allowbreak{}plaquette\allowbreak{}Arithmetic\allowbreak{}Gate}) differs from Model~L1 (plus $\kappa_t$ clock hopping) exclusively by a gauge-link-independent non-negative site-gradient counterterm $\beta\,\mathrm{amgmCoherenceGap}(\rho_t,\psi) = \beta\sum_p(\bar{f}_p - f_\square) \ge 0$ (\texttt{sit\allowbreak{}AMGMGated\allowbreak{}OSAction\_\allowbreak{}eq\_\allowbreak{}sit\allowbreak{}Wilson\allowbreak{}Coherence\allowbreak{}Action\_\allowbreak{}add\_\allowbreak{}gap}; independent of the gauge links $U$, though contributing an $O(a^2)$ per-plaquette scalar-gradient term to dynamical $(\rho_t,\psi)$ fluctuations analyzed below). Consequently, \textbf{within a single action and probability measure} ($0 < Z_{\mathrm{L3}} \le 1$, \texttt{sit\allowbreak{}AMGMGated\allowbreak{}OSPartition\_\allowbreak{}pos\_\allowbreak{}and\_\allowbreak{}is\allowbreak{}Probability}, \texttt{sit\allowbreak{}AMGMGated\allowbreak{}OSAction\_\allowbreak{}and\_\allowbreak{}Law\_\allowbreak{}gauge\allowbreak{}Invariant}), Model~L3 simultaneously achieves: (i) exact $\psi = 0$ gauge-action vanishing, conditional Haar disorder $P_{\mathrm{L3}}(U\mid\rho_t,\psi=0) = \mu_{\mathrm{Haar}}^{\otimes}$, and unconditional marginal $SU(2)$ Haar disorder under the concentrated zero-matter boundary prior $\mu_{\mathrm{site}} = \delta_0 \otimes \mu_{\mathrm{temp}}$ (\texttt{Measure.dirac 0}; \texttt{sit\allowbreak{}AMGMGated\allowbreak{}OSLaw\_\allowbreak{}zero\allowbreak{}Matter\_\allowbreak{}fst\_\allowbreak{}eq\_\allowbreak{}pi\_\allowbreak{}haar}, \texttt{su2\_\allowbreak{}haar\_\allowbreak{}trace\_\allowbreak{}zero\_\allowbreak{}in\_\allowbreak{}amgm\_\allowbreak{}model}, without claiming that a continuous Gaussian matter prior assigns positive measure to $\{\psi \equiv 0\}$); (ii) identical conditional gauge law $P_{\mathrm{L3}}(U\mid\rho_t,\psi) = \mathrm{fixedSiteGaugeLaw}(\rho_t,\psi)$ (\texttt{amgm\allowbreak{}Conditional\allowbreak{}Gauge\allowbreak{}Law\_\allowbreak{}eq\_\allowbreak{}fixed\allowbreak{}Site\allowbreak{}Gauge\allowbreak{}Law}, preserving Eq.~\eqref{eq:haar_48_bound_sec53} and the Kennedy--Pendleton link heatbath kernel verbatim) and exact action equality with Model~L1 on spatially constant $(\rho_t,\|\psi\|^2)$; (iii) coupled $(U,\psi,\rho_t)$ Osterwalder--Schrader reflection positivity on all even-time $2L_t$ tori (\texttt{amgm\allowbreak{}Coupled\allowbreak{}Slab\allowbreak{}Transfer\allowbreak{}Kernel\_\allowbreak{}psd}, \texttt{amgm\_\allowbreak{}gated\_\allowbreak{}os\_\allowbreak{}reflection\_\allowbreak{}positivity\_\allowbreak{}2\allowbreak{}Lt}, \texttt{amgm\_\allowbreak{}gated\_\allowbreak{}os\_\allowbreak{}reflection\_\allowbreak{}positivity\_\allowbreak{}witness\_\allowbreak{}2\allowbreak{}Lt}); and (iv) extension from compact-ball site priors $B_R$ (\texttt{truncateSitePrior}) to arbitrary unbounded/Gaussian unitary-invariant priors via dominated-convergence exhaustion (\texttt{amgm\_\allowbreak{}os\_\allowbreak{}reflection\_\allowbreak{}positivity\_\allowbreak{}of\_\allowbreak{}exhaustion}, \texttt{amgm\_\allowbreak{}gated\_\allowbreak{}os\_\allowbreak{}reflection\_\allowbreak{}positivity\_\allowbreak{}unbounded\_\allowbreak{}prior\_\allowbreak{}2\allowbreak{}Lt}).
\end{enumerate}

Crucially, gauge-link independence of $\beta\,\mathrm{amgmCoherenceGap}(\rho_t,\psi)$ ($P_{\mathrm{L3}}(U\mid\rho_t,\psi) = P_{\mathrm{L1}}(U\mid\rho_t,\psi)$) and exact equality $S_{\mathrm{L3}} = S_{\mathrm{L1}}$ on spatially constant backgrounds do \emph{not} make $\Delta S_{\mathrm{AMGM}} := \beta\,\mathrm{amgmCoherenceGap}(\rho_t,\psi)$ negligible in the continuum for fluctuating $(\rho_t,\psi)$ fields. Around the center $x_c$ of an elementary $(\mu,\nu)$ plaquette, second derivatives cancel identically between $\bar{f}_p$ and $f_\square$ while first gradients survive at $O(a^2)$, $\bar{f}_p - f_\square = \frac{a^2}{8 f(x_c)}[(\partial_\mu f)^2 + (\partial_\nu f)^2] + O(a^4)$ (explicitly $f_0[\cosh(ka/2)-1] = \frac{f_0 k^2}{8}a^2 + \frac{f_0 k^4}{384}a^4 + O(a^6)$ for $f(x)=f_0 e^{kx}$; \texttt{amgm\_\allowbreak{}exp\_\allowbreak{}profile\_\allowbreak{}cosh\_\allowbreak{}minus\_\allowbreak{}one}, \texttt{amgm\_\allowbreak{}exp\_\allowbreak{}profile\_\allowbreak{}quadratic\_\allowbreak{}limit}, \texttt{amgm\_\allowbreak{}centered\_\allowbreak{}2jet\_\allowbreak{}second\_\allowbreak{}deriv\_\allowbreak{}cancels} in \texttt{Stage5\allowbreak{}Extensions.\allowbreak{}Matter\allowbreak{}And\allowbreak{}ADMDomain\allowbreak{}Refinements}). Summing over all $6$ hypercubic plaquette planes ($\sum_{\mu<\nu}[(\partial_\mu f)^2+(\partial_\nu f)^2] = 3(\partial_\mu f)^2$, \texttt{amgm\_\allowbreak{}hypercubic\_\allowbreak{}six\_\allowbreak{}plane\_\allowbreak{}sum}) with $\sum_x = a^{-4}\int d^4x$ yields $\Delta S_{\mathrm{AMGM}} = \frac{3\beta}{8 a^2}\int d^4x\,\frac{(\partial_\mu f)^2}{f} + O(a^0)$, which leaves the Goldstone/angular stiffness $2\kappa_c^{\mathrm{bare}}v_0^2 = \kappa_c^{\mathrm{EFT}}v_0^2$ unshifted while adding a positive-semidefinite rank-$1$ $(h,\delta\rho_t)$ kinetic shift $K_{\mathrm{eff}}(a)$ (\texttt{amgm\_\allowbreak{}induced\_\allowbreak{}kinetic\_\allowbreak{}matrix\_\allowbreak{}properties}), giving $K_{hh}^{\mathrm{eff}}(a) = \kappa_c^{\mathrm{EFT}} + \Delta K_{hh}(a)$ with $\Delta K_{hh}(a) = \frac{3\beta v_0^2}{a^2 \Lambda_{\mathrm{EM}}^4 f_0} > 0$ in \texttt{M\_AUTH\_SIT\_EFT\_v1}. Machine-checked theorem \texttt{amgm\_\allowbreak{}three\_\allowbreak{}eft\_\allowbreak{}scaling\_\allowbreak{}regimes} (verified numerically on $L\in\{8,16,32,64\}$ 4D periodic tori in Figure~\ref{fig:su2-phase-diagram}d) distinguishes three continuum/EFT regimes: \textbf{(Regime~A)} at fixed physical $(v_0,\Lambda_{\mathrm{EM}},\beta>0)$, $\Delta K_{hh}(a)\propto a^{-2}\to +\infty$ as $a\to 0$ (quadratically stiffening/freezing inhomogeneous radial modes rather than vanishing), whereas on the finite Wilsonian UV-cutoff lattice $a = \Lambda_{\mathrm{UV}}^{-1}$ of \texttt{BM\_PHYS\_EFT\_1} the radial kinetic anisotropy is tiny, $\Delta K_{hh}/\kappa_c^{\mathrm{EFT}} = 3\beta v_0^2\Lambda_{\mathrm{UV}}^2/(\kappa_c^{\mathrm{EFT}}\Lambda_{\mathrm{EM}}^4 f_0) = \frac{3}{200{,}000{,}002}\beta \approx 1.50\times 10^{-8}\beta$; \textbf{(Regime~B)} for the pure quadratic \texttt{C09} gate $f = \tilde{c}\,a^2 R_{\mathrm{phys}}^2$, exact $a^2/a^2$ cancellation produces the finite $a$-independent marginal radial kinetic operator $\frac{3}{2}\beta\tilde{c}\int d^4x\,(\partial_\mu R_{\mathrm{phys}})^2$; and \textbf{(Regime~C)} under Wilsonian dimension-$6$ scaling $\Lambda_{\mathrm{EM}}(a)=\xi/a$ ($f = 1 + a^2\|\psi_{\mathrm{phys}}\|^2/\xi^2$), $\Delta K_{hh}(a) = O(a^2)\to 0$ while $f_{\mathrm{EM}}(a)\to 1$ decouples the dimension-$6$ gauge susceptibility in the strict $a\to 0$ limit.

\subsection{Targets (5)--(6): Haar Coherence-Gating Activation Floor, $\mathbb{Z}_N$ Center Symmetry, and First-Order Polyakov Barrier}

\paragraph{Target (5) --- Haar Coherence-Gating Floor (\texttt{haar\_coherence\_gating\_floor\_holds}).}
In $SU(2)$, left multiplication by $-I_2 \in SU(2)$ flips trace signs ($\mathrm{Tr}(-I_2 u) = -\mathrm{Tr}\,u$), yielding $\int_{SU(2)}\mathrm{Tr}\,u\,dh(u) = 0$. At $\psi = 0$ with $f_{\mathrm{gate}}(\rho,0) = 0$, distinct-time spatial plaquettes ($x_{0,3} \ne x_{1,3}$, $l \ge 2$) share zero links and have identically zero connected covariance. For general $\psi$ under Lipschitz gating $0 \le f_{\mathrm{gate}}(r,s) \le L_f s$ and $|O(U)| \le M$, interpolating $\beta' \in [0,\beta]$ across all $6$ plaquette orientations and $4$ corner sites ($6 \times 4 \times 2 = 48$ via $E_\square \le 2$) proves the explicit global activation bound:
\begin{equation}
\left|\int O(U)\,d\mu_{\beta,\psi}(U) - \int O(U)\,d\mu_{0,\psi}(U)\right|
\le 48\,M\,\beta\,L_f\sum_{x\in\Lambda_l}\|\psi(x)\|^2.
\label{eq:haar_48_bound_sec53}
\end{equation}

\paragraph{Target (6) --- Center Symmetry and Heavy-Coherence First-Order Polyakov Transition (\texttt{compact\_polyakov\_first\_order\_transition\_holds}).}
Around the periodic Euclidean thermal circle of circumference $l$, the Polyakov holonomy loop $P(\mathbf{x}) := \prod_{t=0}^{l-1} U((\mathbf{x},t), 3) \in G$ transforms under center multiplication on a single time slice by $z \in Z(G)$ ($Z(U(1)) = U(1)$, $Z(SU(2)) = \{\pm I_2\}$) as $\mathrm{Tr}\,P(\mathbf{x}) \mapsto z\,\mathrm{Tr}\,P(\mathbf{x})$, proving $\langle\mathrm{Tr}\,P(\mathbf{x})\rangle_{\beta} = 0$ at every finite volume while the normalized mean plaquette action interpolates monotonically from strong-coupling disorder $u_{G,l}(0) = 1$ to weak-coupling order $\lim_{\beta\to\infty} u_{G,l}(\beta) = 0$.

Furthermore, adapting the thermal Polyakov effective potential of \emph{The Hawking--Page transition at weak coupling} (\texttt{hexagon:2611.00175v1}, Eq.~18, \url{https://hexagonmath.org/pdf/2611.00175v1}) coupled linearly to a heavy coherence amplitude $\Rcoh$ of mass $m_{\mathrm{coh}} > 0$ ($c > 0$),
\begin{align}
V_{\mathrm{joint}}(a_T, u_1, \Rcoh)
&= (1 - a_T)u_1 + b_0 u_1^2 + c u_1^3 + \frac{1}{2}m_{\mathrm{coh}}^2\Rcoh^2 - \kappa u_1\Rcoh \nonumber\\
&= V_{\mathrm{eff}}(a_T, u_1) + \frac{1}{2}m_{\mathrm{coh}}^2\left(\Rcoh - \frac{\kappa}{m_{\mathrm{coh}}^2}u_1\right)^2,
\label{eq:polyakov_square_completion}
\end{align}
where $V_{\mathrm{eff}}(a_T, u_1) = (1 - a_T)u_1 + b_{\mathrm{eff}} u_1^2 + c u_1^3$ with $b_{\mathrm{eff}} := b_0 - \frac{\kappa^2}{2 m_{\mathrm{coh}}^2}$, whenever coherence-holonomy coupling satisfies $\kappa^2 > 2 b_0 m_{\mathrm{coh}}^2$ we obtain $b_{\mathrm{eff}} < 0$! Consequently, before the Hagedorn point $a_T = 1$, at $a_{\mathrm{crit}} := 1 - b_{\mathrm{eff}}^2/(4c) < 1$, the potential factors identically as
\begin{equation}
V_{\mathrm{eff}}(a_{\mathrm{crit}}, u_1) = c\,u_1\,(u_1 - u_*)^2,
\qquad
u_* := -\frac{b_{\mathrm{eff}}}{2c} > 0,
\qquad
V_{\mathrm{eff}}\!\left(a_{\mathrm{crit}},\,\frac{u_*}{3}\right) = \frac{4}{27}\,c\,u_*^3 > 0,
\label{eq:polyakov_barrier_exact_sec53}
\end{equation}
proving degenerate two-phase coexistence ($u_1 = 0$ confined vs.\ $u_1 = u_* > 0$ deconfined) separated by an exact positive cubic barrier (e.g.\ at $b_0 = 0.5, c = 1, m_{\mathrm{coh}} = 1, \kappa = 2$, $b_{\mathrm{eff}} = -1.5$, $u_* = 0.75$, $a_{\mathrm{crit}} = 0.4375$, and $\Delta V_{\mathrm{eff}} = 1/16 = 0.0625$). In numerical benchmarking (Figure~\ref{fig:polyakov-hysteresis}), we strictly separate microscopic finite-volume fixed-$N_t = 4$ 4D $SU(2)$ thermal lattice scans on $4^3\times 4$ and $6^3\times 4$ at $\kappa = 1.58$ from an independent 2-variable zero-mode effective-potential sampler cross-check of \texttt{heavy\allowbreak{}Coherence\allowbreak{}Polyakov\allowbreak{}Joint}~\eqref{eq:polyakov_square_completion} at $\kappa = 2.0$ (\texttt{is\_\allowbreak{}microscopic\_\allowbreak{}4d\_\allowbreak{}lattice\_\allowbreak{}extraction = false}, $\Delta V_{\mathrm{MCMC}} = 0.0610 \pm 0.0021$ vs.\ $0.0625$, $0.71\sigma$); the quantitative microscopic-to-effective parameter map $(\beta, \kappa_{\mathrm{micro}}, L_s, N_t) \mapsto (a_T, b_0, c, \kappa_{\mathrm{eff}}, V_{\mathrm{eff}})$ and multicanonical $L_s \to \infty$ 4D lattice free-energy extraction remain analytical/open.

\SITConditionalFigure{fig2_polyakov_susceptibility_hysteresis}{Exploratory fixed-$N_t = 4$ 4D $SU(2)$ thermal lattice scans and separate 2-variable zero-mode effective-potential sampler cross-check: (a) Polyakov loop modulus $\langle|L|\rangle$ and intra-branch susceptibility $\chi_L$ across $\beta \in [1.45, 2.60]$ on $4^3\times 4$ ($V_3 = 64$) and $6^3\times 4$ ($V_3 = 216$) using independent Cold ($U=I$) and Hot ($U\sim\mathrm{Haar}$) chains from scratch at every $\beta$ ($300$ thermalization + $200\times 2$ strided = $700$ sweeps/chain, zero mid-scan initialization switching), showing $\chi_L$ at $\kappa = 1.58, \beta = 1.92$ growing $2.12\times$ with spatial volume from $0.959 \pm 0.063$ ($4^3\times 4$) to $2.029 \pm 0.218$ ($6^3\times 4$) relative to the smooth $\kappa = 0$ crossover; (b) independent 2-variable $(u_1, \Rcoh)$ zero-mode effective-potential sampler cross-check of \texttt{heavy\allowbreak{}Coherence\allowbreak{}Polyakov\allowbreak{}Joint} at $(b_0, c, m_{\mathrm{coh}}, \kappa) = (0.5, 1, 1, 2)$, $V_{\mathrm{eff}} = 56$ (\texttt{is\_\allowbreak{}microscopic\_\allowbreak{}4d\_\allowbreak{}lattice\_\allowbreak{}extraction = false}, \emph{not} extracted from the 4D lattice), reconstructing $\Delta V_{\mathrm{MCMC}} = 0.0610 \pm 0.0021$ in $0.71\sigma$ agreement with exact $\Delta V_{\mathrm{eff}} = \frac{4}{27} c u_*^3 = \frac{1}{16} = 0.0625$ (microscopic-to-effective parameter map open).}{fig:polyakov-hysteresis}{Separates fixed-$N_t=4$ exploratory $4^3\times 4$ and $6^3\times 4$ 4D $SU(2)$ susceptibility growth ($2.12\times$ at $\kappa=1.58$) from the 2-variable zero-mode sampler cross-check ($\Delta V_{\mathrm{MCMC}} = 0.0610 \pm 0.0021$ vs.\ $0.0625$, $0.71\sigma$).}

\SITConditionalFigure{fig3_timeslice_correlator_c09}{Exploratory Euclidean time-slice connected plaquette correlator benchmarks $C(t) = \langle P(0)P(t)\rangle - \langle P\rangle^2$ on anisotropic $4^3\times 8$ $SU(2)$ lattices ($6$ independent replicas: $3$ Cold + $3$ Hot, $240$ thermalization + $200\times 3$ strided = $840$ sweeps/replica, $N_{\mathrm{meas}} = 1{,}200$ per regime, Madras--Sokal $\tau_{\mathrm{int}} \le 2.26$): (a) normalized correlator $C(t)/C(0)$ demonstrating exact ultra-local vanishing $C(t \ge 1) = 0$ within $\tau_{\mathrm{int}}$-validated statistical error on the $\psi = 0$ ($f_{\mathrm{gate}}=0$, \texttt{C09}, Models~L1 \& L3 identical) branch alongside coherent $v_0 = 1, \kappa_c = 0.25$ decay for Model~L1 ($\beta = 2.2, 3.0$) and AM--GM OS-positive Model~L3 ($\beta = 2.2$); (b) semilog slice correlator confirming equal-time Haar variance $C(0) = (1.249 \pm 0.020)\times 10^{-3}$ vs.\ exact $1/768 = 1.302\times 10^{-3}$, statistically positive coherent inter-slice correlation $C(1) > 0$ at $t/a = 1$, and coarse finite-volume exploratory two-slice log-ratio estimates $\ln(C(0)/C(1)) = 1.90 \pm 0.07$ (Model~L1, $\beta = 2.2$), $1.99 \pm 0.11$ (Model~L3, $\beta = 2.2$; $\ln(C(1)/C(2)) = 1.13 \pm 0.31$), and $3.20 \pm 0.16$ (Model~L1, $\beta = 3.0$)---an empirical consistency check, \emph{not} a standalone proof of OS reflection positivity (which is proved analytically for all positive-time functionals in Model~L3).}{fig:timeslice-correlator-c09}{Exploratory $4^3\times 8$ $SU(2)$ correlator consistency check confirming $\psi=0$ ultra-local $t/a \ge 1$ decorrelation and coherent inter-slice decay across Model~L1 and OS-positive Model~L3.}

\subsection{Stage-6 Classical $C^4$ Continuum Limit, Periodic Fourier Spectrum, Nonnegative Coercivity, and Exact Stage-5 Embedding}
\label{subsec:stage6-classical-continuum-embedding}

Advancing Stage~5.2's local 2-jet Taylor analysis, Stage~6 (\texttt{Stage6FunctionalResponse.\allowbreak{}ExactPlaquetteHessianAndFourier} [$6$ theorems], \texttt{SmoothTorusAMGMFunctionalConvergence} [$6$ theorems], \texttt{ContinuumRegularityAndLatticeEmbedding} [$7$ theorems]) together with the November~2026 lattice addenda (\texttt{Stage6Addendum.\allowbreak{}PeriodicLatticeFourierEigenvalue} [$3$ theorems] and \texttt{Stage6Addendum.\allowbreak{}NonnegativeFieldCoercivity} [$10$ theorems], $32$ kernel-compiled Lean~4.30 lattice theorems in total under strict \texttt{[propext, Classical.\allowbreak{}choice, Quot.\allowbreak{}sound]}) establishes the complete \emph{classical} continuum limit, periodic Fourier spectrum, and sharp nonnegative-field coercivity bound for the four-corner AM--GM lattice energy directly on $\mathrm{Site}\,l = (\mathrm{Fin}\,4 \to \mathrm{ZMod}\,l)$ (also presented independently of SIT physical interpretations in the standalone mathematical draft \emph{A Classical Continuum Limit for a Four-Corner Lattice Energy}):
\begin{enumerate}[leftmargin=2em]
  \item \textbf{Exact Finite-Lattice Corner Hessian, Phase Invariance, and Periodic Fourier Eigenvalue Equation:} For any smooth positive two-parameter four-corner family $c_i(\varepsilon, t)$ with $c_i(0,0) = f_0 > 0$, all mixed second derivatives $j_i = \partial_\varepsilon\partial_t c_i(0,0)$ cancel identically between $A_p$ and $G_p$ at finite lattice spacing, leaving $H_{\mathrm{corner}}(f_0) = \frac{1}{4f_0}I_4 - \frac{1}{16f_0}\mathbf{1}\mathbf{1}^\top$ (\texttt{mixed\_\allowbreak{}second\_\allowbreak{}deriv\_\allowbreak{}amgm4\_\allowbreak{}eq\_\allowbreak{}pairwise\_\allowbreak{}diff}, \texttt{cornerHessianMat\_\allowbreak{}properties}), strict invariance under local unitary phase rotations with zero phase variations (and zero energy on a uniform gate background, \texttt{amgmCoherenceGap\_\allowbreak{}global\_\allowbreak{}second\_\allowbreak{}variation}), and the exact $L^2(\mathrm{Site}\,l)$ stencil representation $B_{\beta,f_0}(g,h) = \sum_x g(x)(H_{\mathrm{lat}} h)(x)$ (\texttt{amgmLatticeHessianBilin\_\allowbreak{}eq\_\allowbreak{}sum\_\allowbreak{}mul\_\allowbreak{}op}). Every periodic lattice Fourier mode $u(x) = \cos(\varphi_0 + \sum_\mu \frac{2\pi m_\mu}{l} x_\mu)$ satisfies the exact pointwise and bilinear eigenvalue equations $(H_{\mathrm{lat}} u)(x) = a^2\lambda_Q(2\pi m; a)\,u(x)$ ($a = 1/l$; every nonzero mode is an eigenvector, \texttt{periodicLatticeFourierMode\_\allowbreak{}is\_\allowbreak{}hessian\_\allowbreak{}eigenvector}, \texttt{periodicLatticeFourierMode\_\allowbreak{}bilinear\_\allowbreak{}eigenvalue}), obeying $0 \le \frac{3\beta|k|^2}{4f_0} - \lambda_Q(k;a) \le \frac{\beta a^2|k|^4}{16f_0}$ (\texttt{sin\_\allowbreak{}sq\_\allowbreak{}half\_\allowbreak{}sandwich}, \texttt{amgm\_\allowbreak{}fourier\_\allowbreak{}symbol\_\allowbreak{}error\_\allowbreak{}bound\_\allowbreak{}and\_\allowbreak{}limit}).
  \item \textbf{Conventional $C^4$ Regularity and $O(1/l^2)$ Classical Continuum Convergence:} Removing all custom derivative-witness record assumptions, \texttt{smoothPeriodicPositiveTorusField\_\allowbreak{}of\_\allowbreak{}contDiff4} and \texttt{amgm\_\allowbreak{}contDiff4\_\allowbreak{}periodic\_\allowbreak{}torus\_\allowbreak{}convergence} prove from standard Mathlib \texttt{ContDiff} $\mathbb{R}\;4$, strict positivity, and $1$-periodicity on $[0,1]^4$ that for every $\beta \ge 0$ and $n \ge 1$, the actual sampled four-corner functional $Q_n(f)$ satisfies
  \begin{equation}
  \left|Q_n(f) - \frac{3\beta}{8}\int_{[0,1]^4}\frac{|\nabla f(u)|^2}{f(u)}\,du\right| \le \frac{\beta\,(B_4/4 + C_2)}{n^2},
  \label{eq:monograph_stage6_torus_conv}
  \end{equation}
  where the continuum integral is nonnegative (strictly positive when $\beta > 0$ and $f$ is nonconstant), confirmed numerically at order $1.990123$ on a 4D all-6-plane profile ($n=8..32$) and $1.999907$ on a 1D cosine profile ($n=32..256$).
  \item \textbf{Sharp $\frac{3\beta}{4}$ Zero-Division Nonnegative-Field Coercivity Allowing Zero Corners (\texttt{Stage6Addendum.\allowbreak{}NonnegativeFieldCoercivity}, $10$ theorems):} For arbitrary nonnegative corners $c_i \ge 0$ (including zero values, with zero division by $f$) and $u_i := \sqrt{c_i}$, averaging the exact pairing identity $A_4 - G_4 = \frac{1}{4}[(u_i-u_j)^2 + (u_k-u_l)^2] + \frac{1}{2}(\sqrt{u_i u_j}-\sqrt{u_k u_l})^2$ over the two boundary-edge pairings (omitting the diagonal pairing) proves the stronger local bound $A_4(c) - G_4(c) \ge \frac{1}{8}[(u_0-u_1)^2+(u_0-u_2)^2+(u_3-u_1)^2+(u_3-u_2)^2]$ (\texttt{amgm4\_\allowbreak{}ge\_\allowbreak{}one\_\allowbreak{}eighth\_\allowbreak{}boundary\_\allowbreak{}edges}) and the sharp global 4D periodic lattice coercivity bound $a^2\beta\,\mathrm{amgmCoherenceGap} \ge \frac{3\beta}{4}\,a^2\sum_{x,\mu}(\sqrt{F(x+e_\mu)}-\sqrt{F(x)})^2$ (\texttt{amgm\_\allowbreak{}nonnegative\_\allowbreak{}field\_\allowbreak{}lattice\_\allowbreak{}coercivity\_\allowbreak{}three\_\allowbreak{}fourths}, \texttt{stage5\_\allowbreak{}amgmCoherenceGap\_\allowbreak{}nonnegative\_\allowbreak{}coercivity\_\allowbreak{}three\_\allowbreak{}fourths}), strictly improving the valid weaker $\frac{\beta}{2}$ global bound obtained by dropping diagonals from the sharp $\frac{1}{12}$ six-pair corner inequality. Exact sharpness of $\frac{3\beta}{4}$ is witnessed both by a single-site spike $F(0)=s \ge 0$ on every $l \ge 2$ (nontrivial ratio $6s/(8s) = 3/4$ for $s > 0$, \texttt{amgm\_\allowbreak{}coercivity\_\allowbreak{}three\_\allowbreak{}fourths\_\allowbreak{}sharp\_\allowbreak{}single\_\allowbreak{}site\_\allowbreak{}spike}) and by a strictly positive parity-alternating checkerboard $F(x)\in\{A^2, B^2\}$ with $A,B>0$ on every even $l$ (nontrivial ratio $3/4$ for $A \ne B$, \texttt{amgm\_\allowbreak{}coercivity\_\allowbreak{}three\_\allowbreak{}fourths\_\allowbreak{}sharp\_\allowbreak{}positive\_\allowbreak{}checkerboard}); on smooth profiles, the two nonnegative cross-product remainders in the boundary-edge pairing identity supply the additional continuum contribution $\frac{3\beta}{4}\int |\nabla\sqrt{f}|^2\,du$ in aggregate (elevating the sharp finite-lattice constant $3\beta/4$ to the smooth continuum constant $3\beta/2$, not necessarily equally individually). Paired with a uniform bound on the discrete $L^1$ norm $a^4\sum_x F(x)$, the coercivity bound yields strong $L^2(\mathbb{T}^4)$ compactness of the square roots $\sqrt{F_l}$ with weak-derivative $H^1(\mathbb{T}^4)$ control, and hence strong $L^1(\mathbb{T}^4)$ compactness of the fields $F_l$ for future variational $\Gamma$-convergence to $\frac{3\beta}{2}\int |\nabla\sqrt{f}|^2\,du$.
  \item \textbf{Exact Finite-Lattice Stage-5 Embedding and Mandatory Qualifications:} Resolving $\mathrm{ZMod}\,l$ boundary wraparound and $\mathrm{Fin}\,6 \simeq \mathrm{Directions}$ proves exact finite-lattice equality $l^{-2}\beta\,\mathrm{amgmCoherenceGap} = l^{-2}[S_{\mathrm{L3}}(\beta,\kappa_c,0) - S_{\mathrm{L1}}(\beta,\kappa_c)] = Q_l(f)$ with non-vacuous linear ($f \ge 0$) and affine ($f \ge 1$, $f_{\mathrm{gate}}(r,v) = 1 + c v$) matter-field realizations and end-to-end $O(1/l^2)$ convergence (\texttt{sitePlaquetteCorners\_\allowbreak{}eq\_\allowbreak{}torusSampledCorners}, \texttt{rescaled\_\allowbreak{}stage5\_\allowbreak{}amgm\_\allowbreak{}action\_\allowbreak{}gap\_\allowbreak{}eq\_\allowbreak{}torusFunctional}, \texttt{canonicalEmbeddedMatterField\_\allowbreak{}realizes\_\allowbreak{}linear\_\allowbreak{}and\_\allowbreak{}affine\_\allowbreak{}gates}, \texttt{stage5\_\allowbreak{}amgmCoherenceGap\_\allowbreak{}contDiff4\_\allowbreak{}continuum\_\allowbreak{}limit}). Three qualifications remain strict: setting $\kappa_t = 0$ isolates the AM--GM plaquette gap separately from the temporal-field kinetic sum; the positive affine gate $1 + c|\psi|^2$ does not vanish at $\psi = 0$ and must be distinguished from vanishing gates ($f_{\mathrm{gate}}(\rho,0) = 0$) in the zero-matter Haar-disorder theorem; and Eq.~\eqref{eq:monograph_stage6_torus_conv} is a \emph{classical} positive-background theorem leaving variational $\Gamma$-convergence and the interacting \emph{quantum} continuum measure and OS transfer-operator Hamiltonian strictly open (\texttt{OPEN-3}).
\end{enumerate}

\section{Comparative Synthesis: Dated 2025 SIT Gauge-Holonomy Anchors vs.\ October--November 2026 Constructive Yang--Mills and Chiral Lattice Milestones}
\label{sec:2025-vs-2026-ym}

Between February--December 2025 and October--November 2026, two independent research streams converged on compact-group holonomy, lattice admissibility, and machine-checked Lean~4 verification in four-dimensional gauge physics. Adhering strictly to \texttt{QG49} candor, we record both the documented 2025 SIT archival anchors and the sharp axiomatic distinctions separating modified-action coherence gating from unmodified pure 4D Yang--Mills Wightman/OS existence.

\subsection{Chronological Ledger of Primary 2025 SIT and 2026 External Milestones}

\begin{enumerate}[leftmargin=2em]
  \item \textbf{Dated Feb--Dec 2025 SIT Gauge-Holonomy Corpus:}
  \begin{itemize}
    \item \texttt{SIT-49} (Original Pub.\ Feb~9, 2025; v49 Aug~5, 2025) and \texttt{SIT-60} (Aug~8, 2025): introduced coherence-gated gauge dynamics and soliton/barrier interpretations of gauge spectral gaps.
    \item \texttt{SIT-79} (Dec~18, 2025) and Zenodo anchors \texttt{16922512}, \texttt{17311157} (\emph{Coherence-Gated Yang--Mills Gap}, Oct~10, 2025), \texttt{17307956} (\emph{Informational Holonomy and the Electron $g$-Factor}, Oct~4, 2025), \texttt{17308645} (Oct~10, 2025), plus December~24, 2025 technical companions \texttt{Coherence\_\allowbreak{}Gated\_\allowbreak{}Yang\_\allowbreak{}Mills\_\allowbreak{}Gap 3.\allowbreak{}txt}, \texttt{SIT\_\allowbreak{}Hydrogen\_\allowbreak{}Part\allowbreak{}I/II.\allowbreak{}txt}, and \texttt{SIT\_SPIN 6.txt} (ORCID \texttt{0009-0004-5175-9532}): formulated electromagnetic and non-Abelian interactions via principal-bundle Wilson loop holonomies $U[\gamma] = \mathcal{P}\exp\bigl(\frac{iq}{\hbar}\oint_\gamma A_\mu dx^\mu\bigr)$, phase-only spinor minimal coupling ($g_{\mathrm{tree}} = 2$), and toroidal coherence-gated actions (explicitly footnoted in \texttt{Gap 3} fn.~1 as \emph{not} claiming a proof under unmodified Wightman/OS axioms).
  \end{itemize}
  \item \textbf{October--November 2026 Constructive QFT and Formalization Wave:}
  \begin{itemize}
    \item \textbf{Nathaniel Craig (\texttt{arXiv:2610.09093}, Oct 2026):} Constructed a candidate gauge-invariant integration measure for overlap chiral fermions on 4D admissible lattice gauge fields by comparing chiral projectors across successive lattice spacings to establish exponential locality for L\"uscher reconstruction.
    \item \textbf{\texttt{OpenAI-YM-20261028} \& \texttt{yang-mills-group/YM} \texttt{assessment-1.pdf} (Oct~28 / Nov~10, 2026):} Formalized in Lean~4 the existence of a non-trivial continuum 4D relativistic pure Yang--Mills QFT (\texttt{JointExistence}) for any compact simple matrix group $G \subset U(m)$ via an ordered scaling limit ($a_N \to 0$ at fixed physical box $L_p$ with asymptotic freedom $\log(1/r)g_W(r)^2 \to 1/(2b_0)$ and $r^8\log^2(1/r)$ trace-anomaly nontriviality, followed by subsequential $L_p \to \infty$ compactness), while explicitly leaving infinite-volume uniqueness and the continuum mass gap \textbf{unproved} (\texttt{assessment-1.pdf} Sec.~8).
    \item \textbf{Hexagon \texttt{hexagon:2611.00174v1} \& \texttt{hexagon:2611.00175v1} (Nov 2026):} Established optimal 2D CFT spinless modular bootstrap slope $\Lambda_V = 9.1266$ in \emph{The spinless modular bootstrap at large central charge} (\texttt{hexagon:2611.00174v1}, \url{https://hexagonmath.org/pdf/2611.00174v1}) and resolved the weak-coupling $\mathcal{N}=4$ SYM Hawking--Page transition on $S^3$ via three-loop negative Polyakov quartic $b = -1.48\times 10^{-3}\lambda^2 < 0$ (Eq.~18) while diagnosing sharp spatial cutoff anomalies in \emph{The Hawking--Page transition at weak coupling} (\texttt{hexagon:2611.00175v1}, Secs.~6.1 \&~9, \url{https://hexagonmath.org/pdf/2611.00175v1}).
  \end{itemize}
\end{enumerate}

\begin{table}[htbp]
\centering
\caption{Axiomatic comparison across \texttt{OpenAI-YM-20261028}, 2025 SIT \texttt{C09}, canonical continuum \texttt{M\_AUTH\_SIT\_EFT\_v1}, and the Stage-4/5/5.1 4D Haar Wilson-coherence lattice.}
\label{tab:2025-vs-2026-axiomatic}
\vspace{0.4em}
\footnotesize
\begin{tabular}{|p{2.2cm}|p{2.9cm}|p{2.7cm}|p{2.9cm}|p{3.0cm}|}
\hline
\textbf{Axis} & \textbf{\texttt{OpenAI-YM-20261028}} & \textbf{2025 SIT \texttt{Gap 3} (\texttt{C09})} & \textbf{\texttt{M\_AUTH\_SIT\_EFT\_v1}} & \textbf{Stage-4/5.1 SIT 4D Lattice} \\
\hline
\textbf{Axiom System} & Unmodified continuum OS \& Wightman axioms & Non-OS modified continuum action (fn.~1) & $(-,+,+,+)$ Wilsonian EFT ($E\ll\Lambda_{\mathrm{UV}}$) + ADM & 4D Euclidean Wilson lattice + AM--GM coupled $2L_t$ OS (Model~L3) \\
\hline
\textbf{Field Content} & Pure compact simple $G \subset U(m)$ & $A_\mu + (\rho_t,\psi)$ with $f(\rho_t,0)=0$ & $g_{\mu\nu},\rho_t,\psi,A_\mu,\mathrm{SM}$, $f_{\mathrm{EM}},f_{\mathrm{SM}}\ge 1$ & Links $U_e\in G$ + site $(\psi_x,\rho_x)$, general $f_{\mathrm{gate}}\ge 0$ \\
\hline
\textbf{Zero-Field Behavior} & Pure $1/g^2$ YM; $r^8\log^2(1/r)$ trace anomaly & $\Pi^2/(2f)$ singular at $|\psi|=0$ (Tier-4 quarantined) & Nonsingular floor $f_{\mathrm{EM}}\ge 1$; Maxwell at $v_0=0$ & Nonsingular Haar weight $0 < Z_{\mathrm{L3}} \le 1$ even at $f=0$ \\
\hline
\textbf{Continuum Pure-YM Gap} & \textbf{Unproved} (\texttt{assessment-1.pdf} Sec.~8) & Modified-action only; not Clay pure YM & Proca ($v_0>0$) in EFT; not Clay pure YM & Isolated in Target~(7); Clay pure-YM \texttt{NOT\_SOLVED} \\
\hline
\textbf{Verification} & Lean~4 + Comparator & Informal 2025 draft & Lean~4.30 ($2{,}368 + 218 + 227 + 66$) + $74$ SymPy & Lean~4.30 Targets~(1)--(10) + L2/L3 OS ($153+98$) + exploratory MCMC \\
\hline
\end{tabular}
\end{table}

\subsection{Target (7) Conditional UV Bridge and the Five Open Continuum Gauge Frontiers}

In \texttt{Stage4\allowbreak{}Lattice/Open\allowbreak{}AIUVBridge.\allowbreak{}lean} (theorem \texttt{conditional\_\allowbreak{}openai\_\allowbreak{}uv\_\allowbreak{}bridge\_\allowbreak{}holds}), Stage~4 packages the open constructive continuum UV fixed-point and OS-reconstruction hypotheses into the explicit structure \texttt{Open\allowbreak{}AIUVFixed\allowbreak{}Point\allowbreak{}Assumptions} (existence of an asymptotically free \texttt{WilsonFamily} with subsequential \texttt{Hierarchy} convergence, \texttt{Reconstructed} Wightman axioms, \texttt{Operator\allowbreak{}Product\allowbreak{}Expansion}, and nontrivial \texttt{TraceTest}) and proves the honest implication $\mathrm{OpenAIUVFixedPointAssumptions} \to \mathrm{YangMills.JointExistence}$ alongside unconditional normalized Haar existence for every compact simple $G$ and $U(1) \ne \mathrm{simple}$, $SU(2) = \mathrm{simple}$. Per \texttt{assessment-1.pdf} Section~8, five pure-gauge continuum questions remain open across the entire mathematical-physics community: (1) full uniqueness of the $L_p \to \infty$ infinite-volume limit; (2) explicit continuum probability measure on tempered gauge distributions; (3) uniform continuum OS Hamiltonian mass gap $\Delta > 0$; (4) continuum area law for non-local Wilson loops; and (5) asymptotic single-particle glueball scattering states.

\section{Informational Bundle Holonomy for Matter, Chiral Overlap Fermions, and Pure-Import Lean~4.30 Closure}
\label{sec:stage4-matter-holonomy}

Phase~4 of Stage~4 (\texttt{Stage4Matter/*}, audited in \texttt{Stage4\allowbreak{}Phase4\allowbreak{}Matter\allowbreak{}Audit.\allowbreak{}lean}) and Stage~5.1/5.2 (\texttt{Stage5\allowbreak{}Extensions.\allowbreak{}Matter\allowbreak{}And\allowbreak{}ADMDomain\allowbreak{}Refinements}, $66$ declarations) formalize the matter-sector consequences of SIT's bundle-holonomy postulate---spinor magnetic moment $g_{\mathrm{tree}} = 2$, multi-sheet Hydrogen Bohr quantization, and finite-matrix chiral overlap Standard Model coupling---while explicitly separating three epistemic layers: (a) frozen Stage-4 algebraic/conditional specs, (b) Stage-5.1 constructive operator/semiclassical/Neuberger Lean derivations, and (c) analytical or open continuum PDE and Brillouin-zone results.

\subsection{Target (8) and Stage-5.1 Noncommutative Clifford/Pauli Derivation of $g_{\mathrm{tree}} = 2$ with $O(\Lambda_{\mathrm{UV}}^{-2})$ Suppression}

Formalizing the bundle-connection argument of \texttt{SIT\_SPIN 6.txt} across \texttt{Stage4\allowbreak{}Matter/Spinor\allowbreak{}G2\allowbreak{}Phase\allowbreak{}Only.\allowbreak{}lean} (\texttt{spinor\_\allowbreak{}g2\_\allowbreak{}phase\_\allowbreak{}only\_\allowbreak{}holds}) and \texttt{Stage5\allowbreak{}Extensions.\allowbreak{}Matter\allowbreak{}And\allowbreak{}ADMDomain\allowbreak{}Refinements}:
\begin{enumerate}[leftmargin=2em]
  \item \textbf{Frozen Stage-4 Pauli Spin Identity and Algebraic Zeeman Interface:} For all $\mathbf{v}, \mathbf{w} \in \mathbb{R}^3$ with $\mathbf{v}\cdot\boldsymbol{\sigma} = v_1\sigma_1 + v_2\sigma_2 + v_3\sigma_3$, \texttt{SpinorG2PhaseOnly} proves
  \begin{equation}
  (\mathbf{v}\cdot\boldsymbol{\sigma})(\mathbf{w}\cdot\boldsymbol{\sigma})
  = (\mathbf{v}\cdot\mathbf{w})\,I_2 + i\,(\mathbf{v}\times\mathbf{w})\cdot\boldsymbol{\sigma},
  \label{eq:pauli_product_identity}
  \end{equation}
  and checks the algebraic Zeeman consequence for \texttt{let gTree := 2} ($a_{\mathrm{tree}} = (g_{\mathrm{tree}}-2)/2 = 0$).
  \item \textbf{Stage-5.1 Noncommutative Dirac/Pauli Derivation of $g = 2$:} Removing the Stage-4 algebraic assignment \texttt{let gTree := 2}, \texttt{Matter\allowbreak{}And\allowbreak{}ADMDomain\allowbreak{}Refinements} constructs the $4\times 4$ Dirac Clifford matrices (\texttt{diracBeta}, \texttt{diracAlpha}) together with a noncommutative magnetic kinetic-momentum algebra obeying $[\Pi_i, \Pi_j] = i q\hbar\varepsilon_{ijk}B_k$ (\texttt{noncomm\allowbreak{}Momentum\allowbreak{}Witness\_\allowbreak{}commutator}), proves the Pauli--Lichnerowicz identity $(\boldsymbol{\sigma}\cdot\boldsymbol{\Pi})^2 = \boldsymbol{\Pi}^2 I_2 - q\hbar\,\boldsymbol{\sigma}\cdot\mathbf{B}$ (\texttt{noncomm\_\allowbreak{}pauli\_\allowbreak{}lichnerowicz}), and proves that minimal covariant phase coupling with zero bare dimension-5 Pauli term uniquely forces $g_{\mathrm{tree}} = 2$ (\texttt{noncomm\_\allowbreak{}pauli\_\allowbreak{}operator\_\allowbreak{}forces\_\allowbreak{}g\_\allowbreak{}eq\_\allowbreak{}two}, \texttt{dirac\_\allowbreak{}clifford\_\allowbreak{}pauli\_\allowbreak{}derives\_\allowbreak{}g\_\allowbreak{}eq\_\allowbreak{}two}), while demarcating continuum unbounded Foldy--Wouthuysen perturbation theory on $L^2(\mathbb{R}^3,\mathbb{C}^4)$ as an analytical Tier-2 derivation (\texttt{OPEN-3B} for unbounded operator domains).
  \item \textbf{Universal $\Lambda_{\mathrm{UV}}^{-2}$ Dipole Scaling:} Beyond standard QED Schwinger/multi-loop radiative corrections ($\alpha/(2\pi) + \dots$), any SIT scalar-induced dipole shift is suppressed by two powers of the unified EFT ceiling $\Lambda_{\mathrm{UV}}$:
  \begin{equation}
  |\delta a_\ell| = \left|\frac{C_{\mathrm{dip}}(g_\rho + g_R)}{\Lambda_{\mathrm{UV}}^2}\right|
  \le \frac{C_{\mathrm{dip}}(|g_\rho| + |g_R|)}{\Lambda_{\mathrm{UV}}^2},
  \label{eq:g2_uv_suppression}
  \end{equation}
  which for \texttt{BM\_PHYS\_EFT\_1} ($\Lambda_{\mathrm{UV}} = 1.414\times 10^{18}\,\mathrm{GeV}$, $g_{\rho,R} \sim m_\ell / M_{\mathrm{Pl}}$) lies more than $40$ orders of magnitude below current electron and muon $g-2$ experimental uncertainties ($\sim 10^{-13}$ and $\sim 10^{-10}$), guaranteeing complete electroweak precision compatibility.
\end{enumerate}

\subsection{Target (9) and Stage-5.1 Semiclassical Circular Coulomb Derivation of $E_n = E_1/n^2$}

Formalizing \texttt{SIT\_\allowbreak{}Hydrogen\_\allowbreak{}Part\allowbreak{}I.\allowbreak{}txt} across \texttt{Stage4\allowbreak{}Matter/Hydrogen\allowbreak{}Holonomy\allowbreak{}Tower.\allowbreak{}lean} (\texttt{hydrogen\_\allowbreak{}phase\_\allowbreak{}holonomy\_\allowbreak{}tower\_\allowbreak{}holds}) and \texttt{Stage5\allowbreak{}Extensions.\allowbreak{}Matter\allowbreak{}And\allowbreak{}ADMDomain\allowbreak{}Refinements}:
\begin{enumerate}[leftmargin=2em]
  \item \textbf{Single-Valuedness Holonomy Equivalence and Stage-4 Tower Monotonicity:} Around any closed orbital loop $\gamma$, single-valuedness $\exp(i\Phi_{\mathrm{loop}}) = 1$ holds in $\mathbb{C}$ if and only if $\Phi_{\mathrm{loop}} = \oint_\gamma(\nabla\theta - \frac{q}{\hbar}\mathbf{A})\cdot d\boldsymbol{\ell} = 2\pi n$ for some $n \in \mathbb{Z}$, while frozen Stage~4 proves $0 \le \alpha_{\mathrm{eff}}(\Rcoh) := \alpha_{\mathrm{bare}}/f_{\mathrm{EM}}(\Rcoh) \le \alpha_{\mathrm{bare}}$ ($f_{\mathrm{EM}0} \ge 1$) and strict monotonicity $E_m < E_n < 0$ ($1 \le m < n$) for the algebraically supplied sequence $E_n := E_1/n^2$.
  \item \textbf{Stage-5.1 Constructive Coulomb \texttt{HasDerivAt} + $2\pi n$ Holonomy Derivation vs.\ Continuum 3D PDE Scope:} Differentiating the radial circular Coulomb energy via \texttt{HasDerivAt} (\texttt{has\allowbreak{}Deriv\allowbreak{}At\_\allowbreak{}coulomb\allowbreak{}Circular\allowbreak{}Energy}) enforces centripetal force balance $m_e v^2/r = Z e^2/(4\pi\varepsilon_0 f_{\mathrm{EM}0} r^2)$, which coupled to $2\pi r m_e v/\hbar = 2\pi n$ ($n \in \mathbb{N}_+$) constructively derives both the quantized Bohr radii $r_n = n^2 a_B(f_{\mathrm{EM}0})$ and the dressed bound-state spectrum (\texttt{coulomb\_\allowbreak{}holonomy\_\allowbreak{}derives\_\allowbreak{}bohr\_\allowbreak{}spectrum}):
  \begin{equation}
  E_n = \frac{E_1}{n^2} < 0,
  \qquad
  E_1 = -\frac{Z^2 e^4 m_e}{32\pi^2\varepsilon_0^2 f_{\mathrm{EM}0}^2 \hbar^2},
  \qquad
  1 \le m < n \implies E_m < E_n < 0,
  \label{eq:hydrogen_bohr_tower}
  \end{equation}
  recovering $-13.6057\,\mathrm{eV}$ for $Z=1$, $f_{\mathrm{EM}0} \to 1$ and linking optical-clock spectroscopy to variations in $f_{\mathrm{EM}}(\Rcoh)$, while demarcating the full 3D self-adjoint Schr\"odinger/Dirac PDE eigenvalue problem on $H^2(\mathbb{R}^3) \subset L^2(\mathbb{R}^3)$ ($(\ell,m_\ell)$ degeneracy and Laguerre radial wavefunctions) as analytical Tier-2 physics (\texttt{OPEN-3B}).
\end{enumerate}

\begin{remark}[Quarantine of Legacy Clock/BEC Figure (\texttt{fig5}) and Conditional Metrology Status]
\label{rem:monograph-fig5-quarantine}
Legacy Version~80.2 Figure~5 (\texttt{fig5\_clock\_bec\_sensitivity}) plotted optical/nuclear clock and $^{87}\mathrm{Rb}$ BEC sensitivity lines anchored to the legacy $|\alpha_\Phi| \le 3\times 10^{-8}$ Skytree attribution and the downstream $29.4\,\mu\mathrm{Gal}$ BEC ceiling. Because Takamoto~\emph{et~al.}\ (2020) measured $\alpha_{\mathrm{RS}} = (+1.4 \pm 9.1)\times 10^{-5}$ and linking $f_{\mathrm{EM}}(\Rcoh) = 1 + c_{\mathrm{EM}}\Rcoh^2$ to optical clocks or BEC interferometers requires resolving the visible $U(1)$ photon dictionary (\texttt{OPEN-6B}) and state-preparation source $J_{\mathrm{prep}}$ (\texttt{OPEN-6C}), \texttt{fig5} remains retired in Version~81 and both $3\times 10^{-8}$ and $29.4\,\mu\mathrm{Gal}$ are quarantined strictly as hypothetical sensitivity targets (Section~\ref{sec:stage6-action-diagnostics} and Appendices~\ref{app:BEC-benchmark}--\ref{appendix:ClockConstraints}).
\end{remark}

\subsection{Target (10), Constructive Neuberger Overlap Weight, and Four-Layer Stage 1--6 Verification}

Combining Target~\texttt{(10)} (\texttt{Chiral\allowbreak{}Overlap\allowbreak{}SMInterface.\allowbreak{}lean}: finite-matrix Ginsparg--Wilson chiral involution $D\hat{\gamma}_5 + \gamma_5 D = 0$, $\hat{\gamma}_5^2 = I$, normality $D D^\dagger = D^\dagger D$ for any GW matrix $D$, block-diagonal homomorphism $U(3)\times U(2)\times U(1) \hookrightarrow U(6)$, and conditional bounded-weight measure $d\mu_{\mathrm{SIT+SM}}$) with Targets~\texttt{(1)}--\texttt{(9)}, \texttt{SITStage4\allowbreak{}Solution.\allowbreak{}lean} proves the bundled Stage-4 target:
\begin{gather*}
\texttt{SITStage4Proofs.sit\_stage4\_joint\_target\_holds :}\\
\texttt{SITStage4Spec.SITStage4JointTarget},
\end{gather*}
while Stage~5.1 (\texttt{Matter\allowbreak{}And\allowbreak{}ADMDomain\allowbreak{}Refinements}) rules out the trivial $D = 0$ GW solution by constructing a strictly non-zero Neuberger overlap operator $D_{\mathrm{ov}} = \frac{1}{a}(I - \gamma_5 H_W) \ne 0$ (\texttt{neuberger\allowbreak{}Overlap\_\allowbreak{}satisfies\_\allowbreak{}ginsparg\_\allowbreak{}wilson}, \texttt{neuberger\allowbreak{}Overlap\_\allowbreak{}witness\_\allowbreak{}nontrivial}) and proving the Hermitian eigenvalue determinant floor $0 < m_f^{2N} \le \det(D^\dagger D + m_f^2 I) \le ((2/a)^2 + m_f^2)^N$ supplying $W_{\mathrm{ferm}}$ constructively (\texttt{det\_\allowbreak{}conj\allowbreak{}Transpose\_\allowbreak{}mul\_\allowbreak{}self\_\allowbreak{}add\_\allowbreak{}mass\allowbreak{}Sq\_\allowbreak{}ge}, \texttt{massive\allowbreak{}Fermion\allowbreak{}Det\allowbreak{}Weight\_\allowbreak{}constructs\_\allowbreak{}sm\_\allowbreak{}interface}), leaving Brillouin-zone single-pole doubler dispersion ($0 < m_0 < 2r$), admissible-field exponential locality, and massless chiral determinant phases (\texttt{OPEN-5A}) as analytical/open outside Lean.

\paragraph{Eradication of Transitive \texttt{native\_decide}, Stage-5.1/5.2 Counts, and Four-Layer Verification (\texttt{Finding (f)} \& \texttt{Finding 7}).}
To ensure uncompromising proof hygiene without modifying any of the $580$ frozen Stage~1--3 \texttt{.lean} files, Stage~4 Phase~0 factored all declarative data types out of the $10$ legacy Stage-1 files that contained $42$ computational \texttt{native\_decide} calls into $8$ tactic-free \texttt{Stage4Data/*} modules and re-proved all $218$ certificate, normal-order, sparse-phase, deficit-4, and 5-case \texttt{E1C4A} split-projector declarations in $34$ pure \texttt{Stage4Pure/*} modules (\texttt{Stage4PureAudit.lean}). Consequently:
\begin{itemize}
  \item \texttt{SITStage4\allowbreak{}Pure\allowbreak{}Import\allowbreak{}Clean\allowbreak{}Audit.\allowbreak{}lean} ($93$ Stage-4 modules) and \texttt{SITStage5\allowbreak{}Pure\allowbreak{}Import\allowbreak{}Clean\allowbreak{}Audit.\allowbreak{}lean} (\texttt{005-Compile}) verify across all \textbf{$98$ Stage-5.1/5.2 pure-import transitive local \texttt{.lean} modules} ($97$ strict dependencies + audit root, incorporating $544$ Stage-5 extension declarations across \texttt{Coupled\allowbreak{}OSReflection\allowbreak{}Positivity} [$153$], \texttt{Off\allowbreak{}Diagonal\allowbreak{}Curved\allowbreak{}ADMClosure} [$227$], \texttt{Unified\allowbreak{}AMGMGated\allowbreak{}OSPositivity} [$98$], and \texttt{Matter\allowbreak{}And\allowbreak{}ADMDomain\allowbreak{}Refinements} [$66$], packaged in \texttt{standalone-reproducibility-repo/}): \texttt{0} \texttt{sorry}, \texttt{0} \texttt{admit}, \texttt{0} custom axioms, \texttt{0} computational \texttt{native\_decide}, and \texttt{0} bare \texttt{import Mathlib}, with all $64/64$ audited \texttt{\#print axioms} checks depending strictly on \texttt{[propext, Classical.\allowbreak{}choice, Quot.\allowbreak{}sound]}.
  \item \texttt{SITStage4\allowbreak{}Master\allowbreak{}Audit.\allowbreak{}lean} and \texttt{SITStage5\allowbreak{}Master\allowbreak{}Audit.\allowbreak{}lean} (\texttt{006-Compile}) unite Stage~1 (\texttt{SITF49\allowbreak{}Master\allowbreak{}Theorem\allowbreak{}Stack}), Stage~2 + Stage~3 (\texttt{SITStage3MasterAudit}), Stage~4 (\texttt{SITStage4MasterAudit}), and all $6$ compiled Stage-5 modules ($4$ extension modules + $2$ audit roots) across \textbf{$640$ unified local modules} (within the $660$-file frozen \texttt{.\allowbreak{}sitproof-stage5-compiler-20261116} package) across $23/23$ master \texttt{\#print axioms} checks under \texttt{[propext, Classical.\allowbreak{}choice, Quot.\allowbreak{}sound]} with \texttt{0} warnings and \texttt{0} errors.
  \item \texttt{SITStage6Audit.\allowbreak{}lean} (\texttt{007-Compile}, $6$ frozen Stage-6 \texttt{.lean} files in \texttt{Stage6FunctionalResponse/*}, auditing $30/30$ theorems) together with the $3$ publication-branch addendum modules (\texttt{Stage6Addendum.\allowbreak{}PeriodicLatticeFourierEigenvalue} [$3$], \texttt{NonnegativeFieldCoercivity} [$10$], and \texttt{ActionResponsePropagatorIdentity} [$5$]) brings the complete audited Lean~4.30 corpus across Stages~1--6 plus addenda to \textbf{$669$ total \texttt{.lean} files} ($48/48$ Stage-6/Addendum \texttt{\#print axioms} checks strictly within \texttt{[propext, Classical.\allowbreak{}choice, Quot.\allowbreak{}sound]}, \texttt{0} warnings, \texttt{0} errors).
  \item Across the reproducibility suite, verification is explicitly separated into \textbf{four distinct layers}: (V1) static source, topological \texttt{lakefile.toml} parity, and forbidden-token scanning, (V2) dual-engine SymPy/NumPy/$\mathbb{Q}(\sqrt{\det q})$ execution (including all three Stage-6 Python verifiers), (V3) Lean~4.30.0 kernel type-checking and \texttt{\#print axioms}, which mechanically certifies the exact formal \texttt{Prop} theorem types as written, and (V4) human mathematical-physics specification-to-prose fidelity review verifying that every formal predicate (\texttt{IsSPD3}, Model~L3 \texttt{sitAMGMGatedOSAction}, Stage-6 regularity/embedding/observable boundaries, and quarantined provenance rules) matches the physical claims transparently.
\end{itemize}
% --- END INPUT sections/09_part4_sec53_to_55.tex ---
% --- BEGIN INPUT sections/10_appendices_A_to_M.tex ---
\appendix

\section{Glossary of Key Notation and Symbols}
\label{sec:notation}
\label{appendix:A}
\label{appendix:B}

This appendix summarizes the canonical notation for the primary fields, coupling parameters, cutoff scales, and operational observables used throughout Version~81 (\texttt{M\_AUTH\_SIT\_EFT\_v1}, explicitly separated from the July~2026 Version~80 conformal benchmark of Paper~002 in Table~\ref{tab:july-v80-vs-nov-v81}). Spacetime signature follows the mostly-plus relativistic convention $(-,+,+,+)$ with Minkowski metric $\eta_{\mu\nu} = \mathrm{diag}(-1,+1,+1,+1)$ and natural units $\hbar = c = 1$ ($[M] = \mathrm{GeV}$) unless SI factors are explicitly restored.

\subsection*{Core SIT Fields, Canonical Couplings, and Cutoff Scales}

\begin{description}[leftmargin=2.4cm, style=nextline]
  \item[$\psi(x)$] Primitive complex coherence scalar order-parameter field, $\psi(x) = |\psi(x)| e^{i\theta(x)}$ (mass dimension $[M]^1$).
  \item[$\Rcoh(x)$] Dimensionless gauge-invariant coherence ratio, $\Rcoh(x) := |\psi(x)|/\psi_0 \ge 0$, with vacuum value $R_0 = v_0/\psi_0$.
  \item[$\rho_t(x)$] Real gauge-singlet temporal-density scalar field governing operational local clock rates (physical dimension $[T^{-1}]$, natural dimension $[M]^1$), with vacuum value $\rho_0 > 0$.
  \item[$g_{\mu\nu},\,\kappa_G,\,M_{\mathrm{Pl}}$] Lorentzian spacetime metric (dimensionless), gravitational coupling $\kappa_G := 8\pi G = M_{\mathrm{Pl}}^{-2} > 0$, and canonical reduced Planck mass $M_{\mathrm{Pl}} = 2.435323459\times 10^{18}\,\mathrm{GeV}$.
  \item[$A_\mu,\,F_{\mu\nu},\,\hat{q}$] Abelian $U(1)$ gauge connection, curvature $F_{\mu\nu} = \partial_\mu A_\nu - \partial_\nu A_\mu$, and coherence $U(1)$ charge coupling $\hat{q} \ge 0$ in $D_\mu\psi = (\nabla_\mu - i\hat{q}A_\mu)\psi$.
  \item[$f_{\mathrm{EM}}(\Rcoh),\,f_{\mathrm{SM}}(\Rcoh)$] Vacuum-positive constitutive dressings $f_{\mathrm{EM}}(\Rcoh) = 1 + c_{\mathrm{EM}}\Rcoh^2 \ge 1$ and $f_{\mathrm{SM}}(\Rcoh) = 1 + \Rcoh^2/(2 M_{\mathrm{sym}}^2) \ge 1$.
  \item[$\Lambda_{\mathrm{EM}},\,\Lambda_{\mathrm{sym}},\,\Lambda_{\mathrm{UV}}$] Dimension-6 EFT scales $\Lambda_{\mathrm{EM}} = \psi_0/\sqrt{c_{\mathrm{EM}}}$, $\Lambda_{\mathrm{sym}} = \sqrt{2}\,\psi_0 M_{\mathrm{sym}}$, and unified ultraviolet ceiling $\Lambda_{\mathrm{UV}} = \min(M_{\mathrm{Pl}}, \Lambda_{\mathrm{EM}}, \Lambda_{\mathrm{sym}})$.
  \item[$\zeta$] Phenomenological scale parameter entering the separately postulated exploratory energy-density ansatz $\varepsilon_{\mathrm{SIT}}^{\mathrm{phen}} = \zeta\,|\psi|\,\rho_t^2$ (not derived from unchanged \texttt{M\_AUTH\_SIT\_EFT\_v1}, Section~\ref{sec:informational_energy_equivalence}).
  \item[{$\mathcal{C}_{\mathrm{state}}[\rho]$}] Declared nonlinear quantum state functional (replacing the impossible linear operator $\hat{\mathcal{R}}_{\mathrm{coh}}$ of earlier drafts) in the exploratory SIT~3.0 proper-time hypothesis (Section~\ref{sec:SIT3}).
\end{description}

\subsection*{Neuroscientific, Lattice, and Experimental Parameters}

\begin{description}[leftmargin=2.4cm, style=nextline]
  \item[$\phi_{\mathrm{neuron}},\,\phi_{\mathrm{group}},\,\Phi_{\mathrm{token}}$] Instantaneous single-neuron oscillation phase, ensemble mean phase, and local phase-wave differential token $\Phi_{\mathrm{token}} = \Delta(\phi_{\mathrm{neuron}} - \phi_{\mathrm{group}})$ (radians).
  \item[$U(x,\mu),\,U_\square,\,P(\mathbf{x})$] Compact matrix gauge link variable $U(x,\mu) \in G \subset U(n+1)$ on the 4D hypercubic Euclidean lattice $(\mathbb{Z}/l\mathbb{Z})^4$, elementary plaquette holonomy $U_\square$, and thermal Polyakov loop $P(\mathbf{x})$.
  \item[$\alpha_\rho,\,\mathcal{T},\,\alpha_\Phi,\,\chi$] Disambiguated optical-clock response symbols (repairing the Version~79/80.2 reuse of $\alpha_{\mathrm{eff}}$): local field-response coefficient $\alpha_\rho$ ($\Delta\nu/\nu = \alpha_\rho\,\delta\rho_t/\rho_0$), potential-to-field transfer factor $\mathcal{T} := (\rho_0^{-1}\partial\rho_t/\partial(\Phi/c^2))_{\mathrm{lab}}$, gravitational-potential redshift slope $\alpha_\Phi := d\ln\nu / d(\Phi/c^2) = \alpha_\rho\mathcal{T}$ (Tokyo Skytree 2020 measured $\alpha_{\mathrm{RS}} = (+1.4\pm 9.1)\times 10^{-5}$; legacy $3\times 10^{-8}$ and $29.4\,\mu\mathrm{Gal}$ BEC targets are quarantined hypothetical sensitivities), and geometry transfer factor $\chi$ ($\alpha_g = \chi\alpha_\Phi$).
\end{description}

\subsection{Addressing Potential Criticisms and Commitment Levels}

To prevent conflating core laboratory predictions with exploratory extensions, Super Information Theory is organized into distinct levels of empirical and mathematical commitment (refined further into the five-tier formal ledger \texttt{QG49} in Appendix~\ref{app:epistemic-ledger-lean}):
\begin{itemize}
  \item \textbf{Level 0 (Core EFT and Non-Perturbative Lattice Dynamics):} The canonical action \texttt{M\_AUTH\_SIT\_EFT\_v1} ($E \ll \Lambda_{\mathrm{UV}}$), its curved $3+1$D ADM constraint closure, its 4D Haar Wilson-coherence lattice realization, and its laboratory falsifiers in $^{87}\mathrm{Rb}$ Bose--Einstein condensates, optical/nuclear clocks, and atom interferometers. Empirical exclusion of all viable screened windows at Level~0 falsifies the core physical theory.
  \item \textbf{Level 1 (Cosmological and Interdisciplinary Regimes):} Applications of $(\Rcoh, \rho_t)$ gradients to galactic rotation curves, weak-lensing shear residuals, Hubble-tension sampling, and mesoscopic cortical phase-synchronization dynamics. These models are constrained by Level~0 bounds but are modular: revising a cosmological or neural parameter fit does not invalidate Level~0.
  \item \textbf{Level 2 (Conditionally Isolated Arithmetic Spectral Explorations):} Operator-theoretic comparisons with zeta-function spectra, formally decoupled from the physical SIT core under \texttt{ODQ-7} (\texttt{HELD\_PER\_ODQ7}, \texttt{rh\_solved = false}).
\end{itemize}

\subsection{Summary}

Variations of $S_{\mathrm{AUTH}}$ in $(-,+,+,+)$ generate ghost-free, positive-energy scalar, gauge, and metric dynamics that reduce smoothly to General Relativity, Maxwell electrodynamics, quantum mechanics, and Boltzmann kinetic theory in tested limits while yielding crisp laboratory falsification targets.

\section{Roadmap to Formal Proofs and Technical Companion Papers}
\label{app:roadmap}

Alongside this unified Version-80 master monograph, specialized derivations, Monte Carlo codes, and machine-checked Lean~4.30 proof developments are organized across four focused technical companion works:

\subsection{Non-Autonomous Balance Laws, Informational Hysteresis, and Singular Fronts}

Formulates open-system SIT balance laws with time-dependent environmental sources, analyzing path-dependent informational hysteresis ($\Phi_{\mathrm{geom}} \approx \oint C_{\mathrm{rec}}[p]\,dt$) and steep coherence-saturation fronts without violating global conservation of system-plus-bath $J^\mu_{\mathrm{coh}}$.

\subsection{Companion Paper I: Canonical $(-,+,+,+)$ Effective Field Theory, Stage-6 Action-Response Diagnostics, Conditional E\"ot-Wash Non-Exclusion, and Curved $3+1$D ADM Closure}

Presents the complete classical, linearized, and one-loop heat-kernel structure of \texttt{M\_AUTH\_SIT\_EFT\_v1}, Stage-6 multiplicative-vs-conformal and $U(1)$ kinetic-mixing diagnostics (\texttt{ActionIdentityDiagnostics}, $6$ theorems), the exact $(p^2 K + H)^{-1} = V\,\mathrm{diag}(1/(p^2+m_a^2))\,V^\top$ source--detector propagator residue identity, its coordinate-rescaling equivalence $D = \mathrm{diag}(1,\psi_0)$, and its direct canonical instantiation on \texttt{AuthModelParams} proving $0 < m_0^2 \le m_1^2$ and automatic $p^2 + m_a^2 \ne 0$ for $p^2 \ge 0$ (\texttt{Stage6Addendum.\allowbreak{}ActionResponsePropagatorIdentity}, $5$ theorems), the conditional weak-field E\"ot-Wash parallel-slab non-exclusion calculation (\texttt{UnchangedActionWeakFieldObservable}, $5$ theorems) paired with six locked primary experimental provenance records in \texttt{primary\_measurements\_manifest\_v2.json} and quarantine of legacy approximations, infinite-dimensional $\ell^2$ Sobolev Galerkin convergence, and nonlinear curved $3+1$D ADM pointwise formal 3-jet Dirac--Bergmann hypersurface closure modulo explicit spatial divergences $\partial_k B^k$ restricted to the physical Sylvester Riemannian SPD cone \texttt{IsSPD3} ($181$ Stage-3 + $227$ \texttt{Off\allowbreak{}Diagonal\allowbreak{}Curved\allowbreak{}ADMClosure} + \texttt{Matter\allowbreak{}And\allowbreak{}ADMDomain\allowbreak{}Refinements} Lean~4.30 declarations + $74/74$ SymPy checks, demarcating infinite-dimensional Sobolev/Banach phase-space domains with Regge--Teitelboim asymptotic boundary conditions under \texttt{OPEN-2B} and action-to-apparatus gaps under \texttt{OPEN-6}).

\subsection{Companion Paper II \& Standalone Lattice Mathematics Paper: Non-Perturbative 4D Hypercubic Wilson-Coherence Lattice Gauge Theory and Classical $C^4$ Continuum Limit}

Establishes the compact-group 4D Haar Wilson-coherence probability measure $0 < Z \le 1$, exact local non-Abelian gauge invariance, the three-model Osterwalder--Schrader reflection-positivity hierarchy (Models~L1/L2/L3), Stage-6's exact finite-lattice corner Hessian, periodic lattice Fourier-mode eigenvalue equations (\texttt{Stage6Addendum.\allowbreak{}PeriodicLatticeFourierEigenvalue}, $3$ theorems), non-asymptotic 4D Fourier symbol bounds, conventional $C^4$ (\texttt{ContDiff} $\mathbb{R}\;4$) $O(1/l^2)$ convergence to $\frac{3\beta}{8}\int_{[0,1]^4}\frac{|\nabla f|^2}{f}\,du$, exact Stage-5 finite-lattice embedding at $\kappa_t = 0$, and the sharp zero-division nonnegative-field $\frac{3\beta}{4}$ $H^1(\sqrt{F})$ coercivity bound allowing zero corners together with both exact equality witnesses (single-site spike for $l \ge 2$ and strictly positive parity-alternating checkerboard for even $l$) while retaining the valid weaker $\beta/2$ bound and sharp $1/12$ six-pair bound (\texttt{Stage6Addendum.\allowbreak{}NonnegativeFieldCoercivity}, $10$ theorems; $32$ Stage-6/Addendum Lean~4.30 lattice theorems, also presented independently of physical SIT interpretations in the standalone mathematical paper \emph{A Classical Continuum Limit for a Four-Corner Lattice Energy}), alongside heavy-coherence Polyakov coexistence and exploratory finite-volume 4D $SU(2)$ MCMC benchmarks.

\subsection{Companion Paper III: Informational Bundle Holonomy for Matter, Chiral Overlap Fermions, and Pure-Import Lean~4.30 Formalization}

Details noncommutative Clifford/Pauli operator derivation of $g_{\mathrm{tree}} = 2$ with $O(\Lambda_{\mathrm{UV}}^{-2})$ dipole suppression, semiclassical circular Coulomb \texttt{HasDerivAt} plus $2\pi n$ phase-holonomy derivation of the Bohr spectrum $E_n = E_1/n^2$, constructive non-zero finite-matrix Neuberger--Ginsparg--Wilson overlap chirality $D_{\mathrm{ov}} \ne 0$ with $SU(3)_c\times SU(2)_L\times U(1)_Y \hookrightarrow U(6)$ embedding and massive determinant weight $\det(D^\dagger D + m_f^2 I) \ge m_f^{2N}$, and the four-layer verification architecture spanning the $98$-module pure-import closure (\texttt{SITStage5\allowbreak{}Pure\allowbreak{}Import\allowbreak{}Clean\allowbreak{}Audit}), $640$-module unified closure (\texttt{SITStage5MasterAudit}), $6$-module Stage-6 audit (\texttt{SITStage6Audit}, $30/30$ theorems), and $3$-module publication addenda ($18/18$ theorems; $48$ Stage-6/Addendum audited theorems) across $669$ total \texttt{.lean} files.

\section{Chronology of Core Concepts and Bounded Comparison with Causal Fermion Systems}
\label{app:cfs-chronology}

This appendix records the dated public timeline of high-level conceptual framing elements in the SIT precursor corpus (2017--2022) alongside later expositions in Causal Fermion Systems (CFS, 2024--2025). Consistent with Paper~048 (\emph{Causal Fermion Systems and Super Information Theory}, July~28, 2026) and Appendix~\ref{app:chronology}, this chronological comparison documents \emph{narrative-level conceptual parallels only}; it does \emph{not} claim mathematical equivalence, semantic embedding, or independent empirical validation of SIT.

\subsection{Phase 1: Origination of Conceptual Framing in SIT's Precursors (2017--2022)}
The precursor work for SIT explored relational descriptions of information, observation, and time density across neuroscience and physics:

\begin{description}
    \item[2017 (Neural Lace Podcast):] Public expositions introduced two high-level conceptual motifs:
    \begin{itemize}
        \item \textbf{Relational Information:} Framing operational bits around relational \textit{``coincidence patterns''} of events rather than isolated static states.
        \item \textbf{Distributed Internal Observation:} The NAPOT (Neural Array Projection Oscillation Tomography) narrative framing internal observation via distributed internal synchronization without an external observer.
    \end{itemize}
    \item[2022 (Quantum Gradient Time Crystal Dilation - QGTCD):] Precursor physical drafts introduced a scalar time-density concept ($\rho_t$) and explored gradient-driven gravitational analogies.
\end{description}

\subsection{Phase 2: The State of Causal Fermion Systems (Pre-2024)}
In its earlier technical expositions (e.g.\ 2015 overviews), CFS formulated spacetime primarily as the support of a universal measure $\rho$ on abstract operator spaces, while conjecturing that the causal action principle might account for state reduction.

\subsection{Phase 3: Later CFS Expositions and Bounded Comparison Status (2024--2025)}
During 2024--2025, CFS publications developed operator-algebraic accounts of ``effective dynamical collapse'' (2024) and expository descriptions of spacetime as a ``web of correlations'' of a quantum many-body system (2025). While these motifs echo the 2017 relational vocabulary at a high-level narrative plane, Paper~048 establishes that the only controlled mathematical overlap verified to date is a shared regulated Einstein--Dirac effective-equation limit against July~2026 Version~80 Draft~13---which matches an ordinary Einstein--Dirac null control with \emph{zero discriminatory value}---and extending that truncation check to November~2026 \texttt{M\_AUTH\_SIT\_EFT\_v1} remains a separate open calculation (Appendix~\ref{app:chronology}).



\section{Technical Derivations, Proofs, and Mathematical Formalism}
\label{app:tech-derivations}

This appendix records explicit step-by-step variational and post-Newtonian derivations for \texttt{M\_AUTH\_SIT\_EFT\_v1} in signature $(-,+,+,+)$, supporting Sections~\ref{sec:gravitational-sector}--\ref{sec:action-principle}.

\subsection{Full Covariant Variation of the Canonical \texttt{M\_AUTH\_SIT\_EFT\_v1} Action}

Varying $S_{\mathrm{AUTH}}$ (Eq.~\eqref{eq:canonical_action_sec5}) with respect to $\delta g^{\mu\nu}$, $\delta\rho_t$, $\delta\psi^*$, and $\delta A_\nu$ using $\delta\sqrt{-g} = -\frac{1}{2}\sqrt{-g}\,g_{\mu\nu}\delta g^{\mu\nu}$ and $\delta(F_{\alpha\beta}F^{\alpha\beta}) = 2 F_{\mu\lambda}F_\nu^{\ \lambda}\delta g^{\mu\nu} + 4 F^{\mu\nu}\nabla_\mu\delta A_\nu$:
\begin{align}
\delta S_{\mathrm{AUTH}}
&= \int d^4x\,\sqrt{-g}\Biggl\{
  \frac{1}{2\kappa_G}\bigl(G_{\mu\nu} - \kappa_G T_{\mu\nu}^{\mathrm{tot}}\bigr)\delta g^{\mu\nu}
  + \Bigl(\kappa_t\Box_g\rho_t - \frac{\partial V_{\mathrm{tot}}}{\partial\rho_t}\Bigr)\delta\rho_t
  \nonumber\\
&\quad + \left[
    \frac{\kappa_c}{2}D_\mu D^\mu\psi
    - \frac{\partial V_{\mathrm{tot}}}{\partial\psi^*}
    - \frac{c_{\mathrm{EM}}}{4\psi_0^2}\psi F_{\alpha\beta}F^{\alpha\beta}
    + \frac{1}{2\psi_0^2 M_{\mathrm{sym}}^2}\psi\,\mathcal{L}_{\mathrm{matter}}
  \right]\delta\psi^* + \mathrm{c.c.}
  \nonumber\\
&\quad + \Bigl[\nabla_\mu\bigl(f_{\mathrm{EM}}(\Rcoh)F^{\nu\mu}\bigr) - \hat{q}J^\nu_{\mathrm{coh}} - J^\nu_{\mathrm{EM,SM}}\Bigr]\delta A_\nu
\Biggr\},
\label{eq:appD_full_variation}
\end{align}
reproducing Eqs.~\eqref{eq:covariant_einstein_sec5}--\eqref{eq:maxwell_proca_eom_canonical} identically upon setting $\delta S_{\mathrm{AUTH}} = 0$.

\subsection{Algebraic Metric Reduction vs.\ Full Coupled Solution Conditions at Vacuum Equilibrium}

Substituting constant vacuum fields $\rho_t(x) \equiv \rho_0$ and $|\psi(x)| \equiv v_0 = \psi_0 R_0$ with pure-gauge phase $\nabla_\mu\theta - \hat{q}A_\mu = 0$ gives $\partial_\mu\rho_t = 0$, $D_\mu\psi = 0$, $V(\rho_0) = U(v_0) = U_{\mathrm{link}}(\rho_0, v_0) = 0$, $f_{\mathrm{EM}} \equiv f_{\mathrm{EM}}(R_0) \ge 1$, and $f_{\mathrm{SM}} \equiv f_{\mathrm{SM}}(R_0) \ge 1$. Consequently $T_{\mu\nu}^{(\rho_t)} = T_{\mu\nu}^{(\psi)} = T_{\mu\nu}^{(\mathrm{pot})} = 0$, reducing the \emph{metric} equation~\eqref{eq:covariant_einstein_sec5} \emph{algebraically} to $G_{\mu\nu} = 8\pi G_{\mathrm{eff}} T_{\mu\nu}^{\mathrm{SM}} + 8\pi G f_{\mathrm{EM}}(R_0) T_{\mu\nu}^{\mathrm{EM},0}$ with $G_{\mathrm{eff}} = G f_{\mathrm{SM}}(R_0)$. Crucially, this constant-field configuration is an \emph{exact solution of the full coupled system}~\eqref{eq:appD_full_variation} only when the scalar, gauge, and matter equations are simultaneously satisfied:
\begin{enumerate}[leftmargin=2em]
  \item \textbf{Nonzero Dust Scalar Residual at $v_0 > 0$:} In Eq.~\eqref{eq:appD_full_variation} (or doubled radial form~\eqref{eq:Rcoh-eom}), setting $\rho_t \equiv \rho_0$, $|\psi| \equiv v_0 > 0$, $F_{\mu\nu} = 0$, and neutral dust $\mathcal{L}_{\mathrm{matter}} = -\rho_m < 0$ leaves a nonzero scalar-equation residual $\mathcal{E}_{|\psi|} = -\frac{v_0\rho_m}{\psi_0^2 M_{\mathrm{sym}}^2} \ne 0$ (or $-\frac{v_0\rho_m}{2\psi_0^2 M_{\mathrm{sym}}^2}$ in $\delta S/\delta\psi^*$). Thus arbitrary dust-sourced Einstein solutions cannot be embedded as exact coupled solutions merely by freezing $(\rho_t, |\psi|) = (\rho_0, v_0)$ at $v_0 > 0$.
  \item \textbf{Covariant-Constant Gauge Restriction ($\hat{q}v_0 \ne 0$):} A covariantly constant charged vacuum $D_\mu\psi = 0$ with $|\psi| = v_0 > 0$ obeys $[D_\mu, D_\nu]\psi = -i\hat{q}F_{\mu\nu}\psi = 0$, which forces $F_{\mu\nu} = 0$ whenever $\hat{q}v_0 \ne 0$ (and $F_{\alpha\beta}F^{\alpha\beta} \ne 0$ also sources $|\psi|$ whenever $c_{\mathrm{EM}} v_0 \ne 0$).
  \item \textbf{Exact Regimes vs.\ Controlled Decoupling:} Source-free vacuum ($\mathcal{L}_{\mathrm{matter}} = 0, F_{\mu\nu} = 0$) is an exact solution at $(\rho_0, v_0)$; broader exact sourced GR-plus-Maxwell reduction ($G_{\mathrm{eff}} = G$, $f_{\mathrm{EM}}(0) = f_{\mathrm{SM}}(0) = 1$, $m_A = 0$, with zero scalar residual for arbitrary $\mathcal{L}_{\mathrm{matter}}$ and $F_{\mu\nu}$) holds in the separate unbroken parameter limit $v_0 = 0$; and heavy-scalar decoupling ($m_1, m_2 \gg E_{\mathrm{obs}}$) provides a controlled asymptotic approximation.
\end{enumerate}

\subsection{Consistently Normalized Conditional Conformal-Frame Weak-Field Potentials, Shapiro Invariance, and PPN Bound}
\label{app:ppn-derivation}

Around Minkowski space in the static weak-field regime, Einstein-frame metric perturbations $ds_E^2 = -(1 + 2\Phi_E)dt^2 + (1 - 2\Psi_E)d\mathbf{x}^2$ sourced by static Einstein-frame mass $M_E$ satisfy $\nabla^2\Phi_E = \nabla^2\Psi_E = 4\pi G M_E \delta^{(3)}(\mathbf{r})$, giving $\Phi_E(r) = \Psi_E(r) = -G M_E / r$. Simultaneously, linearizing $(\delta\rho_t, \delta|\psi|)$ in canonically normalized variables $\hat{\Phi} = K^{1/2}\delta\Phi$ diagonalizes the positive-definite physical mass matrix $M_{\mathrm{phys}}^2 = K^{-1/2} H K^{-1/2}$ ($H u_a = m_a^2 K u_a$) into physical eigenmasses $m_1, m_2 > 0$ with mixing angle $\vartheta$. \emph{Stage-6 Action-Equivalence Qualification (\texttt{ActionIdentityDiagnostics}):} Stage~6 proves in Lean~4.30 that the unchanged multiplicative portal $+\sqrt{-g}\,f_{\mathrm{SM}}(\Rcoh)\mathcal{L}_{\mathrm{matter}}$ does \emph{not} algebraically match a conformal Jordan metric $g_{\mu\nu}^{(J)} = A^2(\Rcoh)\,g_{\mu\nu}^{(E)}$ at fixed matter variables unless $f_{\mathrm{SM}} = 1$. Conditional on explicitly adopting a universal conformal worldline coupling $S_{\mathrm{pp}} = -m\int A(\Rcoh)\,d\tau_E$ with smooth $A(\Rcoh) > 0$, consistent derivation of both the scalar source and the linearized metric $g_{\mu\nu}^{(J)}/A_0^2 = [1 + 2(A'(v_0)/A(v_0))\delta|\psi|]g_{\mu\nu}^{(E)}$ yields
\begin{equation}
\Phi_J(r) = \Phi_E(r)\bigl[1 + \alpha_{\mathrm{conf}}[A]\,S(r)\bigr],
\qquad
\Psi_J(r) = \Phi_E(r)\bigl[1 - \alpha_{\mathrm{conf}}[A]\,S(r)\bigr],
\label{eq:appD_jordan_potentials}
\end{equation}
with $S(r) = \sin^2\vartheta\,e^{-m_1 r} + \cos^2\vartheta\,e^{-m_2 r} \in (0, 1]$ and $\alpha_{\mathrm{conf}}[A] = \frac{2 M_{\mathrm{Pl}}^2}{\kappa_c}[A'(v_0)/A(v_0)]^2 \ge 0$. Explicitly separating the two normalizations conflated in earlier drafts (Section~\ref{sec:ppn-shapiro-exact}): Model~C1 ($g_{\mu\nu}^{(J)} = f_{\mathrm{SM}}^{-1}g_{\mu\nu}^{(E)}$, $A = f_{\mathrm{SM}}^{-1/2}$) yields $\alpha_{\mathrm{conf}}^{(\mathrm{C1})} = \frac{1}{4}\alpha_{\mathrm{sym}}$, whereas Model~C2 ($g_{\mu\nu}^{(J)} = f_{\mathrm{SM}}^{\pm 2}g_{\mu\nu}^{(E)}$, $A = f_{\mathrm{SM}}^{\pm 1}$) yields $\alpha_{\mathrm{conf}}^{(\mathrm{C2})} = \alpha_{\mathrm{sym}} = \frac{2 M_{\mathrm{Pl}}^2}{\kappa_c}[f'_{\mathrm{SM}}(v_0)/f_{\mathrm{SM}}(R_0)]^2$. Adding the two potentials yields $\Phi_J(r) + \Psi_J(r) = 2\Phi_E(r)$ (conformal-frame Shapiro null-cone invariance), while dividing gives
\begin{equation}
\gamma_{\mathrm{PPN}}(r) = \frac{\Psi_J(r)}{\Phi_J(r)} = \frac{1 - \alpha_{\mathrm{conf}}S(r)}{1 + \alpha_{\mathrm{conf}}S(r)},
\qquad
\bigl|\gamma_{\mathrm{PPN}}(r) - 1\bigr| = \frac{2\alpha_{\mathrm{conf}}S(r)}{1 + \alpha_{\mathrm{conf}}S(r)} < 2.3 \times 10^{-5}
\end{equation}
for all $\alpha_{\mathrm{conf}} S(r) < 1.15 \times 10^{-5}$, proved algebraically in Lean~4.30 (\texttt{ppn\allowbreak{}Gamma\allowbreak{}Bound\_\allowbreak{}of\_\allowbreak{}alpha\_\allowbreak{}and\_\allowbreak{}screening}) for the abstract conformal coupling $\alpha_{\mathrm{conf}}$ subject to the Stage-6 non-equivalence caveat above. For the direct unchanged-action weak-field channel without assuming a conformal Jordan frame, see Section~\ref{sec:stage6-action-diagnostics}.

\subsection{Mutual Information Regulators and Operational Normalization}

For continuous field configurations partitioned into local region $X$ and complement $Y$ at UV coarse-graining scale $\ell_{\mathrm{reg}} \ge \Lambda_{\mathrm{UV}}^{-1}$, regulated mutual information
\[
\mathcal{I}_{\ell_{\mathrm{reg}}}(X : Y) = \int dX\,dY\,p_{\ell_{\mathrm{reg}}}(X,Y)\ln\frac{p_{\ell_{\mathrm{reg}}}(X,Y)}{p_{\ell_{\mathrm{reg}}}(X)p_{\ell_{\mathrm{reg}}}(Y)} \in [0, \mathcal{I}_{\max}(\ell_{\mathrm{reg}})]
\]
remains finite and monotonic under CPTP coarse-graining by the quantum data-processing inequality, ensuring $\Rcoh^{\mathrm{proxy}}(x) = \mathcal{I}_{\ell_{\mathrm{reg}}}/\mathcal{I}_{\max} \in [0,1]$ is ultraviolet-safe across laboratory readouts.

\section{Explicit Recovery of Known Physics in SIT Limits}

\subsection{Gravity Limit: Recovery of General Relativity and Constraints}
\label{app:GRlimit}

To demonstrate the physical viability of SIT, we derive the reduction to Einstein's General Relativity in the appropriate limit, and compute explicit leading corrections.

\paragraph{SIT Action Recap.}
Recall the SIT action:
\begin{align*}
S_{\mathrm{SIT}} = \int d^4x\,\sqrt{-g}\Bigg[
    &\frac{R}{16\pi G}
    - \frac{\kappa_t}{2}\,g^{\mu\nu}\partial_\mu\rho_t\partial_\nu\rho_t - V(\rho_t)
    - f_1(\rho_t, R_{\mathrm{coh}})\sum_{\psi_f} m_{\psi_f}\,\bar\psi_f\psi_f \\
    &- \frac{1}{4}\,f_2(\rho_t, R_{\mathrm{coh}})\,F_{\mu\nu}F^{\mu\nu}
    - \frac{\kappa_c}{2}\,(\partial_\mu |\psi|)(\partial^\mu |\psi|) - U(|\psi|)
    + \mathcal{L}_{\text{int}}(\rho_t, R_{\mathrm{coh}})
\Bigg].
\end{align*}

\paragraph{Assumptions for the Gravity Limit.}
\noindent\textbullet~Assume $R_{\mathrm{coh}}$ is spatially and temporally constant, or its gradients are negligible (strong decoherence regime).
\noindent\textbullet~Take $\rho_t(x) = \rho_0 + \delta \rho_t(x)$, with $\delta \rho_t(x)$ small.
\noindent\textbullet~Potentials $V$ and $U$ have stable minima at $\rho_0$, $R_0$.
\noindent\textbullet~Coupling functions $f_1$, $f_2$ reduce to constants.

\paragraph{Variation with respect to $g_{\mu\nu}$.}
The field equations become:
\[
G_{\mu\nu} = 8\pi G \left[ T_{\mu\nu}^{\text{matter}} + T_{\mu\nu}^{(\rho_t)} \right]
\]
where
\[
T_{\mu\nu}^{(\rho_t)} = \kappa_t \partial_\mu \rho_t \partial_\nu \rho_t - g_{\mu\nu} \left( \frac{\kappa_t}{2} g^{\alpha\beta} \partial_\alpha \rho_t \partial_\beta \rho_t + V(\rho_t) \right)
\]

In vacuum and for constant $\rho_t = \rho_0$, the displayed action term $-\sqrt{-g}\,V(\rho_t)$ gives $T_{\mu\nu}^{(\rho_t)} = -g_{\mu\nu} V(\rho_0)$ (correcting an earlier spurious factor of $1/2$), which vanishes identically at the canonical vacuum $V_{\mathrm{tot}}(\rho_0, v_0) = 0$ of \texttt{M\_AUTH\_SIT\_EFT\_v1} (and otherwise acts as a cosmological-constant shift).

\paragraph{Yukawa Correction in the Weak-Field Limit.}
Linearize about Minkowski space ($g_{\mu\nu} = \eta_{\mu\nu} + h_{\mu\nu}$, $|\delta \rho_t| \ll 1$). In the decoupled single-scalar truncation, the equation of motion for $\delta\rho_t$ reads:
\[
\kappa_t \square \delta\rho_t - V''(\rho_0)\, \delta\rho_t = 0,
\]
whose canonically normalized squared mass is
\[
m_t^2 = \frac{V''(\rho_0)}{\kappa_t}
\]
(reducing to $V''(\rho_0)$ only when $\kappa_t = 1$ is explicitly chosen; in full \texttt{M\_AUTH\_SIT\_EFT\_v1}, the two-mode spectrum is governed by $(p^2 K + H)^{-1}$ in Subsection~E.5). Under a declared matter coupling, a static point mass $M$ acquires a Yukawa modification $\Phi(r) = -\frac{GM}{r}\bigl[1 + \alpha \exp(-m_t r)\bigr]$.

\paragraph{Experimental Constraints and Retirement of Earlier Parameter Box.}
Earlier drafts asserted that SIT is viable only if $\alpha \lesssim 10^{-5}$ and $m_t^{-1} \gtrsim 10^4\,\mathrm{m}$, declaring all deviations outside that box empirically excluded. Version~81 explicitly \emph{retires} that parameter box as an invalid empirical requirement: a sufficiently decoupled or heavy scalar produces no detectable fifth force regardless of interaction range, and canonical benchmark \texttt{BM\_PHYS\_EFT\_1} itself is extraordinarily short-ranged ($\lambda_1 \approx 8.056\times 10^{-31}\,\mathrm{m}$, Section~\ref{sec:stage6-action-diagnostics}). Confronting \texttt{M\_AUTH\_SIT\_EFT\_v1} with torsion-balance and Solar System data requires the conditional response and primary-provenance likelihood program of Section~\ref{sec:stage6-action-diagnostics} and \texttt{OPEN-6A}/\texttt{OPEN-6D} rather than a fixed one-parameter box.

\paragraph{Conclusion.}
Thus, at $(\rho_t, |\psi|) = (\rho_0, v_0)$ where $V_{\mathrm{tot}} = 0$, the metric equation reduces algebraically to General Relativity (forming an exact coupled solution in source-free vacuum, an exact sourced GR+Maxwell reduction in the unbroken limit $v_0 = 0$, and a controlled asymptotic approximation in the heavy-scalar decoupling regime, Appendix~D.2), while any finite-range scalar correction must be tested against primary datasets via the action-to-apparatus likelihood pipeline (\texttt{OPEN-6}).

\subsection{Quantum Field Theory Limit: Flat Spacetime and Decoupled Background}
\label{app:QFTlimit}

\paragraph{Assumptions.}
\noindent\textbullet~Flat metric: $g_{\mu\nu} \to \eta_{\mu\nu}$.
\noindent\textbullet~Unbroken or heavy-scalar decoupled background where $(\rho_t, |\psi|)$ fluctuations decouple and $f_{\mathrm{EM}}, f_{\mathrm{SM}} \to \text{const}$ (rescaled into canonical field normalizations).

\paragraph{Action Simplifies:}
\[
S_{\mathrm{QFT\,limit}} = \int d^4x \left[
    \mathcal{L}_{\text{matter}} + \mathcal{L}_{\text{EM}}
\right],
\]
recovering the standard Minkowski quantum field theory action with canonical equal-time commutation relations $[\hat{\phi}(x), \hat{\pi}(y)] = i \hbar \delta^3(x - y)$ in the decoupled limit. Weak residual $(\rho_t, \psi)$ couplings enter only as cutoff-suppressed perturbative EFT operators below $\Lambda_{\mathrm{UV}}$.

\subsection{Conditional Classical Kinetic and Hydrodynamic Correspondence (Open Microscopic Bridge)}
\label{app:KineticLimit}

Earlier drafts of this subsection identified the classical EFT modulus $\Rcoh(x) = |\psi(x)|/\psi_0$ with local ensemble purity, equated thermodynamic entropy production to $-d\Rcoh/dt$, and asserted that SIT derives the Boltzmann and Navier--Stokes equations directly from $S_{\mathrm{AUTH}}$. Consistent with the operational separation contract (Section~\ref{sec:PrimitiveFieldDefinitions}) and Section~\ref{sec:arrow-of-time}, Version~81 explicitly \emph{re-scopes} this subsection as an external/conditional kinetic correspondence rather than a derivation from $S_{\mathrm{AUTH}}$:
\begin{enumerate}[leftmargin=2em]
  \item \textbf{Standard Kinetic Equations Under Their Actual Assumptions:} In the constant-background decoupled limit where SIT scalar gradients vanish, dilute many-body gases obey the classical Boltzmann equation
  \[
  \frac{\partial f}{\partial t} + \vec{v} \cdot \nabla_x f + \vec{F} \cdot \nabla_p f = \left( \frac{\partial f}{\partial t} \right)_{\mathrm{coll}},
  \]
  and, under Chapman--Enskog hydrodynamic closure, the compressible/incompressible Navier--Stokes equations
  \[
  \rho \left( \frac{\partial \vec{v}}{\partial t} + \vec{v} \cdot \nabla \vec{v} \right)
  = -\nabla p + \eta \nabla^2 \vec{v} + \cdots,
  \]
  strictly by virtue of standard microscopic many-body mechanics (for example, the hard-sphere Boltzmann--Grad propagation-of-chaos theorem of Deng--Hani--Ma~2025).
  \item \textbf{What $S_{\mathrm{AUTH}}$ Does \emph{Not} Derive:} The canonical action \texttt{M\_AUTH\_SIT\_EFT\_v1} does \emph{not} derive a microscopic collision operator $(\partial_t f)_{\mathrm{coll}}$ from $(\rho_t, \psi)$, does \emph{not} identify the classical scalar amplitude $\Rcoh = |\psi|/\psi_0$ with quantum ensemble purity $\mathrm{Tr}(\rho^2)$, and does \emph{not} equate Boltzmann entropy production $-k_B\int C[f]\ln f\,d^3p \ge 0$ to $-d\Rcoh/dt$. Constructing a rigorous Wigner/BBGKY kinetic transport closure for matter coupled to fluctuating $(\rho_t, \psi)$ fields remains an open research problem (see Section~\ref{sec:navier-stokes-boundary} for the phenomenological hyperviscous regularization boundary).
\end{enumerate}

\subsection*{Summary of Reductions and Open Microscopic Bridges}

\begin{description}
    \item[Gravity Reduction] Constant vacuum fields $(\rho_0, v_0)$ reduce the metric equation algebraically to Einstein's equations (exact in source-free vacuum; exact sourced GR+Maxwell in the unbroken limit $v_0 = 0$; controlled asymptotic approximation under heavy-scalar decoupling).
    \item[Decoupled QFT Limit] In flat spacetime with decoupled scalar fluctuations, matter and gauge sectors retain standard Minkowski QFT dynamics up to $\Lambda_{\mathrm{UV}}$-suppressed EFT operators.
    \item[Kinetic Correspondence] Standard Boltzmann and Navier--Stokes equations hold in the decoupled background under their established statistical-mechanical hypotheses, while SIT-specific kinetic collision and transport derivations remain open.
\end{description}

\medskip\hrule\medskip
\noindent\textbf{Scope Note:} Rather than claiming complete first-principles derivation of kinetic theory or quantum measurement from classical scalar fields, Version~81 establishes exact algebraic/unbroken/decoupled consistency with General Relativity and perturbative QFT while keeping microscopic kinetic and measurement bridges explicitly open.

\subsection{Worked Example: Linearized Field Equations and Coupled Dispersion Relations}

We consider small perturbations of the canonical SIT scalar fields of \texttt{M\_AUTH\_SIT\_EFT\_v1} in mostly-plus signature $(-,+,+,+)$ about the broken-phase equilibrium $(\rho_0, v_0)$ (with $\Rcoh = |\psi|/\psi_0$, $R_0 = v_0/\psi_0$, $\psi_0 \ne 0$, and effective radial kinetic normalization $\kappa_R := \kappa_c\psi_0^2 > 0$):
\[
\rho_t(x) = \rho_0 + \delta \rho_t(x), \qquad \Rcoh(x) = R_0 + \delta \Rcoh(x), \qquad h(x) := |\psi(x)| - v_0 = \psi_0\,\delta\Rcoh(x),
\]
packaged either into the dimensionless radial perturbation vector $\delta\Phi_R := (\delta\rho_t, \delta\Rcoh)^\top$ used in this worked example or into the canonical Stage-6 modulus perturbation vector $\delta\Phi_h := (\delta\rho_t, h)^\top = D\,\delta\Phi_R$, where $D := \mathrm{diag}(1, \psi_0)$ is invertible.

\paragraph{Step 1: Expand the canonical action to second order in perturbations}

Starting from the canonical $(-,+,+,+)$ action \texttt{M\_AUTH\_SIT\_EFT\_v1} (Eq.~\eqref{eq:abstract_master_action}, with healthy \emph{negative} scalar kinetic terms under $\eta_{\mu\nu} = \mathrm{diag}(-1,+1,+1,+1)$ so that $+\frac{1}{2}\kappa_t(\partial_t\rho_t)^2 + \frac{1}{2}\kappa_R(\partial_t\Rcoh)^2 \ge 0$):
\[
S_{\mathrm{AUTH}} = \int d^4x \sqrt{-g} \left[
\frac{1}{16\pi G} R
- \frac{\kappa_t}{2} g^{\mu\nu} \partial_\mu \rho_t \partial_\nu \rho_t
- \frac{\kappa_R}{2} g^{\mu\nu} \partial_\mu \Rcoh \partial_\nu \Rcoh
- V_{\mathrm{tot}}(\rho_t, \psi_0 \Rcoh)
+ \cdots \right],
\]
expanding $V_{\mathrm{tot}}$ around its minimum $(\rho_0, \psi_0 R_0)$ yields the positive-definite quadratic potential $\frac{1}{2}\,\delta\Phi_R^\top H_R\,\delta\Phi_R$, where $K_R = \mathrm{diag}(\kappa_t, \kappa_R)$ (or $\mathrm{diag}(\kappa_t, \kappa_c(1+\delta K)\psi_0^2)$ when including the Stage-6 lattice AM--GM shift $\delta K \ge 0$) and $H_R = \partial^2 V_{\mathrm{tot}}/\partial(\rho_t, \Rcoh)^2|_{(\rho_0, R_0)}$ are related by the exact diagonal coordinate congruence $D = \mathrm{diag}(1, \psi_0)$ to Stage~6's $(\delta\rho_t, h)$ kinetic matrix $K_h = \mathrm{diag}(\kappa_t, \kappa_c(1+\delta K))$ and portal Hessian $H_h = \partial^2 V_{\mathrm{tot}}/\partial(\rho_t, |\psi|)^2|_{(\rho_0, v_0)}$:
\[
K_R = D^\top K_h D, \qquad H_R = D^\top H_h D, \qquad J_R = D^\top J_h, \qquad V_R = D^{-1} V_h.
\]
The quadratic scalar Lagrangian is therefore
\[
\mathcal{L}^{(2)}_{\mathrm{scal}} = -\frac{1}{2}\,\eta^{\mu\nu}\,(\partial_\mu\delta\Phi_R)^\top K_R (\partial_\nu\delta\Phi_R) - \frac{1}{2}\,\delta\Phi_R^\top H_R\,\delta\Phi_R + \delta\Phi_R^\top J_R.
\]

\paragraph{Step 2: Derive coupled linearized equations of motion}

Varying $\int d^4x\,\mathcal{L}^{(2)}_{\mathrm{scal}}$ on Minkowski background $\eta_{\mu\nu} = \mathrm{diag}(-1,+1,+1,+1)$ with d'Alembertian $\square := \eta^{\mu\nu}\partial_\mu\partial_\nu = -\partial_t^2 + \nabla^2$ (note $\delta(-\frac{1}{2}\eta^{\mu\nu}\partial_\mu\phi\,\partial_\nu\phi) = +\phi\,\square\phi$ upon integration by parts) yields the exact coupled linearized Klein--Gordon system:
\[
(K_R\,\square - H_R)\,\delta\Phi_R = -J_R,
\qquad\text{i.e.,}\qquad
\begin{pmatrix} \kappa_t\square - (H_R)_{11} & -(H_R)_{12} \\ -(H_R)_{12} & \kappa_R\square - (H_R)_{22} \end{pmatrix}
\begin{pmatrix} \delta\rho_t \\ \delta\Rcoh \end{pmatrix}
= -\begin{pmatrix} J_\rho \\ J_R \end{pmatrix}.
\]
In Euclidean momentum space ($-\square \mapsto p^2 \ge 0$), the canonical Stage-6 instantiation theorem \texttt{Stage6Addendum.\allowbreak{}ActionResponsePropagatorIdentity.\allowbreak{}canonical\_\allowbreak{}sit\_\allowbreak{}two\_\allowbreak{}mode\_\allowbreak{}eigenbasis\_\allowbreak{}and\_\allowbreak{}propagator} constructs directly from \texttt{AuthModelParams} ($v_0 > 0$, $\delta K \ge 0$) a $K_h$-orthonormal generalized eigenbasis $V_h$ with strictly positive squared eigenmasses $0 < m_1^2 \le m_2^2$ (automatically discharging $p^2 + m_a^2 \ne 0$ for every $p^2 \ge 0$), and \texttt{canonical\_\allowbreak{}coord\_\allowbreak{}transform\_\allowbreak{}propagator\_\allowbreak{}equiv} proves that $V_R = D^{-1} V_h$ preserves both the eigenmasses and the source--detector propagator inner product identically:
\begin{gather*}
(p^2 K_R + H_R)^{-1} = V_R\,\mathrm{diag}\!\left(\frac{1}{p^2 + m_1^2},\; \frac{1}{p^2 + m_2^2}\right) V_R^\top,\\[0.3em]
J_{\mathrm{det},R}^\top (p^2 K_R + H_R)^{-1} J_{\mathrm{src},R} = J_{\mathrm{det},h}^\top (p^2 K_h + H_h)^{-1} J_{\mathrm{src},h}.
\end{gather*}

\paragraph{Step 3: Analyze the uncoupled limit ($(H_R)_{12} = 0$, $J_R = 0$)}

When portal mixing vanishes ($(H_R)_{12} = 0$, $J_R = 0$), setting $m_\rho^2 := (H_R)_{11} > 0$ and $m_R^2 := (H_R)_{22} > 0$, the system decouples into two healthy non-tachyonic Klein--Gordon equations under $\square = -\partial_t^2 + \nabla^2$:
\[
\kappa_t \square \delta \rho_t - m_{\rho}^2 \delta \rho_t = 0,
\qquad
\square \delta \Rcoh - \frac{m_{R}^2}{\kappa_R} \delta \Rcoh = 0.
\]
Substituting a plane wave $e^{-i\omega t + i\mathbf{k}\cdot\mathbf{x}}$ ($\square \mapsto +\omega^2 - |\mathbf{k}|^2$) yields the healthy dispersion relations $\omega^2 = |\mathbf{k}|^2 + m_\rho^2/\kappa_t$ and $\omega^2 = |\mathbf{k}|^2 + m_R^2/\kappa_R$, with strictly positive squared masses (invariant under $D = \mathrm{diag}(1, \psi_0)$ since $(H_R)_{22}/\kappa_R = (\psi_0^2 (H_h)_{22})/(\psi_0^2 \kappa_c) = (H_h)_{22}/\kappa_c$).

\paragraph{Step 4: Recover General Relativity limit}

In the decoupled vacuum limit where $\delta\rho_t, \delta\Rcoh \to 0$ (or at energies far below the heavy portal masses $m_1, m_2$), scalar fluctuations freeze at the minimum $V_{\mathrm{tot}}(\rho_0, v_0) = 0$ and the gravitational equations reduce to Einstein's equations $G_{\mu\nu} = 8\pi G_{\mathrm{eff}} T_{\mu\nu}$.

\paragraph{Step 5: Physical interpretation of correlation lengths}

The eigenmasses $m_1, m_2$ set inverse correlation lengths $\lambda_a = \hbar c / m_a$ over which scalar perturbations decay exponentially in static backgrounds (e.g.\ $\lambda_1 \approx 8.056\times 10^{-31}\,\mathrm{m}$ at canonical benchmark \texttt{BM\_PHYS\_EFT\_1}).

\medskip\hrule\medskip
This rederived worked example aligns Appendix~E exactly with the canonical $(-,+,+,+)$ signature of \texttt{M\_AUTH\_SIT\_EFT\_v1} and maps its dimensionless $(\delta\rho_t, \delta\Rcoh)$ matrices $(K_R, H_R)$ to the Stage-6 $(\delta\rho_t, h)$ normalization $(K_h, H_h)$ via the explicit invertible congruence $D = \mathrm{diag}(1, \psi_0)$.

\subsection{Worked Example: Linearized Metric and Canonical Scalar Stress--Energy Tensor}

Consider perturbations of the metric around Minkowski spacetime:
\[
g_{\mu\nu} = \eta_{\mu\nu} + h_{\mu\nu}, \qquad |h_{\mu\nu}| \ll 1,
\]
where $\eta_{\mu\nu} = \mathrm{diag}(-1, +1, +1, +1)$.

\paragraph{Step 1: Expand the canonical action to quadratic order in perturbations}

Expanding $S_{\mathrm{AUTH}}$ to second order in $h_{\mu\nu}$ and $\delta\Phi = (\delta\rho_t, \delta\Rcoh)^\top$ gives:
\begin{equation}
\label{eq:expanded_action}
\begin{split}
S^{(2)} = \int d^4x \Biggl[ &
    \frac{1}{64 \pi G} h^{\mu\nu} \mathcal{E}_{\mu\nu}^{\alpha\beta} h_{\alpha\beta}
    - \frac{\kappa_t}{2} \eta^{\mu\nu} \partial_\mu \delta \rho_t \partial_\nu \delta \rho_t
    - \frac{\kappa_R}{2} \eta^{\mu\nu} \partial_\mu \delta \Rcoh \partial_\nu \delta \Rcoh \\
&   - \frac{1}{2}\,\delta\Phi^\top H\,\delta\Phi
    + \mathcal{L}_{\mathrm{int}} \Biggr],
\end{split}
\end{equation}
where $\mathcal{E}_{\mu\nu}^{\alpha\beta}$ is the Lichnerowicz operator governing linearized gravity and $\frac{1}{2}\delta\Phi^\top H\,\delta\Phi = \frac{1}{2}H_{11}(\delta\rho_t)^2 + H_{12}\,\delta\rho_t\,\delta\Rcoh + \frac{1}{2}H_{22}(\delta\Rcoh)^2$.

\paragraph{Step 2: Derive linearized Einstein equations with canonical scalar stress--energy}

Using $T_{\mu\nu} = -\frac{2}{\sqrt{-g}}\frac{\delta S_{\mathrm{scal}}}{\delta g^{\mu\nu}}$ in $(-,+,+,+)$ signature, varying Eq.~\eqref{eq:expanded_action} with respect to $h_{\mu\nu}$ yields $\mathcal{E}_{\mu\nu}^{\alpha\beta} h_{\alpha\beta} = 16 \pi G \, T_{\mu\nu}^{\mathrm{eff}}$, with $T_{\mu\nu}^{\mathrm{eff}} = T_{\mu\nu}^{\mathrm{matter}} + T_{\mu\nu}^{\mathrm{scal}}$ and exact canonical scalar stress--energy tensor
\begin{align*}
T_{\mu\nu}^{\mathrm{scal}}
&= \kappa_t\,\partial_\mu \delta \rho_t\,\partial_\nu \delta \rho_t
 + \kappa_R\,\partial_\mu \delta \Rcoh\,\partial_\nu \delta \Rcoh \\
&\quad - \eta_{\mu\nu}\!\left(
  \frac{\kappa_t}{2}\eta^{\alpha\beta}\partial_\alpha\delta\rho_t\partial_\beta\delta\rho_t
  + \frac{\kappa_R}{2}\eta^{\alpha\beta}\partial_\alpha\delta\Rcoh\partial_\beta\delta\Rcoh
  + \frac{1}{2}\delta\Phi^\top H\,\delta\Phi
\right).
\end{align*}
In the uncoupled limit ($H_{12} = 0$, $H_{11} = m_\rho^2$, $H_{22} = m_R^2$), this splits cleanly into
\[
T_{\mu\nu}^{\rho_t} = \kappa_t\,\partial_\mu \delta \rho_t\,\partial_\nu \delta \rho_t - \eta_{\mu\nu}\!\left(\frac{\kappa_t}{2}\partial^\alpha \delta \rho_t \partial_\alpha \delta \rho_t + \frac{1}{2} m_{\rho}^2 (\delta \rho_t)^2\right),
\]
\[
T_{\mu\nu}^{\Rcoh} = \kappa_R\,\partial_\mu \delta \Rcoh\,\partial_\nu \delta \Rcoh - \eta_{\mu\nu}\!\left(\frac{\kappa_R}{2}\partial^\alpha \delta \Rcoh \partial_\alpha \delta \Rcoh + \frac{1}{2} m_{R}^2 (\delta \Rcoh)^2\right),
\]
preserving $\kappa_t$ on $\partial_\mu\delta\rho_t\partial_\nu\delta\rho_t$ and keeping the potential density $\frac{1}{2}m_R^2(\delta\Rcoh)^2$ strictly free of any extraneous kinetic factor $\kappa_R$, with positive energy density $T_{00}^{\mathrm{scal}} \ge 0$.

\paragraph{Step 3: Gauge fixing and wave equation}

Choosing harmonic (Lorenz) gauge $\partial^\mu \bar{h}_{\mu\nu} = 0$ for $\bar{h}_{\mu\nu} = h_{\mu\nu} - \frac{1}{2}\eta_{\mu\nu} h$ gives $\square \bar{h}_{\mu\nu} = -16\pi G\,T_{\mu\nu}^{\mathrm{eff}}$ (with $\square = -\partial_t^2 + \nabla^2$).

\paragraph{Step 4: Interpretation and recovery of classical GR}

When $\delta\rho_t, \delta\Rcoh \to 0$, $T_{\mu\nu}^{\mathrm{scal}} \to 0$ and linearized General Relativity is recovered identically.

\paragraph{Step 5: Coupled dynamics with scalar fields}

Together with $\square\bar{h}_{\mu\nu} = -16\pi G\,T_{\mu\nu}^{\mathrm{eff}}$, the scalar perturbations satisfy the canonical coupled Klein--Gordon system $(K\square - H)\delta\Phi = -J$ (reducing for $H_{12}=0, J=0$ to $\kappa_t\square\delta\rho_t - m_\rho^2\delta\rho_t = 0$ and $\square\delta\Rcoh - (m_R^2/\kappa_R)\delta\Rcoh = 0$), with $\partial^\mu T_{\mu\nu}^{\mathrm{scal}} = 0$ on shell when $J = 0$.

\medskip\hrule\medskip
This worked example connects the linearized metric perturbation consistently to the canonical scalar kinetic matrix $K$, potential Hessian $H$, and positive-energy stress--energy tensor of \texttt{M\_AUTH\_SIT\_EFT\_v1}.

\subsection{Retirement of Legacy Parameter Boxes and the \texttt{OPEN-6} Likelihood Program}

Earlier drafts of this appendix listed illustrative phenomenological boxes ($10^{-18}\,\mathrm{eV} \lesssim m_\rho, m_R \lesssim 10^{-3}\,\mathrm{eV}$, $10^{-4} \lesssim \lambda \lesssim 10^{-2}$, and $|\delta R_{\mathrm{coh}}| \lesssim 10^{-5}$) and asserted that those ranges ensured consistency with Solar System, laboratory, cosmological, and quantum-coherence tests. Version~81 explicitly \emph{retires} all such ungrounded parameter boxes as empirical conclusions for three reasons:
\begin{enumerate}[leftmargin=2em]
  \item \textbf{Missing Action-to-Charge-to-Apparatus Map (\texttt{OPEN-6}):} Converting primary Solar System, torsion-balance, optical-clock, or cold-atom datasets into valid exclusion contours on \texttt{M\_AUTH\_SIT\_EFT\_v1} requires deriving microscopic material scalar charges (\texttt{OPEN-6A}), the visible $U(1)_{\mathrm{EM}}$ photon dictionary (\texttt{OPEN-6B}), laboratory preparation sources $J_{\mathrm{prep}}$ (\texttt{OPEN-6C}), and full apparatus transfer functions (\texttt{OPEN-6D}), subject to the mass-calibration and form-factor identifiability controls of Section~\ref{subsec:identifiability-and-qgtcd-guardrails}.
  \item \textbf{Field-Normalization Convention Dependence of Kinetic Coefficients:} Bounding a standalone kinetic coefficient (such as $\lambda$ or $\kappa_R$) without fixing the field normalization and the remaining couplings in the action is convention-dependent, because rescaling a scalar field absorbs kinetic normalization into its mass and interaction couplings (as formalized by $D = \mathrm{diag}(1, \psi_0)$ in Subsection~E.5).
  \item \textbf{Heavy Decoupled Benchmarks Lie Far Outside the Legacy Box:} Canonical benchmark \texttt{BM\_PHYS\_EFT\_1} places the scalar eigenmasses at $m_1 \approx 2.45\times 10^{14}\,\mathrm{GeV}$ and $m_2 \approx 2.83\times 10^{16}\,\mathrm{GeV}$ ($\lambda_1 \approx 8.056\times 10^{-31}\,\mathrm{m}$), demonstrating that heavy, radiatively stable regimes lie over $25$ orders of magnitude above the legacy $10^{-3}\,\mathrm{eV}$ box while remaining strictly non-excluded under the conditional dust closure of Section~\ref{sec:stage6-action-diagnostics}.
\end{enumerate}

Archival exploratory scales from earlier drafts are therefore retained solely as historical context; quantitative empirical bounds belong exclusively to the conditional response and primary-provenance likelihood program documented in Section~\ref{sec:stage6-action-diagnostics} and \texttt{primary\_measurements\_manifest\_v2.json}.

\subsection{Prospective Experimental Arenas (Contingent on \texttt{OPEN-6})}

Once the open action-to-apparatus derivations (\texttt{OPEN-6A}--\texttt{OPEN-6D}) are completed for a specified parameter branch, four experimental arenas provide scoped sensitivity tests: (i)~optical-clock comparisons (with disambiguated $\alpha_\rho$ and $\alpha_\Phi = \alpha_\rho\mathcal{T}$), (ii)~short-range torsion-balance and atom-interferometer inverse-square-law tests (starting with full patterned rotating-mass likelihood reproduction of Lee~2020), (iii)~matched-preparation BEC vs.\ thermal-cloud differential-acceleration tests obeying Eq.~\eqref{eq:bec_matched_residual}, and (iv)~gravitational-wave propagation and astrophysical lensing surveys. No unconditional deviation amplitude is claimed ahead of those derivations.

\section{Holonomy of the Coherence Field and Electromagnetic Tensor}
\label{sec:HolonomyEM}

\subsection{1. Historical Flat-Connection Phase-Holonomy Ansatz and Its $F_{\mu\nu}=0$ Limitation}

To preserve historical continuity while aligning with the corrected gauge structure of Part~II and Section~45, we record the earlier flat-connection phase-holonomy ansatz and state its exact mathematical boundary.

\paragraph{Coherence Field and Phase Fiber:}
In that historical sketch, an auxiliary complex order parameter was written as $\Psi(x) = \sqrt{R_{\mathrm{coh}}(x)}\,e^{i\theta(x)}$ with local phase $\theta(x)$ (distinct from the canonical linear modulus normalization $\Rcoh = |\psi|/\psi_0$ of \texttt{M\_AUTH\_SIT\_EFT\_v1}).

\paragraph{Flat-Connection Ansatz and Holonomy:}
Earlier drafts considered the pure-gradient expression
\[
A_\mu^{\mathrm{flat}}(x) := \frac{\hbar}{e} \partial_\mu \theta(x),
\]
for which transport around a closed loop $C$ in a multiply connected domain yields $\exp\bigl(i\frac{e}{\hbar}\oint_C A_\mu^{\mathrm{flat}}\,dx^\mu\bigr) = \exp(i\Delta\theta)$.

\paragraph{Curvature Limitation ($F_{\mu\nu} = 0$):}
The field strength of any smooth locally defined phase gradient vanishes identically:
\[
F_{\mu\nu}^{\mathrm{flat}} := \partial_\mu A_\nu^{\mathrm{flat}} - \partial_\nu A_\mu^{\mathrm{flat}} = \frac{\hbar}{e} (\partial_\mu \partial_\nu - \partial_\nu \partial_\mu)\theta(x) = 0.
\]
While nontrivial global holonomy can exist around singular vortex cores or multiply connected flux-tube exteriors, a pure gradient does \emph{not} generate ordinary non-flat magnetic fields ($F_{\mu\nu} \ne 0$).

\paragraph{Canonical Interpretation (\texttt{OPEN-6B}):}
Accordingly, Version~81 retains the independent principal-bundle connection $A_\mu$ of \texttt{M\_AUTH\_SIT\_EFT\_v1} (with arbitrary curvature $F_{\mu\nu} = \partial_\mu A_\nu - \partial_\nu A_\mu$), treats $A_\mu^{\mathrm{flat}} = (\hbar/e)\partial_\mu\theta$ solely as a historical pure-gauge ansatz, and keeps the identification of the visible massless photon and visible electromagnetism explicitly open under \texttt{OPEN-6B}.

\subsection{2. Projective Qubit Measurement vs.\ Relative Phase Resetting: Worked Example}

We summarize standard projective measurement for a two-level quantum system while explicitly correcting an error in earlier drafts that equated computational-basis projection with fixing the qubit's relative phase $\theta$ to $0$ or $\pi$.

\paragraph{Qubit Superposition and Basis Coherence:}
Let the pure qubit state be
\[
|\psi_{\mathrm{Q}}\rangle = \alpha |0\rangle + \beta |1\rangle, \qquad |\alpha|^2 + |\beta|^2 = 1,
\]
with density matrix $\rho = |\psi_{\mathrm{Q}}\rangle\langle\psi_{\mathrm{Q}}|$ and computational-basis off-diagonal coherence magnitude $2|\rho_{01}| = 2|\alpha^*\beta|$.

\paragraph{Why Fixing Relative Phase $\theta \in \{0, \pi\}$ Does \emph{Not} Project onto $|0\rangle$ or $|1\rangle$:}
Earlier drafts wrote that measuring in $\{|0\rangle, |1\rangle\}$ ``projects $\theta$ to $0$ or $\pi$'' and thereby yields $|0\rangle$ or $|1\rangle$ with zero superposition. Version~81 explicitly \emph{retracts} that identification: for $|\psi_{\mathrm{Q}}(\theta)\rangle = (|0\rangle + e^{i\theta}|1\rangle)/\sqrt{2}$, setting $\theta = 0$ or $\theta = \pi$ produces the equal-weight superpositions $|+\rangle = (|0\rangle + |1\rangle)/\sqrt{2}$ and $|-\rangle = (|0\rangle - |1\rangle)/\sqrt{2}$, \emph{both} of which retain computational-basis probabilities $(1/2, 1/2)$ and full off-diagonal coherence magnitude $|\rho_{01}| = 1/2 \ne 0$ ($2|\rho_{01}| = 1$)! Resetting a relative phase unitary-rotates a superposition on the Bloch equator; it never collapses the diagonal populations to $(1, 0)$ or $(0, 1)$.

\paragraph{Assumed Projective Quantum Instrument (Born/L\"uders Rule):}
In standard quantum mechanics, selective projective measurement in the $\{|0\rangle, |1\rangle\}$ basis is modeled as an assumed quantum instrument with orthogonal projectors $\Pi_0 = |0\rangle\langle 0|$, $\Pi_1 = |1\rangle\langle 1|$:
\begin{itemize}[leftmargin=2em]
  \item Before measurement: $\rho = \begin{pmatrix} |\alpha|^2 & \alpha\beta^* \\ \alpha^*\beta & |\beta|^2 \end{pmatrix}$ (basis coherence $2|\alpha^*\beta|$).
  \item Outcome $0$ (probability $p_0 = \mathrm{Tr}(\Pi_0\rho) = |\alpha|^2$): $\rho'_0 = \Pi_0\rho\Pi_0 / p_0 = \begin{pmatrix} 1 & 0 \\ 0 & 0 \end{pmatrix}$ ($2|\rho'_{01}| = 0$).
  \item Outcome $1$ (probability $p_1 = \mathrm{Tr}(\Pi_1\rho) = |\beta|^2$): $\rho'_1 = \Pi_1\rho\Pi_1 / p_1 = \begin{pmatrix} 0 & 0 \\ 0 & 1 \end{pmatrix}$ ($2|\rho'_{01}| = 0$).
\end{itemize}

\paragraph{Operational Scope and Open Status:}
Consistent with Sections~\ref{sec:PrimitiveFieldDefinitions} and~\ref{sec:coherence-conservation}, this elimination of basis-dependent off-diagonal matrix elements under an assumed projective instrument is an open-system/measurement update distinct from classical on-shell $U(1)$ Noether charge conservation ($\nabla_\mu J^\mu_{\mathrm{coh}} = 0$) of the EFT scalar $\psi$, and deriving quantum measurement collapse or the Born rule from \texttt{M\_AUTH\_SIT\_EFT\_v1} remains an open foundational problem.

\section{Unital Pure-Dephasing Lindblad Example (Operational Purity Proxy)}
\label{sec:DecoherenceLindblad}

\subsection{Open Quantum Systems and Unital Pure Dephasing}

For a \emph{unital} pure-dephasing Lindblad channel ($L = L^\dagger = \sqrt{\gamma}\,\sigma_z$), off-diagonal coherence and the operational purity proxy $R_{\mathrm{coh}}^{\mathrm{proxy}}$ decay monotonically (consistent with Result~\ref{thm:CoherenceMonotonicity}, which restricts universal purity non-increase to unital CPTP maps and distinguishes $R_{\mathrm{coh}}^{\mathrm{proxy}}$ from the classical EFT amplitude $\Rcoh = |\psi|/\psi_0$).

\paragraph{General Lindblad Equation:}
For a density matrix $\rho$ evolving under Hamiltonian $H$ and decoherence (environmental coupling) described by Lindblad operators $L_k$:
\[
\frac{d\rho}{dt} = -\frac{i}{\hbar}[H, \rho] + \sum_k \left( L_k \rho L_k^\dagger - \frac{1}{2}\{L_k^\dagger L_k, \rho\} \right)
\]

\paragraph{Two-Level System with Pure Dephasing:}
Consider a qubit ($|0\rangle$, $|1\rangle$) with Hamiltonian $H=0$ and a single Lindblad operator $L = \sqrt{\gamma}\, \sigma_z$, where $\gamma$ is the dephasing rate:
\[
\frac{d\rho}{dt} = \gamma \left( \sigma_z \rho \sigma_z - \rho \right)
\]
For initial state
\[
\rho(0) = \begin{pmatrix} p & c \\ c^* & 1-p \end{pmatrix}
\]
the evolution is
\[
\rho(t) = \begin{pmatrix} p & c\, e^{-2\gamma t} \\ c^*\, e^{-2\gamma t} & 1-p \end{pmatrix}
\]
The off-diagonal terms (coherence) decay exponentially with time constant $1/(2\gamma)$.

\paragraph{Decay of the Operational Purity Proxy $R_{\mathrm{coh}}^{\mathrm{proxy}}(t)$:}
Recall from Section~\ref{sec:PrimitiveFieldDefinitions} the finite-dimensional ($d=2$) operational purity proxy (strictly distinct from the classical EFT modulus $\Rcoh = |\psi|/\psi_0$):
\[
R_{\mathrm{coh}}^{\mathrm{proxy}}(t) = \frac{\operatorname{Tr}[\rho^2(t)] - 1/2}{1 - 1/2}.
\]
For the qubit under unital pure dephasing ($L = L^\dagger = \sqrt{\gamma}\,\sigma_z$):
\[
\operatorname{Tr}[\rho^2(t)] = p^2 + (1-p)^2 + 2|c|^2 e^{-4\gamma t},
\qquad
R_{\mathrm{coh}}^{\mathrm{proxy}}(t) = 2\left[ p^2 + (1-p)^2 + 2|c|^2 e^{-4\gamma t} - \frac{1}{2} \right].
\]
As $t\to\infty$, $e^{-4\gamma t}\to 0$ and $R_{\mathrm{coh}}^{\mathrm{proxy}}(t)$ decreases monotonically to $2[p^2 + (1-p)^2 - 1/2]$ (remaining constant if $c = 0$, illustrating that non-unitary dephasing preserves purity on diagonal states).

\paragraph{Mandatory Unitality Scope (Refuting General-Lindblad Purity Monotonicity):}
Monotonic non-increase of $R_{\mathrm{coh}}^{\mathrm{proxy}}(t)$ holds strictly for \emph{unital} CPTP/Lindblad channels ($\sum_k [L_k, L_k^\dagger] = 0$, Result~\ref{thm:CoherenceMonotonicity}). For \emph{non-unital} Lindblad generators such as amplitude damping ($L = \sqrt{\gamma}\,\sigma_-$, Kraus probability $p \in (0,1]$), acting on the maximally mixed state $\rho = I/2$ raises purity to $\mathrm{Tr}(\Phi(\rho)^2) = (1 + p^2)/2$ (for example, at $p = 1/2$, purity increases from $1/2$ to $5/8$). For general (non-unital) CPTP evolution, universal monotonicity is governed by quantum relative-entropy contraction $D(\Phi(\rho)\|\Phi(\omega)) \le D(\rho\|\omega)$ rather than purity decay, and no equivalence with classical closed-system Noether charge conservation ($\nabla_\mu J^\mu_{\mathrm{coh}} = 0$) is claimed.

\section{Computational Methods, Simulations, and Pseudocode}
\label{appendix:tools}

This appendix provides exploratory computational methodologies, pseudocode, and simulation sketches intended for prospective numerical modeling and future empirical testing of SIT hypotheses once microphysical source and detector mappings (\texttt{OPEN-6}) are derived.

\subsection*{Computational methodologies}

Predictive modelling uses three complementary numerical strategies:  

\paragraph{Finite-element integration.}
Equation set \eqref{eq:Einstein-like}--\eqref{eq:Rcoh-eom} is discretised on unstructured tetrahedral meshes with adaptive refinement at steep coherence gradients.  A Crank--Nicolson scheme advances $\rho_{t}$ and $R_{\text{coh}}$ implicitly, conserving the energy to better than $10^{-7}$ per timestep while capturing kilometric oscillations in interferometry test-beds.

\paragraph{Agent-based kinetic Monte Carlo.}
Microscopic ``agents'' carry local phase $\theta_{i}$ and perform collisional updates
$\theta_{i}\!\mapsto\!\theta_{i}+(\theta_{j}-\theta_{i})\exp(-\Delta/\tau_{\rm dec})$,
where $\tau_{\rm dec}$ derives from~\eqref{eq:dissipation}.  Averaging over $10^{9}$ agents reproduces the continuum limit and reveals self-organised synchrony waves analogous to cortical gamma bursts and spiral density waves in protoplanetary disks.

\paragraph{Quantum/classical hybrid circuits.}
High-coherence patches are evolved on NISQ hardware via Suzuki--Trotter factorisation of the $R_{\text{coh}}$ Hamiltonian, while low-coherence regions follow classical Langevin dynamics.  A reversible SWAP gate bridges the two domains every micro-iteration, ensuring phase continuity across the quantum--classical interface.

\subsection{Finite Element SIT Solver Pseudocode}

Finite element methods (FEM) are employed for solving SIT's coupled PDEs numerically. Pseudocode structure:

\begin{verbatim}
initialize_grid(domain_size, resolution)
initialize_fields(rho_t, R_coh, boundary_conditions)

while time < simulation_time:
    compute_gradients(R_coh, rho_t)
    update_fields_via_PDEs(R_coh, rho_t, dt)
    apply_boundary_conditions(fields)
    save_output(fields)
\end{verbatim}

\subsection{Monte Carlo Approaches}

Monte Carlo methods capture informational decoherence processes probabilistically:

\begin{verbatim}
for each simulation_step:
    sample_initial_conditions(R_coh_distribution)
    propagate_information_via_MonteCarlo(fields)
    evaluate_statistics_and_uncertainties(fields)
    record_results()
\end{verbatim}

\subsection{Hybrid Quantum-Classical Pathway Simulations}

Hybrid models integrate classical gravitational dynamics with quantum informational coherence:

\begin{verbatim}
initialize_classical_gravity_fields(g_mu_nu)
initialize_quantum_coherence_fields(R_coh)

for each timestep:
    evolve_quantum_coherence(R_coh)
    feed_coherence_back_to_classical(g_mu_nu, R_coh)
    update_classical_fields_via_GR(g_mu_nu)
    iterate_until_convergence()
\end{verbatim}

These computational approaches facilitate empirical testing and refinement of SIT predictions.

\section{Calibration Pipeline across Metrology and Neuroscience}
\label{sec:calibration-pipeline}
Coefficients of $\Phi_{\mathrm{teleo}}$ and small SIT couplings $(\gamma,\chi)$ are first constrained by optical-clock residuals (Appendix~A), then refined using neural coherence datasets that estimate $\langle R_{\mathrm{coh}}\rangle$ dynamics under volitional tasks, and finally propagated into the BEC benchmark (Appendix~\ref{appendix:E-BEC}) and the three-path prediction (Sec.~\ref{sec:three-path-born}) for parameter continuity across domains.

\section{Prospective Experimental Protocols and Schematic Figures}
\label{appendix:C}

This appendix outlines prospective experimental setups intended to test SIT once the missing action-to-observable derivations (\texttt{OPEN-6}) are completed.

\subsection{Prospective Experimental Protocols}

Candidate experimental venues include:

\begin{itemize}
    \item \textbf{Clock-Cavity Experiments:} High-precision atomic clocks in varied gravitational potentials to probe prospective coherence-time-density couplings once the visible $U(1)$ dictionary and transfer factor $\mathcal{T}$ are derived (\texttt{OPEN-6B}).
    \item \textbf{Aharonov--Bohm (AB) Loops:} Quantum interference setups motivated by bundle-holonomy analogies, conditional on resolving \texttt{OPEN-6B}.
    \item \textbf{SQUID and Cold-Atom Setups:} Ultra-cold atomic ensembles and SQUID magnetometers as prospective venues once $J_{\mathrm{prep}}$ and composite material charges are derived (\texttt{OPEN-6A}, \texttt{OPEN-6C}).
    \item \textbf{Neuroscience Paradigms:} EEG/MEG phase-synchrony paradigms designed to test proposed phenomenological neural-coherence models separately from spacetime EFT.
\end{itemize}

\noindent\textit{Cross-reference:} See Appendix~\ref{appendix:E-BEC} for the $^{87}$Rb BEC hypothetical sensitivity benchmark and scoped falsification protocol.

\subsection{Schematic Figures and Future Sensitivity Requirements}

Conceptual figures illustrate candidate diagnostic observables:

\begin{itemize}
    \item Figure C.1: Schematic phase-holonomy loops illustrating candidate holonomy observables.
    \item Figure C.2: Conceptual neural coherence gradient maps from EEG/MEG.
    \item Figure C.3: Schematic fringe patterns for candidate AB-loop setups.
\end{itemize}

Quantitative sensitivity assessments remain contingent on completing the missing action-to-apparatus derivations (\texttt{OPEN-6}) for each specified parameter state.

\section{Coherence--Gravity Equivalence Prediction for a \texorpdfstring{$^{87}$}{87}Rb BEC}
\label{appendix:E-BEC}
\label{app:BEC-benchmark}

\subsection{Objective}
We provide a numerically specified, falsifiable target for SIT's ``Coherence--Gravity Equivalence Test'' by comparing the free-fall acceleration of a Bose--Einstein condensate (BEC) to that of an otherwise identical thermal cloud.

\subsection{Signal model}
SIT predicts a coherence-dependent refinement of the Weak Equivalence Principle. Let $\alpha_g$ denote the dimensionless coherence--gravity coupling that modifies the ratio $M_g/M_i$. For two ensembles of identical atoms in the same external field $g_E$ but different coherence,
\[
\frac{\delta a}{g_E} \;\equiv\; \frac{a_{\rm BEC}-a_{\rm thermal}}{g_E}
\;=\; \alpha_g\Big(R_{\rm coh}^{\rm (BEC)} - R_{\rm coh}^{\rm (thermal)}\Big).
\]
For a near-pure $^{87}$Rb BEC we set $R_{\rm coh}^{\rm (BEC)}\!\approx\!1$, while the thermal cloud average is $R_{\rm coh}^{\rm (thermal)}\!\approx\!0$, hence
\[
\frac{\delta a}{g_E} \;=\; \alpha_g.
\]

\subsection{Identification with the clock-sector parameter (Hypothetical Sensitivity Benchmark)}
Appendix~\ref{appendix:ClockConstraints} defines the gravitational-potential clock-redshift slope
\(\alpha_\Phi \equiv d\ln\nu/d(\Phi/c^{2}) = \alpha_\rho\mathcal{T}\) (disambiguated from the local field response \(\alpha_\rho\)). While the published Tokyo Skytree optical-lattice clock test (Takamoto~\emph{et~al.}, 2020) reported \(\alpha_{\mathrm{RS}} = (+1.4 \pm 9.1)\times 10^{-5}\), earlier drafts (Versions~79--80.2) adopted an illustrative hypothetical target \(|\alpha_\Phi|_{\mathrm{hyp}} \le 3\times 10^{-8}\). Writing
\[
\alpha_g \;=\; \chi\,\alpha_\Phi,
\]
with a dimensionless transfer factor $\chi$ capturing geometry and sectoral weighting, setting $\chi = 1$ at the hypothetical clock sensitivity $|\alpha_\Phi|_{\mathrm{hyp}} \le 3\times 10^{-8}$ produces strictly a \emph{hypothetical future sensitivity benchmark} (NOT an empirically established ceiling from Tokyo Skytree; deriving $\alpha_g$ from unchanged \texttt{M\_AUTH\_SIT\_EFT\_v1} requires resolving \texttt{OPEN-6A} and \texttt{OPEN-6C}).

\subsection{Hypothetical sensitivity benchmark (NOT an empirical 95\% CL bound)}
With $g_E=9.80665~\mathrm{m\,s^{-2}}$ and $\chi=1$ at the quarantined hypothetical target $|\alpha_\Phi|_{\mathrm{hyp}} \le 3\times 10^{-8}$,
\[
\boxed{\;
\left|\frac{\delta a}{g_E}\right|_{\mathrm{hyp}}
= |\alpha_g|_{\mathrm{hyp}}
= |\alpha_\Phi|_{\mathrm{hyp}}
\le 3\times10^{-8}
\;}
\]
and
\[
\boxed{\;
|\delta a|_{\mathrm{hyp}}
\le 3\times10^{-8}\times 9.80665~\mathrm{m\,s^{-2}}
\approx 2.94\times10^{-7}~\mathrm{m\,s^{-2}}
\approx 29.4~\mu\mathrm{Gal}.
\;}
\]
Version~81 quarantines both $3\times 10^{-8}$ and $29.4\,\mu\mathrm{Gal}$ strictly as hypothetical future sensitivity targets rather than empirical constraints anchored to Tokyo Skytree.

\subsection{Computing $\chi$ and $\delta g$ for a concrete $^{87}$Rb geometry}
To turn this hypothetical sensitivity template into a genuine action-derived prediction (part of \texttt{OPEN-6}), one must compute the microscopic state-preparation source $J_{\mathrm{prep}}$ and composite atomic scalar charges for the specific trap and atom number:
(i) Model the condensate density $n(\mathbf{r})$ (e.g., Thomas--Fermi) and the thermal cloud $n_{\rm th}(\mathbf{r})$ at the experimental temperature.
(ii) Solve the linearized SIT field equations for $(R_{\mathrm{coh}},\rho_t)$ sourced by $n(\mathbf{r})$ and $n_{\rm th}(\mathbf{r})$ to obtain the coherence profiles and the induced $\delta\rho_t$.
(iii) Evaluate the free-fall response functional to extract $\alpha_g$ for each state and set $\chi \equiv \alpha_g/\alpha_\Phi$.
(iv) Insert the calibrated $\chi$ into
\[
\delta a \;=\; \chi\,\alpha_\Phi\,g_E,\qquad
\delta g \equiv \delta a,
\]
using a primary experimental bound on $\alpha_\Phi$ (or a direct action-derived prediction) and report $\delta g$ with geometry, $N$, temperature, and uncertainty.

\paragraph{Falsification protocol.}
Once $\alpha_g$ is derived from the action (\texttt{OPEN-6A}/\texttt{OPEN-6C}), a confirmed nonzero differential acceleration between a pure $^{87}\mathrm{Rb}$ BEC and an equal-mass thermal cloud would provide direct experimental evidence for coherence-dependent coupling, whereas null measurements exclude parameter regions predicting signals above experimental sensitivity; the legacy $29.4\,\mu\mathrm{Gal}$ ($3\times 10^{-8}$) reference line is retained solely as a labeled hypothetical sensitivity scale.

\section{Appendix: Constraints from Optical Clock Metrology (Disambiguated Notation and Provenance)}
\label{appendix:ClockConstraints}

\[
\alpha_\rho \;\text{(local field response)},
\qquad
\mathcal{T} \;\equiv\; \left(\frac{\rho_0^{-1}\,\partial\rho_t}{\partial(\Phi/c^2)}\right)_{\mathrm{lab}},
\qquad
\alpha_\Phi \;\equiv\; \frac{d\ln\nu}{d(\Phi/c^{2})} = \alpha_\rho\,\mathcal{T}.
\]
Version~81 strictly separates $\alpha_\rho$ (response to local field shift $\delta\rho_t/\rho_0$) from $\alpha_\Phi = \alpha_\rho\mathcal{T}$ (response to gravitational-potential shift $\Delta\Phi/c^2$), resolving the Version~79/80.2 notation collision that reused $\alpha_{\mathrm{eff}}$ for both and thereby double-counted $\mathcal{T}$ whenever $\mathcal{T} \ne 1$.

\subsection{Objective}
We clarify the mapping between local time-density perturbations, gravitational-potential redshift residuals, primary optical-clock experimental measurements, and the quarantined hypothetical sensitivity benchmark used in earlier drafts.

\subsection{Mapping SIT to the clock observable}
Conditional on specifying the visible electromagnetic $U(1)$ photon dictionary (\texttt{OPEN-6B}), a local time-density modulation produces a fractional frequency response governed strictly by $\alpha_\rho$:
\[
\frac{\Delta \nu}{\nu}(x) = \alpha_\rho\,\frac{\delta \rho_t(x)}{\rho_0}.
\]
For clocks at locations with differing gravitational potential \(\Phi\), GR gives \((\Delta \nu/\nu)_{\mathrm{GR}} \simeq \Delta \Phi / c^{2}\). Slow spatial variations of \(\rho_t\) induced by \(\Phi\) produce an additional response operationally captured by
\[
\frac{d\ln \nu}{d(\Phi / c^{2})} \equiv \alpha_\Phi .
\]
We use \(\alpha_\Phi\) strictly as the directly measured slope in differential clock comparisons (same species, different potentials) after subtracting the GR redshift.

\subsection{Linearized relation and identification (No Double Counting of $\mathcal{T}$)}
To first order about the laboratory background,
\[
f_2(\rho_t) = 1 + \alpha_\rho \frac{\delta \rho_t}{\rho_0} + \mathcal{O}(\delta \rho_t^{2}), \qquad
\frac{\delta \rho_t}{\rho_0} = \mathcal{T}\,\frac{\Delta \Phi}{c^2} + \cdots
\]
so substituting $\delta\rho_t/\rho_0$ into the local response $\Delta\nu/\nu = \alpha_\rho\,(\delta\rho_t/\rho_0)$ yields
\[
\left(\frac{\Delta\nu}{\nu}\right)_{\mathrm{SIT}}
= \alpha_\rho\,\frac{\delta\rho_t}{\rho_0}
= \alpha_\rho\,\mathcal{T}\,\frac{\Delta\Phi}{c^2}
= \alpha_\Phi\,\frac{\Delta\Phi}{c^2},
\qquad
\alpha_\Phi = \alpha_\rho\,\mathcal{T}.
\]
Crucially, one must \emph{never} multiply $\delta\rho_t/\rho_0$ by $\alpha_\Phi$ for a potential-driven signal, as doing so would yield $\alpha_\rho\mathcal{T}^2(\Delta\Phi/c^2)$ and double-count $\mathcal{T}$ unless $\mathcal{T} = 1$.

\subsection{Primary experimental measurements vs.\ quarantined hypothetical benchmark}
The published Tokyo Skytree transportable $^{87}\mathrm{Sr}$ optical-lattice clock test over a $450\,\mathrm{m}$ elevation difference ($\Delta\Phi/c^2 = 4.9101\times 10^{-14}$; Takamoto~\emph{et~al.}, \emph{Nature Photonics} \textbf{14}, 411--415, 2020) reported
\[
\alpha_{\mathrm{RS}} = (+1.4 \pm 9.1) \times 10^{-5} \quad (1\sigma\text{ combined}),
\]
while NIST/BACON (\emph{Phys.\ Rev.\ Lett.}\ \textbf{137}, 033201, 2026; DOI \texttt{10.1103/g865-9mk1}) achieved optical frequency-ratio uncertainties at or below $3.2\times 10^{-18}$ across $\mathrm{Al}^+/\mathrm{Yb}/\mathrm{Sr}$ clocks linked via phase-stabilized optical fiber networks. By contrast, the legacy value
\[
|\alpha_\Phi|_{\mathrm{hyp}} \lesssim 3 \times 10^{-8} \quad \text{(quarantined hypothetical future sensitivity target; NOT a 95\% CL Skytree bound)}
\]
would correspond to $1.47\times 10^{-21}$ fractional accuracy on the Skytree baseline and is therefore quarantined strictly as a hypothetical sensitivity target in Version~81.

\subsection{Usage in the main text}
Throughout Version~81 (Sections~\ref{sec:map-to-observables}, \ref{sec:atomic-clock-prediction}, \ref{sec:empirical-content}, and \ref{sec:stage6-action-diagnostics}), local field-driven shifts are written strictly with $\alpha_\rho$ and gravitational-potential shifts with $\alpha_\Phi = \alpha_\rho\mathcal{T}$. Converting primary optical-clock bounds into unconditional constraints on \texttt{M\_AUTH\_SIT\_EFT\_v1} remains explicitly blocked under \texttt{OPEN-6B} until the visible $U(1)$ photon dictionary and terrestrial profile $\mathcal{T}$ are derived from the action.

\subsection{Forward look: network clocks and gradiometry}
Clock networks and entangled-state gradiometers can probe spatial derivatives of \(\rho_t\) at improved sensitivity, tightening bounds on \(\alpha_\Phi\) and enabling targeted tests once \texttt{OPEN-6B} is resolved.

\section{Three-Path Interference and Exploratory Sorkin-Term Hypothesis (Open)}
\label{sec:three-path-born}
Consistent with Section~\ref{sec:SIT3} (which retracts the impossible linear operator $\hat{\mathcal{R}}_{\mathrm{coh}}$), any exploratory SIT~3.0 proper-time ansatz must be formulated using a declared \emph{nonlinear} state functional $\mathcal{C}_{\mathrm{state}}[\rho]$ (or restricted linear witness), $d\tau/dt = \rho_g(x)\,\exp[\gamma\,\mathcal{C}_{\mathrm{state}}[\rho]]$, together with an explicit composite-system evolution and measurement prescription ensuring no-signaling. Whether a well-posed nonlinear or open-system extension induces a nonzero Sorkin third-order interference term $I_{123}$ remains an open phenomenological question contingent on specifying that dynamical prescription and deriving the coupling (\texttt{OPEN-6C}), rather than a proved formula of \texttt{M\_AUTH\_SIT\_EFT\_v1}.
% --- END INPUT sections/10_appendices_A_to_M.tex ---
% --- BEGIN INPUT sections/11_appendices_N_to_AA.tex ---
\section{Exploratory Recursion Weight as a Nonnegative Attenuation Cost}
\label{sec:recursion-influence}
\label{app:recursion-influence}
Let $X$ be a real Banach space (with canonical complexification $X_{\mathbb{C}}$) of coarse-grained reduced states $\Sigma_t$ over a compact time interval $t \in [t_i, t_f]$ along a continuous trajectory $p : [t_i, t_f] \to X$, and let $\mathcal{R} : X \to X$ be a Fr\'echet-$C^1$ discrete return map ($\Sigma_{t+\Delta}=\mathcal{R}(\Sigma_t)$) whose complexified Fr\'echet derivative $D\mathcal{R}|_{\Sigma_t} \in \mathcal{B}(X_{\mathbb{C}})$ is a bounded linear operator with Gelfand spectral radius $\rho\bigl[D\mathcal{R}|_{\Sigma_t}\bigr] = \lim_{n\to\infty}\|(D\mathcal{R}|_{\Sigma_t})^n\|^{1/n} \in [0, \|D\mathcal{R}|_{\Sigma_t}\|] < \infty$. Earlier drafts defined an unclipped recursion cost $C_{\mathrm{rec}}[p] = \sup_t \log(\rho[D\mathcal{R}|_{\Sigma_t}])$ and inserted $\exp(-\lambda C_{\mathrm{rec}}[p])$ ($\lambda > 0$). Version~81 repairs both the contraction sign error and the zero-Jacobian singularity in that definition: for any contractive ``clean recursion'' ($0 \le \rho[D\mathcal{R}] \le k < 1$), an unclipped logarithm is negative (making $\exp(-\lambda\log k) = k^{-\lambda} > 1$ amplify rather than attenuate) and diverges at $\rho[D\mathcal{R}] = 0$ ($\log 0 = -\infty$). Defining the nonnegative clipped cost as
\begin{equation}
C_{\mathrm{rec}}^+[p] \;:=\; \sup_{t \in [t_i, t_f]}\,\log\max\!\Bigl(1,\;\rho\bigl[D\mathcal{R}\big|_{\Sigma_t}\bigr]\Bigr) \;\ge\; 0
\label{eq:nonneg_recursion_cost}
\end{equation}
is finite and well-defined for all $\rho\bigl[D\mathcal{R}|_{\Sigma_t}\bigr] \ge 0$ (including $\rho = 0$, where $\log\max(1, 0) = \log 1 = 0$), assigns zero penalty ($C_{\mathrm{rec}}^+ = 0$, weight $1$) to every contractive map ($\rho \le k < 1$), and assigns strictly bounded attenuation $\exp(-\lambda C_{\mathrm{rec}}^+[p]) \in (0, 1]$ to unstable expanding steps ($\rho > 1$):
\[
\mathcal{W}[p] \;=\; \exp\!\left\{\frac{i}{\hbar}\!\int_{t_i}^{t_f}\!dt\,[\mathcal{L}_{\mathrm{SIT}}-\Phi_{\mathrm{teleo}}]\right\}\,
\exp\!\left\{-\lambda\,C_{\mathrm{rec}}^+[p]\right\}.
\]
This single-history effective weight is retained strictly as an exploratory phenomenological ansatz (separate from unchanged \texttt{M\_AUTH\_SIT\_EFT\_v1}); closed-system evolution remains unitary.

\section{Feynman--Vernon Influence-Functional Analogy (Two-History Normalization and Heuristic Scope)}
\label{appendix:FV-recursion}
In open quantum mechanics and Schwinger--Keldysh field theory, integrating out an initially normalized environment state $\rho_{\mathrm{env}}(t_i) \succeq 0$ ($\mathrm{Tr}\,\rho_{\mathrm{env}}(t_i) = 1$) coupled to forward and backward system histories $(p^+, p^-)$ defines the \emph{two-history} Feynman--Vernon influence functional
\[
\mathcal{F}[p^+, p^-] \;:=\; \mathrm{Tr}_{\mathrm{env}}\!\Bigl(U_{\mathrm{env}}[p^+]\,\rho_{\mathrm{env}}(t_i)\,U_{\mathrm{env}}^\dagger[p^-]\Bigr)
\;=\; \exp\!\left\{\frac{i}{\hbar}\Phi_{\mathrm{LS}}[p^+, p^-] \;-\; \Gamma[p^+, p^-]\right\}.
\]
By unitary bath evolution and the Cauchy--Schwarz inequality on $\rho_{\mathrm{env}}(t_i)^{1/2}$, this normalized two-history functional rigorously satisfies $|\mathcal{F}[p^+, p^-]| \le 1$ ($\Gamma[p^+, p^-] \ge 0$) and $\mathcal{F}[p, p] = 1$ ($\Gamma[p, p] = 0$ on identical forward/backward histories), whereas an unnormalized single-history path integral $\int\mathcal{D}q\,e^{(i/\hbar)(S_{\mathrm{env}}+S_{\mathrm{int}})}$ does not by itself obey $|\mathcal{F}| \le 1$. Furthermore, for a linearly coupled Gaussian bath ($S_{\mathrm{int}} = \int dt\,\eta(t)\,G[\Sigma_t]$), exact evaluation yields a \emph{quadratic nonnegative off-diagonal double-time integral} $\Gamma[p^+, p^-] = \frac{1}{2}\iint dt\,dt'\,\Delta G(t)\,\nu(t-t')\,\Delta G(t') \ge 0$ ($\Delta G := G[\Sigma_t^+] - G[\Sigma_t^-]$), not a single-history supremum of a logarithmic Jacobian spectral radius. Writing a single-history damping factor $\exp(-\lambda\,C_{\mathrm{rec}}^+[p])$ via Eq.~\eqref{eq:nonneg_recursion_cost} is therefore an abbreviated phenomenological coarse-graining ansatz motivated by $|\mathcal{F}[p^+, p^-]| \le 1$, rather than an exact Feynman--Vernon path-integral identity.

\section{Existence and Uniqueness of Contractive Fixed Points}
\label{sec:fixed-point}
\label{app:clean-recursion}
Let $X$ be a Banach space of reduced states $\Sigma$ and $\mathcal{R}:X\to X$ a Fr\'echet-$C^1$ coarse-grained return map. If there exists $k\in(0,1)$ with
\[
\|\mathcal{R}(x) - \mathcal{R}(y)\|_X \;\le\; k\,\|x - y\|_X \quad\forall x,y\in X,
\]
then by Banach's fixed-point theorem $\mathcal{R}$ admits a unique fixed point $\Sigma^\star$, with operator norm $\|D\mathcal{R}|_{\Sigma^\star}\| \le k < 1$, spectral radius $\rho[D\mathcal{R}|_{\Sigma^\star}] \le k < 1$, and zero clipped attenuation cost $C_{\mathrm{rec}}^+[\Sigma^\star] = \log\max(1, \rho[D\mathcal{R}|_{\Sigma^\star}]) = 0$ under Eq.~\eqref{eq:nonneg_recursion_cost}.

\section{Driven-Dissipative Topological and Cyclic Attractors (Toy Ring Model)}
\label{appendix:topological-attractors}

We call an \textbf{informational limit cycle} any periodic trajectory $\mathcal{C}$ in a driven-dissipative toy state space $(R_{\mathrm{coh}}, \rho_t, \theta)$ for which $R_{\mathrm{coh}}(t+T) \approx R_{\mathrm{coh}}(t)$, $\rho_t(t+T) \approx \rho_t(t)$, and $\theta(t+T) = \theta(t) + 2\pi n$. Consistent with Section~\ref{sec:arrow-of-time} and Section~\ref{sec:coherence-conservation}, Version~81 explicitly retracts earlier wording asserting an undefined global coherence functional $\mathcal{G}$ satisfying $d\mathcal{G}/dt \le 0$ for all trajectories: in open driven-dissipative oscillator toys, local periodic cycles are sustained by external driving against local damping, while universal open-system monotonicity in finite-dimensional quantum channels holds strictly for unital-CPTP purity and general-CPTP relative entropy (Result~\ref{thm:CoherenceMonotonicity}), distinct from exact closed-system $U(1)$ Noether charge conservation ($\nabla_\mu J^\mu_{\mathrm{coh}} = 0$). On a ring loop $\Gamma$, a nontrivial phase winding has integer topological index $\nu = (1/2\pi)\oint_\Gamma d\theta \in \mathbb{Z}$, and any circulation-versus-dissipation inequality holds only for an explicitly specified dissipative energy balance rather than as a universal theorem of \texttt{M\_AUTH\_SIT\_EFT\_v1}.

\paragraph{Minimal driven-dissipative toy model.} On a ring geometry parameterized by $x \in [0,L)$, write
\begin{align*}
    \partial_t R_{\mathrm{coh}} &= D \partial_x^2 R_{\mathrm{coh}} - \lambda R_{\mathrm{coh}} + \eta \cos(\Omega t) + \chi \sin(\partial_x\theta), \\
    \partial_t \theta &= \omega_0 + \alpha \rho_t + \beta \partial_x^2 \theta, \\
    \partial_t \rho_t &= c^2 \partial_x^2 \rho_t - \mu \rho_t + \sigma R_{\mathrm{coh}},
\end{align*}
with $D,\lambda,\eta,\Omega,\chi,\omega_0,\alpha,\beta,c,\mu,\sigma > 0$. This driven-dissipative PDE system is a phenomenological oscillator toy model (not derived from the reversible canonical action \texttt{M\_AUTH\_SIT\_EFT\_v1}); for drive $\eta$ above threshold and moderate $\chi$, it illustrates how periodic forcing sustains phase winding $\nu$ with loop integral $H = \oint_0^L \partial_x\theta \, dx = 2\pi\nu$ without invoking retrocausality.

\paragraph{Scope.} This appendix does not modify \texttt{M\_AUTH\_SIT\_EFT\_v1}; it documents an exploratory driven-dissipative toy system separate from the conservative relativistic field equations.

\section{Phase Slips and Topological Limit Cycles}

We model a driven local loop as a complex order parameter
\[
\psi(x,t)=A(x,t)\,e^{i\theta(x,t)}
\]
on a periodic ring of length \(L\), evolving under a damped, driven, gauge--covariant TDGL equation
\[
\partial_t \psi
= D(\partial_x - iA)^2 \psi + \big(\mu - g|\psi|^2\big)\psi
+ i\omega_0 \psi + \eta e^{i\Omega t} + \xi(x,t),
\]
with \(A(t)=\Phi(t)/L\) a uniform gauge potential set by a slowly ramped flux \(\Phi(t)\), and \(\xi\) mean-zero complex noise.

The loop's integer winding is
\[
\nu(t)=\frac{1}{2\pi}\oint_0^L \partial_x \theta\,dx,
\]
which is topologically conserved unless the amplitude vanishes at a slip site, \(A(x_s,t_s)\to 0\). At such events \(\theta\) can jump by \(\pm 2\pi\), changing \(\nu\) without retrocausality. The gauge--covariant current is
\[
j(x,t)=\Im\!\Big(\psi^*(\partial_x - iA)\psi\Big), \qquad
\mathcal{C}(t)=\oint j\,dx,
\]
and exhibits discrete jumps synchronized with \(\nu\) transitions. In this driven, damped TDGL toy model, local limit cycles are sustained by external flux ramping and forcing while dissipation damps autonomous relaxation in the undriven limit (distinct from any universal global monotonicity law for reversible canonical fields, Section~\ref{sec:arrow-of-time}).

Numerics with a linear \(\Phi(t)\) ramp show quantized steps in \(\nu\), dips in \(\min_x|\psi|\) preceding each step, and concurrent jumps in \(\mathcal{C}(t)\) (Figs. U1--U3). Final phase and amplitude profiles illustrate slip locations and total phase \(2\pi\nu\) (Figs. U4--U5).

Minimal pseudocode for reproducibility:

\begin{verbatim}
Initialize grid x in [0, L), time step dt, total steps.
Set parameters D, \mu, g, \omega_0, \eta, \Omega, noise_amp.
Initialize \psi(x,0) = small-amplitude complex noise.

For each step n:
    t \leftarrow n\cdot dt
    A \leftarrow \Phi(t)/L   (with \Phi ramped from 0 \to 2\pi)
    Compute \psi_x via centered difference; \psi_xx likewise.
    cov_lap \leftarrow \psi_xx - 2iA \psi_x - A^2 \psi
    nonlin \leftarrow (\mu - g|\psi|^2) \psi + i\omega_0 \psi
    drive \leftarrow \eta exp(i\Omega t)
    noise \leftarrow complex Gaussian, var \propto noise_amp/sqrt(dt)
    \psi \leftarrow \psi + dt*[D*cov_lap + nonlin + drive] + noise*dt
    \theta_unwrap \leftarrow unwrap(arg \psi along x)
    \nu(t) \leftarrow round((\theta_unwrap(L) - \theta_unwrap(0))/2\pi)
    amin(t) \leftarrow min_x |\psi|
    j \leftarrow Im( \psi^* (\psi_x - iA \psi) )
    C(t) \leftarrow sum j\cdot dx
Save \nu(t), amin(t), C(t), and snapshots of \theta_unwrap, |\psi|.
\end{verbatim}

Here, $C(t) = \sum j\, dx$ is the integrated current, $\nu(t)$ is the winding number, and $amin(t)$ is the minimum amplitude.

The generated figures, labeled U1--U5, are available as follows:
\begin{itemize}
    \item \textbf{Fig. U1:} Winding number \(\nu(t)\) from \url{fig1_winding_vs_time.png}.
    \item \textbf{Fig. U2:} Minimum amplitude \(\min|\psi|(t)\) from \url{fig2_min_amplitude.png}.
    \item \textbf{Fig. U3:} Circulation vs. flux from \url{fig3_circulation_vs_flux.png}.
    \item \textbf{Fig. U4:} Final phase from \url{fig4_final_phase.png}.
    \item \textbf{Fig. U5:} Final amplitude from \url{fig5_final_amplitude.png}.
\end{itemize}
The raw simulation data are available in the file \url{phase_slip_winding_sim.npz}.

In plain terms, this toy simulation illustrates how a driven damped 1D ring can execute discrete phase-slip resets when $|\psi|$ dips transiently to zero and the winding number jumps by $\pm 1$. This is a standard driven-TDGL phenomenological model (analogous to superconducting or superfluid rings) rather than a derivation from \texttt{M\_AUTH\_SIT\_EFT\_v1}.

\section{Integration with Prior Work and Conceptual Lineage}
\label{app:prior_work}

\subsection{Evolution from \emph{Super Dark Time} and Related Work}

SIT draws conceptual lineage from \emph{Super Dark Time}, \emph{SuperTimePosition}, and Micah's New Law of Thermodynamics, while separating their respective mathematical models cleanly (Table~\ref{tab:july-v80-vs-nov-v81}). Whereas earlier drafts tried to derive magnetism directly from a phase gradient, the canonical action \texttt{M\_AUTH\_SIT\_EFT\_v1} carries an independent $U(1)$ gauge connection $A_\mu$ with curvature $F_{\mu\nu} = \partial_\mu A_\nu - \partial_\nu A_\mu$ (keeping the visible electromagnetic dictionary open under \texttt{OPEN-6B}).


\subsection*{Teleonomic Hysteresis and a Geometric Phase}
\label{sec:geom-phase}
In a separately modeled dissipative control loop, path-dependent recursion costs motivate exploring an effective geometric phase shift $\Phi_{\mathrm{geom}} \approx \kappa_{\mathrm{geo}}\oint C_{\mathrm{rec}}[p]\,dt$; deriving any observable interferometric signature from \texttt{M\_AUTH\_SIT\_EFT\_v1} remains open under \texttt{OPEN-6C}.

\section{Exploratory Cosmological Sketches and Quantum-Gravity Analogies}
\label{appendix:D}

This appendix preserves exploratory Level-1 cosmological sketches and high-level crosswalk analogies with Loop Quantum Gravity (LQG), causal sets, and string theory. None of the schematic formulas below is derived from unchanged \texttt{M\_AUTH\_SIT\_EFT\_v1} (whose cosmological perturbation theory and UV completion remain strictly open under \texttt{OPEN-6D} and \texttt{OPEN-1}).

\subsection{Exploratory Cosmological Parameterizations}

As phenomenological Level-1 templates (distinct from derived predictions and requiring explicit dimensional transfer functions and mass-calibration checks, Section~\ref{sec:black-holes-voids}), earlier drafts sketched modified lensing ($\Delta\theta_{\mathrm{lens}} \sim \frac{4GM}{c^2 R}[1 + \tilde{\beta}\,\delta\rho_t/\rho_0]$), an effective coherence vacuum-energy term, and an exploratory coupling ansatz $G_{\mathrm{eff}}(\Rcoh) \sim G(1 + \gamma\Rcoh^2)$. Consistent with Section~\ref{sec:darkmatter}, linearizing a flat Friedmann model with an additive coherence density yields $\delta H_0/H_0 \approx \alpha_{\mathrm{cos}}\,\delta\langle\Rcoh\rangle/(2\rho_{\mathrm{tot}})$ (not $6\rho_{\mathrm{tot}}$), and whether any specified SIT cosmological extension is viable against CMB, BAO, and lensing data remains open (\texttt{OPEN-6D}).

\subsection{Connections to LQG, Causal Sets, and String Theory}

We clarify crosswalks to existing quantum gravity theories:

\textbf{Loop Quantum Gravity (LQG):} Spin networks and quantum geometries naturally correspond to quantized informational coherence networks within SIT.

\textbf{Causal Sets:} SIT's time-density fields align directly with causal structure discretization, offering rigorous information-based interpretations of causal set elements.

\textbf{String Theory:} String excitations reflect quantized coherence vibrations, positioning SIT coherence states as possible foundational structures underpinning string dynamics.

\section{Speculative, Outreach, and Metaphorical Extensions}
\label{app:speculative}

This section preserves conceptual metaphors, analogies, and speculative insights enriching broader communication and outreach.

\subsection{Metaphors and Analogies}

SIT conceptual metaphors:

\begin{itemize}
    \item \textbf{Cloth-Twist Spiral:} Visualizing informational coherence distortion analogous to fabric torsion.
    \item \textbf{Electron Pebble:} Quantum state perturbations analogized as ripples from a pebble.
    \item \textbf{Quasicrystal Brain:} Neural phase states mapped onto quasicrystal patterns as high-dimensional embedding analogs.
    \item \textbf{Halfway Universe:} Universal informational equilibrium maintaining overall zero-energy balance.
\end{itemize}

\subsubsection{Visual Outreach Metaphors for Informational Torque}
To help convey the idea of informational torque (Section~\ref{sec:informational_torque}) to non-specialists, three images are useful:

\begin{enumerate}[label=(\alph*)]
\item \textbf{Twisting Fabric.}  
      Aligned phase waves ``wring'' spacetime like cloth; the tighter the twist, the deeper the curvature.

\item \textbf{Choreographed Dancers.}  
      Synchronised rotation pulls the dancers inward; loss of rhythm lets the circle slacken. Coherence does the same for geodesics.

\item \textbf{Water Vortex.}  
      Co-phased surface ripples amplify into a whirlpool; decoherence breaks the vortex apart.
\end{enumerate}

\subsection{Speculative Extensions}

Extended speculative ideas preserved for outreach or future theoretical exploration:

\begin{itemize}
    \item Continuous Cosmic Microwave Background (CMB) coherence interpretations.
    \item Quantum-biological coherence applications in neuroscience and evolution.
    \item Informational attractors linking consciousness, quantum coherence, and gravitational effects.
\end{itemize}

\subsection{Outreach Graphics and Storyboards}

Storyboard and graphics resources developed for broader dissemination:

\begin{itemize}
    \item Infographics illustrating SIT's coherence-based reality.
    \item Animations visualizing SIT coherence dynamics from quantum to cosmic scales.
    \item Public-friendly narratives illustrating SIT philosophical implications (e.g., reality as active information rather than passive substance).
\end{itemize}

This appendix allows SIT's rich conceptual metaphors and speculative possibilities to engage public imagination and interdisciplinary dialogue without diluting the rigorous empirical core of the theory.



\section{Super Information Theory (SIT) Is Constructed Atop a Physical Mechanism---SuperTimePosition}
\label{app:stp-bell-mechanism}

SuperTimePosition (STP) supplies the physical mechanism that renders a local, deterministic, coherence-based field ontology compatible with quantum correlations inside Super Information Theory. Here we state (i) the minimal physical requirements of STP, (ii) the exact manner in which STP interfaces with Bell's theorem, and (iii) its falsifiable experimental signatures.

\subsection{Local Determinism and What Bell Actually Constrains}

Bell's theorem constrains hidden-variable completions under the joint conjunction of two distinct hypotheses:
\begin{enumerate}[leftmargin=2em]
  \item \textbf{Local Response (Parameter Locality):} Outcomes at wing $A$ and wing $B$ depend solely on the local analyzer setting and local hidden state, $A = A(a, \lambda_{\mathrm{loc}})$ and $B = B(b, \lambda_{\mathrm{loc}})$, with zero superluminal causal influence.
  \item \textbf{Measurement Independence (Statistical Setting Independence):} The hidden-state distribution at emission is statistically independent of the analyzer configurations, $\rho(\lambda \mid a, b) = \rho(\lambda)$.
\end{enumerate}

STP introduces a locally real, ultrafast cycling phase variable $\lambda(t)$ whose coarse-grained residue $\lambda \bmod 2\pi$ is undersampled by slower macroscopic detectors, explaining why individual quantum outcomes appear probabilistic despite deterministic phase-gate dynamics. Crucially, however, \emph{determinism alone does not evade Bell}: retaining both local response and measurement independence cannot reproduce the full $-\cos\theta$ singlet correlation curve.

In the SIT+STP framework, singlet correlations arise by physically relaxing \emph{measurement independence} via causal common-past field consistency: detector analyzer configurations are physical matter/electromagnetic distributions whose world-tubes inhabit a shared backward light cone with the emission source, sourcing retarded perturbations in $(\rho_t, \psi)$ across their common causal past so that
\begin{equation}
\rho(\lambda \mid a, b) \ne \rho(\lambda).
\label{eq:bell_common_past_relaxation}
\end{equation}
Strictly retarded hyperbolic propagation ($c_{\mathrm{char}} = 1$) and local outcome response are preserved throughout, with zero superluminal signaling and zero retrocausality.

\subsection{Local Entanglement via Phase-Locking and Operational No-Signaling}

Entangled pairs emerge from common emission events with phase-locked internal cycles, while shared past field history biases $\rho(\lambda \mid a, b)$. Crucially, \emph{operational no-signaling} is enforced by rotational and $\pi$-inversion symmetries of the stationary phase distribution: integrating over the remote outcome leaves single-wing marginals strictly independent of the remote setting, $P(A \mid a, b) = P(A \mid a) = 1/2$ and $P(B \mid a, b) = P(B \mid b) = 1/2$, while joint expectations reproduce $E(a,b) = -\cos(\theta_a - \theta_b)$ in the strong-locking limit.

\subsection{Unified Wave Mechanics from Quantum to Cosmic Scales}

The same dynamical primitives---phase oscillation, synchronization, and relaxation under gradients of $\rho_t$ and $\Rcoh$---operate from microscopic interferometers to collective condensed-matter and astrophysical structures, controlled by coupling scales and environmental decoherence rates.

\subsection{Testable and Falsifiable Predictions of the STP Mechanism}

\begin{itemize}
  \item \textbf{Attosecond and Zeptosecond Phase Substructure:} Probing state transitions at sampling bandwidths approaching internal phase-cycle timescales predicts structured sub-cycle beat modulations rather than featureless white randomness.
  \item \textbf{Gravitational / Time-Density Visibility Tuning:} Entanglement visibility across baselines spanning distinct gravitational potentials is predicted to exhibit small, systematic $\rho_t$-dependent phase-steering shifts.
  \item \textbf{Finite-Locking Visibility Ceiling:} Away from asymptotic strong phase-locking, the steering-to-noise ratio simultaneously fixes both deviations from $-\cos\theta$ and the maximal achievable CHSH parameter $S \le 2\sqrt{2}$.
\end{itemize}

\subsection{Elegance and Parity with Alternative Quantum Formalisms}

Unlike Bohmian mechanics (requiring nonlocal configuration-space pilot waves) or Everettian many-worlds (requiring unobservable ontological branch proliferation), SIT+STP posits only one four-dimensional Lorentzian spacetime carrying $(\psi, \rho_t, g_{\mu\nu}, A_\mu)$ and treats analyzer hardware states physically rather than exempting them from causal field equations.

\subsection{Determinism (Outcome) vs.\ Bell-Evasion (Common-Past Bias)}

Individual detection clicks are determined locally at synchronization ($A = A(a,\lambda_{\mathrm{loc}})$), whereas violations of CHSH bounds up to the Tsirelson ceiling $2\sqrt{2}$ arise exclusively from retarded common-past field correlation~\eqref{eq:bell_common_past_relaxation}.

\subsection{Fast-Switching Bell Tests and the Physical Meaning of ``Setting''}

In loophole-reduced fast-switching Bell experiments, an analyzer ``setting'' $a$ comprises the physical world-tube of the random-number generator, switching optoelectronics, and analyzer crystal. Because those physical degrees of freedom evolve deterministically from the overlap of the backward light cones of wings $A$, $B$, and the source, strictly retarded field propagation through the shared causal past remains causally sufficient without invoking advanced waves, provided local marginals respect operational no-signaling.

\subsection{Integration into the Unified Variational Framework}

SuperTimePosition motivates the exact field content of SIT: $\Rcoh(x)$ measures local phase alignment, $\rho_t(x)$ measures local cycle density, and $S_{\mathrm{AUTH}}$ is the minimal ghost-free covariant EFT governing their coupled evolution.

\section{From Foundational Principles to a Derivational Framework}
\label{app:derivational-framework}

Super Information Theory consolidates four precursor conceptual strands into a single variational framework grounded in $S_{\mathrm{AUTH}}$, distinguishing heuristic correspondences from controlled action-level derivations:
\begin{enumerate}[leftmargin=2em]
  \item \textbf{Thermodynamic Engine (\emph{Micah's New Law} $\to$ Bounded Open-System Lindblad/Relative-Entropy Flow):} Sequential wave-based difference dissipation is aligned strictly with Result~\ref{thm:CoherenceMonotonicity} (Section~\ref{sec:arrow-of-time}): monotonic subsystem-purity non-increase ($d\mathrm{Tr}(\rho^2)/dt \le 0$) holds strictly under \emph{unital} CPTP/Lindblad dynamics ($\Phi(I) = I$, $\sum_k \gamma_k [L_k, L_k^\dagger] = 0$; refuted for general non-unital CPTP maps by amplitude damping taking $I/2$ to purity $(1+p^2)/2 = 5/8$ at $p = 1/2$), whereas \emph{general} CPTP monotonicity is governed by quantum relative-entropy contraction $D(\Phi(\rho)\|\Phi(\omega)) \le D(\rho\|\omega)$, kept distinct from classical closed-system $U(1)$ Noether charge conservation ($\nabla_\mu J^\mu_{\mathrm{coh}} = 0$).
  \item \textbf{Quantum--Gravitational Link (\emph{Super Dark Time} $\to$ Dynamical $\rho_t$ and Portal $U_{\mathrm{link}}$):} The ``mass as localized time density'' concept becomes the covariant scalar $\rho_t(x)$ coupled to $|\psi|^2$ via $U_{\mathrm{link}}$ and screened via $f_{\mathrm{SM}}(\Rcoh)$ (Sections~\ref{sec:framework}--\ref{sec:screening-mechanism}).
  \item \textbf{Deterministic Phase Cycles (\emph{SuperTimePosition} $\to$ Covariant Complex Order Parameter $\psi$):} Rapid phase oscillations and Aharonov--Bohm holonomy motivate the local $U(1)$ bundle geometry of $\psi(x) = |\psi(x)|e^{i\theta(x)}$ coupled to an independent connection $A_\mu$ (Sections~\ref{sec:magnetism-holonomy} and~\ref{sec:stage4-matter-holonomy}).
  \item \textbf{Open-System Feedback (\emph{Self-Aware Networks} $\to$ Exploratory Nonnegative Recursion Cost):} Coarse-graining cyclic environments motivates the heuristic nonnegative clipped attenuation weight $\exp(-\lambda C_{\mathrm{rec}}^+[p])$ ($C_{\mathrm{rec}}^+ \ge 0$, Appendices~\ref{app:recursion-influence}--\ref{app:clean-recursion}) without altering closed-system unitarity.
\end{enumerate}

\section{Bounded Comparison with Causal Fermion Systems and Open Formalization Program}
\label{app:chronology}

Consistent with Paper~048 (\emph{Causal Fermion Systems and Super Information Theory}, July~28, 2026), Version~81 replaces earlier assertions of ``independent validation,'' ``semantic embedding,'' ``formal operational equivalence,'' and ``greater expressive power'' with a strict separation between: (1)~dated narrative-level conceptual parallels, (2)~Paper~048's bounded effective-equation comparison against an ordinary Einstein--Dirac null control, and (3)~an explicitly \emph{unproved} proposed categorical and process-calculus formalization program.

\subsection{Bounded Chronological and Effective-Equation Status (Paper~048)}

Public precursors in the SIT archive (2017 Neural Lace Podcast and NAPOT notes; 2022 QGTCD notes) introduced relational ``coincidence-pattern'' and distributed internal-observation vocabulary prior to 2024--2025 Causal Fermion Systems (CFS) expositions describing ``effective dynamical collapse'' and spacetime as a ``web of correlations.'' However, sharing high-level relational vocabulary is \emph{not} an independent empirical or mathematical validation of SIT.

At the mathematical equation level, Paper~048 establishes a narrow, bounded comparison using frozen sources: under its regulated constant-coupling continuum assumptions, the CFS field equations reduce to Einstein gravity coupled to one on-shell Dirac excitation, matching both a specialized constant-background neutral-Dirac submodel of July~2026 Version~80 Draft~13 and a direct Einstein--Dirac null control. Because all three share the same Einstein--Dirac equation class, that overlap carries \textbf{zero discriminatory value} against the ordinary Einstein--Dirac null and does not map primitive CFS operator measures to SIT fields. Testing whether a consistent constant-background truncation of November~2026 \texttt{M\_AUTH\_SIT\_EFT\_v1} admits the same Einstein--Dirac sector remains a separate open adapter calculation.

\subsection{Proposed Categorical and Process-Calculus Formalization Program (All Constructions Open)}

Earlier drafts sketched a categorical comparison ($\Phi : \mathcal{C}_{\mathrm{CFS}} \to \mathcal{C}_{\mathrm{SIT}}$) and a two-outcome process-calculus comparison over a finite observational fragment $G = (O, \tau, \mathcal{G}, \mathcal{T})$ while prematurely describing the desired conclusions (semantic embedding, weak bisimulation, and non-existence of a reverse functor) as already demonstrated. Version~81 explicitly demotes those statements to \textbf{proposed open mathematical work}:
\begin{enumerate}[leftmargin=2em]
  \item \textbf{Why the Two-Outcome Process Sketch Does Not Establish Operational Equivalence:} Assigning the identical abstract transition expression $P = \tau.(\mathrm{obs}(o_1).0 + \mathrm{obs}(o_2).0)$ to both CFS and SIT proves only equality of those two assigned expressions; it does \emph{not} derive that labeled transition system from either the CFS causal action principle or \texttt{M\_AUTH\_SIT\_EFT\_v1}.
  \item \textbf{Outstanding Obligations Required for Any Future Categorical Theorem (T1--T5):} Any future categorical embedding or separation claim would require explicitly constructing and proving all five currently unproved obligations: (\textbf{T1}) concrete definitions of the open-system categories $\mathcal{C}_{\mathrm{CFS}}$ and $\mathcal{C}_{\mathrm{SIT}}$ directly from the actual physical equations (not postulated labels); (\textbf{T2}) explicit derivation of the black-box functors to $\mathrm{LinRel}$ and proof of natural isomorphism $\alpha$; (\textbf{T3}) explicit construction of $\Phi$ with proved functoriality, monoidal coherence, and current naturality; (\textbf{T4}) a rigorous, invariant-checked separation proof showing why no reverse functor exists (without merely postulating that $\mathcal{C}_{\mathrm{SIT}}$ is closed while $\mathcal{C}_{\mathrm{CFS}}$ is not); and (\textbf{T5}) derivation of operational LTS semantics directly from both full physical theories before verifying Gorla encoding properties. Until T1--T5 are proved, no categorical embedding, operational equivalence, or expressive-power separation between SIT and CFS is claimed.
\end{enumerate}



\section{Historical Development of Super Information Theory (2017--2026)}
\label{app:historical-development}

Super Information Theory developed bottom-up over a continuous nine-year arc spanning systems neuroscience, oscillatory computation, emergent gravitation, effective field theory, and machine-checked formal verification:

\subsection{Phase I (2017): Information as Coincidence, Not Symbol}
Beginning in April 2017 (\emph{Neural Lace Podcast} and SVGN essays), an operational redefinition replaced static symbolic bits with physical \emph{coincidence patterns}---registered alignments of independent events within a temporal coincidence window---privileging phase synchrony and distributed internal observation (\emph{NAPOT}).

\subsection{Phase II (Summer 2022): Oscillatory Computation and Network-Level Agency}
In June 2022, the time-stamped \emph{Self Aware Networks} (SAN) GitHub archive and explainer lectures formalized cognition and control as multi-scale oscillatory phase coordination across distributed networks.

\subsection{Phase III (July 2022): Extending Oscillatory Information into Gravity (QGTCD)}
On July 26, 2022, \emph{Quantum Gradient Time Crystal Dilation} (QGTCD) extended phase-centric coincidence dynamics into gravitation, proposing time density as a structured physical medium whose spatial gradients steer trajectories.

\subsection{Phase IV (2024): Time as a Physical Medium and Volumetric Neural Synthesis}
Throughout 2024, public SVGN physics series and the book \emph{Bridging Molecular Mechanisms and Neural Oscillatory Dynamics} refined field-first intuition for time density alongside dendritic vector embeddings and cortical traveling waves.

\subsection{Phase V (January 2025): Wave-Based Dissipation (\emph{Micah's New Law})}
Early January 2025 formalized equilibration as local iterative dissipation of phase and amplitude differentials across coupled degrees of freedom.

\subsection{Phase VI (January 2025): Operational Time-Density Scalar (\emph{Super Dark Time})}
Late January 2025 promoted time density to an explicit scalar field $\rho_t(x)$ calibrated by local clock rates.

\subsection{Phase VII (February--December 2025): Unification in Two Primitives (Versions 1--79)}
Published originally on February~9, 2025 and refined through Version~79 (December~18, 2025) plus December~24, 2025 holonomy companions, SIT unified $\psi(x)$ and $\rho_t(x)$ inside a single covariant action with symmetron screening, $U(1)$ bundle holonomy, optical-clock bounds, and $^{87}\mathrm{Rb}$ BEC falsifiers.

\subsection{Phase VIII (July--October 2026): July Version~80 Compact Corrective Lineage and Specialized Intervening Papers}
Between Version~79 (December~18, 2025) and the November~2026 formalization series, the public \emph{Version~80} slot was established by the July~2026 compact corrective lineage: Paper~048 (\emph{Causal Fermion Systems and Super Information Theory}, July~28, 2026) cites its frozen comparison baseline as \emph{Version~80 Draft~13} (July~26, 2026), and Paper~002 (\emph{Super Information Theory: Coherence, Time Density, Informational Dissipation, and Gravity}, 26~pages, published July~29, 2026, Zenodo DOI \texttt{10.5281/zenodo.21693867}) published the correction-integrated Version~80 conformal two-real-scalar benchmark, introducing unital-CPTP vs.\ relative-entropy corrections, operational amplitude-vs-purity separation, and mass-calibration/form-factor identifiability tests alongside Papers~003, 007, 014, 028, 033, 034, 055, 062, 063, and \emph{Learning to Feel} (Paper~114: initial edition October~17, 2026; v1.2 October~20, 2026; technical supplement v1.1).

\subsection{Phase IX (November 2026): Canonical \texttt{M\_AUTH\_SIT\_EFT\_v1}, 4D Haar Wilson Lattice, Stage-6 $C^4$ Continuum Limit, Lean~4.30 Verification, and Bidirectional Integration (Version~81)}
Advancing internal November working drafts v80.2 (Stages~1--5.2, November~16, 2026, commit \texttt{8955eab}) and v80.3 (Stage~6 plus three Lean~4.30 publication addenda, Hydrogen house-style alignment, and magnetism/Part~III/lattice consistency repairs, November~18--19, 2026, commit \texttt{30a71e4}) into \textbf{Version~81} (so Version~80 remains cleanly assigned to the July~2026 compact corrective lineage), SIT locked the ghost-free $(-,+,+,+)$ effective action \texttt{M\_AUTH\_SIT\_EFT\_v1} (\texttt{v1\_1}), established the explicit July-vs-November model-continuity contract (Table~\ref{tab:july-v80-vs-nov-v81}), proved infinite-dimensional $\ell^2$ Sobolev convergence and pointwise 3-jet nonlinear curved $3+1$D ADM Dirac--Bergmann closure modulo $\partial_k B^k$ on the Sylvester SPD cone \texttt{IsSPD3}, resolved the \texttt{C09} zero-coherence gauge singularity non-perturbatively on a 4D hypercubic Euclidean Haar Wilson lattice via Model~L3 (\texttt{sitAMGMGatedOSLaw}), proved in Stage~6 and its addenda the exact finite-lattice corner Hessian, periodic lattice Fourier eigenvalue equations, conventional $C^4$ $O(1/l^2)$ classical continuum limit to $\frac{3\beta}{8}\int_{[0,1]^4}\frac{|\nabla f|^2}{f}\,du$, sharp $\frac{3\beta}{4}$ zero-division nonnegative-field coercivity allowing zero corners with both exact equality witnesses, exact Stage-5 finite-lattice embedding at $\kappa_t = 0$, and the canonical propagator inverse $(p^2 K + H)^{-1} = V\,\mathrm{diag}(1/(p^2+m_a^2))\,V^\top$ instantiated directly on \texttt{AuthModelParams}, diagnosed multiplicative-vs-conformal matter couplings and $U(1)$ kinetic mixing, executed a conditional weak-field E\"ot-Wash uniform-slab non-exclusion calculation paired with six corrected primary experimental provenance records (\texttt{primary\_measurements\_manifest\_v2.json}), imported the July~2026 unital-CPTP/relative-entropy repair, operational amplitude-vs-purity separation, Penrose--Onsager/matched-preparation BEC controls, and Paper~002/007/055 identifiability and QGTCD counterexample guardrails, and sealed all $669$ Lean~4.30 files ($48/48$ Stage-6/Addendum theorems) under the five-tier epistemic ledger \texttt{QG49} (\texttt{OPEN-1}--\texttt{OPEN-6} strictly \texttt{NOT\_SOLVED}).

\section{Exhaustive Documentary Record and Independent Timestamp Evidence (2017--2025)}
\label{app:documentary-record}

This appendix preserves the archival record of publicly available artifacts documenting the development of SAN, Neural Lace, NAPOT, QGTCD, Super Dark Time, SuperTimePosition, and Super Information Theory across GitHub, Zenodo, Figshare, Medium/SVGN, Internet Archive Wayback Machine, YouTube, and SoundCloud.

\subsection{Neural Lace Podcast: Complete Episode Record (2017--2019)}
\begin{itemize}
  \item \textbf{nlp1} --- Apr 11, 2017 (monologue). YouTube: \url{https://youtu.be/YrX_68oKuVs}
  \item \textbf{nlp2} --- Apr 13, 2017. YouTube: \url{https://youtu.be/UpOzHjKGa0g}
  \item \textbf{nlp3} --- Apr 22, 2017. YouTube: \url{https://youtu.be/_yKjtTVoVlU}
  \item \textbf{nlp4} --- May 5, 2017. YouTube: \url{https://youtu.be/RRN91RwiYiw}
  \item \textbf{nlp5} --- Recorded May 9, 2017 (Guest: Jules Urbach, OTOY). SoundCloud: \url{https://soundcloud.com/user-899513447/the-neural-lace-podcast-5-jules-urbach-ceo-otoy-inc}
  \item \textbf{nlp7} --- June 7, 2017 (Guest: Android Jones). SoundCloud: \url{https://soundcloud.com/user-899513447/the-neural-lace-podcast-7-guest-android-jones}
  \item \textbf{NerveGear Special} --- Nov 2, 2017. YouTube: \url{https://youtu.be/p9yTPBrrES4}
  \item \textbf{Season 2, Episode 1} --- Aug 2, 2018. YouTube: \url{https://www.youtube.com/watch?v=aexQwTpOwYc}
  \item \textbf{Season 2, Episode 2} --- Sep 11, 2018. YouTube: \url{https://youtu.be/wgMKMn7srHM}
  \item \textbf{Season 2, Episode 4} --- May 11, 2019 (fNIRS). YouTube: \url{https://youtu.be/DlfAOM4dx90}
\end{itemize}
Complete transcripts are preserved at \url{https://github.com/v5ma/selfawarenetworks/blob/main/guide/NLP}.

\subsection{Medium / SVGN Articles with Wayback Timestamp Proof (2017--2021)}
\begin{itemize}
  \item \emph{The Brain as a Special Kind of Hard Drive} (2017): \url{https://web.archive.org/web/20210919212933/https://medium.com/silicon-valley-global-news/the-brain-as-a-special-kind-of-hard-drive-4bd3a99dba4c}
  \item \emph{Humans Are Metal Robots in a Valid Sense} (2017): \url{https://web.archive.org/web/20201020015247/https://medium.com/silicon-valley-global-news/breaking-news-humans-are-metal-robots-literally-a7bd78bb333d}
  \item \emph{Addressing Criticism for ``Humans Are Metal Robots''} (2017): \url{https://web.archive.org/web/20210919191705/https://medium.com/silicon-valley-global-news/addressing-criticism-for-my-humans-are-metal-robot-story-86f1c99d72b6}
  \item \emph{Synaptic Unreliability Undermined by MVP Proof} (2017): \url{https://web.archive.org/web/20210209003006/https://medium.com/silicon-valley-global-news/synaptic-unreliability-a-foundational-concept-found-in-deep-learning-and-in-computational-a97c5dca9e90}
  \item \emph{Mind Code / Brain Code: Go and AlphaGo} (2017): \url{https://web.archive.org/web/20210919213836/https://medium.com/silicon-valley-global-news/the-game-of-go-speculation-about-the-code-of-our-minds-in-our-brains-5a39d7cb6366}
  \item \emph{Neural Lace Journal / Talk Show Index} (Oct 2017): \url{https://web.archive.org/web/20190317101602/https://medium.com/silicon-valley-global-news/the-neural-lace-journal-talk-show-a-compilation-of-articles-and-links-ae3fbb0d5d5a}
  \item \emph{Neural Lace and Deep Learning: Polina Anikeeva} (2018): \url{https://web.archive.org/web/20180530045930/https://medium.com/silicon-valley-global-news/neural-lace-and-deep-learning-6ed70db4e3a7}
  \item \emph{3D Cross-Hair Convolutional Neural Networks} (2018): \url{https://web.archive.org/web/20210919202220/https://medium.com/silicon-valley-global-news/3d-cross-hair-convolutional-neural-networks-5d39e2b565ca}
  \item \emph{Interview with Jules Urbach at GTC 2018} (2018): \url{https://web.archive.org/web/20180528155036/https://medium.com/silicon-valley-global-news/my-interview-with-jules-urbach-the-ceo-of-otoy-at-gtc-2018-a7cadbaa6ec3}
  \item \emph{Question to Jensen Huang on Self-Aware Networks} (2018): \url{https://web.archive.org/web/20251112031520/https://medium.com/silicon-valley-global-news/after-a-mind-exploding-keynote-at-gtc-2018-i-got-the-chance-to-ask-jensen-huang-the-ceo-of-nvidia-a-f1cca7ce04a6}
\end{itemize}

\subsection{2022 GitHub Repositories: SAN and QGTCD}
\begin{itemize}
  \item \textbf{2022-06-17} --- \emph{Self Aware Networks} (SAN) GitHub repository documenting oscillatory computational agency, predictive synchronization, and phase-based network control.
  \item \textbf{2022-07-26} --- \emph{Quantum Gradient Time Crystal Dilation (QGTCD)} initial public GitHub release extending phase-centric coincidence dynamics into gravitational time-density gradients.
\end{itemize}

\subsection{Summer 2022 YouTube Explainer Videos and 2024 Book Release}
\begin{enumerate}[leftmargin=2em]
  \item \emph{1 hour video that explains NAPOT the central thesis of Self Aware Networks}: \url{https://www.youtube.com/watch?v=fLp-yTQ6pSM} (ID: \texttt{fLp-yTQ6pSM}).
  \item \emph{Neural Array Projection Oscillation Tomography Theory 5th Revision}: \url{https://www.youtube.com/watch?v=vixhppNAKPs} (ID: \texttt{vixhppNAKPs}).
  \item \emph{The Self Aware Networks Institute for Neurophysics, Artificial Neurology \& Bio-Synthetic Interfaces} (3 min): \url{https://www.youtube.com/watch?v=VTBNyUM47Zg} (ID: \texttt{VTBNyUM47Zg}).
  \item \emph{Self Aware Networks: A longer explanation of my research, my past work, and what I have built} (13 min): \url{https://www.youtube.com/watch?v=IKbl0ryKRoY} (ID: \texttt{IKbl0ryKRoY}).
  \item Blumberg, Micah (2024). \emph{Bridging Molecular Mechanisms and Neural Oscillatory Dynamics: Explore how synaptic modulation and pattern generation create the brain's seamless volumetric three-dimensional conscious experience}. Amazon ASIN: \texttt{B0DL4701875}, \url{https://www.amazon.com/dp/B0DL4701875}.
\end{enumerate}

\subsection{SVGN / Substack Physics Series (2024--2025) and Zenodo Records}
\begin{itemize}
  \item \emph{A New Unified Field Theory Called QGTCD (Part I)} --- Jan 28, 2024: \url{https://web.archive.org/web/20240128135920/https://www.svgn.io/p/a-new-unified-field-theory-called-qgtcd}
  \item \emph{QGTCD (Part II)} --- Feb 15, 2024: \url{https://web.archive.org/web/20240215141003/https://www.svgn.io/p/qgtcd-part-ii}
  \item \emph{Explain QGTCD Like I Am Six} --- Mar 28, 2024: \url{https://web.archive.org/web/20240328132046/https://www.svgn.io/p/explain-qgtcd-like-i-am-six}
  \item \emph{Quantum Gravity's New Frontier: Time as Medium} --- Oct 21, 2024: \url{https://web.archive.org/web/20241021151234/https://www.svgn.io/p/quantum-gravitys-new-frontier-time-as-medium}
  \item \emph{Dark Time Theory: A Conversation} --- Oct 22, 2024: \url{https://web.archive.org/web/20241022145519/https://www.svgn.io/p/dark-time-theory-a-conversation}
  \item \emph{Micah's New Law of Thermodynamics} --- Jan 1, 2025: \url{https://web.archive.org/web/20250101084219/https://www.svgn.io/p/micahs-new-law-of-thermodynamics}
  \item \emph{Introducing Quantum SuperTimePosition} --- Jan 5, 2025: \url{https://web.archive.org/web/20250115135920/https://www.svgn.io/p/introducing-quantum-supertimeposition}
  \item \emph{SuperTimePosition Measured} --- Jan 17, 2025: \url{https://web.archive.org/web/20250117214702/https://www.svgn.io/p/supertimeposition-measured}
\end{itemize}

\paragraph{Zenodo DOI Records (Chronological by Original Publication Date):}
\begin{itemize}
  \item \emph{Coincidence as a Bit of Information} (2017; consolidated June 21, 2025; Zenodo Aug 21, 2025): \url{https://zenodo.org/records/16922510}
  \item \emph{Micah's New Law of Thermodynamics} (Jan 23, 2025; Figshare DOI \texttt{10.\allowbreak{}6084/m9.\allowbreak{}figshare.\allowbreak{}28264340}): \url{https://zenodo.org/records/16922506}
  \item \emph{Super Dark Time: Gravity Computed from Local Quantum Mechanics} (Jan 27, 2025; Figshare DOI \texttt{10.\allowbreak{}6084/m9.\allowbreak{}figshare.\allowbreak{}28284545}): \url{https://zenodo.org/records/16922502}
  \item \emph{Super Information Theory: The Coherence Conservation Law Unifying the Wave Function, Gravity, and Time} (Feb 9, 2025; Figshare DOI \texttt{10.\allowbreak{}6084/m9.\allowbreak{}figshare.\allowbreak{}28379318}): \url{https://zenodo.org/records/16922512}
  \item \emph{Self Aware Networks: OCA (First Draft)} (May 16, 2025; Figshare DOI \texttt{10.\allowbreak{}6084/m9.\allowbreak{}figshare.\allowbreak{}29085134}): \url{https://zenodo.org/records/16922401}
  \item \emph{Neuroscience in Review: Mapping ``Cortical traveling waves in time and space'' (2025) to Self Aware Networks (2022)} (June 18, 2025; Figshare DOI \texttt{10.\allowbreak{}6084/m9.\allowbreak{}figshare.\allowbreak{}29817383}): \url{https://zenodo.org/records/16922844}
  \item \emph{Neuroscience in Review: Brain Rhythms in Cognition (2024--25) vs.\ Blumberg's Self-Aware Networks (2017--25)} (July 2025; Figshare DOI \texttt{10.\allowbreak{}6084/m9.\allowbreak{}figshare.\allowbreak{}29650250}): \url{https://zenodo.org/records/16922874}
  \item \emph{Timeline Review of Kletetschka's 3D Time Theory \& Blumberg's SIT, SDT, QGTCD} (July 25, 2025; Figshare DOI \texttt{10.\allowbreak{}6084/m9.\allowbreak{}figshare.\allowbreak{}29647640}): \url{https://zenodo.org/records/16922866}
  \item \emph{Informational Holonomy and the Electron g-Factor: A Super Information Theory (SIT) Spinoff} (v4, Oct 4, 2025): \url{https://zenodo.org/records/17307956}
  \item \emph{Coherence-Gated Yang--Mills Gap: A Self-Contained Derivation in Super Information Theory} (Oct 10, 2025; DOI \texttt{10.\allowbreak{}5281/zenodo.\allowbreak{}17311157}): \url{https://doi.org/10.5281/zenodo.17311157}
  \item \emph{Neuroscience in Review: Earl K. Miller's MIT Corpus vs.\ Micah Blumberg's SAN and SIT} (Oct 10, 2025): \url{https://zenodo.org/records/17308645}
  \item \emph{Building Sentient Beings} (Micah Blumberg \& Michael S.\ P.\ Miller, Dec 17, 2025): \url{https://zenodo.org/records/47834429}
\end{itemize}

\setcounter{subsection}{0}
\renewcommand{\thesection}{AA}
\renewcommand{\thesubsection}{AA.\arabic{subsection}}
\section*{Appendix AA: Timestamp Evidence (Wayback Machine Mirrors and ORCID)}
\addcontentsline{toc}{section}{Appendix AA: Timestamp Evidence (Wayback Machine Mirrors and ORCID)}
\label{app:timestamp-aa}

This appendix records independent third-party Internet Archive (Wayback Machine) and ORCID identity anchors corroborating January--July 2025 Figshare and repository availability:
\begin{itemize}
  \item \textbf{ORCID Identity Anchor (Micah Blumberg, \texttt{0009-0004-5175-9532}):} \url{https://web.archive.org/web/20250826004755/https://orcid.org/0009-0004-5175-9532}
  \item \textbf{Super Information Theory (Figshare Wayback Index):} \url{https://web.archive.org/web/*/https://figshare.com/articles/journal_contribution/Super_Information_Theory/*}
  \item \textbf{Super Dark Time (Figshare Wayback Index):} \url{https://web.archive.org/web/*/https://figshare.com/articles/journal_contribution/Super_Dark_Time_Gravity_Computed_from_Local_Quantum_Mechanics/*}
  \item \textbf{Micah's New Law of Thermodynamics (Figshare Wayback Index):} \url{https://web.archive.org/web/*/https://figshare.com/articles/journal_contribution/_b_Micah_s_New_Law_of_Thermodynamics_A_Signal-Dissipation_Framework_for_Equilibrium_and_Consciousness_b_/*}
  \item \textbf{Self Aware Networks: OCA (First Draft, May 16, 2025; S3 PDF Snapshot 2025-07-25):} \url{https://web.archive.org/web/20250725002526/https://s3-eu-west-1.amazonaws.com/pfigshare-u-files/54589349/SelfAwareNetworksOCAFirstDraft.pdf}
\end{itemize}
Note that Wayback capture timestamps record archive crawl dates, whereas original publication dates are independently attested by ORCID metadata, Figshare/Zenodo DOIs, and GitHub commit histories.
% --- END INPUT sections/11_appendices_N_to_AA.tex ---
% --- BEGIN INPUT sections/12_appendix_AB_and_backmatter.tex ---
\setcounter{subsection}{0}
\renewcommand{\thesection}{AB}
\renewcommand{\thesubsection}{AB.\arabic{subsection}}
\section*{Appendix AB: Five-Tier Epistemic Ledger (\texttt{QG49}), \texttt{M\_AUTH\_SIT\_EFT\_v1} Benchmarks, and Stage 1--6 Lean~4.30 Theorem Index}
\addcontentsline{toc}{section}{Appendix AB: Five-Tier Epistemic Ledger (QG49), M\_AUTH\_SIT\_EFT\_v1 Benchmarks, and Stage 1--6 Lean 4.30 Theorem Index}
\label{app:epistemic-ledger-lean}

To enforce uncompromising scientific transparency across mathematical proofs, perturbative effective field theory, numerical lattice simulations, retired historical hypotheses, and open mathematical and observational problems, this appendix records the complete Version~81 five-tier epistemic ledger (\texttt{QG49}), the machine-precision numerical benchmark suite (\texttt{M\_AUTH\_SIT\_EFT\_v1} canon \texttt{v1\_1}), and the complete Stage~1--6 (plus three publication addenda) Lean~4.30 module and capstone theorem index. Each formal theorem addresses a specific bounded mathematical claim; aggregate module or theorem counts never substitute for unresolved microphysical derivations or quantum continuum limits (\texttt{OPEN-1}--\texttt{OPEN-6}).

\subsection{Five-Tier Epistemic Standing Ledger (\texttt{QG49})}

\begin{longtable}{|p{2.3cm}|p{4.3cm}|p{8.4cm}|}
\caption{Complete Five-Tier Epistemic Classification (\texttt{QG49}) for Super Information Theory Version~81.}\label{tab:qg49-ledger}\\
\hline
\textbf{Tier} & \textbf{Classification \& Standard} & \textbf{Verified Scope and Canonical Items} \\
\hline
\endfirsthead
\hline
\textbf{Tier} & \textbf{Classification \& Standard} & \textbf{Verified Scope and Canonical Items} \\
\hline
\endhead
\hline
\endfoot
\textbf{Tier 1} & \texttt{LEAN\_\allowbreak{}KERNEL\_\allowbreak{}VERIFIED\_\allowbreak{}} \texttt{STRICT\_\allowbreak{}STANDARD\_\allowbreak{}AXIOMS} (\texttt{propext}, \texttt{Classical.choice}, \texttt{Quot.sound}; \texttt{0 sorry}, \texttt{0 admit}, \texttt{0} custom axioms, \texttt{0 native\_decide} in clean Stage~2--6 closures) & Lean~4.30.0 kernel-compiled formal theorem types (audited individually by claim rather than treated as an aggregate proof of physical completeness): $2{,}368$ strict Stage~1--3 roots + $218$ \texttt{Stage4PureAudit} declarations + all $10$ Stage~4 lattice/holonomy targets + $544$ Stage~5 extension declarations across $6$ compiled Stage~5 modules ($660$ frozen Stage~1--5 \texttt{.lean} files) + $6$ frozen Stage-6 modules (\texttt{Stage6FunctionalResponse/*}, $30/30$ audited theorems in \texttt{SITStage6Audit}) + $3$ publication-branch addendum modules (\texttt{Stage6Addendum.\allowbreak{}PeriodicLatticeFourierEigenvalue} [$3$], \texttt{NonnegativeFieldCoercivity} [$10$], and \texttt{ActionResponsePropagatorIdentity} [$5$]) = \textbf{$669$ total \texttt{.lean} files} and $48/48$ Stage-6/Addendum strict-axiom theorems. \\
\hline
\textbf{Tier 2} & \texttt{ANALYTICAL\_\allowbreak{}AND\_\allowbreak{}SYMPY\_\allowbreak{}} \texttt{PERTURBATIVE\_EFT} ($\mathrm{rtol} \le 10^{-12}$ SymPy~1.14.0 vs.\ NumPy~2.4.6 + exploratory MCMC) & Sub-Planckian \texttt{M\_AUTH\_SIT\_EFT\_v1} ($E \ll \Lambda_{\mathrm{UV}}$), $74/74$ curved $3+1$D ADM symbolic \& reverse-AD polynomial jet checks, Seeley--DeWitt $(a_0,a_1,a_2)$, continuum 3D Foldy--Wouthuysen / Schr\"odinger Coulomb / Brillouin-zone overlap derivations, deterministic Stage-6 smooth-torus $O(1/n^2)$ convergence regressions ($1.990123$ 4D / $1.999907$ 1D), Stage-6 conditional weak-field uniform-slab E\"ot-Wash non-exclusion calculation under declared dust ($q_{\mathrm{eff}}=1$) vs.\ free-scalar ($q_{\mathrm{eff}}=0$) closures, 2-variable zero-mode effective-potential barrier sampler cross-check ($\Delta V_{\mathrm{MCMC}} = 0.0610 \pm 0.0021$, $0.71\sigma$), and exploratory finite-volume 4D $SU(2)$ MCMC consistency benchmarks. \\
\hline
\textbf{Tier 3} & \texttt{CONDITIONAL\_\allowbreak{}INTERFACES\_\allowbreak{}} \texttt{AND\_BOUNDARIES} (hypotheses explicit in theorem types, \texttt{0 let \_ :=} suppressors) & Stage-4 Target~\texttt{(7)} (\texttt{OpenAI\allowbreak{}UV\allowbreak{}Fixed\allowbreak{}Point\allowbreak{}Assumptions} $\to$ \texttt{YangMills.\allowbreak{}Joint\allowbreak{}Existence}), held arithmetic bridge \texttt{SITF43\allowbreak{}To\allowbreak{}F46\allowbreak{}Riemann\allowbreak{}Bridge} (\texttt{HELD\_PER\_ODQ7}), \texttt{CoherenceTowerSpec}, single-chart Sylvester SPD cone \texttt{IsSPD3} and formal 3-jet identities modulo $\partial_k B^k$, conditional Jordan-frame Cassini identities (contingent on conformal response derivation), conditional E\"ot-Wash uniform-slab closure ($q_{\mathrm{eff}} \in \{0,1\}$), $\kappa_t = 0$ Stage-6 action-gap isolation, and 3D continuum Ladyzhenskaya hypothesis \texttt{hLady}. \\
\hline
\textbf{Tier 4} & \texttt{REFUTED\_\allowbreak{}PURGED\_\allowbreak{}RETIRED\_\allowbreak{}} \texttt{AND\_HISTORICAL} (forbidden in canonical claims) & Wrong-sign/multitemporal/Ostrogradsky ghosts \texttt{C04}, \texttt{C05}, \texttt{C07}; sign/factor typos \texttt{T1}--\texttt{T5}, \texttt{R1}--\texttt{R5}; unconstructed SUSY \texttt{C06}; row-normalized clock collapse \texttt{C08}; continuum $f(\rho_t,0)=0$ $\Pi^2/(2f)$ singularity \texttt{C09}; 1-loop Sakharov induced gravity \texttt{C12}; sharp spatial $P_N$ in boosted 4D YM; $10$ frozen legacy Stage-1 \texttt{native\_decide} files; and Stage-6 quarantined observational shorthands (legacy $3\times 10^{-8}$ Skytree attribution, $29.4\,\mu\mathrm{Gal}$ BEC ceiling reclassified as hypothetical sensitivity, retired \texttt{fig4}/\texttt{fig5}, and hand-coded \texttt{np.piecewise} Yukawa sketch). \\
\hline
\textbf{Tier 5} & \texttt{EXPLICIT\_\allowbreak{}UNSOLVED\_\allowbreak{}OPEN\_\allowbreak{}} \texttt{PROBLEMS} (all strictly \texttt{NOT\_SOLVED}) & (1) \texttt{OPEN-1}: Non-perturbative 4D quantum-gravity UV fixed point at $E \ge \Lambda_{\mathrm{UV}}$ (\texttt{NOT\_SOLVED}); (2) \texttt{OPEN-2}/\texttt{OPEN-2B}: Quantum Wheeler--DeWitt ordering, Sobolev/Banach constraint algebras with Regge--Teitelboim boundary integrals, and global multi-chart Cauchy PDE (\texttt{NOT\_SOLVED}); (3) \texttt{OPEN-3}: Interacting 4D \emph{quantum} continuum-measure convergence and OS transfer-operator Hamiltonian for Model~L3, pure 4D continuum Yang--Mills existence/uniqueness/mass gap ($f\equiv 1$), unmodified 3D Clay Navier--Stokes regularity, and formal unbounded 3D $L^2(\mathbb{R}^3)$ PDEs (\texttt{NOT\_SOLVED}); (4) \texttt{OPEN-4}: Riemann Hypothesis (\texttt{HELD\_PER\_ODQ7}, \texttt{NOT\_SOLVED}); (5) \texttt{OPEN-5}: Massless Weyl chiral determinant phase, three-generation soliton homotopy $\pi_3(\mathcal{M}_{\mathrm{coh}})$, and $P \ne NP$ (\texttt{NOT\_SOLVED}); (6) \texttt{OPEN-6}: Microscopic composite-material scalar charges $\langle\mathcal{L}_m\rangle/\rho_m$, visible $U(1)$ photon dictionary, laboratory $J_{\mathrm{prep}}$, and full patterned rotating-mass E\"ot-Wash torque likelihoods from Lee's dissertation Appendix~A (\texttt{NOT\_SOLVED}). \\
\hline
\end{longtable}

\subsection{Canonical \texttt{M\_AUTH\_SIT\_EFT\_v1} Benchmark Suite (\texttt{BM\_PHYS\_EFT\_1} and Companion Benchmarks)}

Table~\ref{tab:bm-phys-eft-1} records the verified inputs and derived physical observables of canonical benchmark \texttt{BM\_PHYS\_EFT\_1} ($\max\text{SymPy/NumPy rtol} = 2.22\times 10^{-16}$), followed by Table~\ref{tab:six-companion-benchmarks} summarizing the six structural and disambiguation benchmarks from \texttt{SIT-EFT-CANONICAL-LEDGER-v1\_\allowbreak{}1.\allowbreak{}json}.

\begin{table}[htbp]
\centering
\caption{Canonical Physical Benchmark \texttt{BM\_PHYS\_EFT\_1} (CODATA 2022 constants, natural units $[M] = \mathrm{GeV}$).}
\label{tab:bm-phys-eft-1}
\vspace{0.4em}
\footnotesize
\begin{tabular}{|>{\raggedright\arraybackslash}p{5.5cm}|>{\raggedright\arraybackslash}p{4.3cm}|>{\raggedright\arraybackslash}p{5.0cm}|}
\hline
\textbf{Parameter / Derived Observable} & \textbf{Exact Formula / Value} & \textbf{Physical Interpretation / Status} \\
\hline
$M_{\mathrm{Pl}} = (8\pi G)^{-1/2}$ & $2.4353234593382036 \times 10^{18}\,\mathrm{GeV}$ & Canonical reduced Planck scale \\
$\kappa_G = 8\pi G = M_{\mathrm{Pl}}^{-2}$ & $1.68611307195616 \times 10^{-37}\,\mathrm{GeV}^{-2}$ & Einstein coupling \\
$\kappa_t,\;\kappa_c$ & $1.0,\;1.0$ & Canonical scalar kinetic norms \\
$\psi_0 = \rho_0 = m_{\mathrm{link}}$ & $2.0 \times 10^{16}\,\mathrm{GeV}$ & GUT-scale informational baseline \\
$R_0,\;v_0 = \psi_0 R_0$ & $0.01,\;2.0 \times 10^{14}\,\mathrm{GeV}$ & Vacuum coherence ratio \& VEV \\
$c_{\mathrm{EM}},\;M_{\mathrm{sym}}$ & $1.0 \times 10^{-4},\;50.0$ & Dimension-6 gauge \& symmetron couplings \\
$\lambda_\rho,\;\lambda_\psi,\;\lambda_{\mathrm{portal}},\;\hat{q}$ & $0.5,\;0.5,\;0.25,\;0.303$ & Quartic, portal, and $U(1)$ couplings \\
\hline
$\Lambda_{\mathrm{EM}} = \psi_0 / \sqrt{c_{\mathrm{EM}}}$ & $2.0 \times 10^{18}\,\mathrm{GeV}$ & Dimension-6 gauge-gating cutoff \\
$\Lambda_{\mathrm{sym}} = \sqrt{2}\,\psi_0 M_{\mathrm{sym}}$ & $1.4142135623730952 \times 10^{18}\,\mathrm{GeV}$ & Dimension-6 symmetron cutoff \\
$\Lambda_{\mathrm{UV}} = \min(M_{\mathrm{Pl}},\Lambda_{\mathrm{EM}},\Lambda_{\mathrm{sym}})$ & $1.4142135623730952 \times 10^{18}\,\mathrm{GeV}$ & Unified EFT ultraviolet ceiling \\
$f_{\mathrm{EM}0},\;f_{\mathrm{SM}0}$ & $1.00000001,\;1.00000002$ & Vacuum-positive dressings ($\ge 1$) \\
$(\Lambda_{\mathrm{sym}}/\Lambda_{\mathrm{EM}})^2 = (f_{\mathrm{EM}0}-1)/(f_{\mathrm{SM}0}-1)$ & $0.5$ & Exact cutoff-ratio identity \\
$\det(H)$ & $4.8 \times 10^{61}\,\mathrm{GeV}^4$ & Positive-definite radial Hessian \\
Portal eigenmasses $[m_1,\;m_2]$ & $[2.449459\times 10^{14},\;2.828462\times 10^{16}]\,\mathrm{GeV}$ & Non-tachyonic scalar spectrum \\
Proca mass $m_A$ & $6.05999997 \times 10^{13}\,\mathrm{GeV}$ & Screened $U(1)$ longitudinal mode \\
$\alpha_{\mathrm{sym}}$ & $4.74464009 \times 10^{-7}$ & Conditional dust ($q_{\mathrm{eff}}=1$) fifth-force coupling (matches Model~C2 $A=f_{\mathrm{SM}}^{\pm 1}$; Model~C1 $A=f_{\mathrm{SM}}^{-1/2}$ gives $\frac{1}{4}\alpha_{\mathrm{sym}} = 1.18616002\times 10^{-7}$) \\
Light-mode range $\lambda_1$ (dust $q_{\mathrm{eff}}=1$) & $8.055941 \times 10^{-31}\,\mathrm{m}$ & $\log_{10} R_{\mathrm{slab}}(52\,\mu\mathrm{m}) = -2.80331\times 10^{25}$ (Stage~6) \\
$|\gamma_{\mathrm{PPN}} - 1|_{S=1}$ (Jordan-frame) & $9.48927568 \times 10^{-7}$ (C2) / $2.37231976 \times 10^{-7}$ (C1) & Strictly conditional on chosen conformal extension C2 ($A=f_{\mathrm{SM}}^{\pm 1}$) or C1 ($A=f_{\mathrm{SM}}^{-1/2}$); not derived from unchanged $+f_{\mathrm{SM}}\mathcal{L}_m$ \\
Loop factors $(E/\Lambda_{\mathrm{UV}})^{2L+2}$ ($E/\Lambda_{\mathrm{UV}}=10^{-2}$) & $[10^{-4},\;10^{-8},\;10^{-12}]$ ($L=0,1,2$) & Monotone Weinberg suppression \\
\hline
\end{tabular}
\end{table}

\begin{table}[htbp]
\centering
\caption{Companion Structural, Limit, and Disambiguation Benchmarks (\texttt{SIT-EFT-CANONICAL-LEDGER-v1\_\allowbreak{}1.\allowbreak{}json}).}
\label{tab:six-companion-benchmarks}
\vspace{0.4em}
\small
\begin{tabular}{|p{4.6cm}|p{9.8cm}|}
\hline
\textbf{Benchmark ID} & \textbf{Parameters and Verified Identity} \\
\hline
\texttt{BM\_\allowbreak{}QG39\_\allowbreak{}RECONCILED\_\allowbreak{}ADV\_\allowbreak{}D1\_\allowbreak{}01} & $(\psi_0, M_{\mathrm{sym}}, \Rcoh, c_{\mathrm{EM}}) = (1, 100, 0.1, 0.01) \implies \Lambda_{\mathrm{sym}} = 141.4214$, $(\Lambda_{\mathrm{sym}}/\Lambda_{\mathrm{EM}})^2 = 200 \implies \Lambda_{\mathrm{EM}} = 10.0$; setting $c_{\mathrm{EM}} = 5\times 10^{-5}$ aligns $\Lambda_{\mathrm{EM}} = \Lambda_{\mathrm{sym}} = \Lambda_{\mathrm{UV}} = 141.4214$. \\
\hline
\texttt{BM\_\allowbreak{}QG31\_\allowbreak{}KAPPA05\_\allowbreak{}DISAMBIGUATION} & At $\kappa_G = 0.5$, canonical $M_{\mathrm{Pl}} = \kappa_G^{-1/2} = \sqrt{2} = 1.41421356$ while \texttt{QG31} convention $M_{\mathrm{Pl,QG31}} = (2\kappa_G)^{-1/2} = 1.0$. \\
\hline
\texttt{BM\_\allowbreak{}QG16\_\allowbreak{}UNIT\_\allowbreak{}MASS\_\allowbreak{}EIGENVALUES} & $(\rho_0,v_0,\lambda_\rho,\lambda_\psi,m_{\mathrm{link}},\lambda_{\mathrm{portal}}) = (1,1,0.5,1,1,1) \implies H = \bigl(\begin{smallmatrix}2 & -2\\ -2 & 6\end{smallmatrix}\bigr)$, $\mathrm{tr}(H)=8$, $\det(H)=8$, $m_\mp^2 = 4 \mp 2\sqrt{2} = [1.171573, 6.828427]$. \\
\hline
\texttt{BM\_\allowbreak{}CASSINI\_\allowbreak{}YUKAWA\_\allowbreak{}1} & Abstract conformal coupling $\alpha_{\mathrm{conf}} = 10^{-6}$, $S = 1 \implies \gamma_{\mathrm{PPN}} = 0.999998000002$, $|\gamma_{\mathrm{PPN}}-1| = 1.999998\times 10^{-6} < 2.3\times 10^{-5}$ (conditional on conformal response). \\
\hline
\texttt{BM\_V0\_SYMMETRIC\_GR} & Symmetric limit $v_0 = R_0 = 0 \implies f_{\mathrm{EM}0} = f_{\mathrm{SM}0} = 1$, $\alpha_{\mathrm{sym}} = 0$, $\gamma_{\mathrm{PPN}} = \beta_{\mathrm{PPN}} = 1$, $m_A^2 = 0$. \\
\hline
\texttt{BM\_\allowbreak{}LEAN\_\allowbreak{}STRUCTURAL\_\allowbreak{}WITNESS} & Stage-2 Lean \texttt{canonicalAuthParams} (all \texttt{AuthModelParams} $= 1.0$, coupled cutoffs $(10,10,10,1,1)$ with \texttt{OBS-06} closed). \\
\hline
\end{tabular}
\end{table}

All eight Stage~2--3 read-only advisories (\texttt{QG20\_NBE1} theorem-name crosswalk, \texttt{W\_ADV\_TD\_01} zero cosmological constant naming, \texttt{W\_ADV\_TD\_02} Sobolev/Ladyzhenskaya call sites, \texttt{W\_ADV\_TD\_03} sealed tranche tags, \texttt{W\_ADV\_TD\_04} structural witness vs.\ phenomenology separation, \texttt{ADV\_D1\_01} Planck/$\Lambda_{\mathrm{EM}}$ crosswalk, \texttt{ADV\_D1\_02} coupled \texttt{SITAuthoritative\allowbreak{}EFTPackage} with zero \texttt{let \_ :=} suppressors, and \texttt{ADV\_D1\_03} fixed-equilibrium quadratic dust stiffening) are preserved reconciled verbatim.

\subsection{Stage 1--6 (Plus Addendum) Lean~4.30 Module and Capstone Theorem Index}

Table~\ref{tab:lean430-module-index} indexes the complete machine-checked Lean~4.30 formalization (\texttt{Lean v4.30.0}, Mathlib \texttt{c5ea00351c28e24afc9f0f84379aa41082b1188f}) across Stages~1 through~6 plus the three November~2026 publication addenda ($669$ total \texttt{.lean} files).

{\small
\begin{longtable}{|>{\raggedright\arraybackslash}p{2.0cm}|>{\raggedright\arraybackslash}p{5.5cm}|>{\raggedright\arraybackslash}p{7.3cm}|}
\caption{Stage~1--6 (Plus Addendum) Lean~4.30 Module Inventory and Principal Capstone Theorems.}\label{tab:lean430-module-index}\\
\hline
\textbf{Stage} & \textbf{Lean~4.30 Module / Namespace Family} & \textbf{Strict Count \& Principal Proved Theorems} \\
\hline
\endfirsthead
\hline
\textbf{Stage} & \textbf{Lean~4.30 Module / Namespace Family} & \textbf{Strict Count \& Principal Proved Theorems} \\
\hline
\endhead
\hline
\endfoot
\textbf{Stage 1} & \texttt{SITF49\allowbreak{}Master\allowbreak{}Theorem\allowbreak{}Stack}, \texttt{SITTrack\allowbreak{}FCapstone\allowbreak{}Audit\allowbreak{}Purified}, \texttt{SITCore\allowbreak{}Source\allowbreak{}Bridge\allowbreak{}Audit\allowbreak{}Purified}, \texttt{HsitTail.SitUv*} & $2{,}021$ strict standard-axiom roots ($1{,}949 + 53 + 19$): localization recurrence, \texttt{unbounded\_\allowbreak{}fiber\_\allowbreak{}collapses\_\allowbreak{}projected\_\allowbreak{}support}, conditional \texttt{SITF43\allowbreak{}To\allowbreak{}F46\allowbreak{}Riemann\allowbreak{}Bridge}, and \texttt{sit\allowbreak{}F49\_\allowbreak{}master\_\allowbreak{}unconditional\_\allowbreak{}and\_\allowbreak{}boundary\_\allowbreak{}stack}. \\
\hline
\textbf{Stage 2} & \texttt{SITQG09\allowbreak{}Linearized\allowbreak{}Fluctuation\allowbreak{}And\allowbreak{}Obstruction}, \texttt{SITTranche\allowbreak{}BAuthoritative\allowbreak{}Model\allowbreak{}And\allowbreak{}Quantum}, \texttt{SITTranche\allowbreak{}BGRLimits\allowbreak{}And\allowbreak{}EFT}, \texttt{SITTranche\allowbreak{}DCertificates\allowbreak{}And\allowbreak{}Model\allowbreak{}Unification}, \texttt{SITTranche\allowbreak{}DCoherence\allowbreak{}Tower\allowbreak{}And\allowbreak{}Transport} (\texttt{SitQg.Stage2.*}) & $100$ strict roots: \texttt{potentials\allowbreak{}Total\_\allowbreak{}globally\_\allowbreak{}nonneg}, \texttt{vacuum\allowbreak{}Equilibrium\_\allowbreak{}zero\allowbreak{}Cosmological\allowbreak{}Constant}, $7$-mode ADM/Proca positivity, \texttt{ppn\allowbreak{}Gamma\allowbreak{}Bound\_\allowbreak{}of\_\allowbreak{}alpha\_\allowbreak{}and\_\allowbreak{}screening}, \texttt{weinberg\allowbreak{}Loop\allowbreak{}Power\_\allowbreak{}monotone}, \texttt{Jet1Plus1D} Bianchi--Noether conservation, and bidirectional \texttt{ODQ-7} arithmetic decoupling (\texttt{QG31To35}). \\
\hline
\textbf{Stage 3 Phase 3} & \texttt{SITStage3\allowbreak{}Infinite\allowbreak{}Sobolev\allowbreak{}Coherence\allowbreak{}Tower}, \texttt{SITStage3Phase3Audit} (\texttt{Sit\allowbreak{}Qg.\allowbreak{}Stage3.\allowbreak{}Infinite\allowbreak{}Sobolev\allowbreak{}Tower}) & $66$ strict roots: \texttt{l2\allowbreak{}Seq\_\allowbreak{}not\_\allowbreak{}finite\allowbreak{}Dimensional}, \texttt{h2\allowbreak{}Submodule\_\allowbreak{}strict\_\allowbreak{}lt\_\allowbreak{}h1\allowbreak{}Submodule}, \texttt{h1\allowbreak{}Submodule\_\allowbreak{}strict\_\allowbreak{}lt\_\allowbreak{}top}, \texttt{galerkin\allowbreak{}Proj\_\allowbreak{}tendsto\_\allowbreak{}nhds}, \texttt{galerkin\allowbreak{}Proj\_\allowbreak{}l2\_\allowbreak{}defect\_\allowbreak{}bound}, \texttt{galerkin\allowbreak{}Proj\_\allowbreak{}h1\_\allowbreak{}defect\_\allowbreak{}bound}, \texttt{gauge\allowbreak{}Galerkin\allowbreak{}Proj\_\allowbreak{}intertwines}. \\
\hline
\textbf{Stage 3 Phase 4} & \texttt{Stage3ADM.Basic}, \texttt{Pillar1Kinematics}, \texttt{Pillar2Legendre}, \texttt{Pillar3ReducedClosure}, \texttt{Pillar4\allowbreak{}Curved\allowbreak{}Closure\allowbreak{}A.\allowbreak{}.\allowbreak{}E}, \texttt{Pillar5Synthesis}, \texttt{SITStage3\allowbreak{}Curved\allowbreak{}ADMHypersurface\allowbreak{}Closure} & $181$ strict roots: $36$-dim ADM kinematics, $13\times 13$ DeWitt--\texttt{Chol3} SPD Legendre bijection, \texttt{nd\allowbreak{}Htot\_\allowbreak{}cov\_\allowbreak{}Lie\_\allowbreak{}equiv}, curved $\{\mathcal{H}_\perp,\mathcal{H}_\perp\}$, $\{\mathcal{D},\mathcal{H}_\perp\}$, $\{\mathcal{D},\mathcal{D}\}$, $\{\mathcal{H}_\perp,G\}$, $\{\mathcal{D},G\}$ closure, $36 - 2\times 10 = 16$ physical DOF, and \texttt{stage3\_\allowbreak{}phase3\_\allowbreak{}and\_\allowbreak{}phase4\_\allowbreak{}unified\_\allowbreak{}closure}. \\
\hline
\textbf{Stage 4 Phase 0} & $8$ \texttt{Stage4Data/*} modules + $34$ \texttt{Stage4Pure/*} modules + \texttt{Stage4\allowbreak{}Pure.\allowbreak{}Stage4\allowbreak{}Pure\allowbreak{}Audit} (plus isolated \texttt{LegacyDataEquivBridge}) & $218$ audited pure declarations ($216$ unique): tactic-free data extraction eliminating transitive \texttt{native\_decide} and bare \texttt{import Mathlib} across normal-order degrees~1--4, sparse momentum/phase, deficit-4 adjoints, and all $5$ \texttt{E1C4A} split-projector cases. \\
\hline
\textbf{Stage 4 Phase 1} & \texttt{Stage4\allowbreak{}Lattice.\allowbreak{}Public\allowbreak{}Gauge\allowbreak{}And\allowbreak{}Wightman}, \texttt{PublicWilsonLattice}, \texttt{ConcreteGaugeGroups}, \texttt{Concrete\allowbreak{}Gauge\allowbreak{}Compact\allowbreak{}Haar}, \texttt{ConcreteU1Gauge}, \texttt{ConcreteSU2Gauge}, \texttt{SITCoherence\allowbreak{}Lattice\allowbreak{}Defs}, \texttt{SITStage4\allowbreak{}Challenge\allowbreak{}Spec}, \texttt{Stage4Phase1Audit} & Apache-2.0 Wilson/Wightman adaptation to Lean~4.30, normalized Haar probability measures on compact \texttt{GaugeGroup}, concrete $U(1)$ and $SU(2)$ witnesses, coupled lattice action/prior definitions, and SHA-256-frozen comparator target \texttt{SITStage4JointTarget}. \\
\hline
\textbf{Stage 4 Phase 2} & \texttt{Stage4\allowbreak{}Lattice.\allowbreak{}Coupled\allowbreak{}Measure\allowbreak{}Basics}, \texttt{C09\allowbreak{}Finiteness\allowbreak{}And\allowbreak{}Normalization}, \texttt{LocalGaugeInvariance}, \texttt{Broken\allowbreak{}Phase\allowbreak{}EFTUnification}, \texttt{Stage4Phase2Audit} & Targets~\texttt{(1)}--\texttt{(3)}: \texttt{c09\_\allowbreak{}nonperturbative\_\allowbreak{}finiteness\_\allowbreak{}holds} ($0 < Z_{\mathrm{SIT}} \le 1$ even at $f_{\mathrm{gate}}(\rho,0)=0$), \texttt{local\_\allowbreak{}nonabelian\_\allowbreak{}gauge\_\allowbreak{}invariance\_\allowbreak{}holds}, and \texttt{broken\_\allowbreak{}phase\_\allowbreak{}eft\_\allowbreak{}unification\_\allowbreak{}holds} (Elitzur + $O(\varepsilon^2)$ YM trace limit). \\
\hline
\textbf{Stage 4 Phase 3} & \texttt{Stage4\allowbreak{}Lattice.\allowbreak{}Reflection\allowbreak{}Positivity\allowbreak{}Split}, \texttt{Reflection\allowbreak{}Positivity\allowbreak{}Energy\allowbreak{}Split}, \texttt{Reflection\allowbreak{}Positivity\allowbreak{}Kernel}, \texttt{Reflection\allowbreak{}Positivity\allowbreak{}Kernel\allowbreak{}Factor}, \texttt{Reflection\allowbreak{}Positivity\allowbreak{}OS}, \texttt{Haar\allowbreak{}Coherence\allowbreak{}Gating\allowbreak{}Floor}, \texttt{Polyakov\allowbreak{}Center\allowbreak{}Symmetry}, \texttt{Polyakov\allowbreak{}Holonomy\allowbreak{}Transition}, \texttt{OpenAIUVBridge}, \texttt{Stage4Phase3Audit} & Targets~\texttt{(4)}--\texttt{(7)}: \texttt{os\_\allowbreak{}reflection\_\allowbreak{}positivity\_\allowbreak{}holds} ($2L_t$ half-integer time reflection \& PSD \texttt{slabTransferKernel}), \texttt{haar\_\allowbreak{}coherence\_\allowbreak{}gating\_\allowbreak{}floor\_\allowbreak{}holds} ($48 M \beta L_f\sum\|\psi\|^2$), \texttt{compact\_\allowbreak{}polyakov\_\allowbreak{}first\_\allowbreak{}order\_\allowbreak{}transition\_\allowbreak{}holds} ($(4/27)c u_*^3$ cubic barrier), and \texttt{conditional\_\allowbreak{}openai\_\allowbreak{}uv\_\allowbreak{}bridge\_\allowbreak{}holds}. \\
\hline
\textbf{Stage 4 Phase 4} & \texttt{Stage4\allowbreak{}Matter.\allowbreak{}Spinor\allowbreak{}G2\allowbreak{}Phase\allowbreak{}Only}, \texttt{HydrogenHolonomyTower}, \texttt{Chiral\allowbreak{}Overlap\allowbreak{}SMInterface}, \texttt{Stage4\allowbreak{}Phase4\allowbreak{}Matter\allowbreak{}Audit}, \texttt{SITStage4Solution}, \texttt{SITStage4\allowbreak{}Pure\allowbreak{}Import\allowbreak{}Clean\allowbreak{}Audit}, \texttt{SITStage4MasterAudit} & Targets~\texttt{(8)}--\texttt{(10)} \& 4-Stage Closure: \texttt{spinor\_\allowbreak{}g2\_\allowbreak{}phase\_\allowbreak{}only\_\allowbreak{}holds} ($2\times 2$ Pauli product, algebraic $g_{\mathrm{tree}}=2$ Zeeman check, $\Lambda_{\mathrm{UV}}^{-2}$ dipole bound), \texttt{hydrogen\_\allowbreak{}phase\_\allowbreak{}holonomy\_\allowbreak{}tower\_\allowbreak{}holds} ($2\pi n$ iff + supplied $E_1/n^2$ tower monotonicity + $\alpha_{\mathrm{eff}}$ contraction), \texttt{chiral\_\allowbreak{}overlap\_\allowbreak{}sm\_\allowbreak{}interface\_\allowbreak{}holds} (finite-matrix GW chirality premise + $U(6)$ SM + conditional bounded $W_{\mathrm{ferm}}$), \texttt{sit\_\allowbreak{}stage4\_\allowbreak{}joint\_\allowbreak{}target\_\allowbreak{}holds}, and $93$-module Stage-4 pure-import closure. \\
\hline
\textbf{Stage 5, 5.1 \& 5.2} & \texttt{Stage5\allowbreak{}Extensions.\allowbreak{}Coupled\allowbreak{}OSReflection\allowbreak{}Positivity}, \texttt{Stage5\allowbreak{}Extensions.\allowbreak{}Off\allowbreak{}Diagonal\allowbreak{}Curved\allowbreak{}ADMClosure}, \texttt{Stage5\allowbreak{}Extensions.\allowbreak{}Unified\allowbreak{}AMGMGated\allowbreak{}OSPositivity}, \texttt{Stage5\allowbreak{}Extensions.\allowbreak{}Matter\allowbreak{}And\allowbreak{}ADMDomain\allowbreak{}Refinements}, \texttt{SITStage5\allowbreak{}Pure\allowbreak{}Import\allowbreak{}Clean\allowbreak{}Audit}, \texttt{SITStage5MasterAudit} & $153 + 227 + 98 + 66 = 544$ extension declarations across $6$ compiled Stage-5 modules: Model~L2 \texttt{coupled\_\allowbreak{}os\_\allowbreak{}reflection\_\allowbreak{}positivity\_\allowbreak{}2\allowbreak{}Lt} [$153$]; pointwise 3-jet unconstrained 6-comp.\ off-diagonal $\{\mathcal{D},\mathcal{H}_\perp\}$ modulo $\partial_k B^k$ \texttt{offdiag\_\allowbreak{}curved\_\allowbreak{}adm\_\allowbreak{}DLie\_\allowbreak{}Hperp\_\allowbreak{}closure} [$227$]; unified Model~L3 AM--GM-gated zero-matter ($\delta_0\otimes\mu_{\mathrm{temp}}$ and conditional $\psi=0$) Haar + conditional L1 link equivalence + coupled $2L_t$ OS positivity + unbounded/Gaussian prior exhaustion \texttt{amgm\_\allowbreak{}gated\_\allowbreak{}os\_\allowbreak{}reflection\_\allowbreak{}positivity\_\allowbreak{}unbounded\_\allowbreak{}prior\_\allowbreak{}2\allowbreak{}Lt} [$98$]; noncommutative $[\Pi_i,\Pi_j]$ Clifford $g=2$ (\texttt{dirac\_\allowbreak{}clifford\_\allowbreak{}pauli\_\allowbreak{}derives\_\allowbreak{}g\_\allowbreak{}eq\_\allowbreak{}two}), circular Coulomb + $2\pi n$ Bohr spectrum (\texttt{coulomb\_\allowbreak{}holonomy\_\allowbreak{}derives\_\allowbreak{}bohr\_\allowbreak{}spectrum}), non-zero Neuberger $D_{\mathrm{ov}}\ne 0$ + massive $\det(D^\dagger D+m_f^2 I)$ (\texttt{massive\allowbreak{}Fermion\allowbreak{}Det\allowbreak{}Weight\_\allowbreak{}constructs\_\allowbreak{}sm\_\allowbreak{}interface}), Sylvester cone \texttt{IsSPD3} (\texttt{s\_\allowbreak{}lorentz\_\allowbreak{}det\_\allowbreak{}pos\_\allowbreak{}not\_\allowbreak{}spd}, \texttt{isSPD3\_iff\_posDef}, \texttt{spd\_\allowbreak{}offdiag\_\allowbreak{}curved\_\allowbreak{}adm\_\allowbreak{}closure\_\allowbreak{}and\_\allowbreak{}energy\_\allowbreak{}positivity}), and centered-plaquette $O(a^2)$ AM--GM continuum kinetic matching (\texttt{amgm\_\allowbreak{}exp\_\allowbreak{}profile\_\allowbreak{}cosh\_\allowbreak{}minus\_\allowbreak{}one}, \texttt{amgm\_\allowbreak{}exp\_\allowbreak{}profile\_\allowbreak{}quadratic\_\allowbreak{}limit}, \texttt{amgm\_\allowbreak{}centered\_\allowbreak{}2jet\_\allowbreak{}second\_\allowbreak{}deriv\_\allowbreak{}cancels}, \texttt{amgm\_\allowbreak{}hypercubic\_\allowbreak{}six\_\allowbreak{}plane\_\allowbreak{}sum}, \texttt{amgm\_\allowbreak{}induced\_\allowbreak{}kinetic\_\allowbreak{}matrix\_\allowbreak{}properties}, \texttt{amgm\_\allowbreak{}three\_\allowbreak{}eft\_\allowbreak{}scaling\_\allowbreak{}regimes}) [$66$]; $98$-module pure-import closure ($64$ \texttt{\#print axioms} checks) and $640$-module Stage~1--5 master closure ($23$ \texttt{\#print axioms} checks across $660$ \texttt{.lean} files). \\
\hline
\textbf{Stage 6 \& Addenda} & \texttt{Stage6\allowbreak{}Functional\allowbreak{}Response.\allowbreak{}Action\allowbreak{}Identity\allowbreak{}Diagnostics}, \texttt{Exact\allowbreak{}Plaquette\allowbreak{}Hessian\allowbreak{}And\allowbreak{}Fourier}, \texttt{Smooth\allowbreak{}Torus\allowbreak{}AMGM\allowbreak{}Functional\allowbreak{}Convergence}, \texttt{Continuum\allowbreak{}Regularity\allowbreak{}And\allowbreak{}Lattice\allowbreak{}Embedding}, \texttt{Unchanged\allowbreak{}Action\allowbreak{}Weak\allowbreak{}Field\allowbreak{}Observable}, \texttt{SITStage6\allowbreak{}Audit}, \texttt{Stage6\allowbreak{}Addendum.\allowbreak{}Periodic\allowbreak{}Lattice\allowbreak{}Fourier\allowbreak{}Eigenvalue}, \texttt{Nonnegative\allowbreak{}Field\allowbreak{}Coercivity}, \texttt{Action\allowbreak{}Response\allowbreak{}Propagator\allowbreak{}Identity} & $30$ frozen Stage-6 $+$ $18$ addendum theorems ($48/48$ \texttt{\#print axioms} strictly standard; $6 + 3 = 9$ new \texttt{.lean} files $\implies 669$ total across Stages~1--6+): multiplicative-vs-conformal scalar obstruction ($A^2=A^4=f\Rightarrow f=1$), canonical $\chi=e^{-z/2}X$ and $U(1)$ kinetic-mixing residue $1/z_v$ [$6$]; exact 4-corner AM--GM Hessian with $j_i$ cancellation, phase invariance, and 4D Fourier symbol sandwich [$6$]; exact periodic lattice Fourier eigenvalue equations on \texttt{Site $l$} [$3$]; sharp $\frac{3\beta}{4}$ zero-division nonnegative-field coercivity $Q_l(F)\ge \frac{3\beta}{4} a^2\sum_{x,\mu}(\sqrt{F(x+e_\mu)}-\sqrt{F(x)})^2$ allowing zero corners with both exact equality witnesses (single-site spike for $l \ge 2$ and strictly positive parity-alternating checkerboard for even $l$, strictly improving the valid weaker $\frac{\beta}{2}$ bound from the sharp $\frac{1}{12}$ six-pair inequality) [$10$]; canonical propagator inverse $(p^2 K + H)^{-1} = V\,\mathrm{diag}(1/(p^2+m_a^2))\,V^\top$, source--detector residue sum, $D=\mathrm{diag}(1,\psi_0)$ coordinate-rescaling invariance, and direct instantiation on \texttt{Auth\allowbreak{}Model\allowbreak{}Params} proving $0 < m_0^2 \le m_1^2$ and automatic $p^2 + m_a^2 \ne 0$ for $p^2 \ge 0$ [$5$]; explicit $B_4, C_2$ smooth-torus $O(1/n^2)$ convergence [$6$]; \texttt{ContDiff} $\mathbb{R}\;4$ regularity bridge and exact Stage-5 finite-lattice embedding at $\kappa_t=0$ with linear/affine field realizations [$7$]; and unchanged-action two-mode Yukawa and conditional E\"ot-Wash slab non-exclusion certificates [$5$]. \\
\hline
\end{longtable}
}

\setcounter{subsection}{0}
\renewcommand{\thesection}{}
\renewcommand{\thesubsection}{.\arabic{subsection}}
\section*{Further Reading}
\addcontentsline{toc}{section}{Further Reading}

\textbf{Journal Articles}

\noindent Blumberg, Micah (2025). \emph{Super Dark Time}. Figshare Journal Contribution.\\DOI: \href{https://doi.org/10.6084/m9.figshare.28284545}{https://doi.org/10.6084/m9.figshare.28284545}.

\medskip
\noindent Blumberg, Micah (2025). \emph{Micah's New Law of Thermodynamics: A Signal-Dissipation Framework for Equilibrium and Consciousness}. Figshare Journal Contribution.\\DOI: \href{https://doi.org/10.6084/m9.figshare.28264340}{10.6084/m9.figshare.28264340}.

\medskip
\noindent Blumberg, Micah (2025). \emph{Super Information Theory}. Figshare Journal Contribution.\\DOI: \href{https://doi.org/10.6084/m9.figshare.28379318}{https://doi.org/10.6084/m9.figshare.28379318}.

\medskip
\noindent \textbf{Books}

\noindent Blumberg, Micah (2025). \emph{Bridging Molecular Mechanisms and Neural Oscillatory Dynamics: Explore how synaptic modulation and pattern generation create the brain's seamless volumetric three-dimensional conscious experience}. \\Available online at Amazon:\\ \href{https://www.amazon.com/dp/B0DL4701875}{https://www.amazon.com/dp/B0DL4701875}, ASIN: B0DL4701875.

\medskip
\noindent These platforms expand upon the foundational research, offering broader perspectives on consciousness frameworks and computational neuroscience.

\medskip

\noindent \textbf{Related News Stories}
\noindent SVGN.io News features many articles with similar content from the same author as this paper: \href{https://www.svgn.io/p/a-new-book-out-today-bridging-molecular}{https://www.svgn.io/p/a-new-book-out-today-bridging-molecular}.

\medskip

\noindent \textbf{Self Aware Networks Online Archive:} Comprehensive time-stamped notes and original research materials spanning over a decade are available in the Self Aware Networks GitHub repository.

This archive provides detailed documentation of the evolution and refinement of foundational theories, including Super Dark Time (also previously referred to as Quantum Gradient Time Crystal Dilation and Dark Time Theory),

Micah's New Law of Thermodynamics, Neural Array Projection Oscillation Tomography (NAPOT), and Self Aware Networks theory of mind.\\Accessible at: \href{https://github.com/v5ma/selfawarenetworks}{https://github.com/v5ma/selfawarenetworks}.

\medskip
\noindent \textbf{The Neural Lace Podcast:} Explore discussions and analyses regarding consciousness, neuroscience advancements, neural synchronization, EEG-to-WebVR integration, and theoretical physics. The podcast content provides further insight into the conceptual background and implications of the theories presented in the cited works.\\Find episodes of the Neural Lace Podcast via this old link: \href{http://vrma.io}{http://vrma.io}

\medskip

\noindent \textbf{Supplementary Websites and Resources:} Further materials, related projects, and additional context for the research presented can be accessed via the following websites: \href{http://selfawareneuralnetworks.com}{selfawareneuralnetworks.com}, \href{http://selfawarenetworks.com}{selfawarenetworks.com}

\subsection{Influential Voices}
Warm thanks to authors, writers, scientists, mathematicians, or theorists like: Steven Strogatz, 
Peter Ulric Tse, Stephen Wolfram, Gy\"orgy Buzs\'aki, Dario Nardi, Luis Pessoa, Grace Lindsay, 
Oliver Sacks, Michael Graziano, Michael Levin, Michael Gazzaniga, Jeff Hawkins, Nicholas Humphrey, 
Valentino Braitenberg, Douglas Hofstadter, David Eagleman, Olaf Sporns, Jon Lieff, Donald Hebb, 
Netta Engelhardt, Ivette Fuentes, Jacob Bekenstein, Sabine Hossenfelder, Albert Einstein, John Bell, 
David Bohm, Roger Penrose, Earl K. Miller, Eugene Wigner, Basil J. Hiley, John von Neumann, 
Louis de Broglie, Claude Shannon, Alan Turing, Norbert Wiener, Santiago Ramon y Cajal, Eric Kandel, 
Antonio Damasio, Stanislas Dehaene, Patricia Churchland, Christof Koch, Francis Crick, 
Karl Friston, Niels Bohr, Erwin Schr\"odinger, Max Planck, Leonard Susskind, Kip Thorne, 
Freeman Dyson, David Chalmers, Daniel Dennett, Paul Dirac, Isaac Newton, and others for their 
foundational insights into computation, oscillatory synchronization, and higher-order cognition. 
Their work has significantly shaped the wave-computational perspective laid out here or at least 
influenced my thinking on the topic.

\setcounter{subsection}{0}
\section*{Inspirations}
\addcontentsline{toc}{section}{Inspirations}
For the 41st draft of the paper. The author thanks Yu Deng, Zaher Hani, and Xiao Ma for their groundbreaking work rigorously deriving macroscopic irreversibility from classical microscopic dynamics. In addition public discussions on social networks by many individuals have inspired the new clarification on the gauge-theoretic structure of quantum phase, and the new Network Formulation of Super Information Theory. 

The development and empirical framing of Super Information Theory (SIT) has been enriched by interdisciplinary dialogue with the broader communities of mathematical physics, neuroscience, and artificial intelligence, including open discussions on preprint servers and science journalism platforms. The author acknowledges the influence of foundational works in quantum information, gauge theory, thermodynamics, and neural computation, as cited throughout the text.

Constructive feedback from online forums, technical correspondents, and peer reviewers has greatly contributed to the rigor and clarity of this manuscript. Any errors or omissions are the sole responsibility of the author.

\subsection{Influencial Works}
\begin{itemize}
\item Strogatz, S. H. (2003). Sync: The Emerging Science of Spontaneous Order. Hyperion.
\item Strogatz, S. H. (1994). Nonlinear Dynamics and Chaos: With Applications to Physics, Biology, Chemistry, and Engineering. Westview Press.
\item Tse, P. U. (2013). The Neural Basis of Free Will: Criterial Causation. MIT Press
\item Wolfram, S. (2002). A New Kind of Science. Wolfram Media.
\item Buzs\'aki, G. (2006). Rhythms of the Brain. Oxford University Press.
\item Buzs\'aki, G., \& Draguhn, A. (2004). Neuronal oscillations in cortical networks. Science, 304(5679), 1926--1929.
\item Levin, M. (2021). Biological Information and the Problem of Regeneration: Complex Signaling Pathways in Morphogenesis. Nature Reviews Molecular Cell Biology.
\item Hawkins, J., \& Blakeslee, S. (2004). On Intelligence. Times Books.
\item Hawkins, J. (2021). A Thousand Brains: A New Theory of Intelligence. Basic Books.
\item Humphrey, N. (2011). Soul Dust: The Magic of Consciousness. Princeton University Press.
\item Friston, K. (2010). The free-energy principle: A unified brain theory? Nature Reviews Neuroscience, 11(2), 127--138.
\item Engelhardt, N., \& Wall, A. C. (2019). Decoding the black hole interior. Journal of High Energy Physics, 2019(1), 1--20.
\item Von Neumann, J. (1955). Mathematical Foundations of Quantum Mechanics. Princeton University Press.
\item Penrose, R. (1989). The Emperor's New Mind: Concerning Computers, Minds, and the Laws of Physics. Oxford University Press.
\item Shannon, C. E. (1948). A mathematical theory of communication. Bell System Technical Journal, 27(3), 379--423.
\item Crick, F., \& Koch, C. (2003). A framework for consciousness. Nature Neuroscience, 6(2), 119--126.
\item Schr\"odinger, E. (1944). What Is Life? The Physical Aspect of the Living Cell. Cambridge University Press.
\item Hossenfelder, S. (2020). Superdeterminism: A Guide for the Perplexed. arXiv preprint arXiv:2010.01324.
\item Donadi, S., \& Hossenfelder, S. (2020). A Superdeterministic Toy Model. arXiv preprint arXiv:2010.01327.
\item Susskind, L. (2008). The Black Hole War: My Battle with Stephen Hawking to Make the World Safe for Quantum Mechanics. Little, Brown.
\item Hofstadter, D. R. (1979). G\"odel, Escher, Bach: An Eternal Golden Braid. Basic Books.
\item Hofstadter, D. R. (2007). I Am a Strange Loop. Basic Books.

\item[] \textbf{Neuroscience and Free Energy Principle}
\item Adams, R. A., Shipp, S., \& Friston, K. J. (2013). 
\item Predictions not commands: Active inference in the motor system. Brain Structure and Function, 218(3), 611--643.
\item Seth, A. K., \& Friston, K. J. (2016). Active interoceptive inference and the emotional brain. Philosophical Transactions of the Royal Society B: Biological Sciences, 371(1708), 20160007.
\item Rao, R. P., \& Ballard, D. H. (1999). Predictive coding in the visual cortex: A functional interpretation of some extra-classical receptive-field effects. Nature Neuroscience, 2(1), 79--87.
\item[] \textbf{Quantum Mechanics and Thermodynamics}
\item Rovelli, C. (2015). Relational quantum mechanics: A simple explanation. Quantum Physics Letters, 11(3), 5--12.
\item Bekenstein, J. D. (1973). Black holes and entropy. Physical Review D, 7(8), 2333--2346.
\item Zurek, W. H. (2003). Decoherence, einselection, and the quantum origins of the classical. Reviews of Modern Physics, 75(3), 715--775.
\item[] \textbf{Cognitive Science and Artificial Intelligence}
\item Pezzulo, G., \& Cisek, P. (2016). Navigating the affordance landscape: Feedback control as a process model of behavior and cognition. Trends in Cognitive Sciences, 20(6), 414--424.
\item Hassabis, D., Kumaran, D., Summerfield, C., \& Botvinick, M. (2017). Neuroscience-inspired artificial intelligence. Neuron, 95(2), 245--258.
\item Silver, D., Hubert, T., Schrittwieser, J., et al. (2018). A general reinforcement learning algorithm that masters chess, shogi, and Go through self-play. Science, 362(6419), 1140--1144.
\item[] \textbf{Neural Dynamics and Oscillations}
\item Fries, P. (2005). A mechanism for cognitive dynamics: Neuronal communication through neuronal coherence. Trends in Cognitive Sciences, 9(10), 474--480.
\item Buzs\'aki, G., \& Draguhn, A. (2004). Neuronal oscillations in cortical networks. Science, 304(5679), 1926--1929.
\item Singer, W. (1999). Neuronal synchrony: A versatile code for the definition of relations? Neuron, 24(1), 49--65.
\item[] \textbf{Philosophy and Theoretical Perspectives}
\item Dennett, D. C. (1991). Consciousness Explained. Little, Brown and Company.
\item Penrose, R. (1994). Shadows of the Mind: A Search for the Missing Science of Consciousness. Oxford University Press.
\item Humphrey, N. (2011). Soul Dust: The Magic of Consciousness. Princeton University Press.
\item[] \textbf{Predictive Coding and Bayesian Models}
\item Clark, A. (2015). Surfing uncertainty: Prediction, action, and the embodied mind. Oxford University Press.
\item Shipp, S. (2016). Neural elements for predictive coding. Frontiers in Psychology, 7, 1792.
\item Knill, D. C., \& Richards, W. (1996). Perception as Bayesian Inference. Cambridge University Press.
\item[] \textbf{Miscellaneous but Relevant}
\item Thagard, P. (2005). Mind: Introduction to Cognitive Science. MIT Press.
\item Wolfram, S. (2002). A New Kind of Science. Wolfram Media.
\item Nardi, D. (2011). The 16 Personality Types: Descriptions for Self-Discovery. Radiance House.
\item Miller, E. K., \& Cohen, J. D. (2001). An integrative theory of prefrontal cortex function. Annual Review of Neuroscience, 24, 167--202.
\item Levin, M. (2021). Bioelectric signaling: Reprogrammable circuits underlying embryogenesis, regeneration, and cancer. Cell, 184(8), 1971--1989.
\item Braitenberg, V. (1984). Vehicles: Experiments in Synthetic Psychology. MIT Press.
\item Hebb, D. O. (1949). The Organization of Behavior: A Neuropsychological Theory. Wiley.
\item Bohm, D. (1952). A suggested interpretation of the quantum theory in terms of "hidden" variables I. Physical Review, 85(2), 166--179.
\item Kuramoto, Y. ``Self-entrainment of a population of coupled non-linear oscillators,'' in \emph{International Symposium on Mathematical Problems in Theoretical Physics}, Lecture Notes in Physics, 39, Springer (1975).
\item Kuramoto, Y. (1984). Chemical Oscillations, Waves, and Turbulence. Springer.
\item Murray, J. D. \emph{Mathematical Biology: I. An Introduction.} Springer (2002).
\item Arfken, G. B. \& Weber, H. J. \emph{Mathematical Methods for Physicists.} Academic Press (2012).
\item Barab\'asi, A.-L. \emph{Network Science.} Cambridge University Press (2016).

\item[] \textbf{Mathematical Physics, Lattice Gauge Theory, and Formal Verification (2024--2026)}
\item Berry, M. V. (1984). Quantal phase factors accompanying adiabatic changes. \emph{Proceedings of the Royal Society of London A}, 392(1802), 45--57.
\item Simon, B. (1983). Holonomy, the quantum adiabatic theorem, and Berry's phase. \emph{Physical Review Letters}, 51(24), 2167--2170.
\item Deng, Y., Hani, Z., \& Ma, X. (2025). Hilbert's sixth problem: derivation of fluid equations via Boltzmann's kinetic theory. Preprint.
\item Finster, F., et al. (2024--2025). Effective dynamical collapse and spacetime as the web of correlations in Causal Fermion Systems. Preprints.
\item Adelberger, E. G., Heckel, B. R., \& Nelson, A. E. (2009). Torsion-balance experiments: A low-energy frontier of particle physics. \emph{Annual Review of Nuclear and Particle Science}, 59, 117--149.
\item Bertotti, B., Iess, L., \& Tortora, P. (2003). A test of general relativity using radio links with the Cassini spacecraft. \emph{Nature}, 425, 374--376.
\item Craig, N. (2026). A gauge-invariant measure for lattice chiral fermions. \texttt{arXiv:2610.09093}.
\item Yang--Mills Formalization Oversight Repository (\texttt{yang-mills-group/YM}) (2026). \texttt{README.md}, \texttt{checks.md}, and independent assessment (\texttt{assessment-1.pdf}) of \texttt{OpenAI-YM-20261028}: continuum Wightman construction for pure 4D Wilson gauge theory (infinite-volume uniqueness and mass gap unresolved).
\item Hexagon Collaboration (2026a). \emph{The spinless modular bootstrap at large central charge}. \texttt{hexagon:2611.00174v1}, \url{https://hexagonmath.org/pdf/2611.00174v1}.
\item Hexagon Collaboration (2026b). \emph{The Hawking--Page transition at weak coupling}. \texttt{hexagon:2611.00175v1}, \url{https://hexagonmath.org/pdf/2611.00175v1}.
\item Lee, J. G., Adelberger, E. G., Cook, T. S., Fleischer, S. M., \& Heckel, B. R. (2020). New test of the gravitational $1/r^2$ law at separations down to $52\,\mu\mathrm{m}$. \emph{Physical Review Letters}, 124, 101101.
\item Lee, J. G. (2020). \emph{A Fourier-Bessel Test of the Gravitational Inverse-Square Law}. Ph.D.\ Dissertation, University of Washington (Chapter~11, Table~11.4 single-Yukawa fit parameters and Appendix~A corrected run-by-run harmonic torque data, geometries, calibrations, and constrained nuisance parameters).
\item Takamoto, M., Ushijima, I., Ohmae, N., Yahagi, T., Kokado, K., Shinkai, H., \& Katori, H. (2020). Test of general relativity by a pair of transportable optical lattice clocks. \emph{Nature Photonics}, 14, 411--415.
\item Touboul, P., et al.\ (MICROSCOPE Collaboration) (2022). MICROSCOPE mission: Final results of the test of the Equivalence Principle. \emph{Physical Review Letters}, 129, 121102.
\item Asenbaum, P., Overstreet, C., Kim, M., Curti, J., \& Kasevich, M. A. (2020). Atom-interferometric test of the Equivalence Principle at the $10^{-12}$ level. \emph{Physical Review Letters}, 125, 191101.
\item NIST/BACON Collaboration (2026). Optical clock frequency ratios with uncertainties at or below $3.2\times 10^{-18}$ across $\mathrm{Al}^+$, $\mathrm{Yb}$, and $\mathrm{Sr}$ via phase-stabilized optical fiber links. \emph{Physical Review Letters}, 137, 033201 (published July~14, 2026; DOI \texttt{10.1103/g865-9mk1}; \texttt{arXiv:2512.21428}; NIST institutional posting April~10, 2026).
\item Blumberg, M. (2026). \emph{Super Information Theory: Coherence, Time Density, Informational Dissipation, and Gravity} (Paper~002, July~2026 Version~80 Compact Corrective Edition). Self Aware Networks Research Institute, July~29, 2026; Zenodo DOI \texttt{10.5281/zenodo.21693867}.
\item Blumberg, M. (2026). \emph{Level~0 SIT: A Physical-Frame Model and Condensate-Source Falsification Contract} (Paper~003); \emph{Super Dark Time: From Local Time-Rate Gradients to a Testable Scalar-Tensor Bridge} (Paper~007); \emph{Designing a Coherence-Gravity Test with Bose--Einstein Condensates} (Paper~014); \emph{Cross-Disciplinary SIT Translation Guide} (Paper~028); \emph{A Constructive Coherence-Tower Theorem for a $U(1)$ Lattice Toy Model of Super Information Theory} (Paper~033); \emph{SuperTimePosition and Gravitational Phase-Lock} (Paper~034); \emph{Causal Fermion Systems and Super Information Theory} (Paper~048, citing Version~80 Draft~13 dated July~26, 2026); \emph{QGTCD Constraints Note} (Paper~055); \emph{Physics in Review} comparisons (Papers~062--063). Self Aware Networks Research Institute, July~2026.
\item Blumberg, M. (2026). \emph{Learning to Feel} (Paper~114: initial edition October~17, 2026; v1.2 dated October~20, 2026; technical supplement intentionally retaining v1.1; independent neuroscience track). Self Aware Networks Research Institute, October~2026.
\item Blumberg, M. (2026). \emph{A Classical Continuum Limit for a Four-Corner Lattice Energy} (Standalone Mathematical Research Paper, Revised First Draft). Self Aware Networks Research Institute, November~18, 2026.

\end{itemize}

\setcounter{subsection}{0}
\section*{License}
\addcontentsline{toc}{section}{License}
CC BY-SA 4.0: Creative Commons Attribution-ShareAlike
This license allows reusers to distribute, remix, adapt, and build upon the material in any medium or format, so long as attribution is given to the creator. The license allows for commercial use. If you remix, adapt, or build upon the material, you must license the modified material under identical terms.
% --- END INPUT sections/12_appendix_AB_and_backmatter.tex ---

\end{document}
